\documentclass[11pt]{article}
\usepackage[T1]{fontenc}
\usepackage[utf8]{inputenc}
\usepackage{lmodern}
\usepackage[margin=1in]{geometry}
\usepackage{amsmath,amssymb,amsthm,mathrsfs}
\usepackage{enumitem}
\usepackage{microtype}
\usepackage[hidelinks]{hyperref}

\newcommand{\Q}{\mathbb Q}
\newcommand{\C}{\mathbb C}
\newcommand{\Z}{\mathbb Z}
\newcommand{\R}{\mathbb R}
\newcommand{\Qbar}{\overline{\mathbb Q}}
\newcommand{\CH}{\operatorname{CH}}

\newtheorem{theorem}{Theorem}[section]
\newtheorem{lemma}[theorem]{Lemma}
\newtheorem{proposition}[theorem]{Proposition}
\newtheorem{corollary}[theorem]{Corollary}
\numberwithin{equation}{section}
\hypersetup{pdftitle={Weil classes and Hodge classes on abelian powers},pdfauthor={OpenAI}}

\begin{document}
\title{Weil classes and Hodge classes on abelian powers}
\author{OpenAI}
\date{October 6, 2026}
\maketitle

\begin{abstract}
We prove the rational Hodge conjecture in every codimension on every
self-power of a complex abelian sixfold with an imaginary-quadratic
action and a compatible polarization whose rational homological Hermitian
form is hyperbolic of signature \((3,3)\). We also prove it in every
codimension on every self-power of a complex abelian variety of dimension
at most five admitting an imaginary-quadratic action. Both results include
nonsimple varieties and special periods with additional endomorphisms.
The proof uses the companion theorem on the rational Hodge conjecture for
CM abelian varieties.
\end{abstract}

\begingroup
\small
\tableofcontents
\endgroup
\clearpage

\section{Introduction}
\label{p3:sec:introduction}

For a smooth projective complex variety \(X\), the rational Hodge
conjecture asks whether every class in
\[
 H^{2p}(X,\Q)\cap H^{p,p}(X)
\]
is the class of a rational algebraic cycle of codimension \(p\).
For an abelian variety, the question on every self-power is naturally a
question about tensors in its first cohomology: classes can couple several
copies even when the classes on the variety itself are already understood.
This paper proves the conjecture on all self-powers in two families with an
imaginary-quadratic action.

An action of an imaginary quadratic field \(E\) on an abelian variety
\(A\) means a unital embedding \(E\hookrightarrow\operatorname{End}^0(A)\).
Write \(E=\Q(u)\), with \(u^2=-d\) and \(d\in\Z_{>0}\). A rational Riemann form
\(e\) is compatible with the action if
\(e(ax,y)=e(x,\bar a y)\) for every \(a\in E\). Its associated homological
Hermitian form is
\[
 h_E(x,y)=e(x,uy)+u e(x,y),\qquad x,y\in H_1(A,\Q),
\]
conjugate-linear in the first variable and linear in the second. A rank-six
Hermitian form is hyperbolic if it contains a totally isotropic
three-dimensional \(E\)-subspace. The form in the next theorem is the one
coming from the actual polarization, not merely an abstract form of the
same complex signature.

\begin{theorem}[Split Weil sixfolds]
\label{p3ab:split-sixfolds}
Let \(E\) be an imaginary quadratic field and let \(A\) be a complex
abelian sixfold with an \(E\)-action.  Suppose a compatible
polarization has hyperbolic rational Hermitian form of signature
\((3,3)\).  Then the rational Hodge conjecture holds in every
codimension on \(A^m\), for every \(m\geq1\).
\end{theorem}

\begin{theorem}[Quadratic actions in dimension at most five]
\label{p3ab:quadratic-five}
Let \(A\) be a complex abelian variety of dimension at most five
admitting an action of an imaginary quadratic field.  The rational
Hodge conjecture holds in every codimension on \(A^m\), for every
\(m\geq1\).
\end{theorem}

Both statements include nonsimple varieties and periods with additional
endomorphisms. The quadratic action need not be the full endomorphism
algebra or be central on a repeated simple factor. The second theorem
requires no separately chosen compatible polarization: one is obtained by
averaging any rational Riemann form, as shown in
Section~\ref{p3:sec:abelian}.

\subsection{Weil classes and earlier results}

Weil's determinant construction \cite{Weil1979} provides a basic test of the Hodge
conjecture beyond divisor classes. For an abelian variety of dimension
\(2n\) with a compatible imaginary-quadratic action of signature
\((n,n)\), the two embedding determinants form a rational plane of
classes of type \((n,n)\). For a general such variety with \(n>1\),
these classes are not generated by divisors; see
\cite[Sections~4.9--4.11]{vanGeemenWeilNotes}.
Early geometric constructions used automorphisms of curves and Prym
varieties \cite{Schoen1988}. Schoen's addendum corrected the polarization
hypothesis in the earlier fourfold statement and established algebraicity
of the Weil plane for every fourfold with balanced action by
\(\Q(\sqrt{-3})\), without a discriminant restriction
\cite{Schoen1998Addendum}.

Markman's secant-sheaf construction proves algebraicity of the whole Weil
plane for every split Weil sixfold. The preprint states the result for
discriminant \(-1\), while the published account uses the equivalent
split condition on the rational Hermitian form of the compatible
polarization \cite[Theorem~1.5.1]{Markman2025Secant}
\cite[Theorem~1.2]{Markman2026ICM}.
The published account also proves the Hodge conjecture for every
individual abelian variety of dimension at most five, with no
quadratic-action hypothesis \cite[Corollary~1.3]{Markman2026ICM}.
These are all-period statements. The further question in our main
theorems concerns all Hodge classes on every self-power, including the
extra tensors that can occur when the Hodge group drops at a special
period.

The all-powers question also has substantial precedents. Moonen and
Zarhin's low-dimensional classification includes all-powers
divisor-generation results, as well as a fourfold case whose square has
exceptional classes \cite[Theorems~0.1--0.2]{MoonenZarhin1999}.
For every Weil fourfold of discriminant one, Floccari proves the Hodge
conjecture on all powers through the Kuga--Satake correspondence
\cite[Theorems~1.3 and~5.12]{Floccari2026K3}.
Floccari and Fu give another proof using singular varieties of
O'Grady's six-dimensional type \cite[Theorem~1.2]{FloccariFu2026}.
Our argument addresses the two stated families by combining the
split-sixfold Weil plane with a determinant criterion for the actual
Hodge group and the CM algebraicity input specified below.

The geometric proof is a detailed treatment of Markman's construction,
rather than an extension of its period range. It uses Mukai's
Fourier--Mukai duality and the algebraic spin interpretation of
Golyshev--Lunts--Orlov \cite{P3SWMukai1981,GolyshevLuntsOrlov2001}.
We work out the coherent seed, the fixed transform conventions, the full
invariant derived obstruction space, and its passage through quotient
descent and projective-bundle conversion. The full-group comparison and
descent are already part of Markman's method
\cite[Section~7.4.2 and Lemmas~9.3.2 and~9.3.11]{Markman2025Secant}
\cite[Sections~11.3--11.5]{Markman2026ICM}.
The Hochschild--Atiyah comparison follows the character and evaluation
identities of C\u{a}ld\u{a}raru, Buchweitz--Flenner and Huang
\cite{P3HHCaldararuII,P3HHBF,P3HHHuang}, with the abelian specialization
of Todd-corrected compatibility \cite{P3HHCRVdB}.
The deformation argument reconstructs the smooth-projective-germ case
of the coherent-sheaf argument in Buchweitz--Flenner's Theorem~5.1 and
its proof \cite{P3SWBuchweitzFlenner2003}. The needed injectivity is on
the full obstruction group, not merely the image of Hochschild evaluation.

The separate passage from the Weil plane to all powers uses classical
invariant theory to reduce tensors to pairings and determinants
\cite{LehrerZhang2015,LehrerZhang2017,Schrijver2008}.
The two-weight restriction is the standard minuscule argument
\cite[Section~3.2, Lemmas~4--5]{Serre1979HodgeTate}; comparison of product
Hodge groups through graph links has a close predecessor in
Moonen--Zarhin \cite[Section~3.1 and Lemmas~3.3--3.4]{MoonenZarhin1999}.
Adjoining auxiliary factors and extracting lower-dimensional Weil classes
also appears in Schoen's Proposition~10 \cite{Schoen1998Addendum}.
Here CM elliptic factors supply smaller quadratic determinants, while
the actual-group analysis identifies the products selected by the
simultaneous balance equations. Those supply and balance steps, rather
than the invariant spanning theorems alone, give the all-powers conclusions.

\subsection{Proof strategy}

Let \(H=H^1(A,\Q)\). A tensor in \(H^{\otimes 2p}\) is balanced if it
has type \((p,p)\). For \(A^m\), antisymmetrization lifts a Hodge class to
a balanced tensor in \((H^{\oplus m})^{\otimes 2p}\), with the copies
labelled by the factors of the power. Diagonal pullback returns the exterior
class. Section~\ref{p3:sec:tensors} gives a
criterion for algebraicity of these tensors. The multilinear invariant
theorems for symplectic, orthogonal, and special linear groups reduce the
tensors to pairings and top exterior lines
\cite{LehrerZhang2015,LehrerZhang2017,Schrijver2008}. Pairings are divisor
classes. The required products of the remaining lines are supplied by algebraic correspondences
from the full cohomology of CM abelian varieties. A rational Hodge class in
such a supplied subspace lifts to rational Hodge classes on the CM sources;
Theorem~1.1 of \emph{The rational Hodge conjecture for CM abelian
varieties}, dated October 6, 2026, makes those lifts algebraic
\cite{OpenAICM2026}. This is the precise CM input used in this paper.

Section~\ref{p3:sec:abelian} proves the two main theorems from the
Weil-plane input, whose geometric proof follows in
Sections~\ref{p3sw:appendix}--\ref{p3:sec:all-periods}.
Adding CM elliptic curves supplies smaller quadratic determinants whose
signature entries are at most three. The Albert decomposition, the two
Hodge-circle weights, and the small minuscule representations determine
the possible constituents of the actual derived Hodge group. A separate
independence argument identifies their simultaneous product. The balance
equations then cancel opposite determinants or combine them into a
supplied quadratic determinant. This is how the proof accommodates
additional endomorphisms at special periods.

The cycle construction in Section~\ref{p3sw:appendix} begins with a
genus-three Jacobian \(X\). Disjoint
orbits of Abel curves give an ideal sheaf on \(X^2\); a fixed
Fourier--Mukai transform and duality produce a coherent perfect sheaf
\(\mathcal E\) on \(X\times\widehat X\), together with a coherent finite
translation action twisted by line bundles. The normalized Chern character
of this sheaf is the candidate cycle class. Two properties are needed: it
must contain a nonzero Weil component and remain of Hodge type as the Weil
period varies; and, after descent and a projective-bundle construction, an
associated ordinary sheaf must deform so that the normalized character is
algebraic near the initial point.

Section~\ref{p3:sec:obstruction} proves the exact Hochschild--Atiyah trace
identity and its transport for the fixed transform. A direct equivariant
Ext calculation then shows that the evaluation image equals the entire
invariant obstruction space, on which the trace map is injective.
Section~\ref{p3:sec:spin} computes the rational Betti character by a finite
Clifford calculation. It identifies the actual ample polarization, proves
that its homological Hermitian form is hyperbolic, and shows that the
character is Hodge throughout the compatible marked domain.

Section~\ref{p3:sec:deformation} descends the sheaf through the finite
quotient and converts the resulting twisted sheaf into an ordinary sheaf
on a projective bundle without losing any obstruction class. Its
semiregularity theorem produces the required local deformations.
Section~\ref{p3:sec:all-periods} extracts the normalized character by
algebraic projective-bundle operations, spreads the resulting Weil plane through
proper Hilbert parameter spaces, and clears rational lattice denominators
by isogenies. This proves algebraicity of the whole rational Weil plane at
every compatible hyperbolic period.

\section{Algebraic supply and invariant tensors}
\label{p3:sec:tensors}

The cohomology of an abelian variety is generated by its first
cohomology.  We will therefore prove algebraicity first for tensors in
degree one.  Classical invariant theory reduces these tensors to
pairings and top exterior powers.  Pairings are divisor classes.  The
remaining task in each application is to supply the required products
of top exterior powers by algebraic correspondences from CM abelian
varieties.

All varieties are complex, and cohomology is singular cohomology with
the indicated coefficients.  An algebraic correspondence has rational
coefficients unless otherwise stated.  We fix
$k=\Qbar\subset\C$.  An automorphism of $k/\Q$ acts only on
coefficients in the cohomology of the fixed complex variety; it does
not replace that variety by a conjugate.  For a pure Hodge structure
$V$, our convention is
$V(r)^{p,q}=V^{p+r,q+r}$, so the Tate twist $V(r)$ has weight
$\operatorname{wt}(V)-2r$.
A class is \emph{algebraic over $k$} if it belongs to the
$k$-linear span of rational algebraic cycle classes.

\subsection{Correspondences from CM abelian varieties}

For a smooth projective variety $X$, let
\[
 \mathcal C^d(X)=
 \sum_{B,j,\Gamma}
 \operatorname{im}\bigl(
 \Gamma_*:H^{d-2j}(B,\Q)\longrightarrow H^d(X,\Q)\bigr),
 \qquad
 \Gamma\in\CH^{\dim B+j}(B\times X)_\Q ,
\]
where $B$ ranges over CM abelian varieties, including a point, and
$j$ ranges over integers for which the displayed groups and cycles
exist.  Here
and below a CM abelian variety means an abelian variety whose
degree-one Hodge group is a torus.  The displayed map is a morphism of
weight-$d$ Hodge structures from $H^{d-2j}(B,\Q)(-j)$.
We say that a subspace of $H^d(X,k)$ is \emph{CM-supplied} if it is
contained in $\mathcal C^d(X)\otimes_\Q k$.  Thus supply concerns
the full cohomology of the CM source, not only its Hodge classes.

We use the companion theorem that the rational Hodge conjecture holds
for every complex CM abelian variety in every codimension
\cite[Theorem~1.1]{OpenAICM2026}.  This is the only Hodge-conjecture
input in the following lemma.

\begin{lemma}[Operations and Hodge lifts]
\label{p3:cm-operations}
CM supply is preserved by external products, permutations of factors,
cup products, and rational algebraic correspondences.  It is also
preserved by $k$-linear combinations of such correspondences and by
scalar conjugation.  Every rational Hodge class in a CM-supplied
subspace is algebraic.
\end{lemma}

\begin{proof}
External products of correspondence cycles realize products of their
maps, and a product of CM abelian varieties is CM.  Composition is
realized by pulling two cycles to the triple product, intersecting,
and pushing forward.  Graphs give pullback and pushforward.
Diagonal pullback gives cup product; adding integration gives
contraction.
These operations are defined for the smooth projective varieties in
question.  Linearity proves the scalar assertion, and
$\mathcal C^d(X)\otimes k$ is invariant under scalar conjugation
because $\mathcal C^d(X)$ is rational.

Let $\alpha\in H^{2p}(X,\Q)$ have type $(p,p)$ and belong after
scalar extension to $\mathcal C^{2p}(X)\otimes k$.  Membership in
a rational subspace descends under a field extension, so $\alpha$
already belongs to $\mathcal C^{2p}(X)$.  Finitely many of the
defining correspondence maps have an image containing $\alpha$.
They give a morphism
\[
 f:\bigoplus_{\nu=1}^s
 H^{2p-2j_\nu}(B_\nu,\Q)(-j_\nu)
 \longrightarrow H^{2p}(X,\Q)
\]
of polarizable pure Hodge structures.  This morphism splits onto its
image as a rational Hodge morphism: the orthogonal complement of
$\ker f$ for a polarization on the source is a Hodge substructure
and maps isomorphically to $\operatorname{im}f$.  Nondegeneracy on
$\ker f$ follows from the positive Hermitian form obtained using
the Weil operator.  Hence $\alpha$ has a rational Hodge lift.  Its
$\nu$th component before twisting has type
$(p-j_\nu,p-j_\nu)$ on $B_\nu$.  The CM Hodge theorem makes this
component algebraic.  Applying the defining correspondence cycles
and adding proves the assertion.
\end{proof}

\subsection{Degree-one summands and pairings}

We record the coefficient descent needed to pass from complex
eigenspaces to rational algebraic maps.  If $T$ is a finite-dimensional
rational vector space, $P\subset T_\C$ is complex-linear, and
$t\in T_k$ satisfies $\gamma t\in P$ for every
$\gamma\in\operatorname{Gal}(k/\Q)$, then
\begin{equation}
 \label{p3:eq:scalar-descent}
 t\in (T\cap P)\otimes_\Q k.
\end{equation}
Indeed, choose a finite Galois field $K$ containing its coefficients.
For $a\in K$, the trace
$\operatorname{Tr}_{K/\Q}(at)=
\sum_{\tau\in\operatorname{Gal}(K/\Q)}\tau(a)\tau(t)$
is rational and belongs to $P$.  If $b_j$ is a basis of $K$ and
$b_j^*$ is its trace-dual basis, these traces for $a=b_j^*$
are exactly the rational coefficients of $t$ in the basis $b_j$.
Here each $\tau$ is extended to $k$ when applying the hypothesis.
This proves \eqref{p3:eq:scalar-descent}.  Conditions that a tensor
have a specified Hodge type, or that a map preserve Hodge types, are
complex-linear conditions, so the observation applies to them.

Let $X$ be an abelian variety or a smooth connected projective curve.
A subspace $W\subset H^1(X,k)$ is \emph{admissible} if it is the
image of a $k$-linear idempotent $e$ such that every $\gamma e$
preserves the Hodge decomposition after extension to $\C$.
A $k$-linear map between such spaces, or between finite external
direct sums of them, is \emph{admissible} if it preserves the Hodge decomposition
after every scalar conjugation.  We write $\overline W$ for the
scalar complex conjugate of $W$.

A \emph{pair tensor} between $W$ and $W'$ is a tensor
$t\in W\otimes W'$ whose associated map $W^*\to W'$ is an
isomorphism and for which every $\gamma t$ has Hodge type $(1,1)$.
For a tensor on $W$ itself, symmetry or alternation refers to the
ordinary interchange of tensor factors, before the Koszul sign in
product cohomology.

\begin{lemma}[Algebraic maps and contractions]
\label{p3:admissible-operations}
Every admissible map is induced on the indicated summands by a
$k$-linear combination of rational algebraic correspondences.
Every pair tensor is algebraic over $k$, and contraction by its
inverse pairing is an algebraic cohomological operation.
If $W$ is admissible, so is $\overline W$, and there is a pair
tensor between $W$ and $\overline W$.
\end{lemma}

\begin{proof}
For an admissible map, compose with the source projector and the
target inclusion to obtain a map on the ambient degree-one spaces.
Equation~\eqref{p3:eq:scalar-descent} expresses it as a $k$-linear
combination of rational Hodge maps.  A rational Hodge map between
the first cohomologies of abelian varieties is induced by a rational
homomorphism in the contravariant direction: dualize to homology,
multiply to preserve the lattices, and use the resulting complex-linear
map of complex tori.  A homomorphism of abelian varieties is algebraic.
For a curve target use its Abel map to its Jacobian.  For a curve
source, the cup pairing identifies a degree-one Hodge map with a
rational $(1,1)$ class on the product, which is a divisor class by
the Lefschetz $(1,1)$ theorem.  These arguments also apply componentwise
to external direct sums.  In particular an admissible projector
commutes with the scalar extension of the rational Hodge group.

A pair tensor lies in the $H^1\otimes H^1$ summand of degree two
on the product.  Scalar descent and the Lefschetz $(1,1)$ theorem
make it algebraic over $k$.  Let $b:W\otimes W'\to k$ be its
inverse pairing, with its bidegree $(-1,-1)$ understood.  Extend
$b$ by the admissible projectors.  For a rational polarization
pairing $\psi_X$ on the source $H^1(X,\Q)$ there is a unique map
$u:H^1(X',k)\to H^1(X,k)$ satisfying
$b(x,y)=\psi_X(x,u(y))$.  In every scalar conjugate, $b$ pairs
only complementary Hodge degrees; nondegeneracy of $\psi_X$
therefore makes $u$ an admissible map.  The first part algebraizes
$u$.  The remaining contraction by $\psi_X$ is cup product, an
ample power if needed, and integration.  It is algebraic.

The conjugate projector shows that $\overline W$ is admissible.
The restriction of $\psi_X$ to $W\times\overline W$ is perfect.
After extension to $\C$, the positive Hermitian form associated
with the polarization detects every nonzero vector of the Hodge-stable
space $W_\C$ by pairing it with the conjugate of a vector of
$W_\C$.  Its inverse gives a tensor in $W\otimes\overline W$.
For an arbitrary $\gamma$, conjugate the nonzero determinant of
this restricted pairing.  The resulting pairing on
$\gamma W\times\gamma\overline W$ is still perfect and pairs
complementary Hodge degrees, since both spaces are Hodge-stable
and $\psi_X$ is rational.  Its inverse has type $(1,1)$.
This argument does not assume that $\gamma$ commutes with complex
conjugation.
\end{proof}

For an admissible space $W$ of dimension $r$, write
\[
 D(W)=\bigwedge^r W\subset W^{\otimes r}
       \subset H^r(X^r,k)
\]
for its top alternating line on separate tensor slots; put $D(0)=k$
on a point.  Ordinary permutations differ from geometric permutations
only by rational Koszul signs.  Direct-sum factorization and
cancellation of these determinant lines can therefore be performed
algebraically.  More precisely, if an admissible isomorphism identifies
$W$ with $\bigoplus_j W_j$, alternation and the admissible
inclusions and projections induce nonzero maps in both directions
between $D(W)$ and $\bigotimes_jD(W_j)$.  If $W,W'$ have a
pair tensor, the determinant of the inverse pairing contracts
$D(W)\otimes D(W')$ nontrivially to the unit line.  Conversely,
alternating the appropriate tensor power of the pair tensor supplies
$D(W)\otimes D(W')$ algebraically.  For an alternating self-pair
on a $2r$-space, the $r$th exterior power supplies $D(W)$
algebraically.  These statements follow in bases from the nonzero
determinant, or Pfaffian, of the pair tensor and from
Lemma~\ref{p3:admissible-operations}.

\subsection{The tensor criterion}

For a polarizable rational Hodge structure $H$ of weight one
whose only types are $(1,0)$ and $(0,1)$, let
$J=\operatorname{Hdg}(H)$ be the smallest rational algebraic
subgroup of $\operatorname{GL}(H)$ containing the Hodge circle,
which acts on $H^{1,0}$ and $H^{0,1}$ by $z$ and $\bar z$
for $|z|=1$.  It lies in the symplectic group of the polarization.
A tensor in $H^{\otimes 2p}$ is called \emph{balanced} if it
has Hodge type $(p,p)$.
The group $J$ is connected: its identity component is rational and contains
the connected circle, so minimality applies.  It is reductive:
a rational invariant subspace is a Hodge substructure, and
polarizability supplies a rational Hodge complement, which is
again invariant by minimality.  The faithful representation on
$H$ is thus completely reducible.

\begin{lemma}[Hodge tensors determine the Hodge group]
\label{p3:tensor-stabilizer}
Inside the symplectic group of $H$, the group
$\operatorname{Hdg}(H)$ is the pointwise stabilizer of all rational
tensors of type $(a,a)$ in $H^{\otimes 2a}$, for all $a\ge0$.
\end{lemma}

\begin{proof}
By minimality, the rational $J$-invariant tensors are precisely the
rational tensors fixed by the Hodge circle.  To see that these determine
$J$, use Chevalley's theorem to choose a rational representation $T$
of $\operatorname{GL}(H)$ and a line $L\subset T$ whose stabilizer
is $J$ \cite[Theorem~4.27]{MilneAG2022}.  Reductivity gives a
$J$-equivariant projector $e:T\to L$.  Rational representations
of $\operatorname{GL}(H)$ are subquotients of sums of mixed tensor
powers of $H$ \cite[Theorem~4.14]{MilneAG2022}.  Complete reducibility
for $\operatorname{GL}(H)$ \cite[Theorem~22.42]{MilneAG2022}
gives a $\operatorname{GL}(H)$-equivariant embedding of
$\operatorname{End}(T)$ as a direct summand of a finite sum of
mixed tensor powers.  The element $e$ in this representation
is $J$-invariant.
Use the polarization to replace dual tensor factors by ordinary
factors, equivariantly for the ambient symplectic group.  The components
of $e$ are then rational tensors fixed by the circle, hence balanced
tensors in even degrees.  An element fixing every such tensor fixes
$e$, preserves its image $L$, and therefore belongs to $J$.
The reverse inclusion follows from minimality.
\end{proof}

We also use the standard effective weight-one equivalence between toric Hodge
groups and CM abelian varieties, including rational summands.  The following
form ensures an actual CM source.

\begin{lemma}[Toric degree-one summands]
\label{p3:toric-cm-factor}
Let $A$ be a complex abelian variety and let $H\subset H^1(A,\Q)$ be a
rational Hodge summand.  If $\operatorname{Hdg}(H)$ is a torus, then
$H$ is isomorphic as a rational Hodge structure to $H^1(B,\Q)$ for a
CM isogeny factor $B$ of $A$.
\end{lemma}

\begin{proof}
Choose a rational Hodge projector onto $H$.  By the degree-one
realization in Lemma~\ref{p3:admissible-operations}, it is induced by
a rational endomorphism of $A$.  After clearing denominators, its
image is an abelian isogeny factor $B$ with $H^1(B,\Q)\simeq H$.
Its Hodge group is a torus, so $B$ is CM under our convention.

We also describe the embedding-line decomposition used below.
Put $J(H)=\operatorname{Hdg}(H)$.  Its envelope in
$\operatorname{End}_\Q(H)$ is a commutative
semisimple algebra $E_0$: after splitting the torus it is an algebra
of diagonal scalars.  The polarization adjoint preserves $E_0$
and is positive, since these scalars commute with the Weil operator.
A positive involution on a real finite \'etale algebra fixes each
primitive real summand and is complex conjugation on each complex
summand.  Exchanging idempotents would give $e^*e=0$, and the
identity involution on a complex summand would give $i^*i=-1$,
contradicting positivity.  A real summand is impossible here:
the connected compact circle can act on it only trivially, whereas
weight one has only the characters $z,\bar z$.  Thus the rational
field factors of $E_0$ are CM fields.  Splitting the associated
$E_0$-modules into field lines gives rank-one CM Hodge structures.
The projectors onto these field lines are rational Hodge maps.
The same degree-one realization gives abelian isogeny factors of
$B$ with toric Hodge groups.  For a rank-one module over a CM field
$E$, the idempotents of
$E\otimes_\Q k=\prod_{\sigma:E\hookrightarrow k}k$
give one-dimensional embedding lines.  Scalar conjugation
permutes these lines, and each is pure of type $(1,0)$ or
$(0,1)$.
\end{proof}

For admissible $W$, define its imbalance at $\gamma$ by
\[
 \epsilon_\gamma(W)=
 \dim_\C(\gamma W_\C\cap H^{1,0})
 -\dim_\C(\gamma W_\C\cap H^{0,1}).
\]
The determinant line $\gamma D(W)$ is pure of bidegree given by
these two dimensions.  The perfect complementary pairing in
Lemma~\ref{p3:admissible-operations} gives
\begin{equation}
 \label{p3:eq:opposite-imbalance}
 \epsilon_\gamma(\overline W)=-\epsilon_\gamma(W)
 \quad\text{for every }\gamma.
\end{equation}
A self-pair forces equal Hodge dimensions, hence even rank and
imbalance zero.  Abstract self-duality of an $\operatorname{SL}_2$
module alone does not supply an all-conjugate $(1,1)$ self-pair;
such a block may still need to be treated as linear.  These identities remain valid
without commuting $\gamma$ with complex conjugation.

\begin{theorem}[A determinant criterion for all powers]
\label{p3:tensor-criterion}
Let $A$ be a complex abelian variety, $H=H^1(A,\Q)$,
$J=\operatorname{Hdg}(H)$, and $R=J^{\mathrm{der}}$.
Suppose the following data exist.

\begin{enumerate}[label=\textnormal{(\alph*)}]
\item A rational Hodge summand $H_0\subset H$ has toric Hodge group.
Choose a decomposition of $H_{0,k}$ into CM embedding lines,
with repeated occurrences distinguished, and call this finite list
$\mathcal E_0$.
\item The complement of $H_{0,k}$ is a direct sum of blocks of
nonzero admissible spaces.  In a symplectic or an orthogonal block,
each space is equipped with an admissible isomorphism to a
representative $V_i$ chosen as an actual summand of that block,
with, respectively, an alternating or a
symmetric pair tensor on itself.  A linear block has
$\dim V_i\ge2$, with $V_i$ again an actual summand, and each space is equipped with an admissible
isomorphism to either $V_i$ or $\overline V_i$, paired with
each other as in Lemma~\ref{p3:admissible-operations}.
The pair identifies $\overline V_i$ with $V_i^*(-1)$ as a
Hodge-graded space and with the dual as a group representation.
\item Under these identifications the image of $R_k$ contains
\[
 P=\prod_{i\in I_{\mathrm{Sp}}}\operatorname{Sp}(V_i)
   \times\prod_{i\in I_{\mathrm{SO}}}\operatorname{SO}(V_i)
   \times\prod_{i\in I_{\mathrm{SL}}}\operatorname{SL}(V_i).
\]
Each factor acts by its standard representation on its own copies,
by the dual representation on the opposite copies of a linear
block, and trivially on every other block and on $H_{0,k}$.
The symplectic and orthogonal forms are the inverse pairings.
\item Let $W_1,\ldots,W_s$ be any list, with arbitrary repetitions,
from
\[
 \{V_i:i\in I_{\mathrm{SO}}\}
 \ \cup\
 \{V_i,\overline V_i:i\in I_{\mathrm{SL}}\}.
\]
Whenever there are $\ell_1,\ldots,\ell_r\in\mathcal E_0$, again
with arbitrary repetitions, such that
\begin{equation}
 \label{p3:eq:determinant-balance}
 \sum_{j=1}^s\epsilon_\gamma(W_j)
 +\sum_{a=1}^r\epsilon_\gamma(\ell_a)=0
 \quad\text{for every }\gamma\in\operatorname{Gal}(k/\Q),
\end{equation}
the line $\bigotimes_{j=1}^sD(W_j)$ is CM-supplied on its
separate tensor slots.  Empty lists are allowed.
\end{enumerate}

Then for every $p\ge0$ every rational balanced tensor in $H^{\otimes 2p}$ is the
class of a rational algebraic cycle on its separate $A$-slots.
The rational Hodge conjecture holds in every codimension of $A^m$
for every $m\ge1$.

Call the data \emph{saturated} if every singleton $D(V_i)$ for
an orthogonal or linear block is CM-supplied, without any balance
condition.  For saturated data and any CM abelian variety $M$,
the induced data on $A\times M$ are saturated.  In particular the
same tensor and Hodge-conjecture conclusions hold for every power of
$A\times M$.
\end{theorem}

The list of CM lines in \eqref{p3:eq:determinant-balance} is chosen
once and then tested at every scalar conjugation.  This is weaker
than requiring individual determinant supply: special periods in
Section~\ref{p3:sec:abelian} will require only certain products of
column determinants.  Saturation is useful when arbitrary CM
companions are to be added.

\begin{proof}
We use the multilinear first fundamental theorems in their all-degree
form.  For a nondegenerate bilinear form $b$ on $V$, let
$t_b\in V\otimes V$ be its inverse-form, or coevaluation, tensor.
The $\operatorname{Sp}(V)$ invariants in $V^{\otimes d}$
are spanned by permutations of products of $t_b$ for the
alternating form and vanish for odd $d$
\cite[Theorem~4.1 and Corollary~4.2]{LehrerZhang2015}.
For $\operatorname{SO}(V)$, products of $t_b$ for the symmetric
form and a top alternating tensor span all invariants
\cite[Theorem~2.1 and Corollary~4.2]{LehrerZhang2017}.
For $\operatorname{SL}(V)$ in mixed powers of $V,V^*$, the
generators are the identity coevaluation tensor
$\sum_j e_j\otimes e_j^*\in V\otimes V^*$, independent of
the chosen basis, and the top alternating tensors
of $V$ and $V^*$ \cite[Section~5]{Schrijver2008}.
These are spanning statements in every degree; stable-range bounds
concern independence, not spanning.  They hold over $k$ by scalar
extension, since invariants are kernels of linear maps defined over
$k$.  For the product $P$, decompose each tensor slot into its
block summands and apply these statements separately to each factor.

We need one elementary simultaneous-balance observation.  Suppose a
subspace $E\subset T_k$ of a rational pure weight-$2p$ Hodge
structure is spanned by finitely many lines $L_a$, each pure in
Hodge bidegree after every scalar conjugation.  A rational vector
$\alpha\in E$ of type $(p,p)$ is spanned by those $L_a$ whose
every conjugate has type $(p,p)$.  To prove this, for fixed
$\gamma$ group the original lines according to the bidegree of
$\gamma L_a$.  Their group sums form a direct sum, as is seen by
applying $\gamma$ and using the Hodge decomposition.  Let
$q_\gamma$ be the projection retaining the group of type $(p,p)$.
Since $\gamma\alpha=\alpha$, this projection fixes $\alpha$.
Every $L_a$ is an eigenline of every $q_\gamma$ with eigenvalue
zero or one.  The projections therefore commute even if the spanning
lines are dependent.  Only finitely many distinct projections occur,
because each is determined by the subset of the finite list it retains.
Their product fixes $\alpha$ and has image spanned by exactly the
lines retained by every conjugation.  This proves the observation.

Fix a tensor degree.  A rational balanced Hodge tensor is fixed by
$J$, hence by $P$.  The invariant theorems show that the
$P$-invariant subspace is spanned by finitely many lines formed,
up to slot permutations and the prescribed admissible identifications,
from three types of factors:
\[
 \text{lines in }\mathcal E_0,\qquad
 \text{pair tensors},\qquad
 D(V_i)\text{ for }i\in I_{\mathrm{SO}},
 \quad D(V_i),D(\overline V_i)\text{ for }i\in I_{\mathrm{SL}}.
\]
Every such line is pure after every scalar conjugation.  This is part
of the definition for pairs, follows from the top exterior power of
a Hodge-graded space for determinants, and holds for CM embedding
lines.  The simultaneous-balance observation reduces a rational
balanced tensor to generator lines balanced in all conjugates.
Discarding the pair factors from the bidegree count leaves exactly
\eqref{p3:eq:determinant-balance}, with the actual CM lines in
the chosen generator.  Assumption (d) supplies its determinant
product.  Lemma~\ref{p3:toric-cm-factor} supplies the toric summand
by a CM isogeny factor, and the pair tensors and identifications are algebraic by
Lemma~\ref{p3:admissible-operations}.  Lemma~\ref{p3:cm-operations}
therefore places the entire generator line in the scalar extension
of the rational CM-supplied subspace.  Rational membership descends,
and the same lemma makes the original tensor algebraic.

For $A^m$, first cohomology is $H^{\oplus m}$.  The same argument
works with slots labelled by the copies: it merely repeats the old
blocks, and assumption (d) already permits arbitrary repetitions.
Finally
$H^{2p}(A^m,\Q)=\bigwedge^{2p}H^1(A^m,\Q)$.
Normalized antisymmetrization gives a rational Hodge lift of a Hodge
class into the ordinary tensor power.  The tensor result makes that
lift algebraic on separate slots.  Pullback along the diagonal,
equivalently cup product of the slots, returns the exterior class.

Suppose now that the data are saturated.  Scalar conjugation supplies
the opposite determinants, and products supply every determinant
list, so (d) holds.  Let $T$ be the toric Hodge group of
$H^1(M,\Q)$ and let $J'$ be the Hodge group of
$H\oplus H^1(M,\Q)$.  It is a connected reductive subgroup of
$J\times T$ whose projection onto $J$ is surjective by minimality.
Its derived subgroup consequently projects onto $R$ and trivially
onto $T$.  Thus, inside this representation,
\[
 (J')^{\mathrm{der}}=R\times\{1\}.
\]
The product $P$ still acts on the old blocks and trivially on the
enlarged toric summand $H_0\oplus H^1(M,\Q)$.  The old singleton
determinants remain supplied after pullback along $A\times M\to A$
on their slots.  These are saturated data on
$A\times M$, and the proved assertions apply.  If all of $H$ is
toric, there are no classical blocks and the only determinant list
is empty, so this case is included.
\end{proof}

\subsection{Passing from supplied cohomology to powers}

The next elementary transfer passes from complete cohomological supply
to Hodge classes on arbitrary products.  Its hypothesis concerns products
of the source varieties, not merely each source separately.

\begin{lemma}[Correspondence transfer]
\label{p3:correspondence-transfer}
Let $\mathscr S$ be a collection of smooth projective varieties such
that the rational Hodge conjecture holds on every finite product of
members of $\mathscr S$.  Let $Y_1,\ldots,Y_v$ be smooth projective
varieties.  Suppose that, for each $i$ and each degree $q$,
$H^q(Y_i,\Q)$ is spanned by images of maps
\[
 \Gamma_*:H^{q-2j}(S,\Q)\longrightarrow H^q(Y_i,\Q),
 \qquad
 \Gamma\in\CH^{\dim S+j}(S\times Y_i)_\Q,\quad S\in\mathscr S.
\]
Then the rational Hodge conjecture holds on every
finite product of the $Y_i$, with arbitrary repetitions and with
any additional members of $\mathscr S$.

The same conclusion holds if the spanning is given over $k$ by
restrictions of maps induced by $k$-linear combinations of
rational algebraic correspondences, provided their images span
the full target degree after scalar extension.
\end{lemma}

\begin{proof}
Fix a product $Y$ allowed by the conclusion.
In each degree, finite dimensionality makes a finite sum of the
stipulated maps surjective.  Tensor these maps and the identity maps
of any additional source factors.  The K\"unneth decomposition shows
that, in each degree of the target product, a finite sum of maps from
finite source products is surjective.  One uses the indicated
K\"unneth summands only to see that the images span; no algebraic
K\"unneth projector on a source is required.  External products and
composition make all these maps algebraic correspondences.
For a target degree $2p$, write $j_\nu$ for the sum of the individual
shift parameters of one such external-product map, so its cohomological
shift is $2j_\nu$, and write $S_\nu$ for its source
product.  The surjection is a morphism
\[
 \bigoplus_\nu H^{2p-2j_\nu}(S_\nu,\Q)(-j_\nu)
 \longrightarrow H^{2p}(Y,\Q)
\]
of polarizable pure Hodge structures of weight $2p$.
The splitting argument in Lemma~\ref{p3:cm-operations} lifts a
rational target Hodge class.  Its $\nu$th component before
twisting has type $(p-j_\nu,p-j_\nu)$ on $S_\nu$.
It is algebraic by the source hypothesis, so its image is algebraic.

For the scalar variant, expand each scalar correspondence into
finitely many rational correspondence cycles.  Its image on any
specified source subspace is contained in the scalar extension of
the sum of the full images of those rational maps.  If the scalar
images span the full target degree, the quotient by that rational
image sum becomes zero after scalar extension and was already
zero over $\Q$.  The first part therefore applies.
\end{proof}

For later use, if $A$ has saturated data, then the collection of
varieties $A^a\times M$, where $a\ge0$ and $M$ is any CM abelian
variety, satisfies the source hypothesis of
Lemma~\ref{p3:correspondence-transfer}.  A finite product is
$A^{\sum a}\times\prod M$, an abelian factor of a sufficiently high
power of $A\times\prod M$, and its Hodge classes transfer by the
algebraic factor inclusion and projection.  For a curve $C$ with
$A=\operatorname{Jac}(C)$, the Abel map supplies $H^1(C)$ from
$A$, while a point and the fundamental class supply the even
cohomology.  If a smooth projective variety $F$ has
$\mathcal C^q(F)=H^q(F,\Q)$ in every degree, finite
dimensionality expresses each degree by finitely many rational
correspondence images from CM abelian varieties.  External
products then supply the full cohomology of every product of
copies of $C$ and such varieties $F$ from varieties
$A^a\times M$ in the preceding source collection.
Lemma~\ref{p3:correspondence-transfer} gives the corresponding
Hodge-conjecture conclusion.

\section{Quadratic actions at every period}
\label{p3:sec:abelian}

We prove Theorems~\ref{p3ab:split-sixfolds} and~\ref{p3ab:quadratic-five}
by applying the determinant criterion of
Theorem~\ref{p3:tensor-criterion} to the Hodge group at the given period.
The cycle input is algebraicity of the Weil plane on every hyperbolic
Weil sixfold.  We state it below and defer its geometric proof to
Sections~\ref{p3sw:appendix}--\ref{p3:sec:all-periods}.

\subsection{The Weil-plane input}

An action of an imaginary quadratic field on an abelian variety will
always mean a unital embedding
\[
 E\hookrightarrow\operatorname{End}^0(A).
\]
The action need not be the full endomorphism algebra, and it need not
be central on an isotypic power.  We first specify the polarization
convention in the six-dimensional result.

Write \(E=\Q(u)\), where \(u^2=-d\) and \(d\in\mathbb Z_{>0}\), and put
\(V=H_1(A,\Q)\).  A rational Riemann form \(e\) for a polarization
is \emph{compatible} with \(E\) if
\[
 e(ax,y)=e(x,\bar a y)\qquad(a\in E).
\]
In particular \(e(ax,ay)=\operatorname{Nm}_{E/\Q}(a)e(x,y)\),
the norm-pullback condition for a polarized variety of Weil type.
It then defines the Hermitian form, conjugate-linear in the first
variable,
\begin{equation}
 \label{p3ab:eq:hermitian}
 h_E(x,y)=e(x,uy)+u e(x,y)\in E.
\end{equation}
This is the usual homological convention
\cite[Lemma~5.2]{vanGeemenWeilNotes}.  The adjoint identity and
alternation show directly that \(h_E\) is Hermitian and that
\[
 e(x,y)=\operatorname{Tr}_{E/\Q}\bigl((2u)^{-1}h_E(x,y)\bigr);
\]
in particular it is nondegenerate.  For an embedding
\(\sigma:E\hookrightarrow k\), write \(H^1(A,k)_\sigma\) for
the simultaneous \(\sigma\)-eigenspace of the pullback \(E\)-action.
It is admissible: its spectral projector is
\[
 \frac{u^*-\bar\sigma(u)\operatorname{id}}
      {\sigma(u)-\bar\sigma(u)},
\]
and every scalar conjugate is a Hodge endomorphism.
If \(\sigma(u)=i\sqrt d\), the signature of \(h_E\) is \((p,q)\), where
\[
 p=\dim_\C\bigl((H^1(A,k)_\sigma)_\C\cap H^{1,0}(A)\bigr),
 \qquad
 q=\dim_\C\bigl((H^1(A,k)_{\bar\sigma})_\C\cap H^{1,0}(A)\bigr).
\]
Here \(p+q=\dim A\).  On \(V_{\mathbb R}\) viewed via \(\sigma\)
as an \(E\otimes_\Q\mathbb R\simeq\C\)-space, the
complex structure \(I\) has eigenspaces with eigenvalues \(i,-i\).
On them \(u=\sqrt d\,I\) and \(u=-\sqrt d\,I\), respectively, so
\(h_E(x,x)=e(x,ux)\) has the indicated positive and negative signs
by the Riemann positivity \(e(x,Ix)>0\).  The eigenspaces are
Hermitian-orthogonal because \(h_E\) is \(I\)-invariant.  The corresponding
holomorphic \(\sigma\)- and \(\bar\sigma\)-forms have dimensions
\(p\) and \(q\).  A nondegenerate rank-\(2b\) Hermitian form is called
\emph{hyperbolic} if it has a totally isotropic \(E\)-subspace of
dimension \(b\).

Every quadratic action admits a compatible polarization.  Starting
with a rational Riemann form \(e_0\), set
\begin{equation}
 \label{p3ab:eq:averaging}
 e(x,y)=e_0(x,y)+d^{-1}e_0(ux,uy).
\end{equation}
Both summands satisfy Riemann positivity because \(u\) is an invertible
complex-linear map.  Direct substitution gives
\(e(ux,y)=-e(x,uy)\).  A positive integral multiple clears
denominators and is the Riemann form of an ample polarization.

For a sixfold, \(H^1(A,\Q)\) has dimension twelve and is an
\(E\)-space of rank six.  Thus its two spectral eigenspaces each
have dimension six.  The direct sum
\[
 \bigwedge\nolimits^6 H^1(A,k)_\sigma
 \ \oplus\
 \bigwedge\nolimits^6 H^1(A,k)_{\bar\sigma}
 \ \subset H^6(A,k)
\]
is stable under every scalar conjugation: scalar conjugation merely
permutes the two spectral eigenspaces and their exterior powers.
Finite-dimensional Galois descent therefore gives a unique rational
two-dimensional subspace \(\mathcal W_E(A)\subset H^6(A,\Q)\) whose
scalar extension is this sum.  When \(h_E\) has signature \((3,3)\),
each determinant line has type \((3,3)\).

\begin{theorem}[Algebraicity of the hyperbolic Weil plane]
\label{p3sw:weil-plane}
Let \(E\) be an imaginary quadratic field.  Let \(A\) be a complex
abelian sixfold with a unital \(E\)-action and a compatible polarization
whose actual homological Hermitian form \(h_E\) is hyperbolic of signature
\((3,3)\).  Then
\[
 \mathcal W_E(A)\ \subset\
 \operatorname{im}\bigl(\CH^3(A)_\Q
                    \longrightarrow H^6(A,\Q)\bigr).
\]
The assertion holds at every compatible period, including nonsimple
varieties and periods with additional endomorphisms.
\end{theorem}

Markman states this theorem in the discriminant-\(-1\) formulation
\cite[Theorem~1.5.1]{Markman2025Secant}; his published account also
states it for all split Weil sixfolds
\cite[Theorem~1.2]{Markman2026ICM}.  The later sections give the
secant-sheaf and spin proof with the fixed conventions used here.

The theorem supplies the total \(E\)-determinant of a split sixfold.
At a special period that determinant can split into several column
determinants for the larger endomorphism algebra.  Our task is to show
that the products required by Theorem~\ref{p3:tensor-criterion} still
have CM supply.  We first obtain determinant supply by adding CM
elliptic factors, then determine the actual small-dimensional derived
groups, and finally check the determinant products.

\subsection{Determinants supplied by a split sixfold}

For a Hermitian form \(h\), write \([\det h]\) for its determinant
class in \(\Q^\times/\operatorname{Nm}_{E/\Q}(E^\times)\).
Changing an \(E\)-basis multiplies the determinant by a norm.

\begin{lemma}[A rational hyperbolicity test]
\label{p3ab:hyperbolic}
Let \(E\) be imaginary quadratic.  A nondegenerate \(E\)-Hermitian
form of rank \(2b\) and signature \((b,b)\), \(b\geq1\), is
hyperbolic if and only if
\[
 [\det h]=[(-1)^b].
\]
Every hyperbolic form of this rank is isometric to the sum of \(b\)
planes with matrix
\(\left(\begin{smallmatrix}0&1\\1&0\end{smallmatrix}\right)\).
\end{lemma}

\begin{proof}
For necessity, let \(L\) be a totally isotropic \(b\)-space and
choose \(0\ne v\in L\).  Choose \(w\) with
\(h(v,w)=1\) and replace \(w\) by
\(w-\tfrac12h(w,w)v\).  The two vectors now span a hyperbolic
plane.  The projection of \(L\) to the orthogonal complement of
this plane has isotropic dimension \(b-1\): for \(x\in L\), replace \(x\) by
\(x-h(w,x)v\).  Induction gives the displayed sum of planes,
each with determinant \(-1\).  This proves necessity and the
isometry assertion.

Conversely, \(x\mapsto h(x,x)\) is a nondegenerate rational quadratic
form of dimension \(4b\) and signature \((2b,2b)\).  Nondegeneracy
follows from the trace pairing and trace duality.  When \(b\geq2\),
Meyer's theorem gives a nonzero isotropic vector
\cite[Corollary~3.3.4]{PetersSterk2024QuadraticForms}.
Split off its hyperbolic plane.  The complement has signature
\((b-1,b-1)\) and determinant class \([(-1)^{b-1}]\), so one
may repeat until rank two.  Write the last form diagonally as
\(\langle a,b'\rangle\), with \(a,b'\in\Q^\times\).  Its
determinant condition says \(-ab'=\operatorname{Nm}(t)\) for some
\(t\in E^\times\).  The vector \((t/a,1)\) is isotropic because
\[
 a\operatorname{Nm}(t/a)+b'=0.
\]
It splits the last hyperbolic plane.
\end{proof}

The following construction adjoins auxiliary factors and then removes
their determinant lines by algebraic pairings. Schoen
\cite[Proposition~10]{Schoen1998Addendum} uses a related padding and
extraction argument for Weil fourfolds; here the auxiliary factors are
CM elliptic curves and the output is determinant supply on separate slots.

\begin{lemma}[Quadratic determinant supply]
\label{p3ab:determinant-supply}
Let \(A\) be a complex abelian variety with an action of an imaginary
quadratic field \(E\), and let a compatible rational Riemann form
have associated Hermitian form of signature \((p,q)\).
Put \(n=\dim A=p+q\).  Both lines
\[
 D\bigl(H^1(A,k)_\sigma\bigr),\qquad
 D\bigl(H^1(A,k)_{\bar\sigma}\bigr)
\]
are CM-supplied in either of the following cases:
\begin{enumerate}
\item \(n=6\), \((p,q)=(3,3)\), and the form is hyperbolic;
\item \(n<6\) and \(p,q\leq3\).
\end{enumerate}
The second assertion permits every compatible rational Riemann form.
Both assertions hold at every period.
\end{lemma}

\begin{proof}
Theorem~\ref{p3sw:weil-plane} proves that the rational Weil plane is
algebraic on every polarized imaginary-quadratic sixfold whose actual
homological Hermitian form is hyperbolic of signature \((3,3)\).
Its convention is exactly \eqref{p3ab:eq:hermitian}, so it applies directly
to the first case, including special periods.  It supplies the rational
Weil two-plane, not all Hodge classes at such a period.

For a sixfold \(X\) in the first case, the definition of
\(\mathcal W_E(X)\) shows that algebraicity of this
rational plane makes each of its two embedding determinant lines
algebraic over \(k\).  To pass to the
determinant on separate slots, let \(\mu:X^6\to X\) be addition.
For an integer \(t\geq2\), the interpolation polynomial
\[
 \pi_1^X=
 \prod_{\substack{0\leq r\leq2\dim X\\ r\ne1}}
 \frac{[t]^*-t^r\operatorname{id}}{t-t^r}
\]
is an algebraic cohomological projector to \(H^1(X)\), since
\([t]^*\) acts by \(t^r\) on \(H^r(X)\).  For a basis
\(v_1,\ldots,v_6\) of either embedding space, expanding
\(\mu^*v=\sum_i\operatorname{pr}_i^*v\) shows that
\[
 (\pi_1^X)^{\boxtimes6}\mu^*
 (v_1\wedge\cdots\wedge v_6)
 =
 \sum_{\tau\in S_6}\operatorname{sgn}(\tau)
 v_{\tau(1)}\otimes\cdots\otimes v_{\tau(6)}.
\]
The right side is a nonzero vector in the top alternating line.
Thus both determinant lines in the first assertion are algebraic
over \(k\), and hence are CM-supplied from a point.

Now start with an arbitrary \(A\) in the second assertion.  Add
\(3-p\) rank-one \(E\)-Hodge structures of signature \((1,0)\)
and \(3-q\) of signature \((0,1)\).  They are realized by CM
elliptic curves \(C_1,\ldots,C_{6-n}\) with the chosen or the
conjugate \(E\)-action.  In an \(E\)-basis their Hermitian
coefficients may have any positive rational magnitude and the
prescribed signs: positive rational rescaling of an elliptic
polarization leaves the Hodge structure unchanged.

The orthogonal sum on
\(X=A\times C_1\times\cdots\times C_{6-n}\) has signature
\((3,3)\).  Its determinant \(d_0\) in chosen bases is negative.
At least one line has been added, so multiply one added coefficient
by \(-1/d_0>0\).  The new determinant is \(-1\), and
Lemma~\ref{p3ab:hyperbolic} makes the form hyperbolic.  Clearing
denominators scales the form by a positive rational number; in rank
six this changes the determinant by a norm and preserves the required
discriminant.  The first part therefore supplies the determinant of
this actual product.

Put \(L_j=H^1(C_j,k)_\sigma\).  The \(\sigma\)-space of \(X\) is
\[
 H^1(A,k)_\sigma\oplus L_1\oplus\cdots\oplus L_{6-n}.
\]
Project its determinant to successive blocks of slots belonging to
these summands.  The alternation of a concatenated basis maps to a
nonzero multiple of
\[
 D\bigl(H^1(A,k)_\sigma\bigr)\otimes
 L_1\otimes\cdots\otimes L_{6-n}.
\]
These block projections are induced by the product inclusions and
projections, so are algebraic.  Each opposite line \(\bar L_j\)
is supplied by the identity on the CM curve \(C_j\).
Tensoring with these lines and contracting the perfect elliptic
polarization pairings removes all \(L_j\).  The contractions are
nonzero and algebraic by Lemma~\ref{p3:admissible-operations}.
Lemma~\ref{p3:cm-operations} now supplies the required determinant.
Scalar conjugation supplies its opposite.  The construction began
with the fixed \(A\), so no period-generality condition was used.
\end{proof}

\subsection{Columns and their pairings at the given period}

We first decompose first cohomology into admissible irreducible
summands and determine their pairings.  The endomorphism algebra
in this calculation is the actual full algebra at the given period.
We will then determine the derived image on each summand and prove
that the resulting images act independently, as required by
Theorem~\ref{p3:tensor-criterion}.

Let \(B\) be a simple non-CM abelian variety, \(g=\dim B\), and set
\[
 H_B=H^1(B,\Q),\qquad
 J_B=\operatorname{Hdg}(H_B),\qquad
 R_B=J_B^{\mathrm{der}},\qquad
 D=\operatorname{End}_{\mathrm{Hdg}}(H_B).
\]
The degree-one realization in
Lemma~\ref{p3:admissible-operations} identifies \(D\) with the
opposite of \(\operatorname{End}^0(B)\), a division algebra.
It is also the commutant of \(J_B\): the rational centralizer of
an endomorphism commuting with the circle contains \(J_B\) by
minimality.  The group \(J_B\) is connected reductive as in
Section~\ref{p3:sec:tensors}.  If \(Z=Z(D)\) and \([D:Z]=s^2\),
base change of the commutant and complete reducibility give
\begin{equation}
 \label{p3ab:eq:columns}
 D_k=\prod_{\tau:Z\hookrightarrow k}M_s(k),
 \qquad
 H_{B,k}=\bigoplus_{\tau:Z\hookrightarrow k}W_\tau\otimes U_\tau,
 \qquad \dim U_\tau=s.
\end{equation}
Here \(J_B\) acts irreducibly on \(W_\tau\) and trivially on
\(U_\tau\); distinct \(\tau\) give nonisomorphic full-\(J_B\)
modules.  Choose one primitive matrix idempotent in each central
block and identify \(W_\tau\) with its image in \(H_{B,k}\).
We call this actual summand a \emph{column}.  The other primitive
idempotents select its matrix copies, and matrix units identify
them with the chosen column.  These idempotents and maps belong
to \(D_k\), so the spaces and identifications are admissible.
The pairings below will be transported to these chosen copies;
the restriction of the original polarization to an arbitrary copy
need not be nondegenerate.

A column remains irreducible on \(R_B\): the connected
center acts by scalars, and the product of that center with \(R_B\)
surjects onto \(J_B\).  A column is \emph{mixed} if it has both
circle weights \(1,-1\), and \emph{pure} if it has only one.
These terms refer to the chosen embedding: scalar conjugation can
take a pure column to a mixed one.
A self-pair means a pair tensor on the same column.

\begin{lemma}[Albert columns and their pairings]
\label{p3ab:albert-columns}
In Types I--III let \(F=Z(D)\) be totally real of degree \(f\).
In Type IV let the CM center \(K=Z(D)\) have degree \(2f\), and
write \([D:K]=b^2\).  The columns in
\eqref{p3ab:eq:columns} have the following data:
\[
\begin{array}{c|c|c|c}
 \text{Albert type}&\text{necessary divisibility}&\dim_k W&
       \text{pair tensor}\\ \hline
 \mathrm I& f\mid g&2g/f&\text{alternating on }W\\
 \mathrm{II}&2f\mid g&g/f&\text{alternating on }W\\
 \mathrm{III}&2f\mid g&g/f&\text{symmetric on }W\\
 \mathrm{IV}&fb^2\mid g&g/(fb)&\text{between opposite columns}.
\end{array}
\]
The listed tensors are pair tensors in the sense of
Section~\ref{p3:sec:tensors}, and all copy identifications are
admissible.  In Type IV the opposite column is the dual
\(R_B\)-module.
\end{lemma}

\begin{proof}
We use the structural Albert classification with its positive Rosati
involution \cite[Sections~10.13--10.14]{ChaiOortModuli}.  In Types II and
III, \(D\) is quaternionic over \(F\), so its rational dimension
is \(4f\).  Since \(H_B\) is a vector space over \(D\),
\(2f\mid g\).  In Type IV the rational dimension of \(D\) is
\(2fb^2\), giving \(fb^2\mid g\).  Dividing a central block by
its split matrix multiplicity gives the displayed dimensions.  In
Type I the initial column dimension \(2g/f\) must be even by the
alternating form proved next, giving \(f\mid g\).

Choose a rational polarization pairing \(\psi\) on \(H_B\).  On
a first-kind block \(W\otimes U\), the Rosati involution on
\(\operatorname{End}(U)\) is adjoint to a nondegenerate form \(b_U\).
The real-center idempotents are selfadjoint, so these central blocks
are orthogonal and the restriction of \(\psi\) to each is perfect.
Fixing two \(W\)-arguments in \(\psi\) gives a form \(c\) on \(U\)
satisfying
\[
 c(Tu,v)=c(u,T^\dagger v)\qquad(T\in\operatorname{End}(U)).
\]
Writing \(c(u,v)=b_U(Su,v)\) shows that \(S\) commutes with every
matrix, hence is scalar.  Consequently
\[
 \psi|_{W\otimes U}=\beta_W\otimes b_U
\]
for a nondegenerate \(J_B\)-invariant form \(\beta_W\).  Since
\(\psi\) is alternating, the symmetry signs of the two factors are
opposite.

In Type I, \(U\) is one-dimensional and \(b_U\) is symmetric.
In Type II the positive first-kind involution on the split real
quaternion algebra is orthogonal, so \(b_U\) is symmetric after
splitting.  One can see the sign from positivity: for a
symplectic-adjoint involution on \(M_2(\mathbb R)\), a rank-one idempotent
\(e\) would satisfy \(e^\dagger=1-e\), giving
\(\operatorname{Tr}(ee^\dagger)=0\).  In Type III the positive
involution on the Hamilton quaternions is canonical conjugation,
which becomes symplectic-adjoint on \(M_2(k)\).  The alternative
first-kind involution fixes a nonzero pure quaternion \(v\) and
would give \(\operatorname{Trd}(vv^\dagger)=\operatorname{Trd}(v^2)<0\).
Thus \(b_U\) is alternating in Type III and \(\beta_W\) is
symmetric.

These factorizations yield the required pair tensors on the chosen
columns.  For Types I and II first choose a multiplicity line
nonisotropic for \(b_U\).  Restriction of \(\psi\) to its column
copy is then nondegenerate.  Transport its inverse tensor to the
chosen copy of \(W\) by an admissible matrix-unit isomorphism.
For Type III choose two multiplicity lines with
nonzero \(b_U\)-pairing, restrict \(\psi\) between their column
copies, and identify those copies by a matrix unit.  This gives a
nonzero scalar multiple of the symmetric \(\beta_W\); restricting
twice to one multiplicity line would instead give zero.  Again
transport the inverse tensor to the chosen column by a matrix-unit
isomorphism.  The
projectors and matrix units are admissible, and the rational
pairing \(\psi\) pairs complementary Hodge bidegrees.  Its
factored inverse tensors therefore have type \((1,1)\) in every
scalar conjugate.  The chosen nonzero multiplicity pairing remains
nonzero under those conjugations.

In Type IV positivity makes the involution on \(K\) complex
conjugation.  Thus \(\psi\) pairs only opposite central embeddings.
Factoring the split matrix pairing gives a perfect pairing of their
columns.  For the chosen copy \(W\), use the perfect
\(W\)-to-\(\bar W\) pair of
Lemma~\ref{p3:admissible-operations}, and identify \(\bar W\)
with the chosen copy at the opposite embedding by an admissible
matrix-unit isomorphism.  This supplies the pair on the chosen
columns and identifies the opposite column with the derived dual.
\end{proof}

\subsection{The possible derived projections}

The two Hodge-circle weights impose the standard minuscule
restriction on the derived representations; compare
\cite[Section~3.2, Lemmas~4--5]{Serre1979HodgeTate}.
We give the argument in the circle normalization with weights
\(1,-1\), retaining the scalar-conjugate columns needed at special
periods.

\begin{lemma}[Two circle weights]
\label{p3ab:minuscule}
In the irreducible tensor decomposition of a column \(W_\tau\) for
the simple ideals of \(\operatorname{Lie}(R_B)_k\), every nontrivial
simple factor representation is minuscule.
\end{lemma}

\begin{proof}
The sum of a Galois orbit of simple ideals descends to a rational
ideal.  In each such orbit the complexified circle differential has
a nonzero projection to at least one ideal.  Otherwise the kernel
of the rational adjoint action on that ideal would be a proper
rational algebraic subgroup containing the circle, contrary to
minimality of \(J_B\).

Given an acting ideal, pass to a simultaneous scalar conjugate of
it and \(W_\tau\) for which this projection \(x\) is nonzero.
The conjugated column is still admissible, so it has only circle
weights \(1,-1\).  Fix eigenvectors in all other tensor factors.
The semisimple traceless operator \(x\) then has at most two
eigenvalues on its factor.  It has at least two because a nonzero
representation of a simple Lie algebra is faithful and a scalar
traceless semisimple operator is zero.

Choose a Cartan subalgebra containing \(x\).  For any root
\(\alpha\), a Weyl conjugate \(x'\) has \(\alpha(x')\ne0\):
the Weyl orbit of a root spans the dual Cartan space.  If a highest
weight \(\lambda\) satisfied
\(\langle\lambda,\alpha^\vee\rangle\geq2\) for a positive root,
the root string
\(\lambda,\lambda-\alpha,\lambda-2\alpha\) would yield three
different eigenvalues on \(x'\).  Weyl conjugation does not change
the spectrum, so this is impossible.  Hence
\(\langle\lambda,\alpha^\vee\rangle\leq1\) for every positive
root, the minuscule condition.  Scalar conjugation transports it
back to the original factor.
\end{proof}

The part of the minuscule classification needed here is particularly
small \cite[Sections~27.2, 30.4, and~32.1]{Etingof2024Lie}:
\[
\begin{array}{c|l}
 \text{dimension}&\text{simple nontrivial minuscule modules}\\ \hline
 2&A_1\text{ standard, alternating};\\
 3&A_2\text{ standard and dual};\\
 4&A_3\text{ standard and dual};\quad C_2=B_2
       \text{ standard/spin, alternating};\\
 5&A_4\text{ standard and dual};\\
 6&A_5\text{ standard and dual};\quad C_3\text{ standard, alternating};
       \quad A_3\text{ on }\bigwedge^2k^4,\text{ symmetric}.
\end{array}
\]
The unmarked standard and dual modules are nonselfdual.  The
low-rank identifications \(D_3=A_3\) and \(B_2=C_2\) are included.
Every selfdual module in this table has even dimension.

\begin{lemma}[Small derived projections]
\label{p3ab:small-projections}
An alternating self-paired column of dimension \(2,4\), or \(6\)
has the full symplectic derived image.  A symmetric self-paired
column cannot have dimension two; in dimensions four and six its
derived image is, respectively, the full \(\operatorname{SO}_4\)
and \(\operatorname{SO}_6\).

If \(D=K\) is an imaginary quadratic field, \(B\) is non-CM, and
\(\dim_K H_B\leq6\), then the derived image on either embedding
column is the full special linear group, with dual action on the
opposite column.
\end{lemma}

\begin{proof}
A self-paired column is a selfdual derived module.  Every factor in
its irreducible tensor decomposition is then selfdual, by restricting
an isomorphism with the dual to each simple ideal.  In dimensions at
most six the only nontrivial tensor decompositions are \(2\cdot2\)
and \(2\cdot3\).  The latter is not selfdual.  The former is a
tensor product of two alternating \(A_1\) standards and has a
symmetric form.  The displayed table and
Lemma~\ref{p3ab:minuscule} now give exactly the asserted alternating
images, the symmetric \(2\otimes2\) image in dimension four, and
the symmetric \(\bigwedge^2k^4\) image in dimension six.  These
last two are \(\operatorname{SO}_4\) and
\(\operatorname{SO}_6\).  A connected subgroup of
\(\operatorname{SO}_2\) is a torus and cannot act irreducibly on
a two-dimensional space over \(k\), excluding that dimension.
Connectedness promotes the Lie-algebra images to the stated groups.

For the quadratic assertion write \(W,\bar W\) for the two
columns.  Both are mixed: if one were pure, the circle on
\(W\oplus\bar W\) would lie in the rational \(K\)-scalar torus,
making \(B\) CM.  They are derived duals by
Lemma~\ref{p3ab:albert-columns}.  If \(W\) were derived-selfdual,
\(\operatorname{Hom}_{R_B}(W,\bar W)\) would be one-dimensional and
stable under \(J_B\), since \(R_B\) is normal.  Choose a nonzero
element of this line; it is an isomorphism by Schur's lemma.  This
intertwiner shifts the entire circle
weight spectrum by a scalar character.  Both spectra are
\(\{-1,1\}\), so the shift is zero.  This remains true for its
scalar conjugates, since there are only the two mixed embedding
columns.  Extending by the source projector and applying
\eqref{p3:eq:scalar-descent} gives a full-Hodge-group
endomorphism joining the two distinct central embeddings, contrary
to the commutant \(K_k\).  Thus \(W\) is not derived-selfdual.

Every simple ideal of \(\operatorname{Lie}(R_B)_k\) acts on \(W\):
the action on \(W\oplus\bar W=H_{B,k}\) is faithful, and the two
actions are dual.  Each Galois orbit has an ideal with nonzero
circle differential, as in Lemma~\ref{p3ab:minuscule}.  There
cannot be two such active ideals, since two nonconstant factor
spectra have a sumset of at least three elements, whereas \(W\)
has only two weights.  Hence all acting ideals form one Galois
orbit and have the same Lie type.  The only nonselfdual tensor
candidate in dimensions at most six is \(2\otimes3\), with types
\(A_1\) and \(A_2\), so it cannot occur.  The remaining
nonselfdual modules in the table are exactly the standard and dual
standard special linear modules.  This proves the assertion.
\end{proof}

Lemmas~\ref{p3ab:albert-columns}--\ref{p3ab:small-projections} give the
permitted representations once the column dimensions are known.  The
quadratic action and the total dimension now restrict which columns can
actually occur.

\begin{lemma}[The small column list]
\label{p3ab:small-columns}
Let \(A\) have an imaginary-quadratic \(E\)-action and
\(\dim A\leq6\).  Write \(A\sim\prod_B B^{r_B}\) with pairwise
nonisogenous simple factors, and discard the CM factors.  Put
\(g=\dim B\).  After identifying matrix copies, the remaining
columns have the following possibilities.
\begin{enumerate}
\item In Types I and II they have full symplectic derived images
and dimensions \(2,4,6\).
\item In Type III, \(g=4\) or \(6\), the center is \(\Q\), and
\(r_B=1\).  The unique column has full
\(\operatorname{SO}_g\) image.  Each eigenspace of the given
\(E\)-action is one copy of this column and has equal Hodge
multiplicities \(g/2,g/2\).
\item In Type IV, \(D=K\) is a CM field.  If \(E\) admits no
\(\Q\)-embedding into \(K\), then \(K\) is imaginary quadratic,
\(g=3\), and \(r_B=2\); the \(K\)-columns have the full
\(\operatorname{SL}_3\) and dual images.

If \(E\) admits a \(\Q\)-embedding into \(K\), put
\([K:\Q]=2f\) and \(a=g/f\).  The possibilities are
\[
 f=1,\quad 3\leq a=g\leq6;
 \qquad (f,a,g)=(2,2,4),(2,3,6),(3,2,6).
\]
Every column has the full \(\operatorname{SL}_a\) image,
paired with the dual at the opposite embedding.  In the cases
with \(a=2\), some embedding pairs are mixed and some are pure.
\end{enumerate}
All these are projections of the actual derived Hodge group of \(A\).
\end{lemma}

\begin{proof}
The action restricts to each isotypic factor: the central isotypic
projectors commute with it, and a unital homomorphism from a field
to each matrix factor is injective.

In Type I it extends to an action of the quadratic field
\(F\otimes_\Q E\) on \(F^{r_B}\), since \(F\) is totally real.
Thus \(r_B\) is even and \(g\leq3\).  The table in
Lemma~\ref{p3ab:albert-columns} gives alternating column
dimensions \(2,4,6\).  In Type II the same dimensions follow
from \(2f\mid g\leq6\).  Lemma~\ref{p3ab:small-projections}
gives the full symplectic images.

In Type III the initial symmetric dimensions are also \(2,4,6\).
Lemma~\ref{p3ab:small-projections} excludes two and gives the
orthogonal images for four and six.  Since the column dimension
is \(g/f\) and \(gr_B\leq6\), either remaining dimension forces
\(f=r_B=1\).  The injection \(E\hookrightarrow D\) becomes
\(k\times k\hookrightarrow M_2(k)\); its two nonzero
complementary idempotents each have rank one.  Its eigenspaces
are therefore the two separate matrix copies of the column.
The self-pair matches the two Hodge weights perfectly, proving
the asserted equal multiplicities.

For Type IV first eliminate noncommutative \(D\).  If \(b>1\),
the divisibility \(fb^2\mid g\leq6\) leaves only
\[
 (g,f,b,\dim W)=(4,1,2,2).
\]
More generally, a two-dimensional column cannot have every scalar
conjugate mixed.  If it did, any nonzero tensor in \(D(W)\)
would be a nondegenerate alternating pair tensor of type
\((1,1)\) in every conjugate.  If its isomorphism is
\(t_W:W^*\to W\) and the polarization isomorphism is
\(t_{\mathrm{opp}}:W^*\to\bar W\), then
\(t_{\mathrm{opp}}t_W^{-1}:W\to\bar W\) is an
admissible Hodge isomorphism, since the two Hodge shifts cancel.
By \eqref{p3:eq:scalar-descent} it belongs, after extending by
the source projector, to \(D_k\).  This is impossible because
\(D_k\) preserves the distinct center embeddings.

If every Type IV column is pure, the circle acts by a scalar on
each \(W_\tau\otimes U_\tau\), with opposite scalars at opposite
embeddings.  It therefore lies in the rational torus of center
scalars.  Minimality makes the Hodge group toric, so \(B\) is
CM by Section~\ref{p3:sec:tensors}.  A non-CM rank-two factor
must consequently have both pure and mixed embedding pairs.  For
\(f=1\) there is only one pair, excluding the tuple above and
also the field case \(f=1,a=2\).  A one-dimensional column is
likewise pure at every embedding and gives CM.

We now have \(D=K\).  The \(K\)-algebra \(K\otimes_\Q E\)
is a quadratic field \(L/K\) exactly when \(E\) does not admit
a \(\Q\)-embedding into \(K\); otherwise it is \(K\times K\).
The given action on \(\operatorname{Mat}_{r_B}(K)\) extends to
a unital map from this \(K\)-algebra.  In the field case
\(K^{r_B}\) is a vector space over \(L\), so \(r_B\) is
even.  Hence \(g\leq3\); the exclusions just proved leave only
\(f=1,g=3,r_B=2\), with quadratic \(K\) distinct from \(E\).
In the split case its two idempotents have ranks
\(r_+,r_-\geq0\) summing to \(r_B\).  Thus the action may be
noncentral, as in \(a\mapsto\operatorname{diag}(a,\bar a)\)
on \(K^2\) when \(K=E\).  Diagonalizing the idempotents by a
\(K\)-matrix only changes how the matrix copies in each opposite
pair of \(K\)-columns are assigned to the two \(E\)-eigenspaces;
these matrix identifications are admissible.

In the split case \(fa=g\leq6\), \(a\geq2\), and the exclusion
of \(f=1,a=2\) give exactly the displayed list.  When \(f=1\)
the column is mixed and Lemma~\ref{p3ab:small-projections}
applies.  When \(f>1\) its dimension is two or three.  Its
restriction to \(R_B\) is irreducible by the observation following
\eqref{p3ab:eq:columns}; since its dimension is greater than one,
that action is nontrivial.
Lemma~\ref{p3ab:minuscule} and the displayed small table force
\(\operatorname{SL}_2\) or \(\operatorname{SL}_3\), since two
nontrivial factors already have dimension at least four.  This
holds also at a pure column: the active circle differential may
be at another scalar-conjugate column.  The rank-two pure/mixed
assertion was proved above.

Finally, let \(\rho_B\) be the restriction of
\(\operatorname{Hdg}(H^1(A))\) to \(H_B\).  Its rational algebraic
image contains the circle of \(B\), hence contains \(J_B\) by
minimality.  Conversely, \(\rho_B^{-1}(J_B)\) is a rational
algebraic subgroup of \(\operatorname{Hdg}(H^1(A))\) containing the
circle of \(A\).  Minimality of \(\operatorname{Hdg}(H^1(A))\)
forces this inverse image to be the whole group.  Thus the image is
\(J_B\), and the derived projection is also surjective.  The images
just computed are therefore the projections of the actual group of
\(A\).
\end{proof}

\subsection{Independence of the derived images}

Lemma~\ref{p3ab:small-columns} gives the individual projections, whereas
assumption \textnormal{(c)} of Theorem~\ref{p3:tensor-criterion} requires
their simultaneous product.  A proper subdirect image would identify
two projected simple ideals, so that the same derived factor acts
on different column blocks.  We exclude these graph links using
the actual endomorphism algebra and the circle weights.  This
method is closely related to the product and intertwiner arguments
of Moonen and Zarhin
\cite[(3.1) and Lemmas~3.3--3.4]{MoonenZarhin1999}; the proof
below treats the column representations and partially pure cases
in the present dimension range.

\begin{lemma}[Independence of the small columns]
\label{p3ab:independence}
In Lemma~\ref{p3ab:small-columns}, the listed symplectic,
orthogonal, and special linear derived images are independent,
after identifying matrix copies and pairing opposite CM embeddings.
In other words, they satisfy assumption \textnormal{(c)} of
Theorem~\ref{p3:tensor-criterion}.
\end{lemma}

\begin{proof}
Split each \(\mathfrak{so}_4\) target into its two simple ideals.
The Lie image of the derived group of \(A\) is a semisimple
subdirect algebra of the resulting product of simple target
algebras.  A projection of a simple source ideal onto a simple
target is either zero or an isomorphism; exactly one source ideal
maps nontrivially to each target.  Indeed its image is an ideal,
and the images of different source ideals commute.  Thus failure
of independence is exactly a graph link identifying two
isomorphic target ideals.  The link relation is invariant under
scalar Galois transport.

First consider columns mixed after every scalar conjugation,
apart from the two factors of an \(\mathrm{SO}_4\) target and
the \(\mathfrak{so}_6\simeq\mathfrak{sl}_4\) case treated next.  An
isomorphism of the target Lie algebras in our list identifies
their standard modules up to duality in the special linear case.
Replace one column by its opposite in the dual case.  The intertwiner
equations are defined over \(k\); choose a nonzero derived
intertwiner \(t\) over \(k\).  Its space is a line by
irreducibility, and normality of the derived group makes this line
stable under the full Hodge group.  Writing \(h(z)\) for the
Hodge-circle action, we have
\[
 h(z)t h(z)^{-1}=z^m t
\]
for an integer \(m\).  Because \(t\) is an isomorphism, the
target weight spectrum is the source spectrum translated by \(m\).
Both sets are \(\{-1,1\}\), so \(m=0\).  The same argument applies
to every scalar conjugate of \(t\).  Extending it by its source
projector and using \eqref{p3:eq:scalar-descent} puts it in
\(\operatorname{End}_{\mathrm{Hdg}}(H^1(A,\Q))\otimes k\).
This commutant may mix only matrix copies within one central
embedding and one simple isogeny type.  The putative link
contradicts the choice of different columns.
Self-pairs force mixedness in every conjugate, and a non-CM
quadratic-center column is mixed in both conjugates.  Thus this
argument covers all ordinary links outside the partially pure cases.

The only unequal standard representations of isomorphic simple
targets in this range, outside \(\mathrm{SO}_4\), arise from
\(\mathfrak{so}_6\simeq\mathfrak{sl}_4\).  A Type III
\(\mathrm{SO}_6\) factor has dimension six and exhausts \(A\),
so no other non-CM factor can provide such a link.

Consider next a Type IV factor with both pure and mixed columns.
We first exclude links among its own columns.
Here \(r_B=1\), and the columns over one embedding of \(E\)
are the \(f=2\) or \(3\) embeddings of \(K/E\).  The action
of \(\operatorname{Gal}(k/E)\) on them is transitive.  An
invariant equivalence relation on a prime-sized transitive set is
either discrete or universal, since its equivalence classes have
equal size.  On a column the derived part of the circle
differential is its traceless part: the connected center acts
by scalars.  It vanishes on a pure column and is nonzero on a
mixed column.  A graph link identifies the two projections of the
same circle differential, so it preserves whether that differential
is zero.  The universal relation would link a pure column to a
mixed one and is impossible.  Thus all ideals within this factor
are independent.

We must also exclude links to another isogeny factor.  Only the
quartic-center fourfold can have an external non-CM
partner; the other partially pure cases have dimension six.  Its
two \(A_1\) ideals are interchanged by Galois.  The remainder
has dimension at most two and inherits \(E\).  By
Lemma~\ref{p3ab:small-columns}, a non-CM \(A_1\) target there
can only come from a non-CM elliptic curve repeated twice or
from a Type II surface.  In either case there is only one
absolute \(A_1\) target and it is Galois-stable.  Linking it to
one member of the interchanged pair would link it to the other
as well, contradicting their independence.

It remains to treat \(\mathrm{SO}_4\).  Its column is
\(k^2\otimes k^2\) for its two \(A_1\) ideals.  A nonzero
semisimple traceless differential on \(k^2\) has spectrum
\(\{a,-a\}\), \(a\ne0\).  If both factor differentials were
nonzero, the tensor spectrum \(\{\pm a\pm b\}\) would have at
least three values.  Hence at most one factor differential is
nonzero.  The self-pair makes the column mixed, so at least one is
nonzero.  These are the two absolute ideals
of this rational factor, since its action on its sole column,
repeated twice, is faithful.  If both ideals were separately
Galois-stable, the one with zero differential would descend to a rational
adjoint quotient on which the circle acts trivially, contrary to
minimality.  Hence Galois interchanges them.  The same remainder
argument as in the preceding paragraph excludes an external link
to either one.

All graph links have been excluded.  The Lie image is therefore
the full product.  A connected closed subgroup of the product
of the connected classical groups with this full Lie algebra is
the product itself.  This proves the required independent
derived image; no independence of central characters has been
asserted.
\end{proof}

\subsection{The balanced determinant products}

We now verify assumption \textnormal{(d)} of
Theorem~\ref{p3:tensor-criterion}.  Choose an isogeny
\(f:A'=\prod_B B^{r_B}\to A\) as in
Lemma~\ref{p3ab:small-columns}.  Transport \(E\) to
\(\operatorname{End}^0(A')\) by conjugation with the rational
inverse of \(f\).  Pulling a compatible Riemann form back by
\(f_*:H_1(A',\Q)\to H_1(A,\Q)\) identifies the rational
\(E\)-Hermitian spaces isometrically, so their signatures and
hyperbolicity agree.  Isogenies and factor projections are algebraic
correspondences, and pullback and pushforward along an isogeny
transfer the Hodge conjecture in both directions.  We may therefore
work with this product.

Take \(H_0\) to be the sum of its CM isotypic factors.  Once a
decomposition into CM embedding lines is chosen, this choice and
Lemmas~\ref{p3ab:albert-columns}, \ref{p3ab:small-columns},
and~\ref{p3ab:independence} give assumptions
\textnormal{(a)--(c)}.  In each application below the decomposition
is chosen from the geometry of \(A\), before any determinant or CM
list is tested.

The determinant requirements in the two theorems below have three
forms.  Symplectic blocks contribute only algebraic pair tensors,
and an orthogonal column is admissibly isomorphic to an
\(E\)-eigenspace whose determinant is supplied by
Lemma~\ref{p3ab:determinant-supply}.  For a linear
block with quadratic center, the same lemma supplies the individual
determinants when both signature entries are at most three; in
the remaining odd-rank cases the factor exhausts \(A\), and balance
forces cancellation of opposite determinants.  For higher centers,
balance will force equal net exponents on the columns above one
\(E\)-embedding.  Their product is the full \(E\)-determinant,
to which the supply lemma applies.

For any such fixed decomposition, consider any determinant list and
one CM list satisfying \eqref{p3:eq:determinant-balance}.
An orthogonal column has
imbalance zero because of its self-pair.  For each special linear
representative \(W_j\), let \(n_j\in\mathbb Z\) be the number
of its determinant entries minus the number of opposite entries.
Equation~\eqref{p3:eq:opposite-imbalance} rewrites balance as
\begin{equation}
 \label{p3ab:eq:balance}
 \sum_j n_j\epsilon_\gamma(W_j)+c_\gamma=0
 \qquad\text{for every }\gamma\in\operatorname{Gal}(k/\Q),
\end{equation}
where \(c_\gamma\) is the sum of the imbalances of the actual CM
lines in the fixed list.  Opposite determinant pairs are
algebraically supplied by the determinant of their pair tensor,
as explained after Lemma~\ref{p3:admissible-operations}.  Thus
only the net exponents require further supply.  We write
\(d_j=\epsilon_1(W_j)=2\dim W_j^{1,0}-\dim W_j\), where
\(1\) denotes the identity scalar conjugation and
\(W_j^{1,0}=(W_j)_\C\cap H^{1,0}\).  When considering columns
above one \(E\)-embedding, we identify its image with \(E\subset k\);
then \(\operatorname{Gal}(k/E)\) preserves that set of columns.
More precisely, \(\gamma\in\operatorname{Gal}(k/E)\) sends the
central embedding indexing
\(W_j\) to another one, which we denote by \(\gamma j\).
The image of its primitive idempotent may select a different matrix
copy, but the admissible copy identifications preserve Hodge
dimensions.  Consequently
\[
 \epsilon_\gamma(W_j)=d_{\gamma j}.
\]
The balance equations below therefore use permutations of this
fixed imbalance vector, together with the conjugates of the same
fixed CM list.

\begin{proof}[Proof of Theorem~\ref{p3ab:quadratic-five}]
In dimension at most five, Types I and II require only their
algebraic pair tensors.  A Type III factor has dimension four;
its orthogonal column is admissibly isomorphic to an
\(E\)-eigenspace of signature \((2,2)\).
Lemma~\ref{p3ab:determinant-supply} supplies the determinant of
that eigenspace for every compatible rational form.  The
matrix-unit identification transports this supply to the chosen
column's volume.

In Type IV the case where \(E\) does not embed in \(K\) is
absent, since it requires a threefold with multiplicity two.
In the split \(f=1\) case identify \(K\) with \(E\).  The mixed signatures in dimensions three and four
have both entries at most three, so their determinants are
supplied individually by Lemma~\ref{p3ab:determinant-supply}.
A five-dimensional factor exhausts \(A\).  Its single determinant
imbalance is \(2p-5\ne0\).  There is no CM remainder, so
\eqref{p3ab:eq:balance} gives zero net exponent.  No determinant
of signature \((1,4)\) or \((4,1)\) is required.

The only remaining possibility is \((f,a,g)=(2,2,4)\).  Of
the two columns \(W_1,W_2\) above a chosen \(E\)-embedding,
one is mixed and one pure.  Their imbalance vector is therefore
\((0,2)\) or \((0,-2)\), up to order.  A possible remainder
has dimension one and inherits a unital \(E\)-action.  Its
first cohomology is rank one over \(E\), so the circle acts
through the \(E\)-scalar torus; it is an \(E\)-CM elliptic
curve.  Choose its two \(E\)-embedding lines as the fixed
decomposition of this part of \(H_0\), before testing any list.
Every scalar automorphism fixing \(E\) fixes these lines and
their contribution \(c_\gamma\).

Such an automorphism exchanges the two embeddings of \(K/E\).
Subtracting its equation \eqref{p3ab:eq:balance} from the
identity equation gives
\[
 (n_1-n_2)(d_1-d_2)=0.
\]
Since \(d_1\ne d_2\), \(n_1=n_2\).  After supplying opposite
pairs, the remaining line is a power of
\(D(W_1)\otimes D(W_2)\), or of its opposite.  The direct-sum
determinant maps of Section~\ref{p3:sec:tensors} identify the
basic product \(D(W_1)\otimes D(W_2)\), by nonzero algebraic
operations, with the determinant of
the full \(E\)-embedding space of \(B\).  Its signature is
\((1,3)\) or \((3,1)\), so
Lemma~\ref{p3ab:determinant-supply} supplies it.

Every determinant list satisfying balance is now CM-supplied.
Theorem~\ref{p3:tensor-criterion} gives the conclusion.
\end{proof}

\begin{proof}[Proof of Theorem~\ref{p3ab:split-sixfolds}]
The central isotypic projectors are selfadjoint for the compatible
polarization.  Indeed a polarization identifies each simple
isogeny type with its own dual type, and there are no homomorphisms
between distinct types.  The isotypic summands are therefore
orthogonal for \(e\), and also for \(h_E\) because their
projectors commute with \(E\).  Their signatures add to
\((3,3)\).  We do not assert that a proper summand of the
hyperbolic form is itself hyperbolic.

Types I and II again require only pairs.  In Type III, a
four-dimensional column is admissibly isomorphic to an
\(E\)-eigenspace of signature \((2,2)\), whose determinant is
supplied for its arbitrary restricted form.  A six-dimensional
Type III factor exhausts \(A\), so the determinant of either
\(E\)-eigenspace is supplied for the given hyperbolic form.
Both statements follow from Lemma~\ref{p3ab:determinant-supply};
the admissible matrix-unit identifications transport the supply
to the chosen column's volume.

Suppose first that a Type IV factor has \(E\) not embedding
into \(K\).  It is a quadratic-center threefold with multiplicity
two and exhausts \(A\).  Its \(K\)-column is mixed of rank three,
so its imbalance is nonzero.  There is no CM remainder, and
\eqref{p3ab:eq:balance} forces its net determinant exponent to
vanish.

Now consider the split \(f=1\) case and identify \(K\) with \(E\).
For multiplicity one in dimensions three,
four, and five, orthogonal addition of signatures to \((3,3)\)
puts each signature entry at most three.  The only repeated
non-CM case is \(g=3,r_B=2\).  For it use the central \(K=E\)
action on \(B\), of mixed signature \((1,2)\) or \((2,1)\),
and any compatible polarization on \(B\).  The possibly
noncentral given action on \(B^2\) changes only the multiplicities
of its two column types, as proved in
Lemma~\ref{p3ab:small-columns}.  Thus
Lemma~\ref{p3ab:determinant-supply} supplies the required
determinants in all these ranks.  If \(g=6\), \(B\) exhausts
\(A\), and the given hyperbolic determinant supplies both
orientations.

For the quartic-center fourfold \((f,a,g)=(2,2,4)\), replace
\(u\) by \(-u\) if necessary and retain the convention
\(\sigma(u)=i\sqrt d\), so that its restricted Hermitian form has
signature \((1,3)\).  This replacement negates \(h_E\), preserving
hyperbolicity and the total signature \((3,3)\) while interchanging
the two possible restricted signatures.  The complementary two-dimensional isotypic sum
\(T\) has signature \((2,0)\).  In \(H_1(T,\Q)\), of rank
two over \(E\), the complex structure is the same \(E\)-scalar
on the whole space.  Every \(E\)-line in an \(E\)-basis is
therefore a rational Hodge substructure of the same rank-one CM
type.  The degree-one realization makes \(T\) isogenous to a
square of an \(E\)-CM elliptic curve.  Choose the resulting
\(E\)-embedding lines as the fixed decomposition of its contribution
to \(H_0\), before testing any list.
Their imbalance contribution is fixed by automorphisms fixing
\(E\).

The two column imbalances are \((0,-2)\), up to order.
Exchanging the two embeddings of \(K/E\) and subtracting
\eqref{p3ab:eq:balance} gives \(n_1=n_2\).  The remaining
product is a power of \(D(W_1)\otimes D(W_2)\), or its
opposite.  Algebraic direct-sum determinant maps identify the
basic product \(D(W_1)\otimes D(W_2)\)
with the full \(E\)-determinant of \(B\), of signature
\((1,3)\), which is supplied by
Lemma~\ref{p3ab:determinant-supply}.

It remains to consider the six-dimensional higher centers.
Here \(B=A\), so \(c_\gamma=0\).
For \((f,a)=(2,3)\), the two holomorphic multiplicities above
one \(E\)-embedding sum to three.  The alternatives \((0,3)\)
and \((3,0)\) make every column pure and the factor CM.
Otherwise the vector is \((1,2)\), up to order, with imbalance
\((-1,1)\).  Equation~\eqref{p3ab:eq:balance} gives
\(n_1=n_2\).

For \((f,a)=(3,2)\), the three multiplicities belong to
\(\{0,1,2\}\) and sum to three.  The vector \((1,1,1)\)
is excluded by the rank-two argument in
Lemma~\ref{p3ab:small-columns}.  The remaining vector is a
permutation of \((0,1,2)\), with imbalance a permutation of
\((-2,0,2)\).  The Galois action over \(E\) on the three
embeddings of \(K/E\) is transitive and therefore contains a
three-cycle.  Order the columns so that the imbalance is
\((-2,0,2)\).  Applying \eqref{p3ab:eq:balance} to the three
cyclic conjugates gives, up to reversing their order,
\[
 -n_1+n_3=0,\qquad -n_2+n_1=0,\qquad -n_3+n_2=0.
\]
Thus \(n_1=n_2=n_3\).

In either case, after opposite pairs have been supplied, only a
power of \(\bigotimes_{j=1}^fD(W_j)\), or its opposite,
remains.  The center projectors are admissible algebraic
degree-one maps.  Alternation and those projectors identify
the basic product \(\bigotimes_{j=1}^fD(W_j)\) nontrivially
with the determinant of the full
\(E\)-embedding space of \(A\).  The given hyperbolic form
supplies this determinant by
Lemma~\ref{p3ab:determinant-supply}.

These cases exhaust Lemma~\ref{p3ab:small-columns}.  The
equalities above concern arbitrary integer net exponents, so they
verify determinant supply in every tensor degree.  Products,
opposite pairs, and the direct-sum maps preserve CM supply.
Assumption \textnormal{(d)} of
Theorem~\ref{p3:tensor-criterion} is therefore satisfied, and
that theorem proves the result on every self-power.
\end{proof}

\section{The geometric seed for the Weil-plane theorem}
\label{p3sw:appendix}
\label{p3sw:support-component}

We begin the proof of Theorem~\ref{p3sw:weil-plane}, the algebraicity
input used in Lemma~\ref{p3ab:determinant-supply}.  We follow Markman's
secant-sheaf and spin construction
\cite{Markman2025Secant,Markman2026ICM}.
The argument will retain the actual polarized Hermitian space and the
full derived obstruction group throughout descent and deformation.

We first construct a coherent perfect sheaf on the product of a
genus-three Jacobian and its dual, with a coherent action of a finite
group of translations combined with line tensors.  An exact
Hochschild--Atiyah trace identity and a dimension calculation then show
that the semiregularity map is injective on the \emph{entire} invariant
derived \(\operatorname{Ext}^2\) group.  A separate finite Clifford
calculation identifies the normalized rational Betti character of the
same sheaf: it has a nonzero Weil component, and all its components
remain Hodge on the full compatible marked domain.

We next descend the sheaf to a twisted sheaf on the finite quotient.
A projective-bundle construction turns it into an ordinary coherent
sheaf while retaining the full obstruction group.  We prove the
restricted smooth-projective-germ semiregularity theorem needed to
deform that ordinary sheaf, and recover the normalized character by
algebraic projective-bundle operations.  Finally, proper Hilbert
parameter spaces and the Baire theorem extend the local algebraicity
to every marked period; rational isogenies pass from the seed lattice
to every hyperbolic quadratic lattice.  The proof of
Theorem~\ref{p3sw:weil-plane} is completed in
Subsection~\ref{p3sw:all-periods}.

The geometric input is a pair of finite orbits of Abel curves on a
genus-three Jacobian.  We choose the orbits so that their ideals have no
sections after any degree-zero twist of twice the principal polarization.
This vanishing will control the fiber cohomology of a fixed
Fourier--Mukai transform and allow its shifted derived dual to be an
ordinary coherent sheaf.  We construct that sheaf, its coherent
line-twisted translation action, and the finite quotient used in the
deformation argument.  The source ideal and transform are those of
Markman \cite[(6.2.1), (6.2.3), and (9.2.1)]{Markman2025Secant};
the compact-type degeneration below supplies the required orbits.

All varieties in this section are over \(\C\).  If \(V\) is a group
variety, \(t_v(z)=z+v\), and \(t_vD\) denotes the image translate of a
divisor \(D\).  Thus \(t_v^*\mathcal O(D)=\mathcal O(t_{-v}D)\).

\subsection{The Jacobian and the required Abel orbits}

Let \(C\) be a smooth nonhyperelliptic curve of genus three and choose
\(p\in C\).  Give
\[
 X=\operatorname{Pic}^{2}(C)
\]
the group law with origin \(\mathcal O_C(2p)\); equivalently, tensoring
with \(\mathcal O_C(2p)\) identifies \(\operatorname{Pic}^0(C)\) with
\(X\).  Put
\[
 \Theta=W_2(C),\qquad H=\mathcal O_X(\Theta),\qquad
 \widehat X=\operatorname{Pic}^0(X).
\]
For \(v\in X\) define
\begin{equation}
 \lambda(v)=t_v^*H\otimes H^{-1},\qquad
 a(v)=\mathcal O_X(t_v\Theta-\Theta)
     =t_{-v}^*H\otimes H^{-1}=\lambda(-v).
 \label{p3sw:support-lambda-a}
\end{equation}
The theorem of the square makes these homomorphisms, and the canonical
principal polarization makes both isomorphisms.  We use \(a(v)\) both
for a point of \(\widehat X\) and for its represented line bundle.
More generally, for a line bundle \(D\) on an abelian variety or
scheme, \(\phi_D(v)=t_v^*D\otimes D^{-1}\).

In cohomology, \(\Theta\) will also denote the normalized integral Betti
class \(c_{1,\mathrm B}(H)\), so that
\(\int_X\Theta^3/3!=1\).  It does not denote a raw first-jet or Hodge
class.  Here is the sign convention that will be used for this
polarization.  Set
\begin{equation}
 L=H^1(X,\Q),\qquad U=L^\vee=H_1(X,\Q)=H^1(\widehat X,\Q),
 \qquad
 \vartheta:U\longrightarrow L,\quad y\longmapsto\iota_y\Theta.
 \label{p3sw:support-betti-convention}
\end{equation}
The last identification is the one defined by the normalized Poincar\'e
bundle \(\mathcal P\) on \(X\times\widehat X\), whose restriction to
\(X\times\{N\}\) is \(N\).  If \(e_i\) and \(f_i\) are dual bases of
\(L\) and \(U\), then
\(c_{1,\mathrm B}(\mathcal P)=\sum_i e_i\wedge f_i\).
The normalized biextension identity is
\begin{equation}
 (1\times\lambda)^*\mathcal P
 \simeq m_X^*H\otimes\operatorname{pr}_1^*H^{-1}
              \otimes\operatorname{pr}_2^*H^{-1}
              \otimes\underline{H_0},
 \label{p3sw:support-polarization-biextension}
\end{equation}
where \(m_X\) is addition and the constant line \(\underline{H_0}\)
rigidifies the two zero axes.  Comparing the mixed K\"unneth terms in
\eqref{p3sw:support-polarization-biextension} gives
\begin{equation}
 \lambda^*=\vartheta,\qquad a^*=-\vartheta,\qquad
 \lambda_*=-\vartheta,\qquad a_*=\vartheta.
 \label{p3sw:support-polarization-signs}
\end{equation}
Indeed, if
\(\Theta=\frac12\sum_{i,j}c_{ij}e_i\wedge e_j\), its mixed term
after addition is
\(\sum_{i,j}c_{ij}e_i^{(1)}\wedge e_j^{(2)}\).  Consequently
\(\lambda^*f_i=\sum_jc_{ij}e_j=\iota_{f_i}\Theta\).
The two homology maps are the ordinary transposes of the cohomology
maps, and \(\vartheta^t=-\vartheta\), proving the last two signs.

Define the two Abel curves
\begin{equation}
 C_p=\{\mathcal O_C(p+q):q\in C\},\qquad
 \Sigma_p=\{\omega_C(-p-q):q\in C\}.
 \label{p3sw:support-abel-curves}
\end{equation}
They have the normalized Betti class and degree
\begin{equation}
 [C_p]=[\Sigma_p]=\frac{\Theta^2}{2},\qquad
 \Theta\cdot C_p=\Theta\cdot\Sigma_p=3.
 \label{p3sw:support-abel-class}
\end{equation}
For clarity, these normalizations follow directly from the Jacobian
lattice.  In the analytic description
\(J(C)=H^0(C,\omega_C)^*/H_1(C,\Z)\), the Abel map induces the identity
on the displayed lattice.  Take a symplectic basis
\(e_i,f_i\) of \(H^1(C,\Z)\), identified with \(H^1(X,\Z)\) by Abel
pullback.  The canonical principal class is
\(\sum_{i=1}^3e_i\wedge f_i\).  Pairing \(\Theta^2/2\) with each
basis element of \(\Lambda^2H^1(X,\Z)\) gives the same value as
integration of its Abel pullback on \(C\).  Poincar\'e duality proves
the first formula for \(C_p\).  The negative Abel map has the same
action on degree-two cohomology, giving the formula for \(\Sigma_p\).
The degree in \eqref{p3sw:support-abel-class} is then \(\int_X\Theta^3/2=3\).
This lattice description also shows that restriction
\begin{equation}
 \operatorname{Pic}^0(X)\longrightarrow\operatorname{Pic}^0(C_p)
 \label{p3sw:support-abel-restriction}
\end{equation}
is an isomorphism: its integral lattice map and its map on
\(H^{0,1}\) are the Abel pullback isomorphisms.  For a translate of
\(\Sigma_p\) the analogous map is the negative of this isomorphism.

We will repeatedly use curve Riemann--Roch and duality in the following
form.  For a proper Gorenstein curve \(D\) and an invertible sheaf \(N\),
\begin{equation}
 \chi(D,N)=\deg N+\chi(D,\mathcal O_D),\qquad
 H^1(D,N)^\vee=H^0(D,\omega_D\otimes N^{-1}).
 \label{p3sw:support-curve-duality}
\end{equation}
The numerical Riemann--Roch formula is the rank-one case of
\cite[Lemma 53.5.2, Tag 0BS6]{P3SWStacks}; the canonical duality
isomorphism is stated in
\cite[Remark 53.4.3, Tag 0BS4]{P3SWStacks}.  We will apply them to smooth curves.

The temporary restriction \(d\geq3\) loses no imaginary quadratic
field: the final proof of Theorem~\ref{p3sw:weil-plane} replaces a
purely imaginary generator \(u_0\) by \(u=2u_0\), scaling the associated
Hermitian form without changing its isotropic subspaces.

\begin{proposition}[Simultaneous equivariant seeds]
\label{p3sw:seed}
For every integer \(d\geq3\), put \(n=d+1\).  There exist a smooth
nonhyperelliptic genus-three curve \(C\), a point \(p\in C\), and
cyclic subgroups \(G_1,G_2\subset X=\operatorname{Pic}^2(C)\) of
order \(n\), with \(G_1\cap G_2=0\), as follows.  There are cosets
\(\{s_1,\ldots,s_n\}=s_0+G_2\) and
\(\{t_1,\ldots,t_n\}=t_0+G_1\) such that
\begin{equation}
 C_i=t_{s_i}C_p,\qquad \Sigma_j=t_{t_j}\Sigma_p,\qquad
 A=\coprod_{i=1}^n C_i,\qquad \Sigma=\coprod_{j=1}^n\Sigma_j
 \label{p3sw:support-seed-orbits}
\end{equation}
are disjoint within each of the two indicated unions, and
\begin{equation}
 H^0(X,I_A\otimes H^{\otimes2}\otimes N)
 =
 H^0(X,I_\Sigma\otimes H^{\otimes2}\otimes N)=0
 \quad\text{for every }N\in\widehat X.
 \label{p3sw:support-uniform-vanishing}
\end{equation}
\end{proposition}

The uniform vanishing in the proposition has a specific role in the
sheaf construction.  For every \(x\in X\), the product ideal
\(I_A I_{t_{-x}\Sigma}\) is contained in \(I_A\), so every
degree-zero twist of this ideal by \(H^{\otimes2}\) has zero
\(H^0\).  Its quotient in \(\mathcal O_X\) has one-dimensional
support; the ideal sequence and ample-line vanishing will then give
zero \(H^3\).  For general \(x\), the two orbit unions are disjoint.
The restriction of \(H^{\otimes2}\) to each component has degree
six, greater than \(2g(C)-2=4\), which also gives zero \(H^2\)
for these parameters.

We will prove that these cohomology groups are the derived fibers of
the fixed transform.  The first two vanishings give a local model
\(K^1\to K^2\), and the generic \(H^2\)-vanishing makes its dual
differential generically injective.  Integrality of the parameter space
then makes that dual differential injective everywhere.  Its cokernel
is the coherent sheaf we seek.  We first construct the orbits required
by the proposition.

\subsection{Constructing the seed by a compact-type degeneration}

On a product of three elliptic curves, the central support of the
limiting Abel curve will be the union of the three coordinate axes.
The next two lemmas show that four
suitably translated copies of this star impose more vanishing conditions
than a nonzero section of a product of degree-two lines can satisfy.
We will carry this vanishing to a smooth Jacobian through a polarized
family and the proper relative Picard scheme.

\subsubsection{Vanishing on the elliptic star}

\begin{lemma}[Singular marked points of a trilinear form]
\label{p3sw:support-trilinear}
Let \(V_1,V_2,V_3\) be two-dimensional vector spaces and let
\(0\ne T\in V_1\otimes V_2\otimes V_3\).  Suppose a finite indexed
collection
\[
 (x_h,y_h,z_h)\in
 \mathbb P(V_1^*)\times\mathbb P(V_2^*)\times\mathbb P(V_3^*)
\]
has at most two indices in every fiber of each coordinate projection.
At most three indices satisfy all three equations
\begin{equation}
 T(-,y_h,z_h)=T(x_h,-,z_h)=T(x_h,y_h,-)=0.
 \label{p3sw:support-singular-marks}
\end{equation}
Indices, rather than distinct triples, are counted.
\end{lemma}

\begin{proof}
Choose bases and write \(T(x,y,z)=y^tM(x)z\), where
\(M(x)=x_0A+x_1B\) is a pencil of two-by-two matrices.  Put
\(D(x)=\det M(x)\).  Equations \eqref{p3sw:support-singular-marks} imply
\begin{equation}
 M(x)z=0,\qquad y^tM(x)=0,\qquad y^tAz=y^tBz=0.
 \label{p3sw:support-pencil-equations}
\end{equation}
Rank two for \(M(x)\) is impossible.  If its rank is one, its adjugate
is a nonzero multiple of \(zy^t\), so the determinant derivative
formula and \eqref{p3sw:support-pencil-equations} give
\(\partial_{x_0}D(x)=\partial_{x_1}D(x)=0\).
If \(M(x)=0\), its adjugate is zero and the same conclusion holds.
For a nonzero binary quadratic \(D\), the gradient is a nonzero linear
map in characteristic zero and has at most one projective zero.
All qualifying indices then have the same \(x\), and there are at
most two.

It remains to consider \(D\equiv0\).  Every nonzero matrix in the
pencil has rank one.  A one-dimensional pencil makes \(T\) a product
of three linear forms.  For a two-dimensional pencil write
\(A=uv^t\) and \(B=u'v'^t\).  The equality
\[
 0=\det(A+B)=\det[u,u']\det[v,v']
\]
shows that the pencil has a common left or right factor.  After
permuting the variables, both cases have the form
\(T=l(x)q(y,z)\).  If the bilinear form \(q\) has rank two, the last
two partial contractions force \(l(x)=0\), again allowing at most two
indices.  If \(q\) has rank one, write \(T=l(x)m(y)k(z)\).
Equations \eqref{p3sw:support-singular-marks} then hold precisely when at least two of
\(l(x),m(y),k(z)\) vanish.  Each of the three distinguished
projective coordinate values occurs at most twice.  If \(N\) indices
qualify, counting these incidences gives \(2N\leq6\), hence \(N\leq3\).
\end{proof}

\begin{lemma}[The elliptic star]
\label{p3sw:support-elliptic-star}
Let \(J_0=E_1\times E_2\times E_3\) be a product of elliptic curves
with origins, and let \(Z_0\) be the reduced union of the three
coordinate elliptic axes through zero.  Let \(K\subset J_0\) be a
finite subgroup of order at least four whose projection to every
\(E_i\) is injective.  The translates \(Z(g)=g+Z_0\), \(g\in K\),
are pairwise disjoint.  For line bundles \(D_i\) of degree two on
\(E_i\),
\begin{equation}
 H^0\left(J_0,I_{\bigcup_{g\in K}Z(g)}
       \otimes(D_1\boxtimes D_2\boxtimes D_3)\right)=0.
 \label{p3sw:support-star-vanishing}
\end{equation}
The same vanishing holds for any closed subscheme with that support.
\end{lemma}

\begin{proof}
If axis \(i\) at \(g\) met axis \(j\) at \(h\ne g\), the coordinates
of \(g\) and \(h\) outside \(\{i,j\}\) would agree.  There is at least
one such coordinate, contradicting injectivity there.  This proves
disjointness.

Put \(V_i=H^0(E_i,D_i)\).  Formula \eqref{p3sw:support-curve-duality} gives
\(\dim V_i=2\) and \(h^0(D_i(-q))=1\) for every \(q\in E_i\).
Evaluation is therefore basepoint-free and defines a degree-two map
\(\varphi_i:E_i\to\mathbb P(V_i^*)\), each of whose fibers has at most
two distinct points.  The marks
\[
 \bigl(\varphi_1(g_1),\varphi_2(g_2),\varphi_3(g_3)\bigr),
 \qquad g=(g_1,g_2,g_3)\in K,
\]
have at most two indices in each coordinate fiber, since \(K\to E_i\)
is injective.  The product formula for sections identifies a section
of \(D_1\boxtimes D_2\boxtimes D_3\) with a tensor in
\(V_1\otimes V_2\otimes V_3\).  Its restriction to axis \(i\) at \(g\)
is contraction by evaluation at the other two coordinates.  Vanishing
on \(Z(g)\) is exactly \eqref{p3sw:support-singular-marks}.  Lemma
\ref{p3sw:support-trilinear} forbids this for four or more indices for
a nonzero section.  Finally, the ideal of a closed subscheme with
this support is contained in the radical ideal of the reduced union,
which proves the last assertion.
\end{proof}

\subsubsection{Finite cohomology models and base change}

We need to pass from vanishing on the limiting product to vanishing on
nearby Jacobians, uniformly over their Picard twists.  Later we will
identify every derived fiber of the fixed Fourier--Mukai kernel.
The following finite-model statement supplies both semicontinuity for
the first step and arbitrary base change for the second.  It also
defines the determinant line that identifies the limiting polarization.

\begin{lemma}[Finite models and arbitrary base change]
\label{p3sw:support-finite-model}
Let \(B\) be a regular Noetherian scheme of finite dimension, let
\(f:V\to B\) be projective, and let \(\mathcal F\) be a coherent sheaf
flat over \(B\).  Then \(Rf_*\mathcal F\) is locally represented by a
bounded complex \(K^\bullet\) of finite-rank locally free sheaves.
For every morphism \(h:T\to B\) there is a derived base-change
isomorphism
\begin{equation}
 Lh^*Rf_*\mathcal F
 \simeq Rf_{T,*}(\mathcal F|_{V\times_B T}),
 \label{p3sw:support-derived-base-change}
\end{equation}
where the restriction on the right is ordinary, and the left side is
represented locally by \(K^\bullet\otimes\mathcal O_T\).
The functions
\(b\mapsto\dim_{\kappa(b)}H^i(V_b,\mathcal F_b)\) are upper
semicontinuous.  If all derived fibers have cohomology in degrees
\([c,e]\), the local model can be chosen in precisely those degrees.
\end{lemma}

\begin{proof}
Over an affine \(U=\operatorname{Spec}R\subset B\), choose a finite
affine cover \(V_i\) of \(V_U\).  All finite intersections are affine
because \(V_U\) is separated.  Its \v{C}ech complex has terms
\begin{equation}
 C_R^q=\prod_{i_0<\cdots<i_q}
       \Gamma(V_{i_0}\cap\cdots\cap V_{i_q},\mathcal F).
 \label{p3sw:support-cech}
\end{equation}
It is bounded.  Each term is \(R\)-flat: for a module on an affine
scheme, flatness over \(R\) is checked at the total-space stalks, where
it is the assumed relative flatness.  For any \(R\)-algebra \(R'\),
\(C_R^\bullet\otimes_RR'\) is exactly the affine \v{C}ech complex of the
ordinary pullback of \(\mathcal F\) to \(V_U\times_RR'\).
A bounded complex of flat modules is \(K\)-flat, so this tensor product
also computes the derived tensor product.  Affine \v{C}ech cohomology
therefore proves \eqref{p3sw:support-derived-base-change} for affine base changes, and those isomorphisms
glue for an arbitrary \(T\).

The cohomology modules of \(C_R^\bullet\) are finite over \(R\), by
proper coherence for a projective scheme over a Noetherian affine base
\cite[Lemma 30.16.1(3), Tag 0B5T]{P3SWStacks}.  A finite module over
the regular Noetherian finite-dimensional ring \(R\) has a finite
resolution by finite projective modules
\cite[Section 10.110, Tag 065U]{P3SWStacks}.
The truncation triangles, applied to the finitely many cohomology
modules of \(C_R^\bullet\), consequently make its derived object
perfect.  After localizing, it has the stated finite free model.

For such a model the dimension of the \(i\)-th fiber cohomology is
\begin{equation}
 \operatorname{rank}K^i-
 \operatorname{rank}(d^i\otimes\kappa(b))-
 \operatorname{rank}(d^{i-1}\otimes\kappa(b)).
 \label{p3sw:support-rank-semicontinuity}
\end{equation}
Ranks of matrices are lower semicontinuous by their minors; \eqref{p3sw:support-rank-semicontinuity}
proves upper semicontinuity.  Finally, at a term below \(c\), the first
remaining differential is injective on the fiber.  Invert a maximal
minor, split off this term, and replace the next term by its locally
free cokernel.  At a term above \(e\), do the dual operation using a
fiberwise surjective last differential.  Repeating at a chosen point
and shrinking its neighborhood gives a model in \([c,e]\).
\end{proof}

We also use the perfect determinant in a concrete form.  For a finite
locally free complex put
\(\det K^\bullet=\bigotimes_i(\det K^i)^{(-1)^i}\).
The determinant of an acyclic complex is canonically trivialized by
the exact sequences of its cycles.  Applying these trivializations to
cones makes the determinants independent of the finite model, so they
glue and commute with base change.  In particular, a model
\(K^0\to K^1\) of equal ranks has a regular determinant section of
\((\det K^\bullet)^{-1}\); its zero set is where the fiber complex is
not acyclic.

\subsubsection{From the polarized star to a smooth seed}

\begin{proof}[Proof of Proposition~\ref{p3sw:seed}]
\emph{The polarized elliptic limit.}
We first realize the elliptic product with its product principal
polarization as the central Jacobian of a smoothing.  This identifies
the central twice-theta line with the degree-two product lines tested
in Lemma~\ref{p3sw:support-elliptic-star}.

Start with the stable compact-type curve
\begin{equation}
 C_0=R\cup E_1\cup E_2\cup E_3,\qquad R\simeq\mathbb P^1,
 \label{p3sw:support-compact-type-curve}
\end{equation}
where the elliptic tails \(E_i\) meet \(R\) at three distinct nodes
\(e_i\) and have no other intersections.  We choose a projective
one-parameter smoothing \(\mathcal C\to S\) over a smooth pointed
complex algebraic curve \((S,0)\) with regular total surface and with
smooth nonhyperelliptic noncentral fibers.

Here is the existence check, for the unmarked stable curve.  The local
presentation of \(\Omega_{C_0}\) at a node has length one, its local
\(\mathcal Ext^1\) is supported at the nodes, and the local-to-global
sequence gives
\(\operatorname{Ext}^2(\Omega_{C_0},\mathcal O_{C_0})=0\).
The global tangent contribution is the sum of the \(H^1\) of the
tangent lines on the normalized components, with their node points
removed.  It is zero for
\(T_R(-e_1-e_2-e_3)=\mathcal O_{\mathbb P^1}(-1)\), and has
dimension one for each \(T_{E_i}(-e_i)=\mathcal O_{E_i}(-e_i)\).
The three node-smoothing parameters add three more dimensions.
Thus the unmarked versal germ is smooth of dimension six.  The
algebraic input is the smooth Deligne--Mumford stable-curve stack
and density of its smooth-curve locus
\cite[Theorem 109.25.3, Tag 0E99]{P3SWStacks}; it gives an algebraic
versal chart here.  Smooth hyperelliptic genus-three curves have a
five-dimensional locus, as seen from eight branch points on
\(\mathbb P^1\) modulo \(\operatorname{PGL}_2\).  Its closure is
proper in that chart.  A general algebraic curve through the chosen
point is not contained in this closure and has nonzero first
derivative in each node parameter.  The local total-space equation
at each node is then \(uv=t\) times a unit, so the total surface is
regular.  Shrinking \(S\) gives the asserted noncentral fibers.
Choose a smooth point on \(R\) only now.  Smoothness near it lifts it,
after an \'etale localization of \(S\), to a section \(\sigma\).
This point was not included in the six-dimensional unmarked germ.

Let \(\mathcal J\) be the componentwise degree-zero relative Picard
space of \(\mathcal C/S\).  The exact Picard input used here is
\cite[Section 2, p.~136]{P3SWDeligneGabber}: for a stable family this
open group space is separated, smooth, and of finite type, and its
fibers are semiabelian.  The normalization sequence for a connected
nodal curve with dual graph \(\Delta\) gives
\begin{equation}
 1\longrightarrow\mathbb G_m^{b_1(\Delta)}
 \longrightarrow\operatorname{Pic}^0(C_0)
 \longrightarrow\prod_v\operatorname{Pic}^0(C_v)
 \longrightarrow1.
 \label{p3sw:support-normalization-picard}
\end{equation}
Indeed, its kernel is the set of gluing scalars at the nodes modulo
rescaling on each component; one scalar remains per graph cycle.
Our graph is a tree and \(\operatorname{Pic}^0(R)=0\), so
\(\mathcal J_0=\prod_i\operatorname{Pic}^0(E_i)\).  All noncentral
fibers are smooth Jacobians.  Thus every geometric fiber is proper
and connected.  Deligne's
\cite[Section 1.1, pp.~131--132]{P3SWDeligneGabber} states both that a
separated flat finite-presentation space with proper connected
geometric fibers is proper, and that a generic polarization on an
abelian space over a normal integral base extends uniquely and makes
that space projective.  Consequently \(\mathcal J\) is an abelian
scheme and the generic canonical polarization extends to
\(\lambda_{\mathcal J}:\mathcal J\to\widehat{\mathcal J}\).
Its degree is locally constant and is one generically, so it is
principal.  We use only the componentwise degree-zero Picard space;
no separatedness of the unrestricted relative Picard functor is
asserted.

The decomposition of the central Jacobian as an abelian variety does
not yet identify its polarization.  To obtain the product polarization
needed for the vanishing lemma, we compute a theta line by a determinant
of cohomology.  After an \'etale localization, the relative Picard space
has a universal line bundle;
this is the precise universal-line input from
\cite[Section 2.1, last paragraph]{P3SWCoelhoPacini}.
Restrict it to \(\mathcal C\times_S\mathcal J\) and rigidify it along
\(\sigma\); denote it by \(\mathcal U\).  With \(p_C\) the projection
to \(\mathcal J\), put
\begin{equation}
 \mathcal T=
 \left(\det Rp_{C,*}\bigl(\mathcal U\otimes
                 \mathcal O_{\mathcal C}(2\sigma)\bigr)\right)^{-1}.
 \label{p3sw:support-theta-determinant}
\end{equation}
The base \(\mathcal J\) is regular, and this is a flat projective
curve family with a line bundle.  Lemma
\ref{p3sw:support-finite-model} and the determinant construction
following it define \eqref{p3sw:support-theta-determinant} and its base changes.  Every fiber has
Euler characteristic zero and cohomology in degrees zero and one;
locally it has a two-term equal-rank model, with determinant section.
On a smooth nonhyperelliptic fiber the section vanishes exactly on
the effective degree-two locus \(W_2\).  Its generic multiplicity
there is one: at an effective degree-two line the groups \(H^0\)
and \(H^1\) have dimension one, and the first-order differential is
cup product
\[
 H^1(\mathcal O_C)\longrightarrow
 \operatorname{Hom}(H^0(N),H^1(N)).
\]
Its dual is multiplication
\(H^0(N)\otimes H^0(\omega_C N^{-1})\to H^0(\omega_C)\),
which is nonzero on nonzero sections.  Splitting the invertible
blocks of a local two-term model shows that this nonzero first-order
differential is the linear term of its remaining one-by-one
determinant.  Thus the determinant section has a simple zero there.
It follows that the homomorphism \(\phi_{\mathcal T}\) associated to
\(\mathcal T\) is the generic canonical polarization.  It equals
\(\lambda_{\mathcal J}\) over \(S\), since the two morphisms agree
on the schematically dense generic fiber and the target is separated.

On \(C_0\), the rigidification makes \(\mathcal U\) trivial on
\(R\times\mathcal J_0\); its restrictions to the elliptic components
are the Poincar\'e lines \(\mathcal P_i\), normalized at \(e_i\).
The normalization exact sequence for
\(\mathcal U\otimes\mathcal O(2\sigma)\) factors \eqref{p3sw:support-theta-determinant}, up to constant
lines from \(R\) and the nodes, as
\begin{equation}
 \mathcal T_0\simeq
 \mathop{\boxtimes}_{i=1}^3
 \left(\det R\operatorname{pr}_{\widehat E_i,*}(\mathcal P_i)\right)^{-1}.
 \label{p3sw:support-product-determinant}
\end{equation}
For an elliptic curve the \(i\)-th factor is
\(\mathcal O_{\widehat E_i}(0)\).  Its fiber cohomology vanishes off
the origin.  At the origin the one-dimensional \(H^0\) and \(H^1\)
have the nonzero first-order cup-product differential
\(H^0(\mathcal O_{E_i})\otimes H^1(\mathcal O_{E_i})
\to H^1(\mathcal O_{E_i})\); the same determinant argument gives
a simple zero.  Equation \eqref{p3sw:support-product-determinant} therefore proves
\begin{equation}
 (\mathcal J_0,\lambda_{\mathcal J,0})
   =\prod_{i=1}^3(\operatorname{Pic}^0(E_i),\lambda_i)
 \label{p3sw:support-product-polarization}
\end{equation}
as principally polarized abelian varieties.  In particular, the
polarization on this product is the product polarization.

Let \(\mathcal P_{\mathcal J}\) be the normalized relative Poincar\'e
line and set
\begin{equation}
 \mathcal D=(1,\lambda_{\mathcal J})^*\mathcal P_{\mathcal J}
 \quad\text{on }\mathcal J.
 \label{p3sw:support-relative-double-theta}
\end{equation}
The biextension identity gives
\(\phi_{\mathcal D}=\lambda_{\mathcal J}
 +\lambda_{\mathcal J}^{\vee}=2\lambda_{\mathcal J}\).
Concretely, translating the graph in \eqref{p3sw:support-relative-double-theta} by \(x\) produces the
two variable line factors
\(\mathcal P_{\mathcal J}(-,\lambda_{\mathcal J}x)\) and
\(\mathcal P_{\mathcal J}(x,\lambda_{\mathcal J}(-))\);
the remaining factor is a constant fiber.  These represent
\(\lambda_{\mathcal J}x\) and \(\lambda_{\mathcal J}^{\vee}x\).
On a smooth fiber it follows that
\(\mathcal D_s\otimes\mathcal O(-2\Theta_s)\in\operatorname{Pic}^0\):
their associated homomorphisms, equivalently their integral Riemann
forms, agree.  Centrally, \eqref{p3sw:support-product-polarization} and the product Poincar\'e bundle give
\(\mathcal D_0=D_1\boxtimes D_2\boxtimes D_3\), up to a constant
line, with \(\deg D_i=2\).  Every line in
\(\operatorname{Pic}^0(\mathcal J_0)\) is the external product of
three degree-zero lines.

\emph{The flat Abel limit and the torsion orbits.}
We next construct a flat family of Abel images whose central support is
the elliptic star.  Support is enough for the vanishing lemma, since the
ideal of any thickening of that star is contained in its reduced ideal.

The pointed Abel map on the smooth fibers,
\(q\mapsto\mathcal O(q-\sigma)\), extends to a morphism
\(\alpha:\mathcal C\to\mathcal J\).  We give the extension argument
to record the support of its limit.  The map is a rational
\(S\)-map, defined on the smooth-fiber open.  Elimination of
indeterminacy for a regular surface and a proper target
\cite[Lemma 54.4.3, Tag 0C5H]{P3SWStacks} gives a sequence of blowups
at closed points of the central fiber on which it becomes a
morphism.  Each reduced fiber of this sequence over an indeterminacy point
is a connected tree of rational curves.  Every morphism
\(\mathbb P^1\to\mathcal J_0\) is constant: pullback of the
translation-invariant one-forms is zero, and those forms span every
cotangent space, so the differential is zero in characteristic zero.
The resolved map is therefore constant on each such tree.
Its schematic graph image in
\(\mathcal C\times_S\mathcal J\) is integral, proper and birational
over \(\mathcal C\), with one underlying point in each fiber.
A proper morphism with finite fibers is finite
\cite[Lemma 37.44.1, Tag 02LS]{P3SWStacks}.  A finite integral
birational morphism to the normal surface \(\mathcal C\) is an
isomorphism, by integral closedness in its function field.  This
descends the resolved map to \(\alpha\).
The relative first Abel morphism and its tail twists are also given
by \cite[Section 2.2, (7), (9), and Proposition 2.4]{P3SWCoelhoPacini}
for a regular smoothing of a stable compact-type curve; the principal
component in the present star is \(R\).

We compute the specialization directly.  For a smooth \(q_0\in E_i\),
lift it \'etale-locally to a section \(q\).  The divisor \(E_i\) is
Cartier on the regular total surface and
\begin{equation}
 \mathcal O_{\mathcal C}(E_i)|_{E_i}=\mathcal O_{E_i}(-e_i),
 \quad
 \mathcal O_{\mathcal C}(E_i)|_R=\mathcal O_R(e_i),
 \quad
 \mathcal O_{\mathcal C}(E_i)|_{E_j}=\mathcal O_{E_j}\quad(j\ne i).
 \label{p3sw:support-tail-twist}
\end{equation}
Thus \(\mathcal O_{\mathcal C}(q-\sigma+E_i)\) has componentwise
degree zero and agrees generically with the Abel point.
Separatedness of \(\mathcal J\) identifies its special value with
\(\alpha(q)_0\).  The restriction on \(E_i\) is
\(\mathcal O_{E_i}(q_0-e_i)\), all other elliptic restrictions are
trivial, and \(\mathcal O_R(e_i-\sigma_0)\) is trivial.
There is no gluing torus in \eqref{p3sw:support-normalization-picard}.  A section through \(R\) maps to
zero, and continuity includes the nodes.  Identifying
\(\operatorname{Pic}^0(E_i)\) with \(E_i\) using \(e_i\), we get
\begin{equation}
 \alpha_0(R)=0,\qquad
 \alpha_0(q\in E_i)=(0,\ldots,\mathcal O_{E_i}(q-e_i),\ldots,0),
 \qquad
 \alpha_0(|C_0|)=|Z_0|.
 \label{p3sw:support-abel-limit}
\end{equation}
Here \(Z_0\) is exactly the three-axis star of Lemma
\ref{p3sw:support-elliptic-star}.

Let \(\mathcal A\subset\mathcal J\) be the scheme-theoretic closure of
the generic Abel image.  Its ideal is the contraction of the generic
ideal.  At a stalk over any base DVR its quotient embeds in its
localization obtained by inverting a uniformizer, and is therefore
torsion-free over that DVR, hence flat.  The ideal is coherent by
Noetherianity.  The morphism \(\alpha\) is proper, and
\(\mathcal C_\eta\) is dense in \(\mathcal C\).  Its closed topological image is
therefore the closure of its generic image, so \eqref{p3sw:support-abel-limit} gives
\begin{equation}
 |\mathcal A_0|=|Z_0|.
 \label{p3sw:support-closure-support}
\end{equation}
No reducedness of \(\mathcal A_0\) is claimed.  On the punctured base
the relative Abel map is a closed immersion.  Indeed, it is proper
with finite fibers, since every smooth-fiber Abel map is an
embedding, so it is finite by the cited finite-fiber criterion.
Finite pushforward commutes with base change by its affine description,
and the fiber maps
\(\mathcal O_{\mathcal J_s}\to\alpha_{s,*}\mathcal O_{\mathcal C_s}\)
are surjective.  Nakayama's lemma makes the total map surjective.
Its image is flat there and thus has a saturated ideal, so it agrees
with the restriction of the generic closure.  Hence the noncentral
\(\mathcal A_s\) are the actual Abel curves schemewise.

Put \(H_n=(\Z/n)^2\).  Choose isomorphisms
\(H_n\simeq E_i[n]\) and embed \(H_n\) diagonally in
\(\mathcal J_0[n]=\prod_iE_i[n]\).  Denote its two cyclic direct
factors by \(G_1,G_2\).  Multiplication by \(n\) on the abelian scheme
is proper, has invertible differential in characteristic zero and finite
geometric fibers, so its kernel is finite \'etale.  After an \'etale
localization, the chosen torsion points extend to sections with
their group relations.  They remain distinct, since equality of
two sections of a finite \'etale scheme is open and closed.  Thus the
two cyclic factors retain order \(n\) and trivial intersection.
The entire central \(H_n\) projects injectively into each \(E_i\).

For one cyclic factor, the translates \(t_g\mathcal A\) are flat
over \(S\), and their central supports are the disjoint stars of
Lemma \ref{p3sw:support-elliptic-star}.  A pairwise intersection is
closed in the proper scheme \(\mathcal J/S\), so its image in \(S\)
is closed and misses \(0\).  Removing finitely many such images
makes all these translates disjoint.  Their disjoint union
\(\mathcal Z\) is then a closed flat subscheme of \(\mathcal J\),
with actual Abel orbits on noncentral fibers and
\begin{equation}
 |\mathcal Z_0|=\left|\bigcup_g Z(g)\right|.
 \label{p3sw:support-orbit-limit}
\end{equation}
Apply the same argument to the negative closure \(-\mathcal A\)
for the other factor; its central support is still \(Z_0\).
This reasoning uses the closure and its support, and does not assume
that scheme-theoretic image or union formation commutes with the
closed-fiber base change.

\emph{Vanishing for every Picard twist.}
We now have flat orbit closures and a line \(\mathcal D\) representing
twice the principal polarization on every fiber.  It remains to make the
central vanishing uniform on a neighborhood of the smoothing parameter.
The ideal sequence for \(\mathcal Z\) is exact after every base
change to a fiber: both \(\mathcal O_{\mathcal J}\) and
\(\mathcal O_{\mathcal Z}\) are \(S\)-flat.  It also makes
\(I_{\mathcal Z}\) \(S\)-flat.  On
\(\mathcal J\times_S\widehat{\mathcal J}\) the coherent sheaf
\begin{equation}
 \operatorname{pr}_{\mathcal J}^*
        (I_{\mathcal Z}\otimes\mathcal D)\otimes\mathcal P_{\mathcal J}
 \label{p3sw:support-universal-ideal}
\end{equation}
is flat over \(\widehat{\mathcal J}\): locally it is the base change
of an \(S\)-flat module, followed by tensoring with a line.
The projection to the regular scheme \(\widehat{\mathcal J}\) is
projective.  Lemma \ref{p3sw:support-finite-model} makes the locus
\begin{equation}
 \mathcal V=\{N\in\widehat{\mathcal J}_s\text{ for some }s:
 H^0(\mathcal J_s,I_{\mathcal Z_s}\otimes\mathcal D_s\otimes N)
 \ne0\}
 \label{p3sw:support-bad-picard-locus}
\end{equation}
closed.  On the central fiber the exact ideal sequence embeds
\(I_{\mathcal Z_0}\) in \(\mathcal O_{\mathcal J_0}\), and \eqref{p3sw:support-orbit-limit}
gives
\(I_{\mathcal Z_0}\subset I_{\bigcup_gZ(g)}\).
Lemma \ref{p3sw:support-elliptic-star}, applied to every external
degree-two triple, therefore makes
\(\mathcal V\cap\widehat{\mathcal J}_0\) empty.
Properness of \(\widehat{\mathcal J}/S\) makes the image of
\(\mathcal V\) closed.  Removing it leaves a neighborhood of \(0\)
on which the vanishing holds for every Picard twist on every fiber.
Repeat for the negative orbit.

Choose a noncentral complex point of this neighborhood and put
\(C=\mathcal C_s\), \(p=\sigma_s\).  Under the group identification
\(\operatorname{Pic}^0(C)\simeq X\), the positive Abel curve is
\(C_p\).  The negative Abel point \(\mathcal O_C(3p-q)\) differs
from \(\omega_C(-p-q)\) by the fixed degree-zero line
\(\mathcal O_C(4p)\otimes\omega_C^{-1}\); hence the negative curve
is a translate of \(\Sigma_p\).
The difference between \(\mathcal D_s\) and \(H^{\otimes2}\) is
absorbed by the arbitrary Picard twist.  The two factors thus give
exactly \eqref{p3sw:support-seed-orbits}--\eqref{p3sw:support-uniform-vanishing}.
\end{proof}

Only the two within-part disjointness assertions are required; curves
from the two different parts may meet.  They never coincide.  In fact,
if \(t_sC_p=t_t\Sigma_p\), then the fixed degree-two line
\(\omega_C(-2p-s+t)\) has an effective divisor containing every
varying \(q\in C\).  It must have at least two independent sections,
contrary to nonhyperellipticity.  Hence every \(C_i\) and every
\(t_{-x}\Sigma_j\) have zero-dimensional intersection or are
disjoint, for every \(x\in X\).

\subsection{The relative Fourier--Mukai kernel and its coherent dual}

Fix a seed supplied by Proposition~\ref{p3sw:seed}.  We now realize the
cohomological argument following that proposition: first identify the
fixed kernel as a flat family of product ideals, then use its fiber
amplitude to construct the coherent dual.

\subsubsection{The fixed transform and its kernel}

Put \(M=X\times\widehat X\).  A point \((x,N)\) of this product
will parameterize a translation and a Picard twist in one fixed kernel
on the chosen Jacobian \(X\).  The flatness proved below concerns this
family over \(M\); deformation of the sixfold itself is treated later.
We use the derived category of bounded coherent complexes and the
following fixed equivalence:
\begin{equation}
 \begin{split}
 \Psi(F)&=
 R\operatorname{pr}_{\widehat X,*}
  (\operatorname{pr}_X^*F\otimes\mathcal P^{-1})[3],
 \\
 \mu:X\times X&\longrightarrow X\times X,\qquad
                \mu(x,y)=(x+y,y),
 \\
 \Phi&=(1\times\Psi)\circ\mu^*,\qquad
 \widetilde\Phi=\Phi\circ\bigl((H\boxtimes H)\otimes-\bigr).
 \end{split}
 \label{p3sw:support-fixed-transform}
\end{equation}
Here \(1\times\Psi\) is the relative transform on the second factor.
Mukai's inversion theorem
\cite[Theorem 2.2]{P3SWMukai1981} gives the equivalence for the
normalized Poincar\'e kernel.  Pulling that kernel by inversion on one
factor and shifting by \([3]\) are autoequivalences, so it also gives
the stated equivalence for \(\Psi\).  The shear and tensor factors in
\eqref{p3sw:support-fixed-transform} are equivalences as well.  Thus \eqref{p3sw:support-fixed-transform} specifies the transform
completely; no postcomposition shear is implicit.

Use the source and its transform
\begin{equation}
 I=I_\Sigma\boxtimes I_A,\qquad
 G=\widetilde\Phi(I)[-3].
 \label{p3sw:support-source-object}
\end{equation}
On \(X\times M\) use coordinates \((y,x,N)\), and let
\[
 \pi(y,x,N)=(x,N),\qquad \alpha_X(y,x,N)=y+x.
\]
Define
\begin{equation}
 \begin{split}
 \mathscr I&=\operatorname{pr}_y^*I_A,\qquad
 \mathscr J=\alpha_X^*I_\Sigma,\\
 \mathscr H&=\operatorname{pr}_y^*H\otimes\alpha_X^*H
             \otimes\operatorname{pr}_{y,N}^*\mathcal P^{-1},
 \qquad
 \mathscr F=\mathscr H\otimes
                 (\mathscr I\otimes^L\mathscr J).
 \end{split}
 \label{p3sw:support-kernel-sheaf}
\end{equation}
At first the last tensor product is derived.  Expanding the shear
in \eqref{p3sw:support-fixed-transform} changes \(I_\Sigma H\boxtimes I_AH\) into
\(\alpha_X^*(I_\Sigma H)\otimes^L\operatorname{pr}_y^*(I_AH)\);
the second Fourier--Mukai kernel then contributes
\(\mathcal P^{-1}[3]\).  The shift in \eqref{p3sw:support-source-object} cancels this \([3]\).
Consequently, as globally defined derived objects,
\begin{equation}
 G=R\pi_*\mathscr F.
 \label{p3sw:support-kernel-identity}
\end{equation}
This is the fixed-kernel form of
\cite[(6.2.3) and (9.2.1)]{Markman2025Secant}.

\begin{lemma}[Local Tor and product ideals]
\label{p3sw:support-local-tor}
Let \(V\) be a smooth threefold and let \(D_1,D_2\subset V\) be smooth
curves with finite, possibly empty, intersection.  Then
\begin{equation}
 \operatorname{Tor}^{\mathcal O_V}_k(I_{D_1},I_{D_2})=0
 \quad(k>0),
 \qquad
 I_{D_1}\otimes I_{D_2}\ \lhook\joinrel\longrightarrow\ \mathcal O_V
 \label{p3sw:support-tor-product}
\end{equation}
by multiplication, with image \(I_{D_1}I_{D_2}\).
The quotient from \(\mathcal O_V/I_{D_1}I_{D_2}\) to the structure
sheaf of the reduced union has zero-dimensional kernel.  The same
statements hold when each \(D_i\) is a disjoint union of smooth curves.
\end{lemma}

\begin{proof}
At an intersection point choose a regular sequence \(f_1,f_2\) for
\(D_1\).  Its Koszul resolution restricted to the regular
one-dimensional local ring of \(D_2\) has injective left arrow:
at least one \(\bar f_i\) is nonzero because the intersection is
zero-dimensional, and that local ring is a domain.  Thus
\(\operatorname{Tor}_k(\mathcal O_{D_1},\mathcal O_{D_2})=0\)
for \(k\geq2\).  The two ideal sequences give
\[
 \operatorname{Tor}_k(I_{D_1},I_{D_2})
 =\operatorname{Tor}_{k+2}(\mathcal O_{D_1},\mathcal O_{D_2})=0
 \quad(k>0)
\]
and \(\operatorname{Tor}_1(\mathcal O_{D_1},I_{D_2})=0\).
Tensoring the first ideal sequence with \(I_{D_2}\) proves the
injection and product image in \eqref{p3sw:support-tor-product}.  Away from the intersection the
product ideal equals the reduced-union ideal, proving the assertion
about its kernel.  For disjoint unions, at any point at most one
component from each part is present, so the same local calculation
applies.  This is the local fact behind
\cite[Appendix 11, pp.~91--93]{Markman2025Secant}, in the form needed
here.
\end{proof}

\begin{proposition}[The relative kernel and all its base changes]
\label{p3sw:relative-kernel}
For the seed of Proposition \ref{p3sw:seed}, the derived tensor in
\eqref{p3sw:support-kernel-sheaf} is the ordinary coherent sheaf
\(\mathscr I\mathscr J\subset\mathcal O_{X\times M}\).
Both this ideal and its quotient in \(\mathcal O_{X\times M}\) are
flat over \(M\).  Ideal multiplication and both ideal sequences
commute with every base change \(T\to M\).  In particular,
\(\mathscr F\) is \(M\)-flat, and at \(m=(x,N)\) it has the ordinary
fiber
\begin{equation}
 \begin{split}
 \mathscr F_m&=I_{Z_x}\otimes\mathcal M_{x,N},\\
 I_{Z_x}&=I_A I_{\Sigma_x},\qquad
 \Sigma_x=\coprod_jt_{-x}\Sigma_j,\qquad
 \mathcal M_{x,N}=H^{\otimes2}\otimes a(-x)\otimes N^{-1}.
 \end{split}
 \label{p3sw:support-fiber-formula}
\end{equation}
Moreover \(G=R\pi_*\mathscr F\) has a local finite locally free
model, and for every \(h:T\to M\),
\begin{equation}
 Lh^*G\simeq R\pi_{T,*}(\mathscr F|_{X\times T}).
 \label{p3sw:support-kernel-base-change}
\end{equation}
\end{proposition}

\begin{proof}
The subscheme \(A\times M\) is \(M\)-flat, and the subscheme defined
by \(\mathscr J\) is isomorphic to \(\Sigma\times M\) under
\((y,x,N)\mapsto(y+x,x,N)\).  Their ideal sequences make
\(\mathscr I\) and \(\mathscr J\) \(M\)-flat.
Locally along \(A\times M\), its regular codimension-two embedding
gives the ideal resolution
\begin{equation}
 0\longrightarrow\mathcal O
 \longrightarrow\mathcal O^{\oplus2}
 \longrightarrow\mathscr I\longrightarrow0,
 \label{p3sw:support-ideal-resolution}
\end{equation}
whose restriction to each fiber is the corresponding curve-ideal
resolution.  Tensor \eqref{p3sw:support-ideal-resolution} with \(\mathscr J\).  On every fiber the
left arrow is injective by Lemma \ref{p3sw:support-local-tor}.

We use the following precise local flatness criterion.  For a local
map \(R\to S'\) with \(S'\) Noetherian, a finite \(S'\)-module \(Q\),
an \(R\)-flat module \(P\), and an \(R\)-linear map \(u:Q\to P\),
injectivity modulo the maximal ideal of \(R\) implies injectivity
of \(u\) and \(R\)-flatness of its cokernel
\cite[Lemma 10.99.1, Tag 00MD]{P3SWStacks}.
Its mechanism is useful here: flatness compares the maximal-ideal
filtrations and gives injectivity modulo every power; Krull
intersection gives injectivity.  Repeating modulo every ideal of
\(R\) and using the Tor criterion gives flatness of the cokernel.
Apply it at a total-space stalk to
\(\mathscr J\to\mathscr J^{\oplus2}\), with
\(R=\mathcal O_{M,m}\).  It proves that the tensor of \eqref{p3sw:support-ideal-resolution} is exact
on the left and has \(M\)-flat cokernel.  Where
\(\mathscr I=\mathcal O\) the conclusion is immediate.
Thus the derived tensor is the ordinary \(M\)-flat tensor.

Its multiplication map to \(\mathcal O_{X\times M}\) is injective
on every fiber by Lemma \ref{p3sw:support-local-tor}.  The same
criterion, now with the \(M\)-flat target
\(\mathcal O_{X\times M}\), gives an injection with \(M\)-flat
quotient.  Tensoring either exact sequence with any
\(\mathcal O_T\) stays exact because its quotient is \(M\)-flat.
This proves multiplication and arbitrary base change before
pushforward.  The fiber formula follows from
\(t_x^*H=H\otimes a(-x)\) and
\(\mathcal P|_{X\times\{N\}}=N\).
Finally \(M\) is smooth and regular and \(\pi\) is projective.
Lemma \ref{p3sw:support-finite-model}, applied to the actual flat
\(\mathscr F\), gives the local finite model and \eqref{p3sw:support-kernel-base-change}.
\end{proof}

The scheme \(Z_x\) may have embedded points at intersections.
Lemma \ref{p3sw:support-local-tor} says that its reduction is
\(A\cup\Sigma_x\) and that the difference has zero-dimensional
support.  In particular \(Z_x\) always has one-dimensional support.

\subsubsection{The shifted dual is an ordinary sheaf}

\begin{proposition}[The actual coherent dual sheaf]
\label{p3sw:coherent-sheaf}
Let \(D_M(-)=R\mathcal Hom_M(-,\mathcal O_M)\).  Every derived fiber
of \(G\) has cohomology only in degrees one and two, and its
degree-two cohomology vanishes on a nonempty open set.  Consequently
\begin{equation}
 E=D_M(G)[-1]
 \label{p3sw:support-dual-sheaf}
\end{equation}
is an actual coherent perfect sheaf on \(M\), locally with a
length-one locally free resolution, of rank
\begin{equation}
 r=8d>0.
 \label{p3sw:support-rank}
\end{equation}
No reflexivity or codimension assertion is required for this claim.
\end{proposition}

\begin{proof}
For \(m=(x,N)\), Proposition \ref{p3sw:relative-kernel} gives
\begin{equation}
 0\longrightarrow\mathscr F_m\longrightarrow\mathcal M_{x,N}
 \longrightarrow\mathcal O_{Z_x}\otimes\mathcal M_{x,N}
 \longrightarrow0.
 \label{p3sw:support-fiber-exact}
\end{equation}
The line \(\mathcal M_{x,N}\) is ample: it is \(H^{\otimes2}\)
tensored by a topologically trivial line.  Kodaira vanishing applies
on the smooth projective abelian threefold, whose canonical bundle
is trivial, and gives \(H^k(X,\mathcal M_{x,N})=0\) for \(k>0\)
\cite[Theorem 3, p.~1272]{P3SWKodaira1953}.
Hirzebruch--Riemann--Roch for a line bundle on a smooth projective
variety \cite[Main Theorem (M), p.~111]{P3SWHirzebruch1954}, with
the trivial tangent bundle of \(X\), gives
\begin{equation}
 h^0(X,\mathcal M_{x,N})
 =\chi(X,\mathcal M_{x,N})
 =\int_X\frac{(2\Theta)^3}{3!}=8.
 \label{p3sw:support-ambient-euler}
\end{equation}
Since \(I_{Z_x}\subset I_A\), \eqref{p3sw:support-uniform-vanishing} gives \(H^0(\mathscr F_m)=0\).
Since \(Z_x\) has one-dimensional support, \eqref{p3sw:support-fiber-exact} gives
\begin{equation}
 H^3(\mathscr F_m)=0,\qquad
 H^2(\mathscr F_m)=
 H^1(Z_x,\mathcal M_{x,N}|_{Z_x}).
 \label{p3sw:support-degree-two-fiber}
\end{equation}

Each difference image
\(\Sigma_j-C_i\) is the closed image of the proper surface
\(\Sigma_j\times C_i\), and has dimension at most two.  Hence
\begin{equation}
 M^\circ=
 \left(X\setminus\bigcup_{i,j}(\Sigma_j-C_i)\right)\times\widehat X
 \label{p3sw:support-disjoint-open}
\end{equation}
is nonempty and open.  For \(m\in M^\circ\), \(Z_x\) is the disjoint
union of \(2n\) smooth genus-three curves.  By \eqref{p3sw:support-abel-class},
\(\mathcal M_{x,N}\) has degree six on each.  Formula \eqref{p3sw:support-curve-duality} gives
zero \(H^1\) on each component since \(6>2\cdot3-2\).
Thus \eqref{p3sw:support-degree-two-fiber} vanishes on \(M^\circ\), for every \(N\).
Equation \eqref{p3sw:support-kernel-base-change} identifies all these cohomology groups with the
cohomology of the derived fibers of \(G\).

The last assertion of Lemma \ref{p3sw:support-finite-model} now gives
a local model
\begin{equation}
 K^1\xrightarrow{u}K^2
 \label{p3sw:support-two-term-model}
\end{equation}
for \(G\).  On \(M^\circ\) its fiber map is surjective.  Thus
\(u^*:(K^2)^*\to(K^1)^*\) is generically injective.  The variety
\(M\) is integral, and the source of \(u^*\) is locally free; its
generically zero kernel must therefore be zero.  Dualizing \eqref{p3sw:support-two-term-model} and
shifting as in \eqref{p3sw:support-dual-sheaf} leaves its sole cohomology in degree zero, with
resolution
\begin{equation}
 0\longrightarrow(K^2)^*\xrightarrow{u^*}(K^1)^*
 \longrightarrow E\longrightarrow0.
 \label{p3sw:support-sheaf-resolution}
\end{equation}
These local sheaves are the sole cohomology sheaf of the globally
defined object \(D_M(G)[-1]\), so they glue.  This constructs the
actual sheaf, rather than assuming \(E=(\mathcal H^1G)^*\).

Finally, on \(M^\circ\), each curve restriction in \eqref{p3sw:support-fiber-exact} has Euler
characteristic \(6+1-3=4\).  Equations \eqref{p3sw:support-fiber-exact}--\eqref{p3sw:support-ambient-euler} give
\begin{equation}
 \chi(\mathscr F_m)=8-2n\cdot4=-8d.
 \label{p3sw:support-fiber-euler}
\end{equation}
The Euler characteristic of the finite model is
\(-\operatorname{rank}K^1+\operatorname{rank}K^2\).
The rank of \eqref{p3sw:support-sheaf-resolution} is therefore \(8d\).
\end{proof}

\subsection{The coherent translation action and the quotient}

For the group action, enumerate the source cosets in \eqref{p3sw:support-seed-orbits} by
\(G_2\) and \(G_1\), respectively.  Translation on a closed
invariant subscheme gives its ideal the canonical translation
linearization, so \(I=I_\Sigma\boxtimes I_A\) has a genuine
\(G_1\times G_2\)-linearization for the source translations.
The word ``coherent'' in the next statement includes the composition
isomorphisms, not only one isomorphism for each group element.

\begin{proposition}[Coherent conjugation and the graph quotient]
\label{p3sw:coherent-action}
For \(g=(g_1,g_2)\in G_1\times G_2\), put
\begin{equation}
 \begin{split}
 b_g&=(g_1-g_2,\ a(-g_1-g_2))\in M,\\
 N_g&=\operatorname{pr}_X^*a(-g_1)
       \otimes\operatorname{pr}_{\widehat X}^*\mathcal P_{-g_2},
 \qquad
 \mathcal P_v=\mathcal P|_{\{v\}\times\widehat X}.
 \end{split}
 \label{p3sw:support-conjugated-action}
\end{equation}
The conjugation of the source translations by \(\widetilde\Phi\)
has explicit functors
\begin{equation}
 T_g(-)=N_g\otimes t_{b_g,*}(-).
 \label{p3sw:support-action-functors}
\end{equation}
There are unital associative line isomorphisms
\begin{equation}
 N_g\otimes t_{b_g,*}N_h\simeq N_{g+h}
 \label{p3sw:support-line-compositors}
\end{equation}
and compatible isomorphisms \(T_gG\simeq G\).
Perfect duality gives actual sheaf isomorphisms
\begin{equation}
 N_g^{-1}\otimes t_{b_g,*}E\simeq E
 \label{p3sw:support-dual-action}
\end{equation}
with the corresponding unital associative line compositors.
The homomorphism \(g\mapsto b_g\) is injective.  Its image
\(\bar G\) has order \(n^2\), acts freely by translations, and
\begin{equation}
 q:M\longrightarrow Y_0=M/\bar G
 \label{p3sw:support-quotient}
\end{equation}
is a finite \'etale isogeny of abelian sixfolds.
\end{proposition}

\begin{proof}
We verify the conjugation at kernels, which also accounts for its
compositors.  For \(\ell\in\widehat X\) write
\(L_\ell=\mathcal P|_{X\times\{\ell\}}\).
The normalized biextension has the kernel identities
\begin{equation}
 \begin{split}
 (t_v\times1)^*\mathcal P^{-1}
   &\simeq\mathcal P^{-1}\otimes
                    \operatorname{pr}_{\widehat X}^*\mathcal P_{-v},\\
 (1\times t_{-\ell})^*\mathcal P^{-1}
   &\simeq\mathcal P^{-1}\otimes\operatorname{pr}_X^*L_\ell .
 \end{split}
 \label{p3sw:support-poincare-kernel-identities}
\end{equation}
Projection formula and flat base change for translation, applied
to these kernels, give
\begin{equation}
 \Psi\circ t_{v,*}\simeq(\mathcal P_{-v}\otimes-)\circ\Psi,
 \qquad
 \Psi\circ(L_\ell\otimes-)\simeq t_{\ell,*}\circ\Psi.
 \label{p3sw:support-fourier-conjugation}
\end{equation}
These are the inverse-Poincar\'e versions of the translation--tensor
exchange in \cite[Section 3, (3.1)]{P3SWMukai1981}.
Tensor conjugation on the source is
\begin{equation}
 (H\otimes-)\circ t_{v,*}\circ(H^{-1}\otimes-)
 \simeq(a(-v)\otimes-)\circ t_{v,*}.
 \label{p3sw:support-tensor-conjugation}
\end{equation}
For the shear,
\begin{equation}
 \mu^*(a(u)\boxtimes a(v))
       \simeq a(u)\boxtimes a(u+v),\qquad
 \mu^{-1}t_{(g_1,g_2)}\mu=t_{(g_1-g_2,g_2)}.
 \label{p3sw:support-shear-conjugation}
\end{equation}
Thus tensor conjugation and shear give the line
\(a(-g_1)\boxtimes a(-g_1-g_2)\) and the translation
\((g_1-g_2,g_2)\).  Applying \eqref{p3sw:support-fourier-conjugation} on the second factor gives
\eqref{p3sw:support-conjugated-action}--\eqref{p3sw:support-action-functors}.  Commuting the target translation past
\(\mathcal P_{-g_2}\) introduces only a constant one-dimensional
biextension fiber; choosing its trivialization changes the graph
kernel identification by a scalar.

All the displayed identities are identities of integral kernels:
they are obtained by pullback, projection formula, and the
normalized biextension multiplication.  That multiplication is
associative.  Convolution with the fixed kernel of
\(\widetilde\Phi\) and its inverse therefore transports the genuine
source-translation compositors and their associativity diagrams.
The kernel of each functor in \eqref{p3sw:support-action-functors} is an invertible sheaf on the
graph of \(t_{b_g}\).  The product kernel and the kernel for \(g+h\)
have the same graph, since \(b_{g+h}=b_g+b_h\).  A degree-zero
derived isomorphism between these line sheaves is an actual line
isomorphism.  Any two choices differ by a scalar, since
\(H^0(M,\mathcal O_M)^\times=\C^\times\).  Transporting the
compositors through the chosen graph identifications consequently
gives \eqref{p3sw:support-line-compositors} with its unital associativity, and transports the source
linearization to \(T_gG\simeq G\).  This does not infer a scalar
composition from simplicity of \(G\) or \(E\); the scalars come
from the line kernels and their compositors.

Duality commutes with translation by an isomorphism and changes a
line factor to its inverse.  Dualize the preceding compatible
isomorphisms and invert them.  With the shift in \eqref{p3sw:support-dual-sheaf} this gives
\eqref{p3sw:support-dual-action} and the inherited associative line compositors.  Because \eqref{p3sw:support-dual-sheaf}
has only its degree-zero cohomology sheaf, these are actual sheaf
isomorphisms.

Finally, \(b_g=0\) implies \(g_1=g_2\) and \(2g_1=0\), since \(a\)
is injective.  The intersection \(G_1\cap G_2=0\) therefore makes
the kernel trivial.  Nonzero translations have no fixed points.
The quotient by this finite flat closed subgroup exists
\cite[Section 2, p.~136]{P3SWDeligneGabber}; its quotient map is a
\(\bar G\)-torsor, hence finite \'etale because that holds after the
trivializing base change.  The quotient is a proper smooth connected
group space.  Over \(\C\) it is a quasi-projective group scheme
\cite[Section 1.1(a), p.~131]{P3SWDeligneGabber}, hence is a projective
abelian variety.  This proves \eqref{p3sw:support-quotient}.
\end{proof}

We have constructed the positive-rank coherent perfect sheaf \(E\),
its coherent line-twisted action \eqref{p3sw:support-dual-action}, and
the finite \'etale quotient \eqref{p3sw:support-quotient}.  The action
is defined on the entire sheaf by the kernel construction; neither
simplicity nor reflexivity was needed.  In the subsequent obstruction
and character calculations we write
\[
 \mathcal I=I,\qquad \mathcal G=G,\qquad \mathcal E=E.
\]
The obstruction calculation applies to \(\mathcal G\), while the
character and deformation arguments use its coherent shifted dual
\(\mathcal E\).  The flat family constructed here belongs to the fixed
seed.  Subsequent sections use semiregularity to deform an associated
ordinary sheaf locally, and then propagate the resulting algebraic
classes through the period domain.
\section{Hochschild evaluation and the invariant obstruction space}
\label{p3:sec:obstruction}

The geometric construction has produced a coherent perfect sheaf with a
finite line-twisted action. To use its quotient in a deformation argument,
we must control the entire invariant derived obstruction space. We first
prove the exact Atiyah trace identity and its transport through the fixed
Fourier--Mukai equivalence, and then compute the invariant space for the
Abel-orbit seed.

\subsection{Hochschild evaluation and the Atiyah trace}
\label{p3sw:hh-evaluation}

We will relate the evaluation of a Hochschild class on a perfect
complex to contraction against its Chern character, and then compare
these two maps under a fixed Fourier--Mukai equivalence.  This is the
compatibility used in the obstruction argument of
\cite[Diagram~(8.3.4) and Lemma~8.3.10]{Markman2025Secant}.
The proof below calculates the normalized Hochschild--Kostant--Rosenberg
map on the formal diagonal of an abelian variety.  It then obtains the
trace identity from the diagonal adjunction.  In particular, the
calculation fixes one convention for the bivector sign on both sides.

\subsubsection{Grading, contraction, and the first-jet class}

Let \(A\) be a complex abelian variety of dimension \(d\), let
\(\Delta:A\hookrightarrow A\times A\) be the diagonal, and let
\(Q\in\operatorname{Perf}(A)\).  Write
\(\mathcal O_\Delta=\Delta_*\mathcal O_A\).  We use cohomological shifts, so a
vector bundle shifted by \([1]\) lies in degree \(-1\).  Put
\[
 V=H^0(A,T_A),\qquad
 \mathcal S_A=\operatorname{Sym}_{\mathcal O_A}(\Omega_A^1[1])
       =\bigoplus_{k=0}^d\Omega_A^k[k].
\]
The complex on the right has zero differential.  Translation gives
\(T_A=\mathcal O_A\otimes_{\mathbb C}V\).  Our gradings are
\begin{equation}
\begin{split}
 HT^r(A)&=\bigoplus_{a+b=r}H^a(A,\mathord{\bigwedge}^bT_A),\\
 H\Omega_m(A)&=\bigoplus_{k-h=m}H^h(A,\Omega_A^k),\\
 \mathbb H^r(A,\mathcal S_A)&=H\Omega_{-r}(A).
\end{split}
\label{p3sw:hh-eq:gradings}
\end{equation}
All exterior degrees in these sums are between \(0\) and \(d\).

We specify contraction including its order.  If
\(\iota_\xi:\bigwedge^kV^*\to\bigwedge^{k-1}V^*\) is ordinary
interior multiplication, define
\begin{equation}
 \iota^L_{\xi_1\wedge\cdots\wedge\xi_b}
   =\iota_{\xi_1}\circ\cdots\circ\iota_{\xi_b},
 \qquad s_b=(-1)^{b(b-1)/2}.
\label{p3sw:hh-eq:left-interior}
\end{equation}
Composition is read right to left.  Thus
\(\iota^L_{u\wedge v}=\iota^L_u\iota^L_v\), and
\[
 \iota^L_{\xi_1\wedge\cdots\wedge\xi_b}
       (\vartheta_1\wedge\cdots\wedge\vartheta_b)
   =s_b\det\bigl(\vartheta_j(\xi_i)\bigr).
\]
The factor \(s_b\) is the graded dual pairing of
\(\operatorname{Sym}^b(V[-1])\) with
\(\operatorname{Sym}^b(V^*[1])\): interleaving their odd factors
requires \(b(b-1)/2\) crossings.  The pairing of one vector and one
covector is positive.

Use invariant Dolbeault representatives, with the antiholomorphic
cochain before the holomorphic form.  For
\(\eta\in H^a(A,\mathcal O_A)\), \(u\in\bigwedge^bV\),
\(\omega\in H^h(A,\mathcal O_A)\), and
\(\vartheta\in\bigwedge^kV^*\), define
\begin{equation}
 (\eta\otimes u)\mathbin{\lrcorner}(\omega\otimes\vartheta)
   =(-1)^{bh}(\eta\wedge\omega)\otimes\iota^L_u\vartheta.
\label{p3sw:hh-eq:left-action}
\end{equation}
The products used with this action are
\begin{equation}
\begin{aligned}
 (\eta\otimes u)(\eta'\otimes v)
   &=(-1)^{b a'}(\eta\wedge\eta')\otimes(u\wedge v),
       &&\eta'\in H^{a'}(A,\mathcal O_A),\\
 (\omega\otimes\vartheta)(\omega'\otimes\vartheta')
   &=(-1)^{k h'}(\omega\wedge\omega')
                         \otimes(\vartheta\wedge\vartheta'),
       &&\omega'\in H^{h'}(A,\mathcal O_A).
\end{aligned}
\label{p3sw:hh-eq:products}
\end{equation}
These are the total-complex Koszul rules; in particular
\eqref{p3sw:hh-eq:left-action} is a left action for the first product.
They also determine the meaning of powers of a \((1,1)\)-class.
For example, if \(j(\xi)=\xi\mathbin{\lrcorner}\vartheta\) for such a
class, then
\(\xi\mathbin{\lrcorner}\vartheta^2=2j(\xi)\vartheta\).

Write \(\mathcal I_\Delta\) for the diagonal ideal.  Fix the conormal
identification
\[
 df=f\otimes1-1\otimes f
        \pmod{\mathcal I_\Delta^2}.
\]
Let \(a_\Delta\) be the class of the first infinitesimal diagonal:
\begin{equation}
 0\longrightarrow\Delta_*\Omega_A^1
 \longrightarrow\mathcal O_{A\times A}/\mathcal I_\Delta^2
 \longrightarrow\mathcal O_\Delta\longrightarrow0.
\label{p3sw:hh-eq:jet-extension}
\end{equation}
For a kernel \(M\) from \(X\) to \(Y\), we use
\(\Phi_M(-)=Rp_{Y*}(Lp_X^*(-)\otimes^{\mathbf L}M)\).
Evaluation of \(a_\Delta\) by this operation defines the raw
first-jet class
\[
 a_Q:Q\longrightarrow Q\otimes\Omega_A^1[1].
\]
For a bundle it is the class of its first jet sequence.  For a perfect
complex it is the same derived natural transformation.  The power
\(a_Q^b:Q\to Q\otimes\Omega_A^b[b]\) means successive applications
of this class, with each new one-form appended on the right, followed
by the exterior projection.  This fixes both the Yoneda and exterior
orders; \(a_Q^0\) is the identity.

For a coefficient complex \(C\), the perfect partial trace is
\[
 \operatorname{Tr}_{Q,C}:
 \operatorname{Hom}(Q,Q\otimes C)
       \longrightarrow\operatorname{Hom}(\mathcal O_A,C).
\]
Precisely, it is the composite
\[
 \mathcal O_A\xrightarrow{\mathrm{coev}}Q\otimes Q^\vee
 \xrightarrow{\mathrm{braid}}Q^\vee\otimes Q
 \xrightarrow{1\otimes f}Q^\vee\otimes Q\otimes C
 \xrightarrow{\mathrm{ev}\otimes1}C
\]
for \(f:Q\to Q\otimes C\).  These are the duality and symmetry maps
of perfect complexes, with their usual Koszul signs.  On a bounded
locally free representative this is supertrace.  The definition
applies without change to shifted \(C\), and it does not require a
cohomology sheaf of \(Q\) to be locally free.

\subsubsection{The normalized HKR map and the trace identity}

Set
\[
 H_A=L\Delta^*\mathcal O_\Delta,\qquad
 HH^r(A)=\operatorname{Ext}^r_{A\times A}
                   (\mathcal O_\Delta,\mathcal O_\Delta).
\]
We use the derived adjunction \(L\Delta^*\dashv\Delta_*\), with unit
\[
 \eta_\Delta:\mathcal O_\Delta\longrightarrow\Delta_*H_A.
\]
The usual sheafified local bar map has degree-\(n\) formula
\[
 [f_1|\cdots|f_n]\longmapsto
       \frac{df_1\wedge\cdots\wedge df_n}{n!},
\]
with the two outer coefficients multiplied.  This formula specifies
the normalized bar morphism
\(\mathbf I_A:H_A\to\mathcal S_A\); see
\cite[Section~4 and Corollary~4.2]{P3HHCaldararuII} for the local bar
construction.  An element of \(HT^r(A)\) is a morphism
\(\alpha:\mathcal S_A\to\mathcal O_A[r]\) by the graded dual pairing
in \eqref{p3sw:hh-eq:left-interior}.  Define its adjoint Hochschild
class by
\begin{equation}
 \operatorname{HKR}_A(\alpha)
  =\Delta_*\alpha\circ\Delta_*\mathbf I_A\circ\eta_\Delta:
       \mathcal O_\Delta\longrightarrow\mathcal O_\Delta[r].
\label{p3sw:hh-eq:adjoint-hkr}
\end{equation}

Let \(\chi_Q:HH^r(A)\to\operatorname{Ext}^r_A(Q,Q)\) denote
evaluation of a diagonal morphism on \(Q\).
For \(\alpha^{a,b}\in H^a(A,\bigwedge^bT_A)\), contraction with an
Atiyah power means
\(\alpha^{a,b}\mathbin{\lrcorner}a_Q^b
=(1_Q\otimes\alpha^{a,b})\circ a_Q^b\), using the graded-dual morphism
\(\alpha^{a,b}:\Omega_A^b[b]\to\mathcal O_A[a+b]\).
Define the raw Hodge Chern components by
\begin{equation}
 \operatorname{ch}_{H,j}(Q)
      =\operatorname{Tr}_{Q,\Omega_A^j[j]}
                      \left(\frac{a_Q^j}{j!}\right)
         \in H^j(A,\Omega_A^j),\qquad
 \operatorname{ch}_H(Q)=\sum_{j=0}^d\operatorname{ch}_{H,j}(Q).
\label{p3sw:hh-eq:hodge-character}
\end{equation}

\begin{proposition}[Hochschild--Atiyah trace identity]
\label{p3sw:hh-triangle}
Let \(A\) be a complex abelian variety of dimension \(d\), and let
\(Q\in\operatorname{Perf}(A)\).  Use the raw first-jet class \(a_Q\)
from \eqref{p3sw:hh-eq:jet-extension}, and the left-interior and HKR
conventions \eqref{p3sw:hh-eq:left-action} and
\eqref{p3sw:hh-eq:adjoint-hkr}.  Define
\[
 \operatorname{ob}_Q=\chi_Q\circ\operatorname{HKR}_A:
       HT^2(A)\longrightarrow\operatorname{Ext}^2_A(Q,Q).
\]
For \(\alpha=\sum_{a+b=2}\alpha^{a,b}\), this evaluation is
\[
 \operatorname{ob}_Q(\alpha)
       =\sum_{a+b=2}\alpha^{a,b}\mathbin{\lrcorner}
                                      \frac{a_Q^b}{b!}.
\]
On the full group \(\operatorname{Ext}^2_A(Q,Q)\), define
\begin{equation}
\begin{split}
 \sigma_Q&=(\sigma_{Q,j})_{0\leq j\leq d-2}:
   \operatorname{Ext}^2_A(Q,Q)
       \longrightarrow H\Omega_{-2}(A)
          =\bigoplus_{j=0}^{d-2}H^{j+2}(A,\Omega_A^j),\\
 \sigma_{Q,j}(v)&=
   \operatorname{Tr}_{Q,\Omega_A^j[j+2]}
        \left(\left(\frac{a_Q^j}{j!}\right)[2]\circ v\right).
\end{split}
\label{p3sw:hh-eq:sigma}
\end{equation}
The direct sum is zero if \(d<2\).  Set
\[
 c_Q(\alpha)=\alpha\mathbin{\lrcorner}\operatorname{ch}_H(Q)
       \in H\Omega_{-2}(A).
\]
Then
\begin{equation}
 \sigma_Q\circ\operatorname{ob}_Q=c_Q.
\label{p3sw:hh-eq:triangle}
\end{equation}
More precisely, with
\(\operatorname{ch}_{H,k}(Q)=0\) for \(k\notin[0,d]\), each component
is
\begin{equation}
 \sigma_{Q,j}(\operatorname{ob}_Q(\alpha))
    =\sum_{a+b=2}\alpha^{a,b}\mathbin{\lrcorner}
                          \operatorname{ch}_{H,j+b}(Q)
       \quad\text{in }H^{j+2}(A,\Omega_A^j).
\label{p3sw:hh-eq:triangle-components}
\end{equation}
\end{proposition}

The proof uses two identities on the diagonal.  We first show that
the normalized bar map carries the Hochschild action to the fixed
left contraction.  We then show that the adjunction unit becomes
the exponential of the first-jet class.  Naturality of the unit and
of perfect trace will give the stated triangle.

\subsubsection{The normalized HKR action on the formal diagonal}
We next verify globally the normalization and action of these maps on
an abelian variety.

The difference isomorphism
\[
 A\times A\longrightarrow A\times A,\qquad
 (x,y)\longmapsto(y,x-y)
\]
takes \(\Delta\) to \(A\times\{0\}\), and hence identifies the formal
completion along \(\Delta\) with \(A\times\widehat A_0\).
In characteristic zero a commutative formal group has logarithm
coordinates.  In this case they can be obtained by integrating the
closed invariant formal one-forms term by term, with constant term
zero.  The differential of the difference between the logarithm of
the formal sum and the sum of the two logarithms is zero, and its
constant term is zero; thus the logarithm carries the formal group law
to addition.  Its linear term is invertible.  Choose coordinates
\(z_1,\ldots,z_d\) whose linear terms are a basis of invariant
one-forms.  The completed structure sheaf is then
\[
 \mathcal R=\mathcal O_A\widehat\otimes_{\mathbb C}
                         \mathbb C[[z_1,\ldots,z_d]].
\]
These are formal normal coordinates.  No global nonconstant regular
or holomorphic functions on \(A\times A\) cutting out \(\Delta\)
are asserted.

We explain why this completion computes the derived operations that
will be used.  Noetherian completion is flat.  If \(R\) is a
Noetherian local ring along the diagonal, \(J\) is its diagonal ideal,
and \(F\to R/J\) is a finite free resolution, then
\(\widehat R\otimes_R F\) remains a resolution and
\begin{equation}
 \operatorname{Hom}_{\widehat R}
       (\widehat R\otimes_R F,R/J)
       =\operatorname{Hom}_R(F,R/J).
\label{p3sw:hh-eq:completion-hom}
\end{equation}
The same equality holds with a locally free \(R/J\)-module as target.
Such finite resolutions exist locally because the diagonal of the
smooth variety \(A\) is a regular immersion.  The identifications in
\eqref{p3sw:hh-eq:completion-hom} commute with composition and with
tensoring by \(R/J\).  They therefore identify the derived sheaf Hom
complexes on the diagonal, their Yoneda pairings, and the action on
\(L\Delta^*\mathcal O_\Delta\).  After sheaf cohomology they identify
the corresponding global operations.  In particular, the
diagonal-supported morphisms with locally free diagonal coefficients
below can be compared after completion and then descended by these
identifications.

Let \(e_i\) have cohomological degree \(-1\) and correspond to
\(dz_i\), and let \(\xi_i\) be the dual invariant vector.  A global
finite free resolution on the formal completion is
\begin{equation}
 K=\left(\mathcal R\otimes\mathord{\bigwedge}^{\bullet}V^*,
         \ d_K=\sum_i z_i C_i\right),\qquad
 C_i=\iota_{\xi_i},\qquad
 \varepsilon:K\longrightarrow\mathcal O_A.
\label{p3sw:hh-eq:koszul}
\end{equation}
Here \(K^{-n}=\mathcal R\otimes\bigwedge^nV^*\).
The formal variables form a regular sequence, so \(K\) resolves the
diagonal.  Its pullback and derived Hom are represented by
\begin{equation}
\begin{split}
 \mathcal O_A\otimes_{\mathcal R}K
   &=\mathcal O_A\otimes\mathord{\bigwedge}^{\bullet}V^*
       =\mathcal S_A,\\
 \mathcal{H}om_{\mathcal R}(K,\mathcal O_A)
   &=\mathcal O_A\otimes\mathord{\bigwedge}^{\bullet}V,
\end{split}
\label{p3sw:hh-eq:finite-diagonal-complexes}
\end{equation}
with zero internal differentials and with the graded dual
identification just specified.

It remains to identify this splitting with the normalized bar map,
rather than merely with some splitting of cohomology.  Complete the
sheafified local bar resolution on the small diagonal in
\(A^{n+2}\).  Taking the last point as base, use the formal logarithm
to write \(z^{(0)},\ldots,z^{(n)}\) for the logarithms of the other
points minus the last point.  The outer \(\mathcal R\)-action uses
\(z^{(0)}\).  Each completed bar term is a formal power series
extension of the Noetherian sheaf \(\mathcal R\), hence is flat over
\(\mathcal R\).  The ordinary bar contracting homotopy after
forgetting one outer action inserts a unit.  It is continuous for
these formal-adic topologies, so its identities extend to the
completion.  Thus the augmented completed bar complex is a flat
resolution of the diagonal.

There is a global map from \(K\) to this resolution whose term in
degree \(-n\) is
\begin{equation}
 F_n(e_{i_1}\wedge\cdots\wedge e_{i_n})
   =\sum_{\tau\in S_n}\operatorname{sgn}(\tau)
       z^{(1)}_{i_{\tau(1)}}\cdots z^{(n)}_{i_{\tau(n)}}.
\label{p3sw:hh-eq:koszul-bar}
\end{equation}
An interior bar face identifies two successive middle points, so the
alternating determinant on the right becomes zero.  The last face
sets \(z^{(n)}=0\), again giving zero.  The first face sets
\(z^{(1)}=z^{(0)}\).  Expansion along that column gives, for
\(I=(i_1,\ldots,i_n)\),
\[
 bF_n(e_I)
   =\sum_{\ell=1}^n(-1)^{\ell-1}z^{(0)}_{i_\ell}
          F_{n-1}(e_{I\setminus i_\ell})
   =F_{n-1}(d_K e_I).
\]
In degree zero the map is the identity on the augmentation.  It is
therefore a comparison of resolutions.  It is global on the formal
neighborhood: the logarithm makes every face substitution just used
compatible on that neighborhood.

The normalized local bar formula defining \(\mathbf I_A\) restricts
to this completed resolution.  Indeed, it uses only the first jet in
each middle variable along the small diagonal, so is continuous for
the completion; naturality of differentials identifies the formulas
on overlaps.  After pullback to the diagonal, each term of
\eqref{p3sw:hh-eq:koszul-bar} maps to
\(dz_{i_1}\wedge\cdots\wedge dz_{i_n}/n!\): the permutation sign
and the exterior sign cancel.  The \(n!\) terms sum to \(dz_I\).
Consequently \(\mathbf I_A\) is the identity
\(e_I\mapsto dz_I\) on the first complex in
\eqref{p3sw:hh-eq:finite-diagonal-complexes}.  In particular it is a
quasi-isomorphism.  Because the comparison
\eqref{p3sw:hh-eq:koszul-bar} is global, this conclusion has no
additional off-diagonal or higher-cohomology component.  This is the
restriction of the usual sheafified bar HKR morphism, not a newly
chosen splitting.

We now compute its cap action.  The \(C_i\) have degree \(1\), satisfy
\[
 C_iC_j=-C_jC_i,\qquad d_KC_i+C_i d_K=0,
\]
and, for \(u=\xi_1\wedge\cdots\wedge\xi_b\), give a closed lift
\(C_u^L=C_{\xi_1}\cdots C_{\xi_b}\).
Its augmentation on \(K^{-b}\) is \(s_b\) times determinant
evaluation.  More explicitly, the map
\[
 \mathcal O_A\otimes\mathord{\bigwedge}^{\bullet}V
       \longrightarrow\mathcal{E}nd_{\mathcal R}(K),
 \qquad u\longmapsto C_u^L
\]
is a morphism of differential graded algebras and is a
quasi-isomorphism.  To see the last assertion, postcompose with
\(\varepsilon\).  The resulting map to
\(\mathcal{H}om_{\mathcal R}(K,\mathcal O_A)\) is the basis
isomorphism given by the graded dual pairing.  The map
\(\mathcal{E}nd_{\mathcal R}(K)\to
  \mathcal{H}om_{\mathcal R}(K,\mathcal O_A)\)
is itself a quasi-isomorphism, because \(K\) is bounded free and
\(\varepsilon\) is a quasi-isomorphism.  This gives global closed
lifts for the sheaf-level Yoneda and cap pairings.

For clarity about global cohomology, first perform the derived Hom
and derived pullback with the finite free complex \(K\).  The two
resulting complexes in
\eqref{p3sw:hh-eq:finite-diagonal-complexes} are finite complexes of
trivial \(\mathcal O_A\)-bundles.  The closed lifts \(C_u^L\) compute
their sheaf-level pairings.  Resolve only these
\(\mathcal O_A\)-coefficients by the ordinary Dolbeault resolution
\(\mathcal A_A^{0,\bullet}\), and totalize the pairings with Dolbeault
wedge.  The Dolbeault Poincar\'e lemma computes their sheaf
cohomology.  Projective GAGA for coherent cohomology and Ext identifies
these cohomology groups and pairings with their algebraic counterparts.
Neither a Dolbeault resolution of the infinite formal sheaf nor GAGA
for such a sheaf is used.

Every class in \(H^a(A,\mathcal O_A)\) has an invariant Dolbeault
representative; these representatives form
\(\bigwedge^aH^1(A,\mathcal O_A)\).  This standard computation for
a complex torus, together with the triviality of \(T_A\), shows that
decomposable classes \(\eta\otimes u\) span \(HT^r(A)\).
For such a class the totalized operator on a normal Koszul term \(k\)
and a Dolbeault cochain \(\omega\) is
\begin{equation}
 L_{\eta,u}(\omega\otimes k)
       =(-1)^{b|\omega|}(\eta\wedge\omega)\otimes C_u^L(k).
\label{p3sw:hh-eq:totalized-lift}
\end{equation}
This formula denotes the induced pairing on the finite diagonal
complexes above.  The normal lift is closed by its anticommutation
with \(d_K\); totalization is closed because \(\eta\) is closed.
Its augmentation is the cochain for the graded dual pairing.
The sign moves the degree-\(b\) normal operator past the Dolbeault
cochain.

Take the Yoneda convention
\(\beta\smile\gamma=\beta[|\gamma|]\circ\gamma\).
Composing \eqref{p3sw:hh-eq:totalized-lift} gives the first product
in \eqref{p3sw:hh-eq:products}, because
\(C_u^LC_v^L=C_{u\wedge v}^L\).  Tensoring the same lift with the
diagonal gives precisely \eqref{p3sw:hh-eq:left-action}.
The roof for \(\eta_\Delta\) on \(K\) has denominator
\(\varepsilon:K\to\mathcal O_A\) and numerator the projection
\(K\to\mathcal O_A\otimes_{\mathcal R}K\).  The augmentation just
calculated is therefore the adjoint class
\eqref{p3sw:hh-eq:adjoint-hkr}.  Derived adjunction and the
quasi-isomorphism \(\mathbf I_A\) show that
\(\operatorname{HKR}_A:HT^r(A)\to HH^r(A)\) is an isomorphism.
The calculation proves, in particular,
\begin{equation}
 \mathbf I_A\bigl((L\Delta^*\operatorname{HKR}_A(\alpha))h\bigr)
    =\alpha\mathbin{\lrcorner}\mathbf I_A(h),
 \quad
 \alpha\in HT^2(A),\quad h\in\mathbb H^0(A,H_A).
\label{p3sw:hh-eq:cap-hkr}
\end{equation}
Here a derived endomorphism acts on hypercohomology by composition.
The same calculation of the pairings accounts for every summand of
the target \(H\Omega_{-2}(A)\), since invariant representatives span
the relevant groups.

The cochain normalization can also be read directly from
\eqref{p3sw:hh-eq:koszul-bar}.  In unshifted differential-operator
notation it is
\begin{equation}
 \operatorname{HKR}_A(\xi_1\wedge\cdots\wedge\xi_b)
   =\frac{s_b}{b!}\sum_{\tau\in S_b}\operatorname{sgn}(\tau)
        \xi_{\tau(1)}\otimes\cdots\otimes\xi_{\tau(b)}
       \quad\text{in }HH^b(A).
\label{p3sw:hh-eq:cochain-sign}
\end{equation}
Indeed, evaluation on \(F_b\) gives \(s_b\) times determinant
evaluation, the augmentation of \(C_u^L\).
This is the cochain convention in
\cite[Equation~(3.20)]{P3HHCRVdB}.  In particular \(s_2=-1\);
replacing left interior multiplication by determinant-order
contraction would precompose both evaluation and contraction by
\(\operatorname{diag}(1,1,-1)\) on
\(H^2(\mathcal O_A)\oplus H^1(T_A)\oplus
H^0(\bigwedge^2T_A)\).  We make no such replacement.

The formal coordinates, their comparison, and
\eqref{p3sw:hh-eq:totalized-lift} are fixed by diagonal translation.
Translations also fix invariant Dolbeault representatives.  Thus
translations act trivially on \(HT^r(A)\), and
\(\operatorname{HKR}_A\) identifies this with their action on
\(HH^r(A)\).  Finally, \(T_A\) is holomorphically trivial, so its
positive Chern classes and its Atiyah class vanish and
\(\operatorname{td}(A)=\sqrt{\operatorname{td}(A)}=1\).
The explicit identity \eqref{p3sw:hh-eq:cap-hkr} is the abelian
specialization of the compatibility in
\cite[Theorem~1.4]{P3HHCRVdB}; it was proved here without using the
general formality theorem or a separate convention for dual chains.

\subsubsection{The exponential first-jet class and evaluation}

The preceding comparison fixes the homological map.  We next compute
its adjoint on a perfect complex; this supplies the factorials in the
obstruction and trace maps.

On \(K^{-1}\), the extension
\eqref{p3sw:hh-eq:jet-extension} is represented by
\(e_i\mapsto dz_i\).  The sign can be checked by mapping
\(K^0=\mathcal R\) to \(\mathcal R/(z)^2\) and \(K^{-1}\) to
\((z)/(z)^2\) by \(e_i\mapsto z_i\).  The square with the Koszul
differential commutes, and the conormal identification is exactly the
one in \eqref{p3sw:hh-eq:jet-extension}.  A closed lift is
\[
 \mathcal A(k)=\sum_i C_i(k)\otimes dz_i:
       K\longrightarrow(K\otimes V^*)[1].
\]
Successive application, with the forms appended in the order fixed
above, followed by exterior projection, is the lift
\begin{equation}
 \mathcal A^{(b)}(k)=
 \sum_{i_1,\ldots,i_b}C_{i_b}\cdots C_{i_1}(k)
       \otimes dz_{i_1}\wedge\cdots\wedge dz_{i_b}.
\label{p3sw:hh-eq:jet-power-lift}
\end{equation}
It represents the spliced \(b\)-fold first-jet class: composing the
next closed lift splices the next extension, while the displayed form
order is the appended order.  On
\(e_{j_1}\wedge\cdots\wedge e_{j_b}\), its nonzero augmented terms
are indexed by permutations.  The contraction and exterior factors
have the same permutation sign, which cancels.  The result is
\(b!\,dz_{j_1}\wedge\cdots\wedge dz_{j_b}\).  The augmentation is
zero on other Koszul degrees.

On the other hand, the degree-\(b\) component of
\(\Delta_*\mathbf I_A\circ\eta_\Delta\) is the projection from
\(K^{-b}\) followed by \(e_I\mapsto dz_I\).  The preceding
calculation and \eqref{p3sw:hh-eq:completion-hom} therefore prove,
on the original diagonal, the identity of morphisms
\begin{equation}
 \Delta_*\mathbf I_A\circ\eta_\Delta
    =\sum_{b=0}^d\frac{a_\Delta^b}{b!}:
       \mathcal O_\Delta\longrightarrow
                   \bigoplus_{b=0}^d\Delta_*\Omega_A^b[b].
\label{p3sw:hh-eq:universal-exponential}
\end{equation}
This is a direct abelian proof of the universal exponential identity
of \cite[Proposition~4.4]{P3HHCaldararuII}, with its sign and divided
powers specified by the same Koszul comparison.

Let
\[
 AH_Q=\Phi_{\eta_\Delta}(Q):Q\longrightarrow Q\otimes H_A .
\]
The projection formula for a diagonal kernel identifies its target
as written.  Applying the derived integral-kernel operation to
\eqref{p3sw:hh-eq:universal-exponential} gives
\begin{equation}
 (1_Q\otimes\mathbf I_A)\circ AH_Q
      =\sum_{b=0}^d\frac{a_Q^b}{b!}
          :Q\longrightarrow Q\otimes\mathcal S_A .
\label{p3sw:hh-eq:perfect-exponential}
\end{equation}
For a vector bundle, the first infinitesimal diagonal is finite over
the second factor, so its evaluation is the ordinary first jet
sequence.  For a perfect complex the same equality is the derived
kernel operation, computed locally on a bounded complex of vector
bundles.  It commutes with splicing and exterior projection.
Consequently the powers in
\eqref{p3sw:hh-eq:perfect-exponential} are the raw powers \(a_Q^b\)
already defined, also when \(Q\) is shifted.

Equations
\eqref{p3sw:hh-eq:adjoint-hkr} and
\eqref{p3sw:hh-eq:perfect-exponential} give
\begin{equation}
 \chi_Q(\operatorname{HKR}_A(\alpha))
       =(1_Q\otimes\alpha)\circ
                          \sum_{b=0}^d\frac{a_Q^b}{b!}.
\label{p3sw:hh-eq:evaluation}
\end{equation}
If \(\alpha^{a,b}\in H^a(A,\bigwedge^bT_A)\), its morphism from
\(\mathcal S_A\) is zero on all summands except \(\Omega_A^b[b]\).
Its value in \eqref{p3sw:hh-eq:evaluation} is exactly
\(\alpha^{a,b}\mathbin{\lrcorner}(a_Q^b/b!)\), of Ext degree \(a+b\),
with the graded dual and total-complex signs fixed above.  This proves
the evaluation identity in the present perfect abelian setting;
compare \cite[Theorem~A]{P3HHHuang}.

The components in \eqref{p3sw:hh-eq:hodge-character} are the
unscaled \(d\log\) Hodge components.  To identify them with the ordinary
character, first consider a line bundle.  Its first jet extension is
the \(d\log\) transition cocycle; thus \eqref{p3sw:hh-eq:hodge-character} is the exponential
of that first Hodge Chern class.  Perfect trace is additive for a
bounded locally free exact triangle, as is seen from supertrace on
its cone.  Pullback naturality of the first-jet class and the
projective bundle formula allow a vector bundle to be pulled to its
flag bundle, where its successive line quotients give the line-bundle
calculation.  Pullback on the Hodge cohomology in question is
injective by that formula.  Finally, the resolution property of a
smooth projective variety represents every perfect class by an
alternating sum of vector bundles; additivity treats \(Q\).
This gives the character identification of
\cite[Theorem~4.5]{P3HHCaldararuII} in the present convention.

Let
\(\operatorname{ch}_{B,j}(Q)\in H^{2j}(A,\mathbb Q)\) denote the
rational Betti Chern component.  Under the Dolbeault--de Rham
identification
\(H^j(A,\Omega_A^j)=H^{j,j}(A)\subset H^{2j}(A,\mathbb C)\),
the raw comparison is
\begin{equation}
 \operatorname{ch}_{H,j}(Q)
       \longmapsto (2\pi i)^j\operatorname{ch}_{B,j}(Q).
\label{p3sw:hh-eq:tate-normalization}
\end{equation}
Indeed, the exponential sequence with kernel \(2\pi i\mathbb Z\)
sends the \(d\log\) cocycle to \(2\pi i\,c_1^B\), and the line
splitting and additivity just used give the formula in every weight.
Dividing by \((2\pi i)^j\) is the Tate-normalized comparison.  The trace identity
\eqref{p3sw:hh-eq:triangle} uses \(\operatorname{ch}_H\) throughout.

\subsubsection{Proof of the trace identity}

\begin{proof}[Proof of Proposition~\ref{p3sw:hh-triangle}]
For \(\beta:\mathcal O_\Delta\to\mathcal O_\Delta[2]\), write
\(a_A(\beta)=L\Delta^*\beta:H_A\to H_A[2]\).
Naturality of the adjunction unit, with its canonical shift
isomorphisms, is the equality
\[
 \Delta_*a_A(\beta)\circ\eta_\Delta
       =\eta_\Delta[2]\circ\beta.
\]
Evaluate this equality of diagonal morphisms on \(Q\).  The
projection formula gives
\begin{equation}
 (1_Q\otimes a_A(\beta))\circ AH_Q
       =AH_Q[2]\circ\chi_Q(\beta).
\label{p3sw:hh-eq:unit-module}
\end{equation}
This is an identity of derived morphisms before any trace is taken.

Define
\[
 h_Q=\operatorname{Tr}_{Q,H_A}(AH_Q)\in\mathbb H^0(A,H_A),
 \qquad
 \sigma_Q^H(v)
   =\operatorname{Tr}_{Q,H_A[2]}(AH_Q[2]\circ v)
        \in\mathbb H^2(A,H_A).
\]
The perfect partial trace commutes with a morphism of coefficients:
for \(f:C\to C'\) and \(w:Q\to Q\otimes C\),
\begin{equation}
 \operatorname{Tr}_{Q,C'}((1_Q\otimes f)\circ w)
      =f\circ\operatorname{Tr}_{Q,C}(w).
\label{p3sw:hh-eq:partial-trace}
\end{equation}
This follows directly from its coevaluation--braiding--evaluation
definition: \(f\) can be placed after evaluation since it acts only
on the coefficient factor.  Apply it to
\eqref{p3sw:hh-eq:unit-module}.  We obtain
\[
 \sigma_Q^H(\chi_Q(\beta))=a_A(\beta)\circ h_Q.
\]
There is no odd exchange sign, since the degree here is two.

By \eqref{p3sw:hh-eq:perfect-exponential} and
\eqref{p3sw:hh-eq:partial-trace},
\[
 \mathbf I_A(h_Q)=\operatorname{ch}_H(Q),\qquad
 \mathbf I_A(\sigma_Q^H(v))
     =\operatorname{Tr}_Q
          \left(\left(\sum_{j=0}^d\frac{a_Q^j}{j!}\right)[2]
                         \circ v\right).
\]
In the second equality the component with coefficient
\(\Omega_A^j[j+2]\) is \(\sigma_{Q,j}\); those with \(j>d-2\)
vanish because their sheaf cohomology degree exceeds \(d\).
Take \(\beta=\operatorname{HKR}_A(\alpha)\) and apply
\eqref{p3sw:hh-eq:cap-hkr}.  This proves
\eqref{p3sw:hh-eq:triangle}.
For the component statement, contraction of
\(\operatorname{ch}_{H,j+b}(Q)\in
 H^{j+b}(A,\Omega_A^{j+b})\) removes \(b\) holomorphic forms and
adds cochain degree \(a\).  Its target is
\(H^{j+b+a}(A,\Omega_A^j)=H^{j+2}(A,\Omega_A^j)\).
This gives \eqref{p3sw:hh-eq:triangle-components}.
In particular the factors \(1/b!\) in evaluation, \(1/j!\) in
\(\sigma_{Q,j}\), and \(1/(j+b)!\) in the Chern component have
already been accounted for; there is no further scalar.
\end{proof}

The factorization just proved is the perfect-complex specialization
of the character and module identities in
\cite[Proposition~6.2.1, Corollary~6.3.2, and
Proposition~6.4.4]{P3HHBF}.  Their differential-operator convention
has \(a_Q=-\operatorname{At}^{BF}_Q\).  Thus their terms
\((-1)^j(\operatorname{At}^{BF}_Q)^j/j!\) are exactly \(a_Q^j/j!\).
There is no residual sign in
\eqref{p3sw:hh-eq:triangle-components}.  The proof above establishes
the identity from the diagonal unit and the fixed bar comparison,
rather than importing a compatibility between two unspecified HKR
normalizations.

\subsection{Fourier--Mukai transport of evaluation and contraction}
\label{p3sw:hh-fourier-mukai}

We now show that a specified Fourier--Mukai equivalence preserves
the equality \(\ker\operatorname{ob}_Q=\ker c_Q\).  Evaluation is
transported by convolution of kernels.  For character contraction we
construct a homological map using perfect Serre pairings; this also
records the sign of the shift in the fixed transform.

\subsubsection{Enhanced convolution and evaluation}

Let \(A\) and \(B\) be smooth projective complex abelian varieties
of dimensions \(d_A,d_B\).  Let
\(\Phi=\Phi_K:\operatorname{Perf}(A)\to\operatorname{Perf}(B)\)
be a Fourier--Mukai equivalence with perfect kernel \(K\), equipped
with a fixed enhanced inverse kernel \(L\), inverse adjunction units
and counits, and their triangle identities.  Here \emph{enhanced}
means that convolution, morphisms, and the isomorphisms commuting
with cohomological shifts are those of the derived enhancement, not
just chosen identifications of objects in the triangulated category.

For kernels \(P\) from \(X\) to \(Y\) and \(R\) from \(Y\) to \(Z\),
write
\[
 R\star P
  =Rp_{XZ*}\bigl(Lp_{XY}^*P\otimes^{\mathbf L}Lp_{YZ}^*R\bigr).
\]
Thus \(\Phi_{R\star P}=\Phi_R\circ\Phi_P\).
We use the standard enhanced associativity and projection formula for
this convolution.  The products are smooth proper and the
projections are proper perfect morphisms, so these operations send
perfect kernels to perfect kernels.  The fixed adjunctions give
\[
 K\star L\simeq\mathcal O_{\Delta_B},
 \qquad L\star K\simeq\mathcal O_{\Delta_A}.
\]
Conjugation
\[
 \mathcal C_\Phi(M)=K\star M\star L
       :\operatorname{Perf}(A\times A)
           \longrightarrow\operatorname{Perf}(B\times B)
\]
is an exact enhanced equivalence, with the opposite conjugation as
inverse.  Let
\(t:\mathcal C_\Phi(\mathcal O_{\Delta_A})
       \xrightarrow{\sim}\mathcal O_{\Delta_B}\)
be its counit identification.  For \(\beta\in HH^r(A)\), define
\begin{equation}
 \Phi^{HH}(\beta)
       =t[r]\circ\mathcal C_\Phi(\beta)\circ t^{-1}
       \in HH^r(B).
\label{p3sw:hh-eq:hh-conjugation}
\end{equation}
It is a graded isomorphism and preserves our Yoneda convention,
because every displayed shifted morphism uses the enhanced shift
isomorphism.

For \(Q\in\operatorname{Perf}(A)\), let
\(\Phi_{\mathrm{Ext}}:\operatorname{Ext}^2_A(Q,Q)
\xrightarrow{\sim}\operatorname{Ext}^2_B(\Phi Q,\Phi Q)\)
be the induced map with the same shift structure.  Associativity of
convolution, evaluated on \(Q\), gives
\begin{equation}
 \chi_{\Phi Q}(\Phi^{HH}\beta)
       =\Phi_{\mathrm{Ext}}(\chi_Q(\beta)),
       \qquad \beta\in HH^2(A).
\label{p3sw:hh-eq:evaluation-square}
\end{equation}
Indeed, evaluation on \(\Phi Q\) is the composite
\(K\star\beta\star L\star K\) evaluated on \(Q\).  The inverse unit
and counit cancel the adjacent \(L\star K\) by the adjunction triangle
identities.  The outer \(t\) then gives
exactly the morphism obtained by applying \(\Phi\) to the evaluation
on \(Q\).  This proves the equality on morphisms, including their
shifts, rather than only identifying the dimensions of Ext groups.

Use the fixed HKR maps of
Proposition~\ref{p3sw:hh-triangle} to put
\[
 \Phi^{HT}
   =(\operatorname{HKR}_B)^{-1}\Phi^{HH}\operatorname{HKR}_A:
       HT^2(A)\xrightarrow{\sim}HT^2(B).
\]
The evaluation square is already established.  The contraction map has
the corresponding compatibility, with the homological comparison map
constructed below.

\begin{proposition}[Transport of evaluation and contraction]
\label{p3sw:hh-transport}
For the fixed enhanced Fourier--Mukai equivalence
\(\Phi:\operatorname{Perf}(A)\to\operatorname{Perf}(B)\), there exists an
isomorphism
\[
 \Phi^H_{HH}:H\Omega_{-2}(A)\xrightarrow{\sim}H\Omega_{-2}(B)
\]
depending only on that enhanced equivalence, not on a choice of perfect
complex, such that for every \(Q\in\operatorname{Perf}(A)\),
\begin{equation}
\begin{split}
 \Phi_{\mathrm{Ext}}\circ\operatorname{ob}_Q
     &=\operatorname{ob}_{\Phi Q}\circ\Phi^{HT},\\
 \Phi^H_{HH}\circ c_Q&=c_{\Phi Q}\circ\Phi^{HT}.
\end{split}
\label{p3sw:hh-eq:transport-squares}
\end{equation}
Here \(\operatorname{ob}\) and \(c\) have the conventions of
Proposition~\ref{p3sw:hh-triangle}, and \(\Phi_{\mathrm{Ext}}\) is the
induced enhanced map on degree-two Ext.  Consequently
\begin{equation}
\begin{split}
 \Phi^{HT}(\ker\operatorname{ob}_Q)
       &=\ker\operatorname{ob}_{\Phi Q},\\
 \Phi^{HT}(\ker c_Q)&=\ker c_{\Phi Q}.
\end{split}
\label{p3sw:hh-eq:kernel-transport}
\end{equation}
In particular \(\ker\operatorname{ob}_Q=\ker c_Q\) implies
\(\ker\operatorname{ob}_{\Phi Q}=\ker c_{\Phi Q}\).
\end{proposition}

We construct \(\Phi^H_{HH}\) using perfect Serre pairings and then
prove the two squares.  The construction will also retain the
Chern-character sign of a shift in the equivalence.

\subsubsection{Serre pairings and the homological map}

We state explicitly the duality used to define homological
transport.  On a smooth proper variety \(X\), Serre duality is the
perfect natural pairing
\[
 \operatorname{Hom}(P,R)^*
      \simeq\operatorname{Hom}(R,\mathsf S_XP),\qquad
 \mathsf S_XP=P\otimes\omega_X[\dim X],
\]
for perfect \(P,R\); write
\(\operatorname{Tr}_{X,P}:\operatorname{Hom}(P,\mathsf S_XP)
\to\mathbb C\) for its trace.  We use its standard expression as
the Serre trace on \(\mathcal O_X\) after the perfect partial trace.
In particular it is supertrace on a shifted complex.

Put
\[
 \Sigma_A=\Delta_*(\omega_A[d_A]),\qquad
 \Sigma_A^{-1}=\Delta_!\mathcal O_A
               =\Delta_*(\omega_A^{-1}[-d_A]).
\]
Here
\(\Delta_!=\mathsf S_{A\times A}^{-1}\Delta_*\mathsf S_A\)
is left adjoint to \(L\Delta^*\).  This is the adjunction obtained
by applying Serre duality to \(L\Delta^*\dashv\Delta_*\).
The displayed formula for \(\Delta_!\mathcal O_A\) follows from
\(\Delta^*\omega_{A\times A}=\omega_A^{\otimes2}\).
Define Hochschild homology by
\begin{equation}
 HH_m(A)=\operatorname{Hom}_{A\times A}
              (\Sigma_A^{-1}[m],\mathcal O_{\Delta_A})
          \simeq\mathbb H^{-m}(A,H_A),
\label{p3sw:hh-eq:hh-homology}
\end{equation}
where the isomorphism is this adjunction.  For
\(\beta\in HH^2(A)\) and
\(h:\Sigma_A^{-1}[m]\to\mathcal O_{\Delta_A}\), set
\[
 \beta\cap h=(\beta\circ h)[-2]\in HH_{m-2}(A).
\]
Naturality of the adjunction in
\eqref{p3sw:hh-eq:hh-homology} identifies this with the action of
\(L\Delta^*\beta\) used in \eqref{p3sw:hh-eq:cap-hkr}.

For \(m=0,-2\), let
\[
 \mathscr T^m(A)=\operatorname{Hom}_{A\times A}
                 (\mathcal O_{\Delta_A},\Sigma_A[m]).
\]
The Serre functor of \(A\times A\) takes
\(\Sigma_A^{-1}[m]\) to \(\Sigma_A[m]\).  Therefore Serre duality
gives the perfect pairing
\begin{equation}
 \langle h,\delta\rangle_{A,m}
    =\operatorname{Tr}_{A\times A,\,\Sigma_A^{-1}[m]}
                      (\delta\circ h):
 HH_m(A)\otimes\mathscr T^m(A)\longrightarrow\mathbb C.
\label{p3sw:hh-eq:serre-pairing}
\end{equation}
This kernel description and its categorical character are the ones
studied in \cite[Theorems~6 and~10, Section~6.2]{P3HHCW}.
We give the required transport and its cap verification in these
pairings.

The Grothendieck-duality adjoint formulas for a perfect kernel on
smooth proper varieties are
\[
 K^R=K^\vee\otimes p_A^*\omega_A[d_A],
 \qquad
 K^L=K^\vee\otimes p_B^*\omega_B[d_B],
\]
viewed after exchange of the two factors as kernels from \(B\) to
\(A\).  We use these formulas with their adjunction units and
counits.  They give
\[
 K^R\simeq\Sigma_A\star K^L\star\Sigma_B^{-1}.
\]
For an equivalence the fixed inverse \(L\) carries both inverse
adjunctions.  Taking the mate of the last isomorphism with these
adjunctions gives the kernel Serre isomorphism
\[
 \theta:K\star\Sigma_A\xrightarrow{\sim}\Sigma_B\star K.
\]
Its evaluation is the trace-preserving Serre isomorphism: for
\(f:P\to\mathsf S_AP\),
\begin{equation}
 \operatorname{Tr}_{B,\Phi P}(\theta_P\circ\Phi f)
       =\operatorname{Tr}_{A,P}(f).
\label{p3sw:hh-eq:serre-coherence}
\end{equation}
To see why the coherence matters, transport the natural Serre
duality pairing through the equivalence.  Its uniqueness determines
the isomorphism \(\Phi\mathsf S_A\to\mathsf S_B\Phi\) making those
pairings commute; evaluation at the identity gives
\eqref{p3sw:hh-eq:serre-coherence}.  The adjoint formulas with the
same units and counits realize this isomorphism by \(\theta\).
An isomorphism between \(K^R\) and \(K^L\) as unstructured objects
would not specify it.  Smoothness, properness, and perfection are
exactly the hypotheses for these duality and adjunction statements.

Let \(u:\mathcal C_\Phi(\Sigma_A)\xrightarrow{\sim}\Sigma_B\) be
\(\theta\star L\) followed by \(\Sigma_B\star t\), and set
\[
 \Psi^m(\delta)
   =u[m]\circ\mathcal C_\Phi(\delta)\circ t^{-1}:
       \mathscr T^m(A)\xrightarrow{\sim}\mathscr T^m(B)
       \quad(m=0,-2).
\]
For \(\delta\in\mathscr T^{-2}(A)\) and \(\beta\in HH^2(A)\),
associativity and the middle cancellation \(t^{-1}t\) give
\begin{equation}
 \Psi^0(\delta[2]\circ\beta)
      =(\Psi^{-2}\delta)[2]\circ\Phi^{HH}(\beta).
\label{p3sw:hh-eq:test-composition}
\end{equation}
For \(\delta\in\mathscr T^0(A)\), the same associativity and the
inverse adjunction triangle identities give the evaluation identity
\begin{equation}
 (\Psi^0\delta)_{\Phi Q}
       =\theta_Q\circ\Phi(\delta_Q).
\label{p3sw:hh-eq:serre-evaluation}
\end{equation}
Here \(\delta_Q:Q\to\mathsf S_AQ\) is evaluation of the
Serre-valued diagonal morphism.  In this second equality the
triangles cancel the adjacent \(L\star K\) in the evaluation on \(Q\).
Thus both identities use the same \(t\), \(u\), and inverse
adjunctions.

Define the homological isomorphisms for \(m=0,-2\) by the perfect
pairings:
\begin{equation}
 \bigl\langle\Phi_*^{HH}h,\Psi^m\delta\bigr\rangle_{B,m}
       =\langle h,\delta\rangle_{A,m}
       \quad(h\in HH_m(A),\ \delta\in\mathscr T^m(A)).
\label{p3sw:hh-eq:homology-transport}
\end{equation}
Existence, uniqueness, and invertibility follow from the
nondegeneracy of \eqref{p3sw:hh-eq:serre-pairing} and the
invertibility of \(\Psi^m\).

The homological comparison map asserted in
Proposition~\ref{p3sw:hh-transport} is now defined by
\begin{equation}
 \Phi^H_{HH}
    =\mathbf I_B\,\Phi_*^{HH}\,\mathbf I_A^{-1}:
       H\Omega_{-2}(A)\xrightarrow{\sim}H\Omega_{-2}(B).
\label{p3sw:hh-eq:hodge-transport}
\end{equation}
The adjunctions \eqref{p3sw:hh-eq:hh-homology} are understood.
This map depends only on the fixed enhanced equivalence and may
mix the individual Hodge bidegrees within the displayed spaces.

\begin{proof}[Proof of Proposition~\ref{p3sw:hh-transport}]
The first equality in
\eqref{p3sw:hh-eq:transport-squares} follows from
\eqref{p3sw:hh-eq:evaluation-square},
\eqref{p3sw:hh-eq:evaluation}, and the definition of \(\Phi^{HT}\).

For the second square, we verify compatibility with the cap action
and the categorical character.  If \(\beta\in HH^2(A)\), \(h\in HH_0(A)\),
and \(\delta\in\mathscr T^{-2}(A)\), the definition of cap and
the trace of a shifted perfect object give
\[
 \langle\beta\cap h,\delta\rangle_{A,-2}
       =\langle h,\delta[2]\circ\beta\rangle_{A,0}.
\]
The shift is even, so its supertrace sign is \(+1\).  Applying this
on \(A\) and \(B\), and using
\eqref{p3sw:hh-eq:test-composition} and
\eqref{p3sw:hh-eq:homology-transport}, gives
\[
\begin{split}
 \langle\Phi_*^{HH}(\beta\cap h),\Psi^{-2}\delta\rangle_{B,-2}
  &=\langle\beta\cap h,\delta\rangle_{A,-2}\\
  &=\langle h,\delta[2]\circ\beta\rangle_{A,0}\\
  &=\langle\Phi_*^{HH}h,
             (\Psi^{-2}\delta)[2]\circ\Phi^{HH}\beta\rangle_{B,0}\\
  &=\langle\Phi^{HH}\beta\cap\Phi_*^{HH}h,
                                      \Psi^{-2}\delta\rangle_{B,-2}.
\end{split}
\]
The tests \(\Psi^{-2}\delta\) exhaust \(\mathscr T^{-2}(B)\).
Nondegeneracy proves the needed transport, and only in this degree:
\begin{equation}
 \Phi_*^{HH}(\beta\cap h)
       =\Phi^{HH}(\beta)\cap\Phi_*^{HH}(h),
       \qquad \beta\in HH^2(A),\quad h\in HH_0(A).
\label{p3sw:hh-eq:cap-transport}
\end{equation}

Let \(\mathfrak h_Q\in HH_0(A)\) correspond, under
\eqref{p3sw:hh-eq:hh-homology}, to
\(h_Q=\operatorname{Tr}_{Q,H_A}(AH_Q)\).
We verify its trace characterization using the same \(AH_Q\) as in
Proposition~\ref{p3sw:hh-triangle}.  For
\(\delta:\mathcal O_\Delta\to\Sigma_A\), let
\(\delta':H_A\to\omega_A[d_A]\) be its adjoint, so
\(\delta=\Delta_*\delta'\circ\eta_\Delta\).
The adjunction \(\Delta_!\dashv L\Delta^*\), defined by Serre
duality above, is trace-compatible.  Thus
\eqref{p3sw:hh-eq:partial-trace} and the standard partial-trace
expression of the Serre trace give
\begin{equation}
\begin{split}
 \langle\mathfrak h_Q,\delta\rangle_{A,0}
  &=\operatorname{Tr}_{A,\mathcal O_A}(\delta'\circ h_Q)\\
  &=\operatorname{Tr}_{A,\mathcal O_A}
       \left(\operatorname{Tr}_Q
              ((1_Q\otimes\delta')\circ AH_Q)\right)\\
  &=\operatorname{Tr}_{A,Q}(\delta_Q).
\end{split}
\label{p3sw:hh-eq:character-trace}
\end{equation}
The last equality uses evaluation of
\(\Delta_*\delta'\circ\eta_\Delta\).  It includes the supertrace
when \(Q\) is shifted.  Because the pairing is perfect, this
characterizes \(\mathfrak h_Q\) uniquely; it is the categorical
Chern character in \(HH_0(A)\).  Its image under \(\mathbf I_A\)
is \(\operatorname{ch}_H(Q)\) by
\eqref{p3sw:hh-eq:perfect-exponential}.

For a test \(\delta_B=\Psi^0\delta_A\), equations
\eqref{p3sw:hh-eq:homology-transport},
\eqref{p3sw:hh-eq:character-trace},
\eqref{p3sw:hh-eq:serre-evaluation}, and
\eqref{p3sw:hh-eq:serre-coherence} yield
\[
\begin{split}
 \langle\Phi_*^{HH}\mathfrak h_Q,\delta_B\rangle_{B,0}
  &=\langle\mathfrak h_Q,\delta_A\rangle_{A,0}\\
  &=\operatorname{Tr}_{A,Q}((\delta_A)_Q)\\
  &=\operatorname{Tr}_{B,\Phi Q}((\delta_B)_{\Phi Q})\\
  &=\langle\mathfrak h_{\Phi Q},\delta_B\rangle_{B,0}.
\end{split}
\]
The tests exhaust \(\mathscr T^0(B)\), so uniqueness gives
\begin{equation}
 \Phi_*^{HH}(\mathfrak h_Q)=\mathfrak h_{\Phi Q}.
\label{p3sw:hh-eq:character-transport}
\end{equation}

Now take \(\alpha\in HT^2(A)\).  By
\eqref{p3sw:hh-eq:cap-hkr}, the class \(c_Q(\alpha)\) is the
\(\mathbf I_A\)-image of
\(\operatorname{HKR}_A(\alpha)\cap\mathfrak h_Q\).
Equations \eqref{p3sw:hh-eq:cap-transport} and
\eqref{p3sw:hh-eq:character-transport} take its homological
transport to
\[
 \Phi^{HH}(\operatorname{HKR}_A(\alpha))
                    \cap\mathfrak h_{\Phi Q}.
\]
Apply \(\mathbf I_B\) and
\eqref{p3sw:hh-eq:cap-hkr} on \(B\).  The result is
\(c_{\Phi Q}(\Phi^{HT}\alpha)\), as required.  The three comparison
maps in \eqref{p3sw:hh-eq:transport-squares} are isomorphisms.
Taking kernels proves \eqref{p3sw:hh-eq:kernel-transport}.
\end{proof}

\paragraph{Source identification and shifts.}
The Serre-pairing definition also makes the source identification
of homological transport explicit.  Let \(\theta_{\mathcal C}\) be the
trace-preserving Serre isomorphism of the equivalence \(\mathcal C_\Phi\) between the two
kernel categories.  There is a unique
\[
 v:\mathcal C_\Phi(\Sigma_A^{-1})
                \xrightarrow{\sim}\Sigma_B^{-1}
\]
such that
\[
 u=\mathsf S_{B\times B}(v)\circ
             (\theta_{\mathcal C})_{\Sigma_A^{-1}}:
      \mathcal C_\Phi(\Sigma_A)\longrightarrow\Sigma_B.
\]
Indeed, apply the inverse Serre functor to
\(u\circ(\theta_{\mathcal C})_{\Sigma_A^{-1}}^{-1}\).
Preservation of the Serre trace by \(\theta_{\mathcal C}\) then
shows that \eqref{p3sw:hh-eq:homology-transport} is exactly
\[
 h\longmapsto t\circ\mathcal C_\Phi(h)\circ v[m]^{-1}.
\]
The source identification is obtained from Serre and left--right
adjunction coherence, whereas the target identification is \(t\).

The distinction fixes the sign of a shift.  Test it in
\(\operatorname{Perf}(\mathbb C)\) with \(K=\mathbb C[s]\) and
\(L=\mathbb C[-s]\).  The categorical trace of the identity of
\(\mathbb C[s]\) is
\(\mathrm{ev}\circ\mathrm{braid}\circ\mathrm{coev}=(-1)^s\).
Thus \eqref{p3sw:hh-eq:serre-coherence} forces
\(\theta=(-1)^s\) in this test.  It follows that \(\Psi^0\) is
multiplication by \((-1)^s\), and
\eqref{p3sw:hh-eq:homology-transport} makes
\(\Phi_*^{HH}\) on \(HH_0\) multiplication by \((-1)^s\) as well.
In agreement with perfect trace,
\[
 \operatorname{ch}_H(Q[s])=(-1)^s\operatorname{ch}_H(Q).
\]
In particular a shift \([-3]\) reverses this character.  Cancelling
\(K\) and \(L\) on both ends of a homology morphism without the Serre
source identification would lose this sign.  Degree-two evaluation
has an even shift, so no additional odd exchange sign occurs there;
it still uses the enhanced shift isomorphisms in
\eqref{p3sw:hh-eq:evaluation-square}.

The transported maps have source \(HT^2\); no naturality of
\(\sigma_Q\) on arbitrary Ext classes is asserted.
The map \(\Phi^H_{HH}\) is defined by Hochschild homology and the fixed
HKR maps, not identified with an independently normalized Betti integral
transform.
\subsection{The invariant obstruction space of the Abel-orbit seed}
\label{p3sw:invariant-section}

Retain the seed of Proposition~\ref{p3sw:seed}.  Thus \(C\) is a smooth
nonhyperelliptic curve of genus three, \(X=\operatorname{Pic}^2(C)\) has
the specified group origin and is identified with its Jacobian, and
\(H=\mathcal O_X(\Theta)\) is its principal theta line.  Put
\[
 \theta_{\mathrm B}=c_{1,\mathrm B}(H)\in H^2(X,\mathbb Q),
 \qquad \int_X\theta_{\mathrm B}^3=6,\qquad n=d+1\geq4.
\]
The two subschemes supplied there are disjoint unions
\[
 A=\coprod_{j=1}^n C_j,\qquad
 \Sigma=\coprod_{j=1}^n\Sigma_j
\]
of translates of an Abel curve and of its inverse, respectively.  The
cyclic groups \(G_2,G_1\subset X\), each of order \(n\), act transitively
on the components of \(A,\Sigma\), respectively.  Write
\[
 \mathcal I=I_\Sigma\boxtimes I_A\quad\text{on }X^2,\qquad
 G_{\mathrm{src}}=G_1\times G_2.
\]
Each ideal has its canonical translation linearization as the ideal of
an invariant subscheme, and \(\mathcal I\) has the product linearization.
On \(M=X\times\widehat X\), the fixed transform of
Proposition~\ref{p3sw:coherent-action} is
\[
 \widetilde\Phi=\Phi\circ\bigl((H\boxtimes H)\otimes-\bigr),
 \qquad \mathcal G=\widetilde\Phi(\mathcal I)[-3].
\]
Write \(\bar G\simeq G_{\mathrm{src}}\) for the underlying translation
group of its coherent line-translation action.

We will prove that the semiregularity map of
Proposition~\ref{p3sw:hh-triangle} is injective on the full invariant
derived group \(\operatorname{Ext}^2_M(\mathcal G,\mathcal G)^{\bar G}\).
The proof has two computations: the invariant group for \(\mathcal I\)
has dimension \(48\), and contraction with its raw Hodge Chern character
has rank \(48\).  The exact Hochschild--Atiyah identity then identifies
the entire invariant group with the evaluation image.  This calculation
concerns the Abel-orbit seed in Markman's secant-sheaf construction
\cite[Sections~8.3--8.4]{Markman2025Secant}; we give the invariant
calculation directly, using the independently movable orbit components.

\subsubsection{The Abel class and its normal directions}

An \emph{Abel curve} in \(X\) will mean a translate of the image of an
Abel embedding \(C\hookrightarrow X\), or a translate of its image under
\([-1]\), with the Jacobian and polarization conventions of
Section~\ref{p3sw:support-component}.

\begin{lemma}\label{p3sw:abel-character}
Let \(D\subset X\) be an Abel curve.  Its normalized Betti class and
Chern character are
\[
 [D]_{\mathrm B}=\frac{\theta_{\mathrm B}^2}{2},
 \qquad
 \operatorname{ch}_{\mathrm B}(\mathcal O_D)
   =\frac{\theta_{\mathrm B}^2}{2}
      -\frac{\theta_{\mathrm B}^3}{3}.
\]
Normal restriction of translation-invariant vector fields is an
isomorphism
\[
 H^0(X,T_X)\xrightarrow{\ \sim\ }H^0(D,N_{D/X});
\]
both spaces have dimension three.  If
\(Z=\coprod_{j=1}^nD_j\) is a disjoint union of \(n\) Abel curves, then
its ideal satisfies the full normalized Betti identity
\begin{equation}\label{p3sw:invariant-betti-character}
 \operatorname{ch}_{\mathrm B}(I_Z)
 =1-\frac n2\theta_{\mathrm B}^2+\frac n3\theta_{\mathrm B}^3.
\end{equation}
\end{lemma}

\begin{proof}
The fundamental-class formula is
\eqref{p3sw:support-abel-class}, proved there from the Jacobian lattice.
Translation acts trivially on cohomology, and inversion acts by \(+1\)
on \(H^4\), so it holds for every Abel curve in the stated sense.

For the inclusion \(i:D\hookrightarrow X\), use
Grothendieck--Riemann--Roch in its smooth proper form
\[
 \operatorname{ch}(i_*F)\operatorname{td}(X)
   =i_*\bigl(\operatorname{ch}(F)\operatorname{td}(D)\bigr)
\]
for a closed immersion of smooth projective varieties; this is the
form of the Riemann--Roch theorem recorded by Borel and Serre
\cite{p3swBorelSerreRR}.  The tangent bundle of \(X\) is trivial, and
\(\operatorname{td}(D)=1+\tfrac12c_1(T_D)\).  Since
\(\deg T_D=2-2\cdot3=-4\) and
\([\mathrm{pt}]_{\mathrm B}=\theta_{\mathrm B}^3/6\), it follows that
\[
 \operatorname{ch}_{\mathrm B}(\mathcal O_D)
 =i_*\operatorname{td}(D)
 =[D]_{\mathrm B}-2[\mathrm{pt}]_{\mathrm B}
 =\frac{\theta_{\mathrm B}^2}{2}
     -\frac{\theta_{\mathrm B}^3}{3}.
\]
For disjoint \(Z\), one has
\(\mathcal O_Z=\bigoplus_j\mathcal O_{D_j}\).  Additivity in
\(0\to I_Z\to\mathcal O_X\to\mathcal O_Z\to0\) gives
\eqref{p3sw:invariant-betti-character}.

It remains to calculate the normal sections.  For the original Abel
embedding the normal sequence is
\begin{equation}\label{p3sw:invariant-normal-sequence}
 0\longrightarrow T_C\longrightarrow
 \mathcal O_C\otimes H^1(C,\mathcal O_C)
 \longrightarrow N_{C/X}\longrightarrow0.
\end{equation}
Here the middle bundle is \(T_X|_C\).  The dual of the first arrow is
the evaluation map \(H^0(C,K_C)\otimes\mathcal O_C\to K_C\):
differentiating the Abel integral sends a tangent vector \(v\) to
the functional \(\omega\mapsto\omega(v)\).
Consequently, under Serre duality on \(C\), the dual of the map
\[
 H^1(C,T_C)\longrightarrow
 H^1(C,\mathcal O_C)\otimes H^1(C,\mathcal O_C)
\]
in the long exact sequence of
\eqref{p3sw:invariant-normal-sequence} is the multiplication map
\[
 H^0(C,K_C)\otimes H^0(C,K_C)\longrightarrow H^0(C,K_C^2).
\]
This is the usual cohomological Serre duality
\cite{p3swSerreDuality}, applied to the dual evaluation map.

For completeness, multiplication is onto here by an elementary
genus-three argument.  For an effective divisor \(D_2\) of length two,
curve Riemann--Roch, in the form
\(h^0(L)-h^0(K_C\otimes L^{-1})=\deg L+1-g\), gives
\[
 h^0(K_C-D_2)=h^0(D_2)=1.
\]
The last equality uses nonhyperellipticity.  Two independent sections
of a degree-two effective line, after removing their common base divisor,
would give a map to \(\mathbb P^1\) of degree at most two.  Degree one
would make \(C\) rational, and degree two would make it hyperelliptic.
Thus \(K_C\) separates every length-two subscheme.  It has three
sections and degree four, so it embeds \(C\) as an integral plane
quartic.  Such a curve is not contained in a conic.  Restriction of
quadrics therefore makes
\(\operatorname{Sym}^2H^0(K_C)\to H^0(K_C^2)\) injective.  Both spaces
have dimension six, the second by curve Riemann--Roch and
\(\deg K_C^2=8>2g-2\).  The multiplication map is consequently onto,
and its dual is injective.

Since \(H^0(C,T_C)=0\), the long exact sequence of
\eqref{p3sw:invariant-normal-sequence} now identifies
\(H^0(C,N_{C/X})\) with the middle bundle's constant sections
\(H^1(C,\mathcal O_C)=H^0(X,T_X)\).  This is exactly normal restriction
of translation-invariant vector fields.  Translation and inversion
carry the normal sequence and this isomorphism to every Abel curve.
\end{proof}

\subsubsection{The full invariant derived Ext group}

We next compute the derived deformation space of one orbit ideal.
The calculation keeps the local-to-global contribution of the normal
bundle; it does not replace derived Ext by the cohomology of an
endomorphism sheaf.

\begin{lemma}\label{p3sw:invariant-orbit-ext}
Let \(\Gamma\subset X\) be a finite subgroup of order \(n\), and let
\(Z=\coprod_{j=1}^nD_j\) be a disjoint union of Abel curves permuted
transitively by \(\Gamma\).  Give \(I_Z\) its canonical linearization.
There is an isomorphism of \(\Gamma\)-representations
\begin{equation}\label{p3sw:invariant-orbit-representation}
 \operatorname{Ext}^1_X(I_Z,I_Z)
 \simeq H^1(X,\mathcal O_X)
       \oplus\bigoplus_{j=1}^n H^0(X,T_X).
\end{equation}
The first summand is trivial, and \(\Gamma\) acts on the second by
its regular permutation action on the components and the identity on
each translation space.  In particular,
\[
 \dim_{\mathbb C}\operatorname{Ext}^1_X(I_Z,I_Z)^\Gamma=6,\qquad
 \dim_{\mathbb C}\operatorname{Ext}^2_X(I_Z,I_Z)^\Gamma=6,\qquad
 \operatorname{Hom}_X(I_Z,I_Z)^\Gamma=\mathbb C.
\]
\end{lemma}

\begin{proof}
At a point of \(Z\), write its regular codimension-two ideal in the
regular local ring \(R=\mathcal O_{X,x}\) as \(I=(f,g)\), where \(f,g\)
are part of a regular system of parameters.  It has the resolution
\[
 0\longrightarrow R\xrightarrow{(-g,f)}R^2
   \xrightarrow{(f,g)}I\longrightarrow0.
\]
If an endomorphism sends \(f,g\) to \(u,v\), its relation is
\(gu=fv\).  Regularity of \(f,g\) gives \(u=fr\), \(v=gr\) for a
unique \(r\in R\), so \(\operatorname{Hom}_R(I,I)=R\).  Applying
\(\operatorname{Hom}_R(-,I)\) to the resolution shows that the degree-one
cokernel is \(I/I^2\).  Under a change of generators, the dual of the
leftmost generator acquires the inverse determinant.  Consequently the
global degree-one sheaf is
\[
 (I_Z/I_Z^2)\otimes\det(I_Z/I_Z^2)^{-1}
 \simeq (I_Z/I_Z^2)^\vee=N_{Z/X}.
\]
The rank-two isomorphism sends \(v\otimes\ell^\vee\) to the functional
\(w\mapsto\ell^\vee(v\wedge w)\), so it is compatible with changes of
generators.  We view \(N_{Z/X}\) here as a sheaf on \(X\) supported on
\(Z\).  Away from \(Z\) the ideal is \(\mathcal O_X\), and the resolution
has length one.  Thus, globally,
\begin{equation}\label{p3sw:invariant-local-ext}
 \mathcal{H}om_X(I_Z,I_Z)=\mathcal O_X,\qquad
 \mathcal{E}xt_X^1(I_Z,I_Z)=N_{Z/X},\qquad
 \mathcal{E}xt_X^q(I_Z,I_Z)=0\quad(q\geq2).
\end{equation}

The standard local-to-global spectral sequence
\[
 E_2^{p,q}=H^p\!\left(X,\mathcal{E}xt_X^q(I_Z,I_Z)\right)
 \Longrightarrow \operatorname{Ext}^{p+q}_X(I_Z,I_Z)
\]
is the hypercohomology spectral sequence of
\(R\mathcal{H}om_X(I_Z,I_Z)\).  Its five-term sequence and
\eqref{p3sw:invariant-local-ext} give
\begin{equation}\label{p3sw:invariant-local-global}
 0\longrightarrow H^1(X,\mathcal O_X)
 \longrightarrow\operatorname{Ext}^1_X(I_Z,I_Z)
 \xrightarrow{\,e\,}H^0(Z,N_{Z/X})
 \longrightarrow H^2(X,\mathcal O_X).
\end{equation}
The first arrow is also the map induced by the perfect unit
\(\mathcal O_X\to R\mathcal{H}om(I_Z,I_Z)\): the perfect trace composed
with this unit is multiplication by the rank, which is one.

Translate the \(n\) curves independently.  The locus in \(X^n\) where
these translated proper curves remain disjoint is an open neighborhood
of the original tuple.  Over it their union is a disjoint union of flat
translated families.  Its ideal is flat as well, by the exact sequence
from the structure sheaf of the union.  Restriction of this family to
first-order tangent vectors gives a linear map
\[
 s:\bigoplus_{j=1}^nH^0(X,T_X)
   \longrightarrow\operatorname{Ext}^1_X(I_Z,I_Z).
\]
Indeed a flat deformation over \(\mathbb C[\epsilon]/(\epsilon^2)\)
gives the extension
\(0\to I_Z\to I_{Z,\epsilon}\to I_Z\to0\).
For a vector field \(\xi\), the moving generators are
\(f-\epsilon\xi(f)\) and \(g-\epsilon\xi(g)\).  Lifting the displayed
relation \((-g,f)\) gives
\[
 (-g)\bigl(f-\epsilon\xi(f)\bigr)
 +f\bigl(g-\epsilon\xi(g)\bigr)
 =\epsilon\bigl(g\xi(f)-f\xi(g)\bigr).
\]
Thus the local Ext cocycle is represented by
\(\overline{g\xi(f)-f\xi(g)}\in I/I^2\), up to the single sign in the
chosen Ext boundary convention.  Choose \(\ell=\bar f\wedge\bar g\) and
\(\ell^\vee(\ell)=1\).  Under the written determinant
identification \(v\otimes\ell^\vee\mapsto
(w\mapsto\ell^\vee(v\wedge w))\), this displayed representative sends
\(\bar f,\bar g\) to \(-\overline{\xi(f)},-\overline{\xi(g)}\),
respectively.  Reversing the boundary convention reverses both signs.
Consequently the edge map \(e\) sends the geometric deformation to its
normal restriction up to one fixed sign, independent of the vector field
and the component.  By Lemma~\ref{p3sw:abel-character}, \(e\circ s\)
is therefore the direct sum of the normal-restriction isomorphisms
multiplied by that fixed sign.  It follows that
\eqref{p3sw:invariant-local-global} reduces to the split short exact
sequence
\[
 0\longrightarrow H^1(X,\mathcal O_X)
 \longrightarrow\operatorname{Ext}^1_X(I_Z,I_Z)
 \xrightarrow{e}H^0(Z,N_{Z/X})\longrightarrow0,
\]
with splitting given by \(s\) after this signed normal-restriction
identification.

This splitting is equivariant.  Translation acts trivially on
\(H^1(X,\mathcal O_X)\), as is seen from translation-invariant Dolbeault
representatives, and it acts as the identity on \(H^0(X,T_X)\).
It sends the deformation of \(D_j\) to the deformation of its translated
component.  Since a group of order \(n\) acts transitively on \(n\)
components, that permutation action is regular.  This proves
\eqref{p3sw:invariant-orbit-representation}; its invariant summands have
dimensions three and three.

We also use Serre duality for perfect complexes on a smooth projective
threefold:
\[
 \operatorname{Ext}^i(F,G)^\vee
 \simeq\operatorname{Ext}^{3-i}(G,F\otimes\omega_X).
\]
This is the extension of the usual locally free Serre duality
\cite{p3swSerreDuality} by finite locally free resolutions.  Choose a
nonzero translation-invariant holomorphic three-form on \(X\).  For
\(F=G=I_Z\), the resulting perfect pairing is
\[
 \operatorname{Ext}^1(I_Z,I_Z)\otimes
 \operatorname{Ext}^2(I_Z,I_Z)
 \longrightarrow H^3(X,\mathcal O_X)\simeq\mathbb C.
\]
It is \(\Gamma\)-invariant: translations preserve the chosen form,
composition, and the perfect trace.  Thus these two Ext groups are dual
\(\Gamma\)-representations.  Averaging over the finite group in
characteristic zero identifies invariants of the dual with the dual of
the invariants, so the invariant \(\operatorname{Ext}^2\) also has
dimension six.  Finally \eqref{p3sw:invariant-local-ext} gives
\(\operatorname{Hom}(I_Z,I_Z)=H^0(X,\mathcal O_X)=\mathbb C\), with
trivial action.
\end{proof}

\begin{proposition}\label{p3sw:invariant-ext}
For the source object \(\mathcal I=I_\Sigma\boxtimes I_A\) and its
canonical \(G_{\mathrm{src}}=G_1\times G_2\) linearization,
\[
 \dim_{\mathbb C}
 \operatorname{Ext}^2_{X^2}(\mathcal I,\mathcal I)^{G_{\mathrm{src}}}
 =6+36+6=48.
\]
This is the invariant subspace of the full derived Ext group.
\end{proposition}

\begin{proof}
The ideal sheaves are perfect by their local length-one resolutions.
The perfect K\"unneth isomorphism for their derived global Hom complexes
is
\[
 \mathbf R\operatorname{Hom}_{X^2}(\mathcal I,\mathcal I)
 \simeq
 \mathbf R\operatorname{Hom}_X(I_\Sigma,I_\Sigma)
 \otimes_{\mathbb C}^{\mathbf L}
 \mathbf R\operatorname{Hom}_X(I_A,I_A).
\]
Taking degree two gives
\begin{align}
 \operatorname{Ext}^2_{X^2}(\mathcal I,\mathcal I)
 \simeq{}&
 \operatorname{Ext}^2_X(I_\Sigma,I_\Sigma)
       \otimes\operatorname{Hom}_X(I_A,I_A) \notag\\
 &\oplus\operatorname{Ext}^1_X(I_\Sigma,I_\Sigma)
       \otimes\operatorname{Ext}^1_X(I_A,I_A) \notag\\
 &\oplus\operatorname{Hom}_X(I_\Sigma,I_\Sigma)
       \otimes\operatorname{Ext}^2_X(I_A,I_A).
 \label{p3sw:invariant-product-ext}
\end{align}
Here the isomorphism follows by taking cohomology of the tensor product
over \(\mathbb C\) of the two derived global Hom complexes; over the
field \(\mathbb C\) there are no additional Tor terms.  It is equivariant
for the independent actions of \(G_1,G_2\).  Averaging in each factor
gives \((V_1\otimes V_2)^{G_1\times G_2}
=V_1^{G_1}\otimes V_2^{G_2}\).  Applying
Lemma~\ref{p3sw:invariant-orbit-ext} to the three summands of
\eqref{p3sw:invariant-product-ext} gives \(6\), \(6\cdot6\), and \(6\).
Because \eqref{p3sw:invariant-product-ext} is a decomposition of derived
Ext itself, the count retains every local-to-global contribution.
\end{proof}

\subsubsection{The rank of contraction with the raw Hodge character}

We now compute the map that Proposition~\ref{p3sw:hh-triangle} factors
through this \(48\)-dimensional invariant group.  Use the gradings
\eqref{p3sw:hh-eq:gradings}; in particular, for an abelian variety \(B\),
\[
 H\Omega_{-k}(B)=\bigoplus_{q-p=k}H^q(B,\Omega_B^p)
 \qquad(k\in\mathbb Z).
\]
For a perfect object \(Q\) on \(B\), extend character contraction to every
degree \(k\geq0\) by
\[
 c_Q^k:HT^k(B)\longrightarrow H\Omega_{-k}(B),\quad
 c_Q^k(\alpha)=\alpha\mathbin{\lrcorner}
                    \operatorname{ch}_{\mathrm H}(Q).
\]
Here \(\operatorname{ch}_{\mathrm H}\) is the raw first-jet character
\eqref{p3sw:hh-eq:hodge-character}, and contraction uses the
left-interior action \eqref{p3sw:hh-eq:left-action}, including its
Dolbeault Koszul sign.  Thus \(c_Q^2=c_Q\) in
Proposition~\ref{p3sw:hh-triangle}.

The normalization of the character matters for that identity.  Put
\(\theta_{\mathrm H}=a_H\in H^1(X,\Omega_X^1)\).  Under the raw de
Rham comparison,
\begin{equation}\label{p3sw:invariant-weight-dictionary}
 \theta_{\mathrm H}\longmapsto(2\pi i)\theta_{\mathrm B},\qquad
 \operatorname{ch}_{\mathrm H,k}(Q)
   \longmapsto(2\pi i)^k\operatorname{ch}_{\mathrm B,k}(Q)
 \quad\text{in codimension }k.
\end{equation}
This is the weightwise dictionary of
Proposition~\ref{p3sw:hh-triangle}; the raw first-jet class is the
negative of the Buchweitz--Flenner Atiyah convention.  In particular
Lemma~\ref{p3sw:abel-character} gives for either orbit ideal
\begin{equation}\label{p3sw:invariant-hodge-character}
 f_n:=\operatorname{ch}_{\mathrm H}(I_Z)
 =1-\frac n2\theta_{\mathrm H}^2+\frac n3\theta_{\mathrm H}^3.
\end{equation}
There is no codimension-one component.  Each term in
\eqref{p3sw:invariant-hodge-character} has the raw weight specified by
\eqref{p3sw:invariant-weight-dictionary}.

Use the abelian Hodge-cohomology identifications
\[
 \mathsf U=H^1(X,\mathcal O_X),\qquad
 \mathsf W=H^0(X,\Omega_X^1),\qquad
 \mathsf T=H^0(X,T_X)=\mathsf W^\vee,
 \qquad
 H^q(X,\Omega_X^p)=\wedge^q\mathsf U\otimes\wedge^p\mathsf W.
\]
These follow from the invariant Dolbeault forms on the complex torus.
The spaces \(\mathsf U,\mathsf W,\mathsf T\) each have dimension three.  The
polarization class \(\theta_{\mathrm H}\in\mathsf U\otimes\mathsf W\)
is nondegenerate, so in the fixed contraction convention
\[
 j:\mathsf T\longrightarrow\mathsf U,\qquad
 j(\xi)=\xi\mathbin{\lrcorner}\theta_{\mathrm H}
\]
is an isomorphism.

\begin{lemma}\label{p3sw:invariant-one-factor-rank}
For either orbit ideal \(I_Z\), the map \(c_{I_Z}^1\) is injective of
rank six, and \(c_{I_Z}^2\) is surjective of rank six.
\end{lemma}

\begin{proof}
An element of \(HT^1(X)=\mathsf U\oplus\mathsf T\) is a pair
\((\eta,\xi)\).  Using \eqref{p3sw:invariant-hodge-character}, its
projections under \(c_{I_Z}^1\) to
\(H^1(\mathcal O_X)\) and \(H^2(\Omega_X^1)\) are respectively
\begin{equation}\label{p3sw:invariant-degree-one-projections}
 \eta,\qquad -n\,j(\xi)\wedge\theta_{\mathrm H}.
\end{equation}
The first uses the rank term of the character; the second uses
\(-n\theta_{\mathrm H}^2/2\).  There is no contribution from \(\xi\)
to the first projection or from \(\eta\) to the second, because the
codimension-one character is zero.  The signs in
\eqref{p3sw:invariant-degree-one-projections} are those of the specified
left-interior convention.

The map \(v\mapsto v\wedge\theta_{\mathrm H}\) on \(\mathsf U\) is
injective.  If \(v\ne0\), choose bases with \(v=u_1\) and
\(\theta_{\mathrm H}=\sum_{i=1}^3u_i\otimes w_i\), absorbing the fixed
contraction signs in the dual basis.  The terms of
\(v\wedge\theta_{\mathrm H}\) with \(i=2,3\) are linearly independent
nonzero tensors.  Therefore the two projections in
\eqref{p3sw:invariant-degree-one-projections} determine \(\eta\) and
\(\xi\), proving the first assertion.

The only summands of \(H\Omega_{-2}(X)\) are
\begin{equation}\label{p3sw:invariant-degree-two-target}
 H^2(X,\mathcal O_X)=\wedge^2\mathsf U,\qquad
 H^3(X,\Omega_X^1)=\wedge^3\mathsf U\otimes\mathsf W.
\end{equation}
Each has dimension three.  The
\(H^2(\mathcal O_X)=\wedge^2\mathsf U\) block of \(HT^2(X)\) maps
identically to the first summand, using the rank term \(f_{n,0}=1\).
The \(H^1(T_X)=\mathsf U\otimes\mathsf T\) block has no contribution
there, since \(f_{n,1}=0\), and its projection to the second summand is
\begin{equation}\label{p3sw:invariant-degree-two-map}
 \eta\otimes\xi\longmapsto
 -n\,\eta\wedge j(\xi)\wedge\theta_{\mathrm H}.
\end{equation}
For each \(k\in\{1,2,3\}\), take \(\eta\) and \(j(\xi)\) to be the two
basis elements \(u_i,u_j\) complementary to \(u_k\).  Their wedge with
\(\theta_{\mathrm H}\) is a nonzero multiple of
\((u_1\wedge u_2\wedge u_3)\otimes w_k\).  Hence
\eqref{p3sw:invariant-degree-two-map} spans the second summand.
These two source blocks already give rank six; the target in
\eqref{p3sw:invariant-degree-two-target} has dimension six.  This proves
surjectivity.  The bivector block \(H^0(\wedge^2T_X)\) is not needed
for this lower bound.
\end{proof}

\begin{lemma}\label{p3sw:invariant-product-rank}
For \(\mathcal I=I_\Sigma\boxtimes I_A\), the map
\[
 c_{\mathcal I}^2:HT^2(X^2)\longrightarrow H\Omega_{-2}(X^2)
\]
has rank \(48\).
\end{lemma}

\begin{proof}
K\"unneth gives the domain decomposition
\[
 HT^2(X^2)=
 \bigl(HT^2(X)\otimes HT^0(X)\bigr)
 \oplus\bigl(HT^1(X)\otimes HT^1(X)\bigr)
 \oplus\bigl(HT^0(X)\otimes HT^2(X)\bigr).
\]
The product character is \(f_n\boxtimes f_n\), and \(HT^0(X)=\mathbb C\)
acts by \(c_{I_Z}^0(1)=f_n\ne0\).  The written left-action and product
formulas \eqref{p3sw:hh-eq:left-action} and
\eqref{p3sw:hh-eq:products}, applied to the K\"unneth decomposition,
give the external-product contraction rule with its total-complex
signs.  They therefore send the three blocks onto
\[
 \operatorname{im}(c_{I_\Sigma}^2)\otimes\mathbb C f_n,\qquad
 \operatorname{im}(c_{I_\Sigma}^1)\otimes
       \operatorname{im}(c_{I_A}^1),\qquad
 \mathbb C f_n\otimes\operatorname{im}(c_{I_A}^2).
\]
Their dimensions are \(6,36,6\), by
Lemma~\ref{p3sw:invariant-one-factor-rank}.  The target has the direct
K\"unneth decomposition
\[
 H\Omega_{-2}(X^2)
 =\bigoplus_{r+s=2}H\Omega_{-r}(X)\otimes H\Omega_{-s}(X).
\]
The first image lies in the
factorwise Hodge-imbalance summand \((2,0)\), the second in \((1,1)\),
and the third in \((0,2)\).  Here the imbalance of
\(H^q(X,\Omega_X^p)\) is \(q-p\); the two-factor Hodge K\"unneth
decomposition is a direct sum for these ordered pairs.  Hence the
three images cannot cancel, regardless of their K\"unneth signs.
Their ranks add to \(6+36+6=48\).
\end{proof}

\subsubsection{The exact rank comparison and the fixed transform}

The two calculations now have the same dimension.  We use the exact
local Hochschild results to turn this equality into injectivity on the
whole invariant group.  Write \(o_Q=\operatorname{ob}_Q\) for evaluation
and \(\sigma_Q\) for the perfect-trace map of
Proposition~\ref{p3sw:hh-triangle}.  For a perfect object \(Q\) on an
abelian variety \(B\), their domains and targets are
\[
 HT^2(B)\xrightarrow{\ o_Q\ }\operatorname{Ext}^2_B(Q,Q)
        \xrightarrow{\ \sigma_Q\ }H\Omega_{-2}(B).
\]

\begin{proposition}\label{p3sw:invariant-injectivity}
For the source \(\mathcal I=I_\Sigma\boxtimes I_A\) with its canonical
\(G_{\mathrm{src}}=G_1\times G_2\)-linearization,
\begin{equation}\label{p3sw:invariant-source-conclusion}
 \operatorname{im}(o_{\mathcal I})
   =\operatorname{Ext}^2_{X^2}(\mathcal I,\mathcal I)^{G_{\mathrm{src}}},
 \qquad
 \ker(o_{\mathcal I})=\ker(c_{\mathcal I}^2),
\end{equation}
and \(\sigma_{\mathcal I}\) is injective on this full invariant group.

Let \(M=X\times\widehat X\) and use the fixed enhanced equivalence
\[
 \widetilde\Phi=\Phi\circ\bigl((H\boxtimes H)\otimes-\bigr),
 \qquad
 \mathcal G=\widetilde\Phi(\mathcal I)[-3].
\]
Give \(\mathcal G\) the coherent line-translation action from
Proposition~\ref{p3sw:coherent-action}, and denote its underlying
translation group by \(\bar G\simeq G_{\mathrm{src}}\).  Then
\begin{equation}\label{p3sw:invariant-target-conclusion}
 \operatorname{im}(o_{\mathcal G})
   =\operatorname{Ext}^2_M(\mathcal G,\mathcal G)^{\bar G},
 \qquad
 \ker(o_{\mathcal G})=\ker(c_{\mathcal G}^2),
\end{equation}
and
\[
 \sigma_{\mathcal G}\big|_{
       \operatorname{Ext}^2_M(\mathcal G,\mathcal G)^{\bar G}}
 \quad\text{is injective}.
\]
\end{proposition}

\begin{proof}
Proposition~\ref{p3sw:hh-triangle} proves for every perfect \(Q\)
the exact identity
\[
 \sigma_Q\circ o_Q=c_Q^2
\]
with the definitions just given.  In particular it is an identity for
the perfect trace on derived Ext, with the same contraction convention
as the rank computation.  The components with \(q>\dim B-2\) vanish
because \(H^{q+2}(B,\Omega_B^q)=0\).

Translations act trivially on \(HT^2(X^2)\): they fix the invariant
tangent fields and the invariant Dolbeault representatives of
\(H^a(\mathcal O_{X^2})\).  The translation-invariant diagonal-kernel
HKR construction identifies this with their trivial action on
\(HH^2(X^2)\).  Apply the diagonal-kernel evaluation square
\eqref{p3sw:hh-eq:evaluation-square} to each translation equivalence,
and identify its translate of \(\mathcal I\) with \(\mathcal I\) by
the canonical linearization.  Naturality of evaluation then gives
\(g\cdot o_{\mathcal I}(\alpha)=o_{\mathcal I}(g\cdot\alpha)
=o_{\mathcal I}(\alpha)\) for \(g\in G_{\mathrm{src}}\).  Hence
\[
 \operatorname{im}(o_{\mathcal I})
 \subset
 W_{\mathrm{inv}}
 :=\operatorname{Ext}^2_{X^2}(\mathcal I,\mathcal I)^{G_{\mathrm{src}}}.
\]
Proposition~\ref{p3sw:invariant-ext} gives
\(\dim W_{\mathrm{inv}}=48\), while
Lemma~\ref{p3sw:invariant-product-rank} and the exact identity give
\[
 48=\operatorname{rank}(c_{\mathcal I}^2)
 \leq\operatorname{rank}(o_{\mathcal I})
 \leq\dim W_{\mathrm{inv}}=48.
\]
Thus the inclusion is equality and both maps have rank \(48\).
The identity also gives
\(\ker(o_{\mathcal I})\subset\ker(c_{\mathcal I}^2)\); equality of
ranks on their common domain makes the kernels equal.  If
\(w=o_{\mathcal I}(\alpha)\in W_{\mathrm{inv}}\) and
\(\sigma_{\mathcal I}(w)=0\), then \(c_{\mathcal I}^2(\alpha)=0\),
so the kernel equality gives \(w=0\).  This proves
\eqref{p3sw:invariant-source-conclusion} and source injectivity.

Apply Proposition~\ref{p3sw:hh-transport} to the fixed enhanced
equivalence \(Q\mapsto\widetilde\Phi(Q)[-3]\).  It includes tensoring
by \(H\boxtimes H\), the fixed Fourier--Mukai kernel, and the displayed
shift.  That proposition gives isomorphisms
\[
 u:HT^2(X^2)\xrightarrow{\sim}HT^2(M),\quad
 v:\operatorname{Ext}^2(\mathcal I,\mathcal I)
       \xrightarrow{\sim}\operatorname{Ext}^2(\mathcal G,\mathcal G),
 \quad
 \gamma:H\Omega_{-2}(X^2)\xrightarrow{\sim}H\Omega_{-2}(M)
\]
such that
\begin{equation}\label{p3sw:invariant-transport-squares}
 v\,o_{\mathcal I}=o_{\mathcal G}\,u,\qquad
 \gamma\,c_{\mathcal I}^2=c_{\mathcal G}^2\,u.
\end{equation}
The map \(\gamma\) is the HKR expression of the enhanced
Hochschild-homology isomorphism: its cap and character naturality,
including the shift sign, is exactly what gives the second equality.
No identification of \(\gamma\) with a separately normalized Betti integral
transform is used.

Proposition~\ref{p3sw:coherent-action} constructs the target action by
coherent conjugation of the canonical source action by this same
equivalence.  Hence \(v\) is an isomorphism of the corresponding
representations.  The lines in the target line-translation action
cancel on derived endomorphisms, and the underlying group is the
specified \(\bar G\simeq G_{\mathrm{src}}\).  Therefore
\[
 v(W_{\mathrm{inv}})
 =\operatorname{Ext}^2_M(\mathcal G,\mathcal G)^{\bar G}.
\]
The first equality of \eqref{p3sw:invariant-transport-squares} transports
the source image equality to the image equality in
\eqref{p3sw:invariant-target-conclusion}.  Both \(u\) and \(\gamma\) are
isomorphisms, so its second equality transports the source kernel
equality to the target kernel equality.  Finally apply
\(\sigma_{\mathcal G}o_{\mathcal G}=c_{\mathcal G}^2\), the target
instance of Proposition~\ref{p3sw:hh-triangle}.  An invariant \(z\)
has the form \(o_{\mathcal G}(\alpha)\); if
\(\sigma_{\mathcal G}(z)=0\), the target kernel equality gives \(z=0\).
This proves injectivity on the full target invariant Ext group.
\end{proof}

The fixed support construction identifies
\[
 \mathcal E=R\mathcal{H}om_M(\mathcal G,\mathcal O_M)[-1]
\]
as the actual coherent perfect sheaf of positive rank, with the
dual coherent action on the same underlying translation group.
Proposition~\ref{p3sw:twisted-descent} uses the full injection just
proved, transfers it by perfect duality, and then performs finite
\'etale twisted descent.  The exact dual sign and that derived descent
are left to that proposition.  The output here is injection on
\(\operatorname{Ext}^2(\mathcal G,\mathcal G)^{\bar G}\) itself, not
only on a chosen Hochschild subspace.  No identification of that
derived Ext group with
\(H^2(M,\mathcal H^0(R\mathcal{H}om_M(\mathcal G,\mathcal G)))\)
is used.
\section{The spin character and the polarized quadratic seed}
\label{p3:sec:spin}
\label{p3sw:spin-seed}

The preceding construction gives a coherent perfect sheaf
\(\mathcal E\) of rank \(r=8d\) on \(M=X\times\widehat X\).
Its normalized rational Betti character is
\begin{equation}
\label{p3sw:spin-eq:normalized-character}
 \kappa_M
   =\exp\!\left(-\frac{c_{1,\mathrm B}(\mathcal E)}r\right)
             \operatorname{ch}_{\mathrm B}(\mathcal E),
\end{equation}
where \(\operatorname{ch}_{\mathrm B}\) and \(c_{1,\mathrm B}\)
are the ordinary rational Betti Chern character and first Chern class.
We compute \(\kappa_M\) and identify the polarized quadratic structure
on this same sixfold.
The calculation uses the secant spinors of
Markman~\cite[Sections~2.4 and~6]{Markman2025Secant}.
We give a finite Clifford proof of the required transform identity,
so the identity applies to the algebraic spin group, including its
elements that do not preserve the initial Hodge structure.

\subsection{The fixed transform and its source character}

Put \(\ell_X=6=2g_X\), with \(g_X=3\), and set
\[
 L=H^1(X,\Q),\qquad U=L^\vee=H^1(\widehat X,\Q),\qquad V=L\oplus U.
\]
Choose \(e_1,\ldots,e_{\ell_X}\) with
\(\int_X e_1\wedge\cdots\wedge e_{\ell_X}=1\), and let \(f_1,\ldots,f_{\ell_X}\)
be the dual basis of \(U\).  Our normalized Poincar\'e line has class
\[
 p=c_{1,\mathrm B}(\mathcal P)=\sum_{i=1}^{\ell_X} e_i\wedge f_i.
\]
For an ordered subset \(P\subset\{1,\ldots,\ell_X\}\), write \(e_P\) and
\(f_P\) for the corresponding increasing products.  For disjoint
ordered subsets \(I,J\), let \(\epsilon_{I,J}\) be the sign in
\(e_Ie_J=\epsilon_{I,J}e_{I\cup J}\), and set it to zero if they
overlap.  Define reversal by
\[
 \tau(e_P)=(-1)^{|P|(|P|-1)/2}e_P .
\]

We use Grothendieck--Riemann--Roch in the following standard form:
if \(f:B\to C\) is a proper morphism of smooth projective complex
varieties and \(F\) is perfect, then
\[
 \operatorname{ch}_{\mathrm B}(Rf_*F)\operatorname{td}(C)
   =f_*\bigl(\operatorname{ch}_{\mathrm B}(F)\operatorname{td}(B)\bigr).
\]
This is the Riemann--Roch theorem of Borel and
Serre~\cite{p3swBorelSerreRR}.
In a Fourier transform between abelian varieties the Todd classes
are one.  This gives the usual cohomological Poincar\'e transform;
compare the calculation in
\cite[Section~4.2.3]{GolyshevLuntsOrlov2001}.
The second-factor transform in the fixed functor \(\Phi\) of
\eqref{p3sw:support-fixed-transform} uses the inverse Poincar\'e line
and the shift \([g_X]\).  It is therefore
\begin{equation}
\label{p3sw:spin-eq:fourier}
 F_-(\alpha)=(-1)^{g_X}(\pi_{\widehat X})_*
            \bigl(e^{-p}\pi_X^*\alpha\bigr).
\end{equation}
The factor \((-1)^{g_X}\) is the Chern-character sign of the shift.

\begin{lemma}[The fixed Fourier matrix]
\label{p3sw:fourier-matrix}
For \(P\subset\{1,\ldots,\ell_X\}\), \(k=|P|\), one has
\[
 F_-(e_P)=(-1)^{k(k+3)/2}\epsilon_{P,P^c}f_{P^c}.
\]
In particular \(F_-\) is a rational isomorphism preserving parity.
\end{lemma}

\begin{proof}
In \eqref{p3sw:spin-eq:fourier}, only the complementary indices in
\(e^{-p}\) can contribute to integration on \(X\).  Moving their
\(f\)-factors past the \(e\)-factors gives the coefficient
\[
 (-1)^{\,g_X+(\ell_X-k)(\ell_X-k+1)/2}\epsilon_{P,P^c}.
\]
The difference of the displayed exponent and \(k(k+3)/2\) is
\(2(g_X+1)(g_X-k)\), since \(\ell_X=2g_X\), and hence is even.  Every basis
vector goes to a signed complementary basis vector, proving the
remaining assertions.
\end{proof}

For \(\mu(x,y)=(x+y,y)\) define the graded K\"unneth maps
\begin{equation}
\label{p3sw:spin-eq:fixed-maps}
 \phi=(1\otimes F_-)\mu^*,\qquad
 \nu=(F_-\otimes1)\mu^*,\qquad
 D=F_-^{-1}\otimes F_-,\qquad
 \phi'=\phi(1\otimes\tau)=D\nu(1\otimes\tau).
\end{equation}
The input of \(D\) is written in the order \(U\oplus L\), and its
output in the order \(L\oplus U\); all identifications use the graded
shuffle.  These are fixed rational isomorphisms, determined only by
the oriented lattice, addition, and \(p\).  Equivalently,
\[
 \phi'(s\otimes t)=(-1)^{g_X}\int_{X_y}
          s(x+y)\tau(t)(y)e^{-\sum_i y_i z_i}.
\]
This is finite exterior-algebra coefficient extraction in
\(y_1\cdots y_{\ell_X}\); the output variables \(x,z\) represent \(L,U\).
In particular the maps do not vary with a complex period having the
given marking.

We now identify the tensor to which this transform applies.  Write
\(\Theta=c_{1,\mathrm B}(H)\in\Lambda^2L\) for the principal class
of the seed.  Both orbit ideals have \(n=d+1\) disjoint Abel
components, so Lemma~\ref{p3sw:abel-character} gives their common
Betti character
\[
 f_n=1-\tfrac n2\Theta^2+\tfrac n3\Theta^3.
\]
Thus the sheaves \(F_1=I_A(H)\) and \(F_2=I_\Sigma(H)\) used in
\(\Phi\) have common character
\[
 v=e^\Theta f_{d+1}
   =1+\Theta-\tfrac d2\Theta^2-\tfrac d6\Theta^3.
\]
Fix \(K=\Q(u)\), where \(u^2=-d\) and \(d\geq3\) is the integer
of the seed construction, and put
\[
 \lambda_+=e^{u\Theta},\qquad \lambda_-=e^{-u\Theta}
       \quad\text{in }(\Lambda L)_K.
\]
Expanding the exponentials in powers of \(\Theta\) gives
\begin{equation}
\label{p3sw:spin-eq:source-character}
 v=\tfrac12\bigl[(1+u^{-1})\lambda_+
                    +(1-u^{-1})\lambda_-\bigr].
\end{equation}
The two exponential classes are linearly independent: both have
constant term \(1\), and their degree-two terms are opposite and
nonzero.  Reversal acts by \((-1)^j\) on \(\Theta^j\), so it
interchanges \(\lambda_+\) and \(\lambda_-\).  The rational tensor
\(T=v\otimes\tau v\) therefore lies in the tensor square of their
two-dimensional \(K\)-span.  Grothendieck--Riemann--Roch gives
\begin{equation}
\label{p3sw:spin-eq:sheaf-transform}
 \beta:=\operatorname{ch}_{\mathrm B}\Phi(F_2\boxtimes F_1)
       =\phi(v\otimes v)=\phi'(T).
\end{equation}
To relate this class to the original sheaf, let \(\iota\) be the
ring involution acting by \((-1)^j\) on \(H^{2j}(M,\Q)\).
The shifts in \(\mathcal G=\Phi(F_2\boxtimes F_1)[-3]\) and
\(\mathcal E=D_M(\mathcal G)[-1]\) give
\begin{equation}
\label{p3sw:spin-eq:sheaf-duality}
 \mathcal E=D_M\bigl(\Phi(F_2\boxtimes F_1)\bigr)[2],\qquad
 \operatorname{ch}_{\mathrm B}(\mathcal E)=\iota\beta,\qquad
 \beta_0=r=8d.
\end{equation}
Here the rank is supplied by Proposition~\ref{p3sw:coherent-sheaf}.
We next introduce the Clifford action on \(\Lambda L\) and compare
its transport through \(\phi'\) with the graded action on
\(H^*(M)\).  The exponential correction will disappear under the
normalization defining \(\kappa_M\).

\subsection{The exact Clifford identity}

Give \(V\) the symmetric split form
\[
 Q(e_i,f_j)=\delta_{ij},\qquad Q(e_i,e_j)=Q(f_i,f_j)=0.
\]
Use the Clifford algebra with relation \(vw+wv=Q(v,w)\).
We use the standard Clifford construction of the spin group in this
normalization: the norm-one even Clifford normalizer
\(\operatorname{Spin}(V,Q)\) acts on \(V\) by conjugation through
\(\rho:\operatorname{Spin}(V,Q)\to SO(V,Q)\), and acts on a Clifford
module by multiplication.  The norm is Clifford reversion times the
element.  For this split rational construction, see
\cite[Sections~3.2 and~3.4--3.5]{GolyshevLuntsOrlov2001}.
Only this construction, not a realization of arbitrary spin elements
by derived autoequivalences, is used here.

On \(S=\Lambda L\) the Fock action is
\[
 c(e)=\epsilon_e=e\wedge-,\qquad c(f)=\iota_f.
\]
Write \(m_g=c(g)\).  Conjugation by the fixed \(\phi'\) defines the
rational algebraic representation
\[
 \rho'_g=\phi'(m_g\otimes m_g)(\phi')^{-1}
       \quad\text{on }\Lambda V.
\]
The following calculation is the matrix-unit Clifford realization
underlying \cite[Proposition~3.2.1(e)]{GolyshevLuntsOrlov2001},
with the signs of the fixed transform retained.

\begin{proposition}[Fixed Clifford identity]
\label{p3sw:spin-identity}
For every algebraic spin element \(g\), one has the exact equality
of rational representations
\begin{equation}
\label{p3sw:spin-eq:identity}
 \rho'_g=e^{p/2}\Lambda\rho(g)e^{-p/2}
       =e^{(p-\Lambda^2\rho(g)p)/2}\Lambda\rho(g).
\end{equation}
The exponentials denote exterior multiplication.  In particular
\[
 [\rho'_g(1)]_0=1,\qquad
 [\rho'_g(1)]_2=\tfrac12\bigl(p-\Lambda^2\rho(g)p\bigr).
\]
\end{proposition}

\begin{proof}
Exterior multiplication and contraction give every matrix unit
between the monomial basis vectors of \(S\): first project onto a
specified monomial by products of the commuting occupation
projections \(\epsilon_{e_i}\iota_{f_i}\) and
\(\iota_{f_i}\epsilon_{e_i}\), and then remove and insert the required
generators.  The ordered Clifford monomials span at most
\(2^{2\ell_X}\) dimensions, which equals \(\dim\operatorname{End}(S)\).
Thus
\[
 C(V,Q)\simeq\operatorname{End}(S).
\]
Its even part is the parity-preserving part: the inclusion follows
from parity of the generators, and both sides have dimension
\(2^{2\ell_X-1}\).  Put \(F=f_1\cdots f_{\ell_X}\)
in the Clifford algebra and define
\[
 \varphi:S\otimes S\longrightarrow C(V,Q),\qquad
 \varphi(s\otimes t)=sF\tau(t).
\]
On \(S\) this is the rank-one operator
\begin{equation}
\label{p3sw:spin-eq:rank-one}
 c(\varphi(s\otimes t))(r)
   =(-1)^{g_X}[\tau(t)\wedge r]_{\mathrm{top}}\,s.
\end{equation}
Indeed \(F\) contracts the full oriented monomial with sign
\((-1)^{\ell_X(\ell_X-1)/2}=(-1)^{g_X}\).  The pairing
\(\langle t,r\rangle=[\tau(t)\wedge r]_{\mathrm{top}}\) is perfect.
Exterior multiplication is self-adjoint for it by reversal of
products.  Contraction is self-adjoint as well: apply \(\iota_f\)
to a zero product in exterior degree \(\ell_X+1\) and use
\[
 \tau\iota_f|_{\Lambda^kL}
     =(-1)^{k-1}\iota_f\tau|_{\Lambda^kL}.
\]
Consequently Clifford reversion is adjunction for this pairing.
Every spin element has reversion norm one, and therefore preserves
the pairing.  Formula~\eqref{p3sw:spin-eq:rank-one} now proves that
\(\varphi\) is an isomorphism intertwining the diagonal spin action
with conjugation on \(C(V,Q)\).

For any bilinear form \(B\) on \(V\) with \(B(v,v)=Q(v,v)/2\), set
\[
 c_B(v)=\epsilon_v+\iota_{B(v,-)}
       \quad\text{on }\Lambda V,\qquad
 \psi_B(x)=c_B(x)1 .
\]
The Clifford relations follow directly from the exterior and
contraction relations.  On each ordered Clifford monomial
\(\psi_B\) has that exterior monomial as its leading term, so
\(\psi_B\) is a triangular isomorphism.  Take
\[
 B_0(e_i,f_j)=\delta_{ij},\quad B_0(f_i,e_j)=0,\qquad
 B_s=Q/2,\qquad
 A_0=B_0-B_s=\tfrac12\sum_i e^i\wedge f^i,
\]
with other entries of \(B_0\) zero.  Here \(e^i,f^i\) are the
covectors dual to \(e_i,f_i\).

The fixed map \(\nu\) is exactly the nonsymmetric Clifford symbol:
\begin{equation}
\label{p3sw:spin-eq:symbol}
 \nu(1\otimes\tau)=\psi_{B_0}\varphi .
\end{equation}
To check every sign, let \(I'=P\setminus I\), and let \(j=|J|\).
Both sides on \(e_P\otimes e_J\) equal
\begin{equation}
\label{p3sw:spin-eq:symbol-basis}
 (-1)^{j(j-1)/2}
 \sum_{I\subseteq P}
   \epsilon_{I',I}\epsilon_{I,J}
   (-1)^{\sum_{i\in I'}i-|I'|}
   f_{(I')^c}\wedge e_{I\cup J}.
\end{equation}
For the left side, expand \(\mu^*e_P\) and apply
Lemma~\ref{p3sw:fourier-matrix}.  For the right side use
\[
 c_{B_0}(e_i)=\epsilon_{e_i}+\iota_{f^i},\qquad
 c_{B_0}(f_i)=\epsilon_{f_i},
\]
and
\[
 \iota_{f^{i_1}}\cdots\iota_{f^{i_s}}f_{\{1,\ldots,\ell_X\}}
       =(-1)^{\sum_{i\in I'}i-|I'|}f_{(I')^c}.
\]
Moving the remaining \(e_I\) past the \(f\)'s changes
\(\epsilon_{I,I'}\) into \(\epsilon_{I',I}\).  This proves
\eqref{p3sw:spin-eq:symbol}.

For this finite exterior calculation define
\(\iota_{\alpha\wedge\beta}=\iota_\beta\iota_\alpha\) for covectors
on \(V\).  This fixes the contraction convention locally; the
Hochschild contraction used earlier was specified independently.
The ordinary commutator identity
\[
 [\iota_{\alpha\wedge\beta},\epsilon_v]
   =\alpha(v)\iota_\beta-\beta(v)\iota_\alpha
   =\iota_{(\alpha\wedge\beta)(v,-)}
\]
and commutation of even contractions with contractions give
\[
 c_{B_0}(v)=e^{\iota_{A_0}}c_{B_s}(v)e^{-\iota_{A_0}},
 \qquad
 \psi_{B_0}=e^{\iota_{A_0}}\psi_{B_s}.
\]
The second equality has no scalar, since
\(e^{-\iota_{A_0}}1=1\).  All exponentials are finite.
If \(R=\Lambda\rho(g)\), preservation of \(Q\) gives
\[
 Rc_{B_s}(v)R^{-1}=c_{B_s}(\rho(g)v),
 \qquad
 R\psi_{B_s}(x)=\psi_{B_s}(gxg^{-1}).
\]
Thus the symmetric symbol is spin-equivariant.

It remains to conjugate \(\iota_{A_0}\) by \(D\).  On the input
\(\Lambda(U\oplus L)\) fix the orientation
\(\Omega=f_1\cdots f_{\ell_X} e_1\cdots e_{\ell_X}\).  The exterior determinant
pairing of \(Q\) defines \(PD_Q\) by
\[
 (PD_Q\eta,\zeta)_Q=[\eta\wedge\zeta]_\Omega.
\]
For \(|P|=k\), \(|J|=j\), and \(\delta=k+j\), direct evaluation gives
\[
 PD_Q(f_Je_P)=(-1)^{kj}\epsilon_{J,J^c}\epsilon_{P,P^c}
                      e_{J^c}f_{P^c}.
\]
Applying Lemma~\ref{p3sw:fourier-matrix} to the two factors of \(D\)
therefore yields
\begin{equation}
\label{p3sw:spin-eq:duality}
 D|_{\Lambda^\delta V}
      =(-1)^{\,g_X+\delta(\delta+3)/2}PD_Q .
\end{equation}
For example, before parity reduction the exponent in the coefficient
of \(D(f_Je_P)\) is
\[
 (\ell_X-j)(\ell_X-j+3)/2+k(k+3)/2+j.
\]
Its difference from \(g_X+\delta(\delta+3)/2+kj\) is even; this verifies
the factor order as well as the shift.

Raising a covector by \(Q\) is denoted by a superscript \(\sharp\).
The exterior derivation rule, applied twice, gives
\[
 PD_Q(\iota_{A_0}\eta)=A_0^\sharp\wedge PD_Q(\eta).
\]
The signs in \eqref{p3sw:spin-eq:duality} in degrees
\(\delta-2\) and \(\delta\) have ratio \(-1\).  Moreover \(Q\) raises
\(e^i\) to \(f_i\) and \(f^i\) to \(e_i\), so
\[
 A_0^\sharp=\tfrac12\sum_i f_i\wedge e_i=-p/2,\qquad
 D\iota_{A_0}D^{-1}=\epsilon_{p/2}.
\]
Thus \(De^{\iota_{A_0}}D^{-1}=e^{p/2}\).
The map \(PD_Q\), and hence \(D\), commutes with \(R\): \(R\)
preserves \(Q\), degree, and orientation, since \(\det\rho(g)=1\).
It follows that \(F_s=D\psi_{B_s}\varphi\) intertwines the diagonal
spin action with \(R\).  Combining the displayed identities gives
\[
 \phi'=D\psi_{B_0}\varphi=e^{p/2}F_s,
\]
which proves \eqref{p3sw:spin-eq:identity}.  Every map used is a
rational algebraic map, so the equality is one of algebraic
representations.  Applying it to \(1\) proves the two stated
components.
\end{proof}

In particular the degree-two term
\(a(g)=\frac12(p-\Lambda^2\rho(g)p)\) satisfies
\[
 a(gh)=a(g)+\Lambda^2\rho(g)a(h).
\]
This follows either by direct substitution or by composing
\eqref{p3sw:spin-eq:identity}.  It fixes the normalization of the
character calculation below.

\subsection{The quadratic action and its actual polarization}

To identify the polarization preserved by the spinor stabilizers,
recall the contraction map of the principal class:
\[
 \vartheta:U\longrightarrow L,\qquad
 \vartheta(y)=\iota_y\Theta,\qquad
 \iota_y(e\wedge e')=y(e)e'-y(e')e.
\]
Thus \(y'(\vartheta y)=\Theta(y,y')\) and
\(\vartheta^t=-\vartheta\).  The principal polarization gives an
isomorphism
\[
 \lambda:X\longrightarrow\widehat X,\qquad
 \lambda(z)=t_z^*H\otimes H^{-1}.
\]
The isomorphism in the sheaf construction is
\(a(z)=t_{-z}^*H\otimes H^{-1}=\lambda(-z)\).
For the normalized Poincar\'e class \(p=\sum e_i\wedge f_i\),
the signs proved in \eqref{p3sw:support-polarization-signs} are
\[
 \lambda^*:U\longrightarrow L\ \text{ is }\vartheta,\qquad
 a^*=-\vartheta.
\]
Under the dual identifications on homology,
\(\lambda_*=-\vartheta\) and \(a_*=\vartheta\).
No symmetry of \(H\) is required.

The two exponential spinors \(\lambda_+,\lambda_-\) introduced in
\eqref{p3sw:spin-eq:source-character} are exchanged by scalar
conjugation.  Their span is the scalar extension of the rational
plane generated by
\((\lambda_++\lambda_-)/2\) and
\((\lambda_+-\lambda_-)/(2u)\).  Let \(J\) be the pointwise algebraic
spin stabilizer of this rational plane, and let
\(\widetilde L_K\) be the ordered stabilizer of the two individual
spinor lines over \(K\).

\begin{proposition}[The polarized quadratic seed]
\label{p3sw:polarized-seed}
On the actual sixfold \(M=X\times\widehat X\), define
\[
 \eta(u)(x,N)=\bigl(a^{-1}(N),[-d]a(x)\bigr)
             =\bigl(-\lambda^{-1}(N),[d]\lambda(x)\bigr).
\]
This gives a unital \(K\)-action.  The ample line
\[
 \mathscr L_M=\operatorname{pr}_X^*H^{\otimes d}
          \otimes\operatorname{pr}_{\widehat X}^*(a^{-1})^*H
\]
has a class \(\theta_M=c_{1,\mathrm B}(\mathscr L_M)\) fixed by the
vector images of \(J\) and \(\widetilde L_K\).  The action is of Weil
type \((3,3)\), is compatible with this polarization, and its actual
homological Hermitian form is hyperbolic.

More precisely, the pullback \(f=\eta(u)^*\) on \(V=L\oplus U\) and
its two eigenspaces are
\begin{equation}
\label{p3sw:spin-eq:oriented-action}
 \begin{split}
 f(w,y)&=(d\vartheta y,-\vartheta^{-1}w),\qquad f^2=-d,\\
 W_+&=\operatorname{Ann}(\lambda_+)
       =\{(-u\vartheta y,y):y\in U_K\},\\
 W_-&=\operatorname{Ann}(\lambda_-)
       =\{(u\vartheta y,y):y\in U_K\}.
 \end{split}
\end{equation}
The eigenvalue on \(W_+\) is \(+u\), and that on \(W_-\) is \(-u\).
\end{proposition}

\begin{proof}
The Fock formula
\(\iota_y e^{u\Theta}=u\vartheta(y)\wedge e^{u\Theta}\)
gives the two annihilators in
\eqref{p3sw:spin-eq:oriented-action}.  They are complementary because
\(\vartheta\) is invertible.  Substitution in the displayed \(f\)
gives the asserted eigenvalues and \(f^2=-d\).  It also gives
\[
 Q(fv,w)=-Q(v,fw).
\]

Let \(s_a=1\times a:X^2\to M\) and
\(s_\lambda=1\times\lambda\).  Equation~\eqref{p3sw:support-polarization-signs}
gives the complete pullback dictionary
\begin{equation}
\label{p3sw:spin-eq:action-matrices}
 \begin{array}{c|c|c|c}
 &s^*(w,y)&s^*W_+&s^*fs^{*-1}\ \text{on }L\oplus L\\ \hline
 s_a&(w,-\vartheta y)&\{(uz,z)\}&
       f_a=\begin{pmatrix}0&-d\\1&0\end{pmatrix}\\[2mm]
 s_\lambda&(w,\vartheta y)&\{(-uz,z)\}&
       f_\lambda=\begin{pmatrix}0&d\\-1&0\end{pmatrix}.
 \end{array}
\end{equation}
The homology transpose of \(f_a\) is
\[
 T_a=\begin{pmatrix}0&1\\-d&0\end{pmatrix},
\]
which is induced by the algebraic endomorphism
\((x,z)\mapsto(z,[-d]x)\) of \(X^2\).  Conjugating by \(s_a\)
gives the stated \(\eta(u)\), so its pullback is \(f\).
Since \(\eta(u)^2=[-d]\), the resulting homomorphism
\(\Q[u]/(u^2+d)\to\operatorname{End}^0(M)\) is unital and injective.
The map \(s_a\) is holomorphic at the original seed.  The graph
\(\{(uz,z)\}\) has three-dimensional intersections with each of
\(H^{1,0}(X^2)\) and \(H^{0,1}(X^2)\); the same is true of the
opposite graph.  This proves the Weil multiplicities for the actual
action.

To identify its invariant cohomology class, first form the covariant
alternating form
\[
 \Xi(v,w)=Q(fv,w).
\]
Let \(\xi\in\Lambda^2 V\) be obtained by raising both indices of
\(\Xi\) with \(Q\).  If matrices are written on cohomology, its
coefficient matrix is
\begin{equation}
\label{p3sw:spin-eq:raised}
 C=Q^{-1}(f^tQ)Q^{-1}=Q^{-1}f^t=-fQ^{-1}.
\end{equation}
This raising is necessary: the covariant form \(\Xi\) itself is not
the degree-two cohomology class.

Here is its exact value on the seed.  Put
\(B_{ij}=\Theta(f_i,f_j)\).  The skew invertible matrix \(B\)
represents the bivector \(\Theta\); the contraction convention means
\(\vartheta=-B\) in column coordinates, hence \(a^*=B\).
In \(s_a\)-coordinates,
\[
 Q_a=\begin{pmatrix}0&B^{-1}\\-B^{-1}&0\end{pmatrix},\qquad
 f_a^tQ_a=\begin{pmatrix}-B^{-1}&0\\0&-dB^{-1}\end{pmatrix},\qquad
 Q_a^{-1}=\begin{pmatrix}0&-B\\B&0\end{pmatrix}.
\]
Multiplication gives
\begin{equation}
\label{p3sw:spin-eq:ample-class}
 Q_a^{-1}(f_a^tQ_a)Q_a^{-1}
       =\begin{pmatrix}dB&0\\0&B\end{pmatrix},\qquad
 s_a^*\xi=d\Theta_1+\Theta_2.
\end{equation}
This is exactly the class of \(H^{\otimes d}\boxtimes H\).
Consequently \(\xi=c_{1,\mathrm B}(\mathscr L_M)=\theta_M\), and
\(\mathscr L_M\) is ample.

For a spin element \(g\),
\(\operatorname{Ann}(m_gs)=\rho(g)\operatorname{Ann}(s)\), as follows
by conjugating the Clifford action.  Thus the vector image of the
ordered line stabilizer preserves \(Q\) and both \(W_+\) and \(W_-\).
It commutes with \(f\), and hence
\[
 \Xi(\rho(g)v,\rho(g)w)
  =Q(f\rho(g)v,\rho(g)w)=Q(\rho(g)fv,\rho(g)w)=\Xi(v,w).
\]
Raising by the invariant \(Q\) is equivariant, so it fixes
\(\theta_M\).  The same holds for its subgroup \(J\).

It remains to verify compatibility and hyperbolicity on actual
homology.  Let \(E\) be the principal Riemann form of \(H\) on
\(U=H_1(X,\Q)\).  By \eqref{p3sw:spin-eq:ample-class}, the actual
Riemann form on \(U\oplus U\) is
\[
 e((v,w),(v',w'))=dE(v,v')+E(w,w').
\]
It is positive at the original period because \(H^{\otimes d}\boxtimes H\)
is ample.  On the two factor indices put \(D_0=\operatorname{diag}(d,1)\).
Direct multiplication gives
\[
 D_0^{-1}T_a^tD_0=-T_a,\qquad T_a^tD_0T_a=dD_0.
\]
These scalar matrices commute with \(E\), so
\(e(T_ax,y)=-e(x,T_ay)\).  For \(k=b+cu\in K\) it follows that
\[
 e(\eta(k)x,y)=e(x,\eta(\bar k)y),\qquad
 \eta(k)^*\theta_M=(b^2+dc^2)\theta_M
                   =\operatorname{Nm}_{K/\Q}(k)\theta_M.
\]
These are precisely the compatibility identities.

Use the actual second-linear form
\[
 H_{\mathrm{hom}}(x,y)=e(x,T_ay)+u e(x,y).
\]
The skew-adjoint identity, \(T_a^2=-d\), and alternation of \(e\)
give
\[
 H_{\mathrm{hom}}(T_ax,y)=-uH_{\mathrm{hom}}(x,y),\quad
 H_{\mathrm{hom}}(x,T_ay)=uH_{\mathrm{hom}}(x,y),\quad
 H_{\mathrm{hom}}(y,x)=\overline{H_{\mathrm{hom}}(x,y)}.
\]
Choose a rational symplectic basis \(v_1,\ldots,v_6\) for \(E\)
with \(E(v_{2i-1},v_{2i})=b_i\in\Q^\times\).  The elements
\(s_i=(v_i,0)\) form a \(K\)-basis because
\(T_as_i=(0,-dv_i)\).  Directly on this basis,
\begin{equation}
\label{p3sw:spin-eq:homological-gram}
 H_{\mathrm{hom}}(s_i,s_j)=du\,E(v_i,v_j).
\end{equation}
Its Gram matrix is the orthogonal sum of
\[
 \begin{pmatrix}0&du b_i\\-du b_i&0\end{pmatrix}
     =\begin{pmatrix}0&c_i\\\bar c_i&0\end{pmatrix},
 \qquad c_i=du b_i\ne0.
\]
Every block has an isotropic coordinate line and is a hyperbolic
Hermitian plane.  This proves hyperbolicity of the actual homological
form, and its signature is \((3,3)\).

For comparison, the auxiliary cohomological form
\(H_V(v,w)=dQ(v,w)+uQ(fv,w)\) is also second-linear:
\(H_V(fv,w)=-uH_V(v,w)\) and \(H_V(v,fw)=uH_V(v,w)\).
The homological conclusion above uses
\eqref{p3sw:spin-eq:homological-gram} directly, without a
cohomology-to-homology trace-dual convention.
\end{proof}

\subsection{The ordered stabilizer and its exterior invariants}

The next calculation identifies exactly which invariant part can
carry a Weil class.  It also distinguishes fixing the two spinor
vectors from merely fixing their lines.

\begin{proposition}[Ordered spinor stabilizer]
\label{p3sw:ordered-stabilizer}
Let \(W=W_+\).  Under \(V_K=W\oplus W^*\), the ordered line
stabilizer and the pointwise stabilizer are
\[
 \widetilde L_K
   =\{(g,t)\in GL(W)\times\mathbb G_m:t^2=\det g\},\qquad
 m(g,t)=t^{-1}\Lambda g,\qquad
 J_K=\{(g,1):g\in SL(W)\}.
\]
The characters on \(\lambda_+\) and \(\lambda_-\) are \(t,t^{-1}\).
For \(0\leq j\leq6\),
\begin{equation}
\label{p3sw:spin-eq:invariants}
 H^{2j}(M,\Q)^J=
 \begin{cases}
  \Q\theta_M^j,&j\ne3,\\
  \Q\theta_M^3\oplus\mathcal W_K(M),&j=3,
 \end{cases}
 \qquad H^{\mathrm{odd}}(M,\Q)^J=0 .
\end{equation}
The ordered line stabilizer fixes every \(\theta_M^j\) and acts on
the two scalar Weil lines by \(t^2,t^{-2}\).
\end{proposition}

\begin{proof}
The two annihilators in
\eqref{p3sw:spin-eq:oriented-action} are complementary maximal
isotropic spaces, paired perfectly by \(Q\).  Identify the second
with \(W^*\).  The Clifford algebra on \(W\oplus W^*\) acts on
\(\Lambda W\) by wedge and contraction.  As in the matrix-unit
argument above it is the full matrix algebra on this module.
Hence this module is isomorphic to \(S_K\) as a Clifford module.
The common annihilator of all wedge operators from \(W\) is the top
line \(K\omega\), while the common annihilator of contractions from
\(W^*\) is the vacuum line \(K1\).  The two original spinor lines
therefore correspond to \(K\omega,K1\), in that order.

An ordered stabilizer must preserve the two annihilators, so its
vector action is \(\operatorname{diag}(g,g^{-T})\) for \(g\in GL(W)\).
Conversely a lift of such a block action preserves both one-dimensional
annihilator spaces.  The operator \(P_g=\Lambda g\) on \(\Lambda W\)
implements this block action by direct conjugation of wedge and
contraction.  It is in the even Clifford algebra because it preserves
parity.  For the top reversal pairing,
\[
 \langle P_gs,P_gt\rangle=\det(g)\langle s,t\rangle.
\]
Thus its Clifford reversion norm is \(\widetilde P_gP_g=\det g\).
Any other Clifford lift of the same vector action differs by a scalar:
their quotient commutes with all of the full matrix algebra.
A norm-one lift is therefore \(aP_g\) with \(a^2\det g=1\).
Writing \(t=a^{-1}\) proves the asserted algebraic cover and action.
The top and vacuum characters are \(t^{-1}\det g=t\) and \(t^{-1}\).
Fixing both vectors is exactly \(t=1\), hence exactly \(SL(W)\).
In particular the extra covering element \((1,-1)\), which negates
both spinors, is not in \(J_K\).  This is a statement of algebraic
groups; it does not assert that every \(K\)-rational determinant has
a square root in \(K\).  The map
\(t\mapsto(\operatorname{diag}(t^2,1,\ldots,1),t)\) is an algebraic
section of the character \(t\); in particular \(t^2,t^{-2}\) are
distinct nontrivial characters.

Now
\[
 \Lambda^k(V_K)
   =\bigoplus_{a+b=k}\Lambda^aW\otimes\Lambda^bW^*.
\]
The invariant vectors in a summand are the intertwiners
\(\operatorname{Hom}_{SL(W)}(\Lambda^bW,\Lambda^aW)\).
Here is the elementary calculation of those intertwiners.  In a
basis \(w_1,\ldots,w_6\), the diagonal torus
\(\prod t_i=1\) acts on \(w_I\) with weight \(\prod_{i\in I}t_i\).
Two such weights, for \(I,J\subset\{1,\ldots,6\}\), agree on this
torus exactly when their indicator vectors differ by a constant
multiple of \((1,\ldots,1)\).  Thus either \(I=J\), or
\((I,J)=(\{1,\ldots,6\},\varnothing)\), or the reverse.
For \(1\leq a\leq5\), any invariant subspace of \(\Lambda^aW\)
is a sum of its one-dimensional torus weight spaces.  The elementary
unipotents \(w_i\mapsto w_i+c w_j\) connect every two \(a\)-subsets
by successive replacements, so a nonzero invariant subspace contains
all the weight spaces.  This proves irreducibility.  In an endomorphism
commuting with the torus the coefficients on these weight spaces
are diagonal, and commuting with the same unipotents makes them all
equal.  Consequently the only invariant blocks are
\[
 a=b,\qquad (a,b)=(6,0),(0,6),
\]
and each is one-dimensional.  The balanced identity tensor is the
degree-two class \(\sum_i w_i\wedge w_i^*\); its \(j\)-th power is
nonzero and spans the balanced block \(a=b=j\).  The already
constructed nonzero invariant \(\theta_M\) is a scalar multiple of
that class.  The other two lines are \(\det W,\det W^*\), with
ordered-group characters \(\det g=t^2\) and \((\det g)^{-1}=t^{-2}\).

Invariants commute with extension of scalars: they are the kernel of
the coaction difference, and field extension is flat.  The group
\(J\) and \(\theta_M\) are rational.  The two determinant lines are
exchanged by scalar conjugation and descend to the rational Weil
plane.  They occupy the exterior summands
\((a,b)=(6,0),(0,6)\) in \(\Lambda^aW\otimes\Lambda^bW^*\),
whereas \(\theta_M^3\) occupies the \((a,b)=(3,3)\) summand.
Thus the sum in degree six is direct.  Descent gives
\eqref{p3sw:spin-eq:invariants}, including the vanishing in odd
degrees.
\end{proof}

\subsection{The normalized character has a nonzero Weil component}

For an even class \(\alpha=\sum_{j=0}^6\alpha_j\), with
\(\alpha_j\in H^{2j}(M,\Q)\) and \(\alpha_0\ne0\), define
\[
 N(\alpha)=e^{-\alpha_1/\alpha_0}\alpha .
\]
If \(b\) has degree two, then \(N(e^b\alpha)=N(\alpha)\), because
the degree-two term changes by \(\alpha_0b\).  If \(R\) is a graded
algebra automorphism, then \(N(R\alpha)=RN(\alpha)\).

\begin{proposition}[The exceptional component of the actual sheaf]
\label{p3sw:exceptional-character}
The normalized character \(\kappa_M\) in
\eqref{p3sw:spin-eq:normalized-character} is a rational
\(J\)-invariant even class, and its degree-six
component has a nonzero rational Weil component:
\[
 (\kappa_M)_3=a_M\theta_M^3+w_M,\qquad
 a_M\in\Q,\quad 0\ne w_M\in\mathcal W_K(M).
\]
\end{proposition}

\begin{proof}
We return to the source tensor \(T=v\otimes\tau v\) in
\eqref{p3sw:spin-eq:sheaf-transform}.  Expanding
\eqref{p3sw:spin-eq:source-character} gives
\begin{equation}
\label{p3sw:spin-eq:secant-tensor}
 \begin{split}
 T={}&
 \frac{d+1}{4d}(\lambda_+\otimes\lambda_+
                         +\lambda_-\otimes\lambda_-)\\
 &+\frac{d-1}{4d}(\lambda_+\otimes\lambda_-
                         +\lambda_-\otimes\lambda_+)\\
 &+\frac{u}{2d}(\lambda_-\otimes\lambda_+
                         -\lambda_+\otimes\lambda_-).
 \end{split}
\end{equation}
It is rational and fixed by \(J\).  It is not fixed by
\(\widetilde L_K\): the two square components have the distinct
nontrivial characters \(t^2,t^{-2}\), with nonzero coefficient
\((d+1)/(4d)\), while the mixed components have character one.

Since \(T\) is \(J\)-fixed, the identity \(\beta=\phi'(T)\) makes
\(\beta\) \(J\)-fixed for \(\rho'\).
Proposition~\ref{p3sw:spin-identity} and the two normalization
identities above show that \(N(\beta)\) is \(J\)-fixed for
\(\Lambda\rho\).

Suppose its degree-six component had zero Weil projection.
Proposition~\ref{p3sw:ordered-stabilizer} would then put every
component of \(N(\beta)\) in the power ring \(K[\theta_M]\), which
is fixed by \(\widetilde L_K\).  The class \(e^{-p/2}\beta\) is
\(J\)-fixed for \(\Lambda\rho\), by
\eqref{p3sw:spin-eq:identity}.  Its degree-two term
\[
 b=\beta_1-rp/2
\]
is therefore in \(H^2(M,\Q)^J=\Q\theta_M\), and is also fixed by
\(\widetilde L_K\).  But
\[
 e^{-p/2}\beta=e^{b/r}N(\beta).
\]
Under the supposition the right side is fixed by \(\widetilde L_K\).
The exact Clifford identity would make \(\beta\) fixed by \(\rho'\)
on that group, and the invertibility of \(\phi'\) would make \(T\)
fixed by the diagonal spin action.  This contradicts
\eqref{p3sw:spin-eq:secant-tensor}.  Hence
\[
 N(\beta)_3=a\theta_M^3+w,\qquad
 a\in\Q,\quad 0\ne w\in\mathcal W_K(M).
\]
The coefficients are rational because the direct sum
\eqref{p3sw:spin-eq:invariants} is rational.

The duality identity \eqref{p3sw:spin-eq:sheaf-duality} gives
\(\kappa_M=\iota N(\beta)\).
The involution commutes with the graded vector action and preserves
the rational decomposition in
\eqref{p3sw:spin-eq:invariants}.  In degree six it multiplies both
terms by \(-1\), so the Weil component remains nonzero.  Renaming
the two rational terms proves the proposition.
\end{proof}

\subsection{The full normalized character is Hodge on the Weil domain}

We now relate the vector stabilizer to the actual polarized complex
structures.  This uses the polarization matrix itself, so no
cohomological Hermitian trace-dual convention is needed.

\begin{proposition}[The actual polarized Hodge circle]
\label{p3sw:hodge-circle}
Keep the rational homology of the seed \(M\), the action \(\eta\),
and the Riemann form of \(\mathscr L_M\) fixed.  At every compatible
marked complex structure of Weil type \((3,3)\), every component
\((\kappa_M)_j\) has Hodge type \((j,j)\).
\end{proposition}

\begin{proof}
Let \(C\) be the coefficient matrix of \(\theta_M\) on cohomology.
It is also the matrix of its alternating Riemann form on dual
homology.  Equation~\eqref{p3sw:spin-eq:raised} gives
\[
 C=-fQ^{-1}.
\]
For a cohomology automorphism \(g\) commuting with the invertible
\(f\), this implies
\begin{equation}
\label{p3sw:spin-eq:centralizer}
 \begin{split}
 gCg^t=C
 &\Longleftrightarrow -f(gQ^{-1}g^t)=-fQ^{-1}\\
 &\Longleftrightarrow gQ^{-1}g^t=Q^{-1}
 \Longleftrightarrow g^tQg=Q.
 \end{split}
\end{equation}
Thus among \(K\)-linear cohomology operators, fixing the actual
polarization bivector is exactly orthogonality for \(Q\).

Let \(I\) be a compatible complex structure on homology and let
\(T=\eta(u)_*\).  Put \(h_t=\cos t+\sin t\,I\), and let
\(g_t=h_t^t\) be its cohomology pullback.  Since \(IT=TI\) and
\(T^t=f\), one has \(g_tf=fg_t\).  Compatibility
\(e(Ix,Iy)=e(x,y)\) says \(I^tCI=C\).  Since \(I^2=-1\), it also
gives \(I^tC+CI=0\); expansion of \(h_t^tCh_t\) then gives
\(h_t^tCh_t=C\).  This is \(g_tCg_t^t=C\).
Equation~\eqref{p3sw:spin-eq:centralizer} places \(g_t\) in the
centralizer of \(f\) in \(O(Q)\), and its path from the identity
places it in \(SO(Q)\).

After the usual complex scalar extension of \(W_+\), the two
Hodge-circle weights have multiplicities three and three.  Hence
\[
 \det(g_t|_{W_+})=e^{3it}e^{-3it}=1.
\]
Using the inverse circle convention only interchanges the factors.
By Proposition~\ref{p3sw:ordered-stabilizer}, after extending to
\(K\) the determinant-one vector centralizer is \(SL(W_+)\), and
its unique lift fixing both spinors is \((g,1)\in J_K\).
This identity of algebraic groups descends from \(K\).  The actual
polarized Hodge circle is therefore in the vector image of \(J\).

On \(H^{p,q}\) the circle has weight \(p-q\), up to simultaneous
reversal.  A degree-\(2j\) class fixed by it has only type \((j,j)\).
Proposition~\ref{p3sw:exceptional-character} makes every component
of \(\kappa_M\) \(J\)-invariant, proving the assertion at every
compatible marked period.
\end{proof}

Only a constant rational transport will be needed for the quotient
\(q:M\to Y_0\).  Transporting homology, the \(K\)-action, and the
polarization by \(q_*\) identifies the marked complex structures;
under these identifications \(q^*\) is a constant rational
cohomology map and respects every corresponding Hodge decomposition.
This observation concerns classes on the marked domain, and requires
no relative isogeny or relative descent of the sheaf.
\section{Descent and deformation}
\label{p3:sec:deformation}

We now pass from the equivariant sheaf on the Jacobian product to an
ordinary sheaf whose full obstruction group can be used in deformation.
The first step is coherent twisted descent through the finite quotient and
a projective-bundle conversion. The second is the restricted
smooth-projective-germ semiregularity theorem needed for that ordinary
sheaf.

\subsection{Twisted descent and conversion to an ordinary sheaf}
\label{p3sw:descent-section}

The deformation argument will use an ordinary coherent sheaf on a
projective family.  We first construct a twisted sheaf on the quotient of
the abelian sixfold and compare its full derived obstruction group with
the invariant obstruction group on the covering sixfold.  We then cancel
the twist on a relative projective bundle.  Throughout, the group action
on an endomorphism complex is the conjugation action induced by the given
line-twisted action; an action on the sheaf by ordinary automorphisms is
not required.

\subsubsection{Characters and the duality sign}

For a perfect complex \(P\) on a smooth complex manifold \(V\), write
\[
 a_P=-\operatorname{At}_{\mathrm{BF}}(P)
       \in\operatorname{Ext}^1_V(P,P\otimes\Omega_V^1).
\]
This is the raw first-jet Atiyah class with positive line-bundle Chern
sign.  The opposite sign is the convention of
Buchweitz--Flenner \cite[Definition~3.4 and Definition~4.1]{P3SWBuchweitzFlenner2003}.
Use the Atiyah-power and perfect-trace conventions of
Proposition~\ref{p3sw:hh-triangle}.  Define
\begin{equation}
\label{p3sw:descent-character}
 \begin{aligned}
 \operatorname{ch}_{H,k}(P)&=\frac{\operatorname{tr}(a_P^k)}{k!},
 \qquad
 \operatorname{ch}_H(P)=\sum_k\operatorname{ch}_{H,k}(P),\\
 \sigma_P(v)&=
 \left(\operatorname{Tr}_{P,\Omega_V^j[j+2]}
   \left(\left(\frac{a_P^j}{j!}\right)[2]\circ v\right)\right)_{j\geq0}.
 \end{aligned}
\end{equation}
Here \(\operatorname{tr}\) is the perfect trace, and
\[
 \sigma_P:\operatorname{Ext}^2_V(P,P)\longrightarrow
 \mathcal H_{-2}(V):=\bigoplus_{j\geq0}H^{j+2}(V,\Omega_V^j).
\]
The summands with \(j>\dim V-2\) vanish.  These definitions also apply to
a perfect weight-one complex on a finite-band gerbe: its derived
endomorphisms and their trace have weight zero and hence descend to the
underlying manifold.

On a smooth projective \(V\), let
\(\iota_V:H^q(V,\Omega_V^p)\hookrightarrow H^{p+q}(V,\C)\) be the
Hodge inclusion for the ordinary de Rham differential \(d\).  Our
normalized Betti convention is
\begin{equation}
\label{p3sw:descent-tate}
 \operatorname{ch}_{B,k}(P)
   =(2\pi i)^{-k}\iota_V\bigl(\operatorname{ch}_{H,k}(P)\bigr),
 \qquad
 \operatorname{ch}_B(P)=\sum_k\operatorname{ch}_{B,k}(P).
\end{equation}
For an ordinary perfect complex this is its rational topological Chern
character.  For a twisted complex the expression initially defines a
complex Betti class; rationality in the present construction will follow
from finite descent.  The inclusion \(\iota_V\) is multiplicative and
commutes with pullback.  On the semiregularity summand
\(H^{j+2}(\Omega_V^j)\), the corresponding normalization is
\((2\pi i)^{-j}\iota_V\): its weight is the form degree \(j\), not half
the total cohomological degree.  Thus, if
\(\iota_V(\gamma_H)=(2\pi i)\gamma_B\) for
\(\gamma_H\in H^1(\Omega_V^1)\), multiplication by \(\gamma_H\) on the
raw targets corresponds weightwise to multiplication by \(\gamma_B\).
All identities involving \(\sigma\) will remain in raw Hodge cohomology.

Return to the fixed perfect complex \(\mathcal G\) on
\(M=X\times\widehat X\).  Write
\[
 D(\mathcal G)=R\mathcal{H}om_M(\mathcal G,\mathcal O_M),
 \qquad E=D(\mathcal G)[-1]=\mathcal E.
\]
By Proposition~\ref{p3sw:coherent-sheaf}, \(E\) is a coherent
perfect sheaf of rank \(r=8d>0\).
Proposition~\ref{p3sw:coherent-action} supplies the coherent
line-twisted actions on \(\mathcal G\) and \(E\), both indexed by
\(G_{\mathrm{src}}\) and identified with the translation subgroup
\(\bar G\) through \(g\mapsto b_g\).  Duality inverts the line
factors and preserves the translations and coherent composition law.
Proposition~\ref{p3sw:invariant-injectivity} proves that
\(\sigma_{\mathcal G}\) is injective on the full group
\(\operatorname{Ext}^2_M(\mathcal G,\mathcal G)^{\bar G}\).
The next calculation transfers this injection to \(E\).

\begin{lemma}
\label{p3sw:descent-duality}
Even transpose induces an isomorphism
\[
 d_2:\operatorname{Ext}^2_M(\mathcal G,\mathcal G)^{\bar G}
       \xrightarrow{\ \sim\ }
       \operatorname{Ext}^2_M(E,E)^{\bar G}.
\]
With the convention \eqref{p3sw:descent-character}, its components obey
\begin{equation}
\label{p3sw:descent-dual-sign}
 \sigma_{E,j}(d_2v)=(-1)^{j+1}\sigma_{\mathcal G,j}(v)
 \quad\text{for every }j\geq0.
\end{equation}
Consequently, injectivity of \(\sigma_{\mathcal G}\) on the full
\(\bar G\)-invariant \(\operatorname{Ext}^2\) implies injectivity of
\(\sigma_E\) on the full \(\bar G\)-invariant \(\operatorname{Ext}^2\).
\end{lemma}

\begin{proof}
Perfect duality is an anti-equivalence and turns the given line-twisted
action into its inverse-line action.  The quasi-isomorphism from
\(D(\mathcal G)[-1]\) to the degree-zero sheaf \(E\) is canonical in the
derived category and is equivariant.  Transpose therefore gives the
asserted invariant isomorphism in even degree.

The sign can be computed in a finite complex of free modules.  The dual
Atiyah class is the negative transpose of the original class.
Transposing a product reverses both its Ext-degree-one factors and its
one-form factors.  The two reversal Koszul signs cancel: equivalently,
an \(\Omega^1[1]\)-valued Atiyah operation has even combined parity.
The Ext-degree-two operation \(v\) is also even, so its movement past
these factors gives no further sign.  Supertrace is invariant under
transpose, while the shift \([-1]\) changes supertrace by \(-1\).
The \(j\) negative dual Atiyah factors and this last sign give
\eqref{p3sw:descent-dual-sign}.  Evaluation, transpose, and the
Atiyah construction respect changes of finite free model, so the local
calculation glues.  The diagonal operator with entries
\((-1)^{j+1}\) is invertible, which proves the final assertion.
\end{proof}

\subsubsection{Determinant normalization and derived descent}

Let
\[
 q:M\longrightarrow Y_0:=M/\bar G
\]
be the finite etale translation isogeny of
Proposition~\ref{p3sw:coherent-action}.  We write the action on \(E\)
using pullback translations.  For \(h=-b_g\), set \(L_h=N_g^{-1}\).
Since \(t_{b_g,*}=t_{-b_g}^*=t_h^*\), the action
\eqref{p3sw:support-dual-action} becomes
\[
 L_h\otimes t_h^*E\simeq E\qquad(h\in\bar G).
\]
These isomorphisms retain the unital associative line compositors of
that proposition.  Reindexing by inversion on \(\bar G\) preserves the
invariant subspace of the derived endomorphisms.

\begin{proposition}
\label{p3sw:twisted-descent}
There are a class \(\alpha\in H^2(Y_0,\mu_r)\), a coherent
\(\alpha\)-twisted sheaf \(\mathcal B\) of rank \(r\) with
\(\det\mathcal B\simeq\mathcal O_{Y_0}\), and an invertible
\(q^*\alpha\)-twisted sheaf \(\mathcal L\) on \(M\), such that
\begin{equation}
\label{p3sw:descent-twisted-pullback}
 q^*\mathcal B\simeq E\otimes\mathcal L.
\end{equation}
Let \(\mathfrak Y_0\) be the \(\mu_r\)-gerbe of \(\alpha\).  For every
integer \(i\), this isomorphism induces
\begin{equation}
\label{p3sw:descent-ext}
 \operatorname{Ext}^i_{\mathfrak Y_0}(\mathcal B,\mathcal B)
   \simeq\operatorname{Ext}^i_M(E,E)^{\bar G},\qquad
 \operatorname{Ext}^i_0(\mathcal B,\mathcal B)
   \simeq\operatorname{Ext}^i_0(E,E)^{\bar G},
\end{equation}
where the subscript \(0\) denotes the kernel of the perfect trace to
\(H^i(\mathcal O)\).  Put
\[
 \lambda_H=\frac{a_{\det E}}r=\frac{\operatorname{tr}(a_E)}r.
\]
Under the first isomorphism \(\phi\) in degree two, the raw Hodge square
\begin{equation}
\label{p3sw:descent-square}
\begin{array}{ccc}
 \operatorname{Ext}^2(\mathcal B,\mathcal B)&
       \xrightarrow{\ \phi\ }&\operatorname{Ext}^2(E,E)^{\bar G}\\
 \big\downarrow{\scriptstyle\sigma_{\mathcal B}}&&
       \big\downarrow{\scriptstyle e^{-\lambda_H}\sigma_E}\\
 \mathcal H_{-2}(Y_0)&\xrightarrow{\ q^*\ }&\mathcal H_{-2}(M)
\end{array}
\end{equation}
commutes, both on full Ext and on its trace-free summand.  In particular,
full invariant injectivity for \(E\) makes \(\mathcal B\) semiregular.
By Lemma~\ref{p3sw:descent-duality}, full invariant injectivity for
\(\mathcal G\) is sufficient for this conclusion.
\end{proposition}

\begin{proof}
We construct the twisted sheaf before comparing its derived
endomorphisms.  Choose an evenly covered good analytic cover of \(Y_0\),
refined so that the finitely many transition lines are trivial on the
required sheet intersections.  Take every sheet over every cover member,
so the sheets \(U_i\) cover \(M\), and retain its image \(V_i\) as a
labelled member of the cover.  Transport
\(E|_{U_i}\) to a sheaf \(B_i\) on \(V_i\).  On \(V_i\cap V_j\) for
\(i<j\), the deck transformation comparing the two sheets and a local
trivialization of its line \(L_h\) give an isomorphism
\(\rho_{ij}:B_j\to B_i\).  Use the unit and inverse composition data
to impose \(\rho_{ii}=\operatorname{id}\) and
\(\rho_{ji}=\rho_{ij}^{-1}\).  On a triple intersection, the coherent
composition law cancels the translations; the remaining line
trivializations differ by an invertible scalar function.  Thus
\(\rho_{ij}\rho_{jk}\rho_{ki}=c_{ijk}\operatorname{id}\).

We use the determinant functor for perfect complexes in its elementary
local form: a finite free complex has the alternating tensor product of
the determinants of its terms.  It commutes with restriction and
isomorphism, and the determinant of \(c\operatorname{id}\) on a
rank-\(r\) perfect complex is \(c^r\).  Each \(B_i\) is perfect, so
\(\det B_i\) is a line.  Trivialize these lines and let \(\eta_{ij}\)
be the determinant of \(\rho_{ij}\); then \(\eta_{ii}=1\) and
\(\eta_{ji}=\eta_{ij}^{-1}\).  After a further good-cover refinement,
choose an \(r\)-th root \(\widetilde\eta_{ij}\) of this nowhere-zero
holomorphic function for \(i<j\), and impose
\(\widetilde\eta_{ii}=1\) and
\(\widetilde\eta_{ji}=\widetilde\eta_{ij}^{-1}\).  Set
\(\rho'_{ij}=\widetilde\eta_{ij}^{-1}\rho_{ij}\).  These maps have
determinant one and satisfy \(\rho'_{ii}=\operatorname{id}\) and
\(\rho'_{ji}=(\rho'_{ij})^{-1}\).  The coherent composition law gives
an invertible holomorphic scalar \(\epsilon_{ijk}\) defined by
\[
 \rho'_{ij}\rho'_{jk}=\epsilon_{ijk}\rho'_{ik}.
\]
Taking determinants gives \(\epsilon_{ijk}^r=1\), so these scalars are
locally constant with values in \(\mu_r\).  Associativity on a fourfold
intersection gives the exact Cech cocycle identity
\[
 \epsilon_{ijk}\epsilon_{ikl}=\epsilon_{jkl}\epsilon_{ijl}.
\]
The repeated-index multipliers are one; scalar multiplication is
faithful because \(B_i\) has positive rank on a smooth manifold.
The resulting \(\mu_r\)-cocycle defines \(\alpha\), and the sheaves
\(B_i\) with transitions \(\rho'_{ij}\) define \(\mathcal B\).
Their determinant trivializations glue, so
\(\det\mathcal B\simeq\mathcal O_{Y_0}\).  Rank, coherence, and
perfection are local, and hence pass from \(E\) to \(\mathcal B\).
This is the determinant-normalized form of the local construction in
\cite[Lemma~9.3.6 and equation~(9.3.3)]{Markman2025Secant};
it uses the perfect determinant, not an exterior power of a locally
free sheaf.

Let \(\theta_i:(q^*B_i)|_{U_i}\xrightarrow{\sim}E|_{U_i}\) be the
local identification from transport.  On an actual overlap
\(U_i\cap U_j\), the comparing deck transformation is the identity,
so the coherent line action gives an invertible scalar \(l_{ij}\) with
\[
 \theta_i(q^*\rho'_{ij})\theta_j^{-1}=l_{ij}\operatorname{id}_E.
\]
The normalized transitions give \(l_{ii}=1\), \(l_{ji}=l_{ij}^{-1}\),
and \(l_{ij}l_{jk}=(q^*\epsilon_{ijk})l_{ik}\).  Use \(l_{ij}\) as
the transition from the \(j\)-th trivial line to the \(i\)-th.  These
lines define the weight-one \(q^*\alpha\)-twisted line \(\mathcal L\)
with the same multiplier as \(q^*\mathcal B\), and the local
identifications give \eqref{p3sw:descent-twisted-pullback}.  On translated
intersections, conjugation by these identifications is the original conjugation action
on the derived endomorphisms of \(E\).

For clarity, a twisted sheaf here is a coherent weight-one sheaf on
\(\mathfrak Y_0\), equivalently the just described Cech data.
Weight-zero sheaves and complexes descend uniquely to \(Y_0\).
A coherent twisted sheaf on the smooth manifold is absolutely perfect:
in a local trivialization of the gerbe it is a finite module over a
regular analytic local ring and has a finite free resolution.  Write
\[
 \mathcal C_{\mathcal B}
   =R\mathcal{H}om_{\mathfrak Y_0}(\mathcal B,\mathcal B)
\]
for its weight-zero descent.  It is a perfect complex on \(Y_0\),
locally \(\mathcal B^\vee\otimes^L\mathcal B\).

We use the following derived base-change identity.  For an absolutely
perfect \(P\), any complex \(Q'\), and any morphism \(f:T\to V\),
\begin{equation}
\label{p3sw:descent-rhom-basechange}
 Lf^*R\mathcal{H}om_V(P,Q')
   \simeq R\mathcal{H}om_T(Lf^*P,Lf^*Q').
\end{equation}
Indeed, locally the left side is
\(Lf^*(P^\vee\otimes^L Q')\); derived pullback is symmetric
monoidal and carries the dual of a perfect complex to the dual of its
pullback.  This also proves the identity on gerbe trivializations.
The unit and trace, being coevaluation and evaluation, commute with
the same derived pullbacks.  This is an assertion about derived
pullback; ordinary pullback under a nonflat map is not being substituted.

Apply \eqref{p3sw:descent-rhom-basechange} to the etale map \(q\), where
ordinary and derived pullback agree.  Line cancellation in
\eqref{p3sw:descent-twisted-pullback} gives
\begin{equation}
\label{p3sw:descent-rhom}
 q^*\mathcal C_{\mathcal B}
 \simeq R\mathcal{H}om(q^*\mathcal B,q^*\mathcal B)
 \simeq R\mathcal{H}om_M(E,E).
\end{equation}
Both the line factors and their scalar gluing cancel.  The descent datum
on the right is consequently the honest conjugation action of \(\bar G\),
although the action on \(E\) itself is line-twisted.

For a finite covering, \(q_*\) on sheaves is exact: over an evenly
covered open set it is the direct sum of the finitely many sheet
restrictions.  Etale descent gives
\(\mathcal C_{\mathcal B}\simeq(q_*q^*\mathcal C_{\mathcal B})^{\bar G}\).
Over \(\C\), invariants are the split summand defined by
\(|\bar G|^{-1}\sum_{h\in\bar G}h\).  Applying derived global sections
to that summand and using exactness of \(q_*\) gives
\[
 R\Gamma(Y_0,\mathcal C_{\mathcal B})
   \simeq R\Gamma(M,q^*\mathcal C_{\mathcal B})^{\bar G}.
\]
Taking cohomology and using \eqref{p3sw:descent-rhom} proves the full
Ext comparison in \eqref{p3sw:descent-ext}.

The perfect unit and trace are equivariant maps
\begin{equation}
\label{p3sw:descent-trace}
 \mathcal O_{Y_0}\xrightarrow{u}\mathcal C_{\mathcal B}
       \xrightarrow{\operatorname{tr}}\mathcal O_{Y_0},
 \qquad \operatorname{tr}u=r\operatorname{id}.
\end{equation}
The last equality is the alternating-rank calculation in a finite free
model.  Thus \(u/r\) splits the trace.  If
\(\mathcal C_{\mathcal B,0}=\operatorname{fib}(\operatorname{tr})\),
then
\[
 \mathcal C_{\mathcal B}\simeq
 \mathcal O_{Y_0}\oplus\mathcal C_{\mathcal B,0},\qquad
 H^iR\Gamma(Y_0,\mathcal C_{\mathcal B,0})
 =\ker\bigl(\operatorname{Ext}^i(\mathcal B,\mathcal B)
                   \to H^i(\mathcal O_{Y_0})\bigr).
\]
The splitting commutes with \eqref{p3sw:descent-rhom}, since derived
pullback and tensoring by an invertible line preserve evaluation and
coevaluation.  Descent of this summand proves the trace-free comparison.
In particular, the calculation retains the higher sheaf-Ext terms:
it does not replace \(H^2R\Gamma(\mathcal C_{\mathcal B,0})\) by
\(H^2(\mathcal End_0(\mathcal B))\).

It remains to compare the semiregularity maps.  The cotangent target
on a \(\mu_r\)-gerbe is that of its underlying manifold.  Equivalently,
the scalar triple-gluing functions are locally constant, so their
differentials vanish in first jets.  The pullback and tensor rules for
Atiyah classes, with the differential
\(q^*\Omega_{Y_0}^1\simeq\Omega_M^1\), give
\begin{equation}
\label{p3sw:descent-atiyah}
 q^*a_{\mathcal B}=a_E+a_{\mathcal L}\operatorname{id}_E,
 \qquad
 a_{\mathcal L}=-\frac{a_{\det E}}r=-\lambda_H.
\end{equation}
The second equality follows by tracing the first and using
\(\det\mathcal B=\mathcal O_{Y_0}\); the trace of \(a_E\) is the
Atiyah class of its perfect determinant.  The pullback compatibility is also
\cite[Proposition~3.14]{P3SWBuchweitzFlenner2003}; the tensor rule
follows from the Leibniz rule for first jets on a local perfect model.
The scalar \(-\lambda_H\) commutes with the Atiyah operations in the
combined Ext/form grading.  Expanding the exponential and applying the
natural perfect trace yields
\[
 q^*\sigma_{\mathcal B}(v)
     =e^{-\lambda_H}\sigma_E(\phi v),
\]
which is \eqref{p3sw:descent-square}.  Multiplication by
\(\lambda_H\in H^1(\Omega_M^1)\) raises both Hodge indices and is
nilpotent on \(\mathcal H_{-2}(M)\); \(e^{-\lambda_H}\) therefore
has inverse \(e^{\lambda_H}\).  Full invariant injectivity on the
right implies full injectivity on the left.  The trace-free comparison
also shows that injectivity restricted to the invariant trace-free
group would suffice: the \(j=0\) component of \(\sigma_{\mathcal B}\)
is the perfect trace itself, so its vanishing first puts a class into
the trace-free summand.
\end{proof}

The normalized Betti counterpart of \(\lambda_H\) is
\[
 \lambda_B=\frac{c_1^{B}(\det E)}r,\qquad
 \iota_M(\lambda_H)=(2\pi i)\lambda_B,
\]
where \(c_1^B\) is the usual integral Betti Chern class.  Applying
\eqref{p3sw:descent-atiyah} to the character and then using
\eqref{p3sw:descent-tate} weight by weight gives
\begin{equation}
\label{p3sw:descent-character-pullback}
 \begin{split}
 q^*\operatorname{ch}_H(\mathcal B)
   &=e^{-\lambda_H}\operatorname{ch}_H(E),\\
 q^*\operatorname{ch}_B(\mathcal B)
   &=e^{-\lambda_B}\operatorname{ch}_B(E).
 \end{split}
\end{equation}
The right side of the second equality is rational.  The Betti trace
\(q_*^B\) satisfies \(q_*^Bq^*=(\deg q)\operatorname{id}\), so
\((\deg q)^{-1}q_*^B\) proves that
\(\operatorname{ch}_B(\mathcal B)\) is rational as well.  No
integrality of \(\lambda_B\) is required.

\subsubsection{A relative projective bundle carrying an ordinary sheaf}

Let \(a:Y\to S\) be a holomorphic abelian scheme of relative dimension
six over a sufficiently small contractible representative of a smooth
analytic germ, with central fiber \(Y_0\).  Assume that the family is
polarized.  Locally it admits a line \(\mathcal L_Y\) representing the
polarization and ample relative to \(a\).  More precisely, if only the
polarization homomorphism was specified, the standard rigidified
relative Picard component of representing lines is a torsor under the
smooth dual abelian scheme; a local analytic section through a chosen
central representative supplies \(\mathcal L_Y\).  Relative ampleness
holds after shrinking.  This is the only local choice of a representing
polarization line that will be used.

We realize the twisting class by a finite projective representation
and extend its projective-space bundle over \(Y\).  On the central
fiber, the inverse-twisted tautological line cancels the twist of the
pulled-back sheaf \(\mathcal B\).  The vanishing of higher cohomology
of the structure sheaf on projective space retains the full obstruction
group in this conversion.

\begin{proposition}
\label{p3sw:ordinary-conversion}
For the class \(\alpha\) and the semiregular twisted coherent sheaf
\(\mathcal B\) in Proposition~\ref{p3sw:twisted-descent}, put
\(m=r^{12}\) and \(N=6+m-1\).  There exist a holomorphic
\(\mathbb P^{m-1}\)-bundle \(p:Z\to Y\), a determinant-trivial
\(\alpha\)-twisted vector bundle \(Q_0\) of rank \(m\) on \(Y_0\), and
an inverse-twisted tautological line
\(\ell_0=\mathcal O_{\mathbb P(Q_0)}(1)\), with \(Z_0=\mathbb P(Q_0)\).
The ordinary line
\begin{equation}
\label{p3sw:descent-anticanonical}
 H=\omega_{Z/Y}^{-1}=\det T_{Z/Y}
\end{equation}
satisfies \(H|_{Z_0}\simeq\ell_0^m\) and restricts to
\(\mathcal O_{\mathbb P^{m-1}}(m)\) on every \(p\)-fiber.  The
normalized Betti class and raw Hodge class
\begin{equation}
\label{p3sw:descent-zeta}
 \zeta_{B,s}=\frac{c_1^B(H|_{Z_s})}{m},\qquad
 \zeta_{H,s}=\frac{a_{H|_{Z_s}}}{m},\qquad
 \iota_{Z_s}(\zeta_{H,s})=(2\pi i)\zeta_{B,s}
\end{equation}
have the following properties: \(\zeta_B\) is a horizontal rational
class of type \((1,1)\), and its restriction to every \(p_s\)-fiber
is the normalized hyperplane class.  Every \(Z_s\) and \(p_s\) is
projective algebraic, and \(a\circ p\) has a relative ample line.

The sheaf
\[
 F_0=p_0^*\mathcal B\otimes\ell_0
\]
is ordinary coherent on \(Z_0\), and it is semiregular.  More precisely,
there is a full derived isomorphism
\begin{equation}
\label{p3sw:descent-ordinary-ext}
 \operatorname{Ext}^2_{Z_0}(F_0,F_0)
       \simeq\operatorname{Ext}^2_{\mathfrak Y_0}(\mathcal B,\mathcal B)
\end{equation}
under which
\begin{equation}
\label{p3sw:descent-ordinary-sigma}
 \sigma_{F_0}(p_0^*v)
       =e^{\zeta_{H,0}}p_0^*\sigma_{\mathcal B}(v).
\end{equation}
If a rational horizontal section
\(\kappa_B=\sum_{k=0}^6\kappa_{B,k}\) specializes to
\(\operatorname{ch}_B(\mathcal B)\) and has type \((k,k)\) in degree
\(2k\), then
\begin{equation}
\label{p3sw:descent-horizontal-character}
 A_B=p^*\kappa_B e^{\zeta_B}=\sum_{k=0}^{N}A_{B,k}
\end{equation}
is rational, horizontal, and of type \((k,k)\) in degree \(2k\).
Its raw Hodge section in that degree is the edge of
\((2\pi i)^kA_{B,k}\), and its central value is
\(\operatorname{ch}_{H,k}(F_0)\).  Finally, on every fiber,
\begin{equation}
\label{p3sw:descent-extraction}
 p_{s*}^B\bigl(\zeta_{B,s}^{\,m-1}e^{-\zeta_{B,s}}A_B(s)\bigr)
       =\kappa_B(s),
\end{equation}
where \(p_{s*}^B\) is normalized Betti fiber integration.
\end{proposition}

\begin{proof}
We first realize the exact finite-band class \(\alpha\).  Set
\(\Lambda=H_1(Y_0,\Z)\simeq\Z^{12}\).  The universal coefficient
theorem for the torus gives
\[
 H^2(Y_0,\mu_r)=\operatorname{Hom}(\mathop{\bigwedge}\nolimits^2\Lambda,\mu_r);
\]
there is no Ext term because \(H_1(Y_0,\Z)\) is free.  Let \(b\) be
the alternating pairing corresponding to \(\alpha\), and put
\(G_r=\Lambda/r\Lambda\).  For an ordered basis \(e_1,\ldots,e_{12}\)
of \(\Lambda\), define
\[
 c(x,y)=\prod_{i<j}b(e_i,e_j)^{x_i y_j}\quad(x,y\in G_r).
\]
This is a well-defined bimultiplicative cocycle with
\(c(x,y)c(y,x)^{-1}=b(x,y)\).  On the \(m=|G_r|=r^{12}\)
dimensional vector space with basis \(\{v_y:y\in G_r\}\), set
\[
 T_xv_y=c(x,y)v_{x+y}.
\]
Then \(T_xT_y=c(x,y)T_{x+y}\), so \(x\mapsto[T_x]\) is a
projective representation.  For \(r=1\) this is the trivial
one-dimensional representation.

The multiplication map \([r]:Y\to Y\) is a finite etale torsor
under \(Y[r]\).  Its differential on the relative tangent bundle is
multiplication by \(r\); it is proper, and every abelian fiber has
\(r^{12}\) inverse images of a point.  Thus it is a proper local
biholomorphism of that degree.  Contractibility of \(S\) trivializes
the finite local system \(Y[r]\) as \(G_r\), with its central
identification.  Form \(Z\) as the quotient of
\(Y\times\mathbb P^{m-1}\) by the free diagonal action covering
translation by \(G_r\), taking the inverse action on the projective
coordinate.  If local sections of the torsor satisfy
\(s_j=s_i+x_{ij}\), this convention gives
\([s_j,z]=[s_i,[T_{x_{ij}}]z]\).  Local sections of the finite
covering show that \(p:Z\to Y\) is a holomorphic projective-space
bundle with these projective transitions.

We must check the class in \(H^2(Y_0,\mu_r)\), not only its image in
the analytic Brauer group.  Rescale each \(T_x\) to have determinant
one.  The multiplier of these special-linear lifts is \(\mu_m\)-valued:
it is scalar with determinant one.  It differs from \(c\) by a scalar
coboundary and still has commutator \(b\).  The monodromy of the
central \([r]\)-torsor is \(\Lambda\to\Lambda/r\Lambda\).
For any constant coefficient group \(C'\),
\[
 H^2(Y_0,C')=\operatorname{Hom}(\mathop{\bigwedge}\nolimits^2\Lambda,C').
\]
In the bar complex the cycle
\([e_i|e_j]-[e_j|e_i]\) represents \(e_i\wedge e_j\), and an
inflated multiplier evaluates on it as its alternating commutator.
Thus the special-linear obstruction class has pairing \(b\).  The
inclusion \(\mu_m\hookrightarrow\Q/\Z\) induces an injection on
this \(H^2\), so that class is the image of \(\alpha\) under
\(\mu_r\hookrightarrow\mu_m\).  On a common good cover, its
\(\mu_m\)-cocycle therefore differs from a chosen \(\mu_r\)-cocycle
for \(\alpha\) by the coboundary of a \(\mu_m\)-valued one-cochain.
Multiplying the special-linear lifts by that cochain gives transitions
with exactly the chosen \(\mu_r\)-multiplier and still determinant
one.  They define \(Q_0\), and their projectivizations give \(Z_0\).

Use the convention that \(\mathbb P(Q_0)\) parametrizes lines in
\(Q_0\).  The tautological \(\ell_0=\mathcal O(1)\) then has the
inverse twist.  The twisted factors cancel in the Euler sequence
\[
 0\longrightarrow\mathcal O_{Z_0}
 \longrightarrow p_0^*Q_0\otimes\ell_0
 \longrightarrow T_{Z_0/Y_0}\longrightarrow0.
\]
Since \(\det Q_0=\mathcal O_{Y_0}\), its determinant is
\(H|_{Z_0}=\ell_0^m\).  On any projective-space fiber the same
anticanonical line is \(\mathcal O(m)\).

On each \(s\), the quotient construction is the algebraic quotient
of the projective variety \(Y_s\times\mathbb P^{m-1}\) by a free
finite algebraic action.  It is projective: tensor the translates of
an ample line on the product, use the permutation linearization on
that tensor product, descend it along the free finite quotient, and
use descent of ampleness under a finite surjection.  The map \(p_s\)
is algebraic and projective, for instance by its closed graph in a
projective product.  The global line \(H\) is \(p\)-ample.  The
composition rule for relative ampleness gives positive \(u,v\) for
which \(H^u\otimes p^*\mathcal L_Y^v\) is ample for \(a\circ p\).
The class \(\zeta_{B,s}\) is the restriction of the global class
\(c_1^B(H)/m\).  Restrictions of a total class along a proper smooth
submersion are horizontal, and a holomorphic line gives type \((1,1)\).
Its fiber degree is one by the calculation of \(H\).  The raw equality
in \eqref{p3sw:descent-zeta} is the degree-one comparison
\eqref{p3sw:descent-tate}; centrally it gives
\(a_{\ell_0}=\zeta_{H,0}\).

We now verify ordinary semiregularity.  The twists of \(p_0^*\mathcal B\)
and \(\ell_0\) cancel, so \(F_0\) is ordinary coherent.  It need not
be locally free.  Since \(p_0\) is flat, the derived identity
\eqref{p3sw:descent-rhom-basechange} and line cancellation give
\[
 R\mathcal{H}om_{Z_0}(F_0,F_0)\simeq p_0^*\mathcal C_{\mathcal B}.
\]
Moreover \(Rp_{0*}\mathcal O_{Z_0}=\mathcal O_{Y_0}\).  To check
this equality, pull back along the finite etale cover
\([r]:Y_0\to Y_0\); the bundle becomes
\(Y_0\times\mathbb P^{m-1}\), where the assertion is the vanishing
of the higher cohomology of \(\mathcal O_{\mathbb P^{m-1}}\).
Flat base change for this fixed projective fiber and faithful flatness
give the equality on \(Y_0\).  The perfect projection formula now yields
\[
 \begin{split}
 \operatorname{Ext}^2_{Z_0}(F_0,F_0)
 &\simeq H^2R\Gamma(Z_0,p_0^*\mathcal C_{\mathcal B})\\
 &\simeq H^2R\Gamma(Y_0,\mathcal C_{\mathcal B})
 =\operatorname{Ext}^2_{\mathfrak Y_0}(\mathcal B,\mathcal B).
 \end{split}
\]
This is \eqref{p3sw:descent-ordinary-ext}; it uses the full derived
Hom complex.  Pullback and line tensoring preserve evaluation and
coevaluation, so this identification commutes with the unit and perfect
trace under \(p_0^*:H^2(Y_0,\mathcal O_{Y_0})\xrightarrow{\sim}
H^2(Z_0,\mathcal O_{Z_0})\) and identifies their trace-free kernels.
The tensor rule and
\(a_{\ell_0}=\zeta_{H,0}\) give
\eqref{p3sw:descent-ordinary-sigma}.

The raw Hodge pullback in that equation is injective.  Indeed, for
\(\eta_H\in H^{j+2}(Y_0,\Omega_{Y_0}^j)\), put
\(\eta_B=(2\pi i)^{-j}\iota_{Y_0}(\eta_H)\), a complex Betti class.
Comparison commutes with pullback.  With \(f=m-1\), normalized Betti
integration gives
\[
 p_{0*}^B\bigl(\zeta_{B,0}^{\,f}p_0^*\eta_B\bigr)=\eta_B,
 \qquad
 \int_{\mathbb P^f}c_1^B(\mathcal O(1))^f=1.
\]
Thus \(p_0^*\eta_H=0\) implies \(\eta_B=0\), and then
\(\eta_H=0\).  This avoids any choice of normalization for Hodge
Gysin.  The exponential \(e^{\zeta_{H,0}}\) is invertible on the
finite direct sum \(\mathcal H_{-2}(Z_0)\).  Equation
\eqref{p3sw:descent-ordinary-sigma} and semiregularity of
\(\mathcal B\) prove full semiregularity of \(F_0\).

Finally, pullback is a morphism of cohomology local systems and
\(\zeta_B\) is horizontal, so \eqref{p3sw:descent-horizontal-character}
is horizontal.  Each degree-\(2k\) part is rational of type \((k,k)\).
Its flat lift \((2\pi i)^kA_{B,k}\) lies in \(F^k\), and its edge
in \(R^k(a\circ p)_*\Omega^k_{Z/S}\) is the raw Hodge class to be used
in deformation theory.  Multiplicativity of the comparison and the
one Tate factor contributed by each \(\zeta_H\) give, centrally,
\begin{equation}
\label{p3sw:descent-central-character}
 \begin{split}
 A_H(0)&=p_0^*\operatorname{ch}_H(\mathcal B)e^{\zeta_{H,0}}
          =\operatorname{ch}_H(F_0),\\
 A_B(0)&=p_0^*\operatorname{ch}_B(\mathcal B)e^{\zeta_{B,0}}
          =\operatorname{ch}_B(F_0).
 \end{split}
\end{equation}
Here \(A_{H,k}\) denotes the edge of \((2\pi i)^kA_{B,k}\), and
\(A_H=\sum_kA_{H,k}\).
The normalized projection formula and the hyperplane degree give
\[
 \begin{split}
 p_{s*}^B\bigl(\zeta_{B,s}^{\,m-1}e^{-\zeta_{B,s}}A_B(s)\bigr)
 &=p_{s*}^B\bigl(\zeta_{B,s}^{\,m-1}p_s^*\kappa_B(s)\bigr)\\
 &=\kappa_B(s)\,p_{s*}^B(\zeta_{B,s}^{\,m-1})
 =\kappa_B(s).
 \end{split}
\]
All exponentials are finite polynomials in cohomology.  This proves
\eqref{p3sw:descent-extraction} and the proposition.
\end{proof}
\subsection{Deformation over a smooth projective analytic germ}
\label{p3sw:bf-section}

The ordinary coherent-sheaf deformation statement below is the
projective case of the argument in Buchweitz--Flenner
\cite[Theorem~5.1 and its proof, pp.~174--179]{P3SWBuchweitzFlenner2003},
whose proof constructs a flat analytic sheaf with the prescribed central
fiber.  We reconstruct that argument for the projective bundle in
Proposition~\ref{p3sw:ordinary-conversion}, keeping the full coefficient
obstruction group explicit.  Projectivity permits a comparison of the
Chern characters of all infinitesimal lifts and an effectivity argument
through compatible quotients.

\begin{theorem}
\label{p3sw:projective-germ}
Let \(S=(S,0)\) be a smooth complex analytic germ, and let
\(\pi:\mathcal X\to S\) be a proper holomorphic submersion of relative
dimension \(n\), equipped with a holomorphic \(\pi\)-ample line
\(\mathcal L_{\mathcal X}\).  Let \(F_0\) be an ordinary coherent
sheaf on \(X_0\).  For \(I\subseteq\{1,\ldots,n\}\), suppose the following
raw semiregularity map is injective:
\begin{equation}
\label{p3sw:bf-seed-sigma}
 \begin{split}
 \sigma_{I,0}:\operatorname{Ext}^2_{X_0}(F_0,F_0)&\longrightarrow
       \bigoplus_{k\in I}H^{k+1}(X_0,\Omega_{X_0}^{k-1}),\\
 v&\longmapsto
       \left(\operatorname{Tr}_{F_0,\Omega_{X_0}^{k-1}[k+1]}
       \left(\left(\frac{a_{F_0}^{k-1}}{(k-1)!}\right)[2]\circ v\right)\right)_{k\in I}
 \end{split}
\end{equation}
Suppose there are rational flat sections
\(\beta_{B,k}\) of \(R^{2k}\pi_*\Q\), for \(k\in I\), such that
\[
 \beta_{B,k}(0)=\operatorname{ch}_{B,k}(F_0),\qquad
 \beta_{B,k}\in F^k(R^{2k}\pi_*\C\otimes\mathcal O_S).
\]
The second condition means that the flat de Rham class itself lies in
the indicated Hodge filtration.  Then, after shrinking \(S\), there
is a coherent \(S\)-flat sheaf \(F\) on \(\mathcal X\), with
\(F|_{X_0}\simeq F_0\), for which
\begin{equation}
\label{p3sw:bf-fiber-character}
 \operatorname{ch}_{B,k}(F|_{X_s})=\beta_{B,k}(s)
       \quad(s\text{ near }0,\ k\in I).
\end{equation}
These classes are rational algebraic cycle classes.  No rank,
simplicity, stability, local-freeness, reflexivity, or fixed-support
hypothesis is imposed on \(F_0\).
\end{theorem}

We use the raw convention \(a=-\operatorname{At}_{\mathrm{BF}}\)
from \eqref{p3sw:descent-character}.  With ordinary de Rham
differential \(d\), the raw degree-\(k\) Chern character is the Hodge
edge of \((2\pi i)^k\operatorname{ch}_{B,k}\).
Thus the prescribed raw Hodge section in the theorem is the edge
of \((2\pi i)^k\beta_{B,k}\).  The factor does not change
horizontality or membership in \(F^k\).

\subsubsection{Absolute perfectness and the coefficient obstruction group}

Write
\[
 R=\mathcal O_{S,0}=\C\{t_1,\ldots,t_e\},\qquad
 A_\nu=R/\mathfrak m^{\nu+1},\qquad
 X_\nu=\mathcal X\times_S\operatorname{Specan}A_\nu.
\]
Here \(\mathfrak m\) is the maximal ideal of \(R\).  The integer
\(\nu\geq0\) records the infinitesimal order.

We construct compatible \(A_\nu\)-flat sheaves \(E_\nu\) with
central fiber \(F_0\).  If \(E_\nu\) has been constructed, the
obstruction to the next lift has coefficients in
\(J_\nu=\ker(A_{\nu+1}\to A_\nu)\).  The first lemma identifies
its group with
\(\operatorname{Ext}^2_{X_0}(F_0,F_0)\otimes_\C J_\nu\), on
which central semiregularity remains injective.  For these extensions,
differentiation gives an injection
\(J_\nu\to\Omega^1_{A_{\nu+1}/\C}\otimes_{A_{\nu+1}}A_\nu\).
The comparison with Gauss--Manin then forces the obstruction to vanish
if each selected component of the relative Chern character of
\(E_\nu\) is the edge of its prescribed flat filtered class.
We prove this character identity for every
infinitesimal lift in Lemma~\ref{p3sw:bf-artin-character}; thus any
lift produced at one stage satisfies the hypothesis needed at the next.
Finally, we express the sheaves as quotients of one vector bundle by
lifting the quotient-defining sections compatibly.  The resulting
formal section of the relative Quot space admits an analytic
approximation preserving the central quotient.  Pulling back the universal
quotient gives a flat sheaf with central fiber \(F_0\), whose fiber
characters are determined by flat transport.

\begin{lemma}
\label{p3sw:bf-coefficients}
Let \(A=A_\nu\), \(X=X_\nu\), and let \(E\) be an \(A\)-flat
coherent sheaf on \(X\) whose central fiber is \(F_0\).
Then \(E\) is absolutely perfect as an \(\mathcal O_X\)-complex,
with local projective dimension at most \(n\).  If \(W\) is a
finite \(A\)-module killed by the maximal ideal and
\(W_X=\mathcal O_X\otimes_AW\), there are natural identifications
\begin{equation}
\label{p3sw:bf-coefficient-identities}
 \begin{split}
 \operatorname{Ext}^2_X(E,E\otimes_AW)
   &\simeq\operatorname{Ext}^2_{X_0}(F_0,F_0)\otimes_\C W,\\
 H^q(X,\Omega^j_{X/A}\otimes_AW)
   &\simeq H^q(X_0,\Omega_{X_0}^j)\otimes_\C W.
 \end{split}
\end{equation}
Under them the coefficient semiregularity map
\[
 \sigma_{I,W}(v)=
 \left(\operatorname{Tr}_{E,(\Omega_{X/A}^{k-1}\otimes_AW)[k+1]}
 \left(\left(\frac{a_{X/A}(E)^{k-1}}{(k-1)!}
       \otimes_A^{\mathbf L}1_W\right)[2]\circ v\right)\right)_{k\in I}
\]
is exactly \(\sigma_{I,0}\otimes1_W\), and is therefore injective.
\end{lemma}

\begin{proof}
At \(x\in X_0\), put \(B=\mathcal O_{X,x}\) and \(M_E=E_x\).
Both are \(A\)-flat, while \(B_0=B\otimes_A\C\) is regular local
of dimension \(n\).  Choose a minimal finite-rank free
\(B\)-resolution of \(M_E\), initially allowed to be infinite.
Tensoring it over \(A\) with \(\C\) remains exact because its
terms and \(M_E\) are \(A\)-flat.  It is still minimal over
\(B_0\), since maximal-ideal entries remain in the maximal ideal
after reduction.  The regular-local projective-dimension theorem
gives projective dimension at most \(n\) over \(B_0\), so this
minimal reduced resolution has no terms in higher degree.  Its
ranks are the ranks before reduction, and the original resolution
therefore stops as well.  This proves absolute
\(\mathcal O_X\)-perfectness.  It is separate from relative
Tor-amplitude \([0,0]\), which is precisely the assumed
\(A\)-flatness; the nonreduced total space has not been declared
regular.
For the general fiberwise perfectness criterion, see
\cite[Proposition~4.5 and Lemma~4.6, p.~169]{P3SWBuchweitzFlenner2003}.

Let \(i:X_0\hookrightarrow X\).  Flatness of \(\mathcal O_X\)
and \(E\) over \(A\), together with derived associativity, gives
\begin{equation}
\label{p3sw:bf-derived-coefficient}
 Li^*E\simeq E\otimes_A^L\C=F_0,\qquad
 E\otimes_{\mathcal O_X}^L W_X
   \simeq E\otimes_A^L W
   =E\otimes_AW=i_*(F_0\otimes_\C W).
\end{equation}
In particular the last derived tensor is an ordinary sheaf.  We
do not use the generally false identity
\(Li^*W_X=\mathcal O_{X_0}\otimes_\C W\).
Perfect duality and the projection formula for the closed
immersion give
\[
 R\mathcal{H}om_X(E,E\otimes_AW)
 \simeq E^\vee\otimes^L i_*(F_0\otimes W)
 \simeq i_*(F_0^\vee\otimes^L F_0\otimes W).
\]
Derived global sections prove the first identity in
\eqref{p3sw:bf-coefficient-identities}.  The coefficient sheaf
\(\Omega^j_{X/A}\otimes_AW\) is the closed pushforward of
\(\Omega_{X_0}^j\otimes_\C W\), proving the second.
These are finite direct-sum identities, not higher-direct-image
base change for an endomorphism complex.

The relative first-jet Atiyah class pulls back to that of \(F_0\),
and evaluation \(E^\vee\otimes E\to\mathcal O_X\) pulls back to
central evaluation.  Apply these identities to the displayed
perfect-Hom formula.  Each coefficient semiregularity component
becomes the corresponding central component tensored with \(W\).
This proves the asserted equality and injectivity.
\end{proof}

\subsubsection{A square-zero lift and its semiregularity obstruction}

Fix \(\nu\), put \(A'=A_{\nu+1}\), \(A=A_\nu\), and write
\[
 J=\mathfrak m^{\nu+1}/\mathfrak m^{\nu+2},\qquad
 K=\Omega^1_{A'/\C}\otimes_{A'}A,\qquad
 j:X=X_\nu\hookrightarrow X'=X_{\nu+1}.
\]
The ideal of \(X\) is \(J_X=\mathcal O_X\otimes_AJ\).
We have \(J^2=0\) and \(\mathfrak m_AJ=0\).  Moreover
\begin{equation}
\label{p3sw:bf-jacobi}
 K=\Omega^1_{R/\C}/\mathfrak m^{\nu+1}\Omega^1_{R/\C}
       \simeq A^{\oplus e},\qquad d:J\hookrightarrow K.
\end{equation}
Indeed \(d(\mathfrak m^{\nu+2})\subset
\mathfrak m^{\nu+1}\Omega_R^1\).  On homogeneous terms the last
map is
\[
 \operatorname{Sym}^{\nu+1}(\mathfrak m/\mathfrak m^2)
 \longrightarrow
 \operatorname{Sym}^{\nu}(\mathfrak m/\mathfrak m^2)
          \otimes(\mathfrak m/\mathfrak m^2),
\]
given by differentiation.  Euler's formula
\(\sum_i t_i\partial_i f=(\nu+1)f\) proves injectivity in
characteristic zero.

\begin{lemma}
\label{p3sw:bf-one-step}
Let \(E\) be an \(A\)-flat coherent sheaf on \(X\) with central
fiber \(F_0\).  Suppose, for every \(k\in I\), that its raw relative
character
\(\operatorname{ch}^{\rm rel}_{H,k}(E)
=\operatorname{tr}(a_{X/A}(E)^k)/k!\)
is the restriction of the Hodge edge of a flat \(F^k\) de Rham
class on \(X'/A'\).  Then \(E\) has a coherent \(A'\)-flat
sheaf lift on \(X'\), with its prescribed restriction to \(X\).
\end{lemma}

\begin{proof}
We first record the precise coefficient exactness needed for
\eqref{p3sw:bf-jacobi}.  The relative ample line algebraizes every
\(X_\nu\), and every coherent sheaf on it, over the finite complex
algebra \(A_\nu\).  We use projective analytic GAGA in the form of
an exact equivalence on coherent sheaves and comparison on their
cohomology; local analytification is faithfully flat
\cite{P3SWSerreGAGA1956}.  The algebraic model is projective and
smooth over \(A_\nu\): its analytification is locally a smooth
product, and smoothness for these finite-type models is reflected
by analytification.  Deligne's proper-smooth theorem
\cite[Theorem~5.5]{P3SWDeligne1968} applies in characteristic zero,
including to nonreduced Artinian bases.  It states that
\(H^q(X_\nu,\Omega^j_{X_\nu/A_\nu})\) is finite locally free,
commutes with arbitrary algebra base change, and that the
Hodge-to-de-Rham spectral sequence degenerates.

Here is the finite-module consequence.  For a finite \(A\)-module
\(W\), apply algebra base change to the split square-zero algebra
\(A\oplus W\).  On the same underlying space the relative
differential sheaf is
\(\Omega^j_{X/A}\oplus(\Omega^j_{X/A}\otimes_AW)\).
The base-change isomorphism respects the inclusion and augmentation
of \(A\oplus W\).  Taking the \(W\)-summand gives, naturally in \(W\),
\begin{equation}
\label{p3sw:bf-hodge-coefficients}
 H^q(X,\Omega^j_{X/A}\otimes_AW)
       \simeq H^q(X,\Omega^j_{X/A})\otimes_AW.
\end{equation}
The latter Hodge module is free, so this functor of \(W\) is exact.
In particular \(d:J\hookrightarrow K\) stays injective after
\eqref{p3sw:bf-hodge-coefficients}.  GAGA transfers the same
assertion to analytic sheaves.

For the remaining derived calculation, algebraize \(X'\), reduce
it to obtain \(X\), and algebraize \(E\); use the same symbols for
this projective algebraic pair.  Faithfully flat local
analytification reflects \(A\)-flatness, and the proof of
Lemma~\ref{p3sw:bf-coefficients} gives absolute perfectness on
the algebraic model too.  The derived foundations used here are
the usual tensor and adjunction formalism for modules, cotangent
transitivity, the first conormal homology
\(H^{-1}(L_{X/X'})=J_X/J_X^2\), and the functorial derived
principal-parts construction of the Atiyah class.  For a perfect
complex the trace is evaluation; its cyclicity and its
compatibility with pullback and coefficient maps follow in a
finite free model.  We also use the characteristic-zero shifted
exterior convention, for a cotangent complex \(L\),
\(L\operatorname{Sym}^k(L[1])\simeq
(L\!\bigwedge^k L)[k]\).
These are the standard algebraic cotangent and perfect-duality
constructions; no obstruction theorem for a general perfect
complex is imported.

Only ordinary relative objects have to be compared back to
analysis.  For the smooth morphism \(X/A\), its cotangent
complex is \(\Omega^1_{X/A}[0]\) and its derived self-product
is underived by flatness.  The first-principal-parts sequence is
split for the input \(\mathcal O_X\)-action, so tensoring it with
any coherent \(E\) is exact.  Thus the relative derived Atiyah
triangle is the ordinary first-jet extension.  Ordinary
differentials, conormal sequences, and first jets commute with
analytification, as follows from finite presentations; perfect
trace does too by local free evaluation.  For
\(Q=E\otimes_AJ\), GAGA for the bounded coherent complex
\(E^\vee\otimes^LQ\) compares
\(\operatorname{Ext}^2_X(E,Q)\) with the analytic group, and
ordinary GAGA compares the Hodge targets.  Hence the coefficient
map of Lemma~\ref{p3sw:bf-coefficients} is the same one on the
model.  We do not compare an absolute nonreduced analytic
cotangent complex or an \(\operatorname{Ext}^1\) group on \(X'\);
the extension below will be constructed algebraically and then
analytified.

Define the cotangent extension class
\[
 \delta_X:L_{X/\C}\longrightarrow L_{X/X'}
       \longrightarrow H^{-1}(L_{X/X'})[1]=J_X[1].
\]
The whole complex \(L_{X/X'}\) need not be \(J_X[1]\).
Smoothness of \(X'/A'\), restricted to \(X\), gives the exact
sequence
\begin{equation}
\label{p3sw:bf-differentials}
 0\longrightarrow\mathcal O_X\otimes_AK
 \longrightarrow\Omega^1_{X'/\C}\otimes\mathcal O_X
 \longrightarrow\Omega^1_{X/A}\longrightarrow0.
\end{equation}
The conormal map \(J_X\to\Omega^1_{X'/\C}\otimes\mathcal O_X\)
factors through its first arrow by \(1\otimes d\).  It is
injective by \eqref{p3sw:bf-jacobi} and \(A\)-flatness.
Consequently the ordinary conormal sequence ending in
\(\Omega^1_{X/\C}\) is short exact as well.

Map the transitivity triangle
\[
 Lj^*L_{X'/\C}\longrightarrow L_{X/\C}\longrightarrow L_{X/X'}
\]
to this conormal triangle by the degree-zero truncations.
The induced cone map \(L_{X/X'}\to J_X[1]\) is the canonical
truncation: on \(H^{-1}=J_X/J_X^2\) it is the identity, and,
since \(L_{X/X'}\in D^{\leq-1}\), truncation adjunction
determines a map to \(J_X[1]\) by that map on \(H^{-1}\).
Now map the conormal sequence to
\eqref{p3sw:bf-differentials}, using \(1\otimes d\) on the
left and the identity in the middle.  Naturality of boundaries
and exterior comultiplication give the derived square
\begin{equation}
\label{p3sw:bf-conormal-square}
\begin{array}{ccc}
 L\!\bigwedge^k L_{X/\C}&\xrightarrow{\ \rho_k\ }&\Omega^k_{X/A}\\
 \big\downarrow{\scriptstyle\delta_{X,k}}&&
        \big\downarrow{\scriptstyle\nabla'_k}\\
 (\Omega^{k-1}_{X/A}\otimes_AJ)[1]&
        \xrightarrow{\ 1\otimes d\ }&
        (\Omega^{k-1}_{X/A}\otimes_AK)[1].
\end{array}
\end{equation}
Here \(\rho_k\) projects each factor to relative forms,
\(\delta_{X,k}\) contracts one factor with \(\delta_X\) and
projects the others, and \(\nabla'_k\) is the exterior boundary
of \eqref{p3sw:bf-differentials}.  The exterior signs are those
of the shifted symmetric convention just specified.  This
argument does not require a regular immersion.

We next prove the required sheaf obstruction statement.  By
\eqref{p3sw:bf-derived-coefficient},
\[
 Q=E\otimes_AJ
    \simeq E\otimes_{\mathcal O_X}^LJ_X
\]
is an ordinary coherent sheaf.  Put
\(\mathcal K=Lj^*j_*E\).  The Tor sequence for
\(0\to J_X\to\mathcal O_{X'}\to\mathcal O_X\to0\) gives
\(H^0(\mathcal K)=E\) and \(H^{-1}(\mathcal K)=Q\);
lower cohomology is allowed.  Its truncation triangle is
\begin{equation}
\label{p3sw:bf-self-intersection}
 Q[1]\longrightarrow\tau_{\geq-1}\mathcal K
       \longrightarrow E\xrightarrow{b_E}Q[2].
\end{equation}
Derived adjunction identifies
\(\operatorname{Ext}^1_{X'}(j_*E,j_*Q)\) with
\(\operatorname{Hom}_{D(X)}(\mathcal K,Q[1])\).
Because \(Q\) is a sheaf, this Hom factors through
\(\tau_{\geq-1}\mathcal K\).  Applying
\(\operatorname{Hom}(-,Q[1])\) to
\eqref{p3sw:bf-self-intersection} gives the exact segment
\begin{equation}
\label{p3sw:bf-extension-sequence}
 \operatorname{Ext}^1_{X'}(j_*E,j_*Q)
 \longrightarrow\operatorname{Hom}_X(Q,Q)
 \longrightarrow\operatorname{Ext}^2_X(E,Q).
\end{equation}
Use the cochain convention in which a connecting map is the
differential of a lift.  The last map sends \(\operatorname{id}_Q\)
to \(b_E\).  If the first class is represented by
\(0\to j_*Q\to E'\to j_*E\to0\), the Tor boundary after
restriction sends \(j_0\otimes e\) to \(j_0\widetilde e\),
where \(\widetilde e\) lifts \(e\) in \(E'\).  This follows
directly from the two-term presentation
\(J_X\to\mathcal O_{X'}\) of \(\mathcal O_X\), with the same
connecting-map convention.  It is the map \(Q\to Q\) in
\eqref{p3sw:bf-extension-sequence}.  Thus \(b_E=0\) precisely
when there is an extension
\begin{equation}
\label{p3sw:bf-multiplication-extension}
 0\longrightarrow E\otimes_AJ\longrightarrow E'
       \longrightarrow E\longrightarrow0
\end{equation}
whose kernel identification is the multiplication map
\(J_X\otimes_{\mathcal O_X}E\to E'\).
Yoneda Ext supplies an actual extension; it is coherent because
coherent modules on the noetherian model are closed under
extensions.

We identify \(b_E\) with the Atiyah contraction, including its
orientation.  Keep the two actions on
\[
 \mathcal C=\mathcal O_X^{(1)}
       \otimes^L_{\mathcal O_{X'}}\mathcal O_X^{(2)},\qquad
 \mathcal I=\operatorname{fib}(\mathcal C\to\mathcal O_X).
\]
Tensor over action (1), using action (2) as output.  Then
\(\mathcal C\otimes^L_{(1)}E=\mathcal K\),
\(\mathcal I\in D^{\leq-1}\), and
\(H^{-1}(\mathcal I)=J_X\).
Universal derivation on action (2) gives the diagonal
linearization
\(\ell:\mathcal I\to L_{X/X'}\), oriented by
\(1\otimes f-f\otimes1\mapsto df\).
To check its first homology, resolve factor (2) by a local free
DG algebra over \(\mathcal O_{X'}\).  A degree-\(-1\)
generator \(u\) with \(d_{\rm DG}u=j_0\in J_X\) represents
\(j_0\) in \(H^{-1}(\mathcal I)\); its universal differential
represents the same class in
\(H^{-1}(L_{X/X'})=J_X/J_X^2\).  Thus \(\ell\) induces the
identity under the canonical first-conormal-homology
identifications.  Universal derivation makes this check
compatible on overlaps.

The homotopy pushout of
\(\mathcal I\otimes^L_{(1)}E\to\mathcal K\) along
\(\ell\otimes E\) is the derived first-principal-parts triangle
\[
 E\otimes^L L_{X/X'}\longrightarrow P^1_{X/X'}(E)
 \longrightarrow E
 \xrightarrow{\operatorname{At}_{X/X'}(E)}
 E\otimes^L L_{X/X'}[1].
\]
The boundary is the principal-parts definition of Atiyah.
Its sign relative to Buchweitz--Flenner can be read from first
jets.  For ordinary degree-zero jets put
\(\jmath(p)=(1\otimes1)\otimes p\), with \(p\) on action (1)
and the module action on action (2).  With
\(df=1\otimes f-f\otimes1\),
\[
 \jmath(pf)-\jmath(p)f=-p\,df.
\]
A connection satisfies
\(\nabla(pf)=\nabla(p)f+p\,df\), so
\(\jmath+\nabla\) is the module splitting.  In a free DG model
its failure to commute with the differential is
\([d_{\rm DG},\nabla]\).  This is precisely
\(\operatorname{At}_{\mathrm{BF}}\) in
\cite[Definition~3.4]{P3SWBuchweitzFlenner2003}, not its negative.

Let \(t:L_{X/X'}\to J_X[1]\) be truncation.  Since \(t\ell\)
is the identity on \(H^{-1}\), truncation adjunction identifies
it with \(\mathcal I\to J_X[1]\).  After tensoring with \(E\),
the analogous assertion has target \(Q[1]\): the tensor
cohomology spectral sequence has in degree \(-1\) only
\(H^{-1}(\mathcal I)\otimes E=Q\), and all higher Tor terms
have lower degree.  The derived target is \(Q[1]\) by
\eqref{p3sw:bf-derived-coefficient}.  Naturality of the fiber
and truncation triangles therefore identifies \(b_E\) with the
principal-parts Atiyah class followed by \(t[1]\).  Naturality
under \(L_{X/\C}\to L_{X/X'}\) gives
\[
 b_E=(1_E\otimes\delta_X)[1]\operatorname{At}_{\mathrm{BF},X/\C}(E).
\]
Define \(o_E=-b_E\).  With \(a=-\operatorname{At}_{\mathrm{BF}}\),
we have proved
\begin{equation}
\label{p3sw:bf-obstruction}
 o_E=\langle\delta_X,a_{X/\C}(E)\rangle.
\end{equation}
It vanishes if and only if
\eqref{p3sw:bf-multiplication-extension} exists with its
multiplication identification.
The use of the two highest cohomology degrees of \(\mathcal K\)
depends on \(Q\) being a sheaf.  For a general perfect complex,
the Hom in \eqref{p3sw:bf-extension-sequence} need not factor
through this truncation.

Absolute perfectness supplies the trace in the next identity.
Put \(L=L_{X/\C}\).  Under
\(L\operatorname{Sym}^k(L[1])\simeq(L\!\bigwedge^kL)[k]\),
the Atiyah morphism \(a:E\to E\otimes L[1]\) is a degree-zero
coefficient-valued endomorphism.  The map
\(\delta_X[1]:L[1]\to J_X[2]\) extends as the derivation
contracting each symmetric generator.  Its output coefficient
has even shift two.  Expanding contraction of \(a^k\) gives
\(k\) terms with the same supertrace by graded cyclicity: the
symmetric algebra accounts for the internal Koszul signs, and
rotation of the degree-zero Atiyah operations and the even
coefficient leaves no residual sign.  Trace commutes with
coefficient maps.  Dividing by \(k!\) proves
\begin{equation}
\label{p3sw:bf-trace-identity}
 \sigma_{k,J}(o_E)
 =\delta_{X,k}\left(\frac{\operatorname{tr}(a_{X/\C}(E)^k)}{k!}\right).
\end{equation}
This supplies the particular obstruction and trace identities
corresponding to \cite[Propositions~4.4 and 4.2, pp.~166--169]{P3SWBuchweitzFlenner2003}.

The multiplication identification also gives \(A'\)-flatness.
For completeness, the square-zero flatness criterion says that
an \(A'\)-module \(M'\) is flat if \(M'/JM'\) is \(A\)-flat and
\(J\otimes_A(M'/JM')\to M'\) is injective.  The latter
condition is \(\operatorname{Tor}^{A'}_1(A,M')=0\).
The change-of-rings Tor sequence then gives
\(\operatorname{Tor}^{A'}_1(W,M')=0\) for every \(A\)-module
\(W\).  Filtering an arbitrary \(A'\)-module by \(JW\subset W\)
gives the same vanishing for every \(W\), which is flatness.
For \eqref{p3sw:bf-multiplication-extension}, the kernel
identification gives \(E'/JE'=E\) and the required injection.
Apply the criterion to the algebraic stalks.  Flat local
analytification preserves this \(A'\)-flatness, coherence, and
the extension.

We now relate the obstruction to the assumed horizontal class.
The exterior boundary \(\nabla'_k\) in
\eqref{p3sw:bf-conormal-square} is the restricted Hodge-graded
Gauss--Manin map.  Indeed, quotient the absolute de Rham
complex of \(X'/\C\) by forms with at least two base
differentials.  The connecting map of the resulting sequence
is base differentiation in product coordinates.  Filtering
by relative form degree gives on the graded term exactly the
exterior boundary of \eqref{p3sw:bf-differentials}.  Hodge
degeneration identifies the degree-\(2k\) graded term with
\(H^k(\Omega^k)\), and
\eqref{p3sw:bf-hodge-coefficients} identifies its restriction
to \(A\).  Therefore the restriction of the edge of a flat
\(F^k\) class on \(X'/A'\) is killed by \(\nabla'_k\).

The image under \(\rho_k\) of the absolute character in
\eqref{p3sw:bf-trace-identity} is the relative character of
\(E\), by Atiyah naturality.  Evaluate
\eqref{p3sw:bf-conormal-square} on that absolute character
and use \eqref{p3sw:bf-trace-identity}.  The result is
\begin{equation}
\label{p3sw:bf-semiregularity-square}
 (1\otimes d)_*\sigma_{k,J}(o_E)
       =\nabla'_k\bigl(\operatorname{ch}^{\rm rel}_{H,k}(E)\bigr).
\end{equation}
This is the restricted Gauss--Manin square used in
\cite[Lemma~5.8 and Proposition~5.9]{P3SWBuchweitzFlenner2003}.
The hypothesis of the lemma kills its right side for \(k\in I\).
The left coefficient map is injective by
\eqref{p3sw:bf-hodge-coefficients}, and
\(\sigma_{I,J}\) is injective by
Lemma~\ref{p3sw:bf-coefficients}, since \(\mathfrak m_AJ=0\).
Thus \(o_E=0\).  Equations
\eqref{p3sw:bf-obstruction} and
\eqref{p3sw:bf-multiplication-extension}, followed by the
flatness criterion, give the required analytic flat sheaf lift.
\end{proof}

\subsubsection{The Chern character of every infinitesimal lift}

The one-step lemma can be iterated only after checking that a chosen
lift has the prescribed relative character, rather than merely the
prescribed central restriction.  The projectivity assumption supplies
a direct proof for every stage.

\begin{lemma}
\label{p3sw:bf-artin-character}
For every \(A_\nu\)-flat coherent sheaf \(E_\nu\) on \(X_\nu\)
with central fiber \(F_0\), and every \(0\leq k\leq n\), the ordinary flat comparison
\[
 H^l_{\rm dR}(X_\nu/A_\nu)
       \simeq H^l(X_0,\C)\otimes_\C A_\nu
\]
satisfies
\begin{equation}
\label{p3sw:bf-artin-character-formula}
 \begin{split}
 \operatorname{ch}^{\rm rel}_{H,k}(E_\nu)
 &=\operatorname{edge}_k\bigl((2\pi i)^k
                 \operatorname{ch}_{B,k}(F_0)\otimes1\bigr),\\
 (2\pi i)^k\operatorname{ch}_{B,k}(F_0)\otimes1
 &\in F^kH^{2k}_{\rm dR}(X_\nu/A_\nu).
 \end{split}
\end{equation}
The assertion holds independently of the choice of the lift.
\end{lemma}

\begin{proof}
For \(k=0\) the assertion is constancy of rank on each connected
component.  Assume \(k\geq1\).
The projective model of \(X_\nu\) has an ample line.  Repeated
Serre global-generation surjections from finite sums of negative
powers of that line resolve \(E_\nu\).  The local
projective-dimension bound in Lemma~\ref{p3sw:bf-coefficients}
makes the \(n\)-th syzygy locally free.  GAGA gives a global finite holomorphic vector-bundle
resolution \(V^\bullet\to E_\nu\).  No compatibility of these
resolutions as \(\nu\) varies is asserted or needed.

Choose a differentiable trivialization of the proper submersion
\(\mathcal X\to S\).  Holomorphic \(t\)-jets through order \(\nu\)
of relative smooth forms define a morphism
\[
 \Omega^\bullet_{X_\nu/A_\nu}
       \longrightarrow \mathcal A^\bullet_{X_0}\otimes_\C A_\nu,
\]
where \(\mathcal A^\bullet_{X_0}\) is the smooth de Rham complex.
Explicitly, after the trivialization take
\[
 \sum_{|\mathbf a|\leq\nu}
 \frac{(\partial_t^{\mathbf a}\omega)|_{t=\bar t=0}}
      {\mathbf a!}\,t^{\mathbf a}.
\]
Leibniz's rule makes this compatible with products, and jets
commute with the fiber differential.  The jets of the fiberwise
type projections give a bigrading on the target; their
integrability identities persist under taking jets.  In a local
holomorphic product chart the resulting Dolbeault complex is the
ordinary one with coefficients \(A_\nu\), up to an invertible
formal smooth coordinate change.  The coefficientwise Dolbeault
lemma therefore makes it a fine resolution of each
\(\Omega^p_{X_\nu/A_\nu}\), respecting the form filtration.

Both total complexes resolve the constant sheaf \(A_\nu\).
For the holomorphic complex, filter by powers of \(\mathfrak m\);
its associated graded complex is
\(\Omega^\bullet_{X_0}\otimes_\C\operatorname{gr}A_\nu\),
to which the holomorphic Poincare lemma applies.  The jet map
is the identity on constants.  It therefore realizes the
ordinary flat de Rham comparison in the statement.

A holomorphic vector bundle \(V\) on \(X_\nu\) becomes a smooth
\(A_\nu\)-bundle under the same jet map, and is smoothly
isomorphic to \(V_0^\infty\otimes A_\nu\).  Indeed, modulo
successive powers of \(\mathfrak m\), the discrepancy between its
transition matrices and the central ones is a Cech one-cocycle
with values in smooth sections of
\(\operatorname{End}(V_0^\infty)\otimes
\mathfrak m^a/\mathfrak m^{a+1}\).  This sheaf is fine, so a
smooth gauge change removes that discrepancy.  Induction through
the finitely many powers of \(\mathfrak m\) gives the isomorphism.

Choose a smooth right-module connection
\(D=\bar\partial_V+\nabla^{1,0}\) on this bundle, and put
\(F_D=D^2\).  Its curvature has no \((0,2)\)-part.  In the
Dolbeault resolution the BF commutator
\([\bar\partial,\nabla]\), with its form coefficient written
after the \((0,q)\)-form, maps under exterior multiplication to
\(F_D^{1,1}\).  The sign also follows from the first-jet
calculation above: for a line with frames \(e_j=e_i g_{ij}\),
the BF cocycle is \(-d\log g_{ij}\), while the connection
matrices differ by \(d\log g_{ij}\); the Cech--Dolbeault
boundary carries this cocycle to \(+\bar\partial A_i\).
Thus \(a=-\operatorname{At}_{\mathrm{BF}}\) is represented by
\(-F_D^{1,1}\).  It follows that the closed form
\[
 \frac{(-1)^k}{k!}\operatorname{tr}(F_D^k)
\]
belongs to \(F^k\) and that its edge is the raw degree-\(k\)
Atiyah character.

This de Rham class is the constant central class.  Use a constant
connection on \(V_0^\infty\otimes A_\nu\) and the smooth
isomorphism just constructed.  For
\(D_u=D_0+u(D_1-D_0)\), the Chern--Simons identity is
\[
 \partial_u\frac{\operatorname{tr}(F_{D_u}^k)}{k!}
 =d\,\frac{\operatorname{tr}((D_1-D_0)F_{D_u}^{k-1})}{(k-1)!}.
\]
Integration in \(u\) proves independence of the connection;
the polynomial identity is valid over the finite complex algebra
\(A_\nu\).  For this right-connection convention, the normalized
topological character is
\(\operatorname{tr}\exp(-F_D/(2\pi i))\).  Hence the signed
unscaled form above represents
\((2\pi i)^k\operatorname{ch}_{B,k}(V_0)\) in the flat comparison.

Apply the calculation termwise to \(V^\bullet\).  A
grading-preserving connection represents
\(\operatorname{At}_{\mathrm{BF}}\) by its commutator with
\(d_{V^\bullet}+\bar\partial\), and \(a\) by the negative of
that commutator.  In a power, every term involving
\(d_{V^\bullet}\) has positive internal degree and zero
supertrace.  The remaining terms are the alternating sum of the
signed bundle curvature traces.  Their edge is therefore
\(\operatorname{ch}^{\rm rel}_{H,k}(E_\nu)\).  By
\eqref{p3sw:bf-derived-coefficient}, the central restriction of
the resolution resolves \(F_0\).  Its alternating central
character is \(\operatorname{ch}_{B,k}(F_0)\), which proves
\eqref{p3sw:bf-artin-character-formula}.
\end{proof}

\subsubsection{Compatible quotients and analytic effectivity}

\begin{proof}[Proof of Theorem~\ref{p3sw:projective-germ}]
Let \(\beta_{H,k,\nu}\) be the edge of the restriction of
\((2\pi i)^k\beta_{B,k}\) to \(X_\nu/A_\nu\).  A flat section
is uniquely determined near the origin by its central value.
Thus Lemma~\ref{p3sw:bf-artin-character} identifies the
relative character of every \(E_\nu\) with
\(\beta_{H,k,\nu}\).  Starting with \(E_0=F_0\), the
assumption that \(\beta_{B,k}\) lies in \(F^k\) provides the
flat lift to order \(\nu+1\).  Lemma~\ref{p3sw:bf-one-step}
then gives \(E_{\nu+1}\) with its identification on \(X_\nu\).
This constructs a compatible tower of flat coherent sheaves.
Only the extensions \(A_{\nu+1}\to A_\nu\) and their modules
\(J_\nu\) were used, not arbitrary Artinian extensions.

We make the classifying maps compatible before applying analytic
approximation.  Choose \(c\) large enough that
\(F_0(c)=F_0\otimes\mathcal L_{\mathcal X,0}^{c}\) is generated
by \(b\) global sections and has \(H^1(X_0,F_0(c))=0\).
Let
\(\mathcal V=(\mathcal L_{\mathcal X}^{-c})^{\oplus b}\)
and fix the central quotient
\(\mathcal V_0\twoheadrightarrow F_0\).  Given a quotient
\(\mathcal V_\nu\twoheadrightarrow E_\nu\), tensor the
extension for \(E_{\nu+1}\) by
\(\mathcal L_{\mathcal X}^{c}\).  Its kernel is
\(i_{0*}(F_0(c)\otimes_\C J_\nu)\), whose first cohomology
vanishes.  The global-section map to order \(\nu\) is therefore
surjective.  Lift the \(b\) sections defining the quotient.
Their map to \(E_{\nu+1}\) is surjective by Nakayama's lemma,
because its central restriction is the fixed central quotient.
We obtain compatible quotients
\(\mathcal V_\nu\twoheadrightarrow E_\nu\) for every \(\nu\).

We use the relative analytic Quot theorem of Pourcin
\cite[Theorem~2, pp.~451--452]{P3SWPourcin1969}: proper flat
quotients of a coherent sheaf on a separated finite-dimensional
complex space over a separated finite-dimensional complex base,
which may be nonreduced, are represented by a complex space
with a universal flat quotient.  These hypotheses hold here;
in particular quotient supports are proper because
\(\mathcal X\to S\) is proper.  Let \(T\) be the germ of this
Quot space at the central quotient.  The compatible quotients
give unique compatible \(S\)-maps
\(\operatorname{Specan}A_\nu\to T\).  Passing to the inverse
limit gives a local \(R\)-algebra map
\[
 \mathcal O_{T,0}\longrightarrow\widehat R
          =\C[[t_1,\ldots,t_e]],
\]
that is, a formal section of \(T\to S\).

A graph embedding presents this analytic \(S\)-germ as
\[
 \mathcal O_{T,0}
 =R\{z_1,\ldots,z_h\}/(g_1,\ldots,g_v).
\]
The formal section gives zero-constant formal series \(z_j(t)\)
solving these finitely many convergent equations.  Artin's
analytic approximation theorem
\cite[Theorem~1.2, p.~277]{P3SWArtin1968} says that such a
formal solution over \(\C\) has an exact convergent solution
with any prescribed finite-order congruence.  Take congruence
modulo \(\mathfrak m\).  The solution still maps the origin to
the central quotient, and, with \(t\) fixed as the parameters,
defines an analytic section \(S\to T\).  Pullback of the
universal quotient gives a coherent \(S\)-flat sheaf \(F\)
with central fiber \(F_0\).  The convergent solution need not
realize the entire chosen formal tower; retaining the central
fiber is sufficient.

It remains to identify characters on actual fibers.  The total
space \(\mathcal X\) is smooth, so \(F\) is absolutely perfect.
Separately, \(S\)-flatness and derived associativity give, for
every \(s\),
\begin{equation}
\label{p3sw:bf-derived-fiber}
 Li_s^*F
 \simeq F\otimes^L_{\mathcal O_{\mathcal X}}
       (\mathcal O_{\mathcal X}\otimes^L_{\mathcal O_S}k(s))
 \simeq F\otimes^L_{\mathcal O_S}k(s)
 =F|_{X_s}[0].
\end{equation}
Near \(X_0\), \(F\) has a finite resolution by total holomorphic
vector bundles: use relative Serre generation repeatedly and
the regular-local projective-dimension bound on the smooth
total space, then shrink by properness at the final locally
free syzygy.  The alternating normalized topological Chern
character of these bundles restricts by
\eqref{p3sw:bf-derived-fiber} to
\(\operatorname{ch}_B(F|_{X_s})\).  Restrictions of a total
cohomology class along a proper smooth submersion are flat.
Their central value is \(\operatorname{ch}_B(F_0)\), so
uniqueness of flat transport proves
\eqref{p3sw:bf-fiber-character}.

Every \(X_s\) is projective algebraic, and projective GAGA
algebraizes \(F|_{X_s}\).  On the smooth projective fiber a
finite locally free resolution identifies its normalized
topological character with the cycle class of its rational
algebraic Chern character.  Its raw degree-\(k\) Atiyah
character is the weightwise \((2\pi i)^k\) image of that class.
This proves algebraicity and completes the theorem.
\end{proof}

For \(p:Z\to Y\) in Proposition~\ref{p3sw:ordinary-conversion},
the morphism \(a\circ p\) satisfies the theorem's projectivity
hypothesis, \(F_0\) is ordinary coherent and fully semiregular,
and its relative dimension is \(N=6+m-1\).  Taking
\(I=\{1,\ldots,N\}\), the raw input in degree \(k\) is exactly
the edge of \((2\pi i)^kA_{B,k}\).  The theorem uses neither a
relative version of the finite isogeny \(q\) nor a deformation
of the twisted sheaf \(\mathcal B\).  Its proof also does not
assert arbitrary classical base change for global Ext groups
or higher direct images.  The smooth projective germ and the
ordinary coherent-sheaf condition are the stated limits of the
deformation result.
\section{Algebraicity at every compatible period}
\label{p3:sec:all-periods}

We apply the descent and deformation results to the normalized character
computed for the polarized seed. Projective-bundle operations first give
local algebraicity. A marked period family and proper Hilbert parameter
spaces then extend the Weil plane to every compatible period, and rational
isogenies remove the dependence on the initial lattice.

\subsection{The descended character on the marked Weil domain}
\label{p3sw:endpoint-marked-domain}

Fix \(K=\Q(u)\), with \(u^2=-d\) and \(d\geq3\), and use the seed
\(M=X\times\widehat X\) constructed above.  Its coherent perfect sheaf
\(\mathcal E\) has rank \(r=8d\).  We write \(\operatorname{cl}_B\) for
the usual rational Betti cycle-class map, normalized so that the first
Chern class \(c_1^B(L)\) of a line bundle is integral.  The normalized
character on the seed is
\begin{equation}
\label{p3sw:endpoint-seed-character}
 \kappa_{M,B}
   =\exp\!\left(-\frac{c_1^B(\mathcal E)}r\right)
       \operatorname{ch}_B(\mathcal E)
   =\sum_{j=0}^6\kappa_{M,B,j}.
\end{equation}
Here \(c_1^B(\mathcal E)=c_1^B(\det\mathcal E)\), so this is the
class \(\kappa_M\) of Section~\ref{p3:sec:spin} in normalized Betti
notation.  Proposition~\ref{p3sw:polarized-seed} supplies the unital
\(K\)-action on \(M\) and its ample class \(\theta_M\), and
Proposition~\ref{p3sw:exceptional-character} gives
\begin{equation}
\label{p3sw:endpoint-exceptional-component}
 \kappa_{M,B,3}=c_\theta\theta_M^3+w_M,\qquad
 c_\theta\in\Q,\qquad 0\ne w_M\in\mathcal W_K(M).
\end{equation}
Proposition~\ref{p3sw:hodge-circle} proves that every
\(\kappa_{M,B,j}\) has type \((j,j)\) at every marked complex
structure of Weil type \((3,3)\) compatible with this action and
polarization.  We will transport these properties to the quotient
on which the descended sheaf is defined.

\subsubsection{The actual homological Hermitian space}

Let \(e_M\) be the Riemann form of the ample class \(\theta_M\).
The actual homological Hermitian form is
\[
 h_M(x,y)=e_M(x,\eta_M(u)_*y)+u e_M(x,y).
\]
In the second-linear convention of \eqref{p3ab:eq:hermitian},
Proposition~\ref{p3sw:polarized-seed} proves that this form is
hyperbolic of signature \((3,3)\); its three hyperbolic blocks are
computed in \eqref{p3sw:spin-eq:homological-gram}.

Let \(q:M\longrightarrow Y_0\) be the isogeny in
Proposition~\ref{p3sw:twisted-descent}.  Put
\[
 V_0=H_1(Y_0,\Q),\qquad \Lambda_0=H_1(Y_0,\Z),\qquad
 \theta_0=(q^*)^{-1}\theta_M.
\]
Transport the rational action through \(q\) in the category up to
isogeny, and write \(T_0=\eta_0(u)_*\) for the transported action
on \(V_0\).  The alternating form \(e_0\) associated to \(\theta_0\)
is characterized by
\[
 e_0(q_*x,q_*y)=e_M(x,y).
\]
It is rational and is a positive Riemann form at \(Y_0\); a positive
integer multiple is integral on \(\Lambda_0\) and is an ample
polarization.  The second-linear Hermitian form
\begin{equation}
\label{p3sw:endpoint-reference-hermitian}
 h_0(x,y)=e_0(x,T_0y)+u e_0(x,y)
\end{equation}
is carried to \(h_M\) by \(q_*\), so it too is hyperbolic of
signature \((3,3)\).

\subsubsection{The full compatible domain}

Let \(\mathscr D\) be the set of real complex structures \(I\) on
\(V_{0,\R}\) satisfying
\begin{equation}
\label{p3sw:endpoint-compatible-structures}
 I^2=-1,\qquad IT_0=T_0I,\qquad
 e_0(Ix,Iy)=e_0(x,y),\qquad e_0(x,Ix)>0\quad(x\ne0).
\end{equation}
This is the full marked compatible domain for the fixed rational
space, form, and lattice.  We verify both its connectedness and the
complex structure used below.

Choose the embedding \(\sigma:K\hookrightarrow\C\) for which
\(\sigma(u)=i\sqrt d\), and view \(V_{0,\sigma}=V_0\otimes_{K,\sigma}\C\)
with its Hermitian form \(h_{0,\sigma}\).  A real \(I\) commuting with
\(K\) is complex-linear on this space.  Commutation with \(T_0\)
together with preservation of \(e_0\) implies that \(I\) is
unitary for \(h_{0,\sigma}\).  Its \(i\)- and \(-i\)-eigenspaces are
therefore orthogonal.  On the first of them \(T_0=\sqrt d\,I\), and
on the second \(T_0=-\sqrt d\,I\).  Since
\(h_0(x,x)=e_0(x,T_0x)\), Riemann positivity makes the first
eigenspace positive and the second negative.  Each has dimension
three because \(h_0\) has signature \((3,3)\).

Conversely, if \(P\subset V_{0,\sigma}\) is a positive three-plane,
then \(P^\perp\) is negative of dimension three.  Set \(I=i\) on
\(P\) and \(I=-i\) on \(P^\perp\).  This is a real complex structure
commuting with \(K\) and preserving \(h_0\), hence also \(e_0\),
which is the coefficient of \(u\) in \eqref{p3sw:endpoint-reference-hermitian}.
For \(x=x_++x_-\) in this orthogonal decomposition,
\[
 e_0(x,Ix)
   =d^{-1/2}\bigl(h_{0,\sigma}(x_+,x_+)
                       -h_{0,\sigma}(x_-,x_-)\bigr)>0
 \quad(x\ne0).
\]
Thus positive three-planes parameterize exactly \(\mathscr D\).
Fix a positive/negative orthogonal splitting
\(V_{0,\sigma}=V_+\oplus V_-\).  Projection of a positive
three-plane to \(V_+\) is an isomorphism, so that plane is the graph
of a unique map \(Z:V_+\to V_-\).  Positivity is precisely
\begin{equation}
\label{p3sw:endpoint-matrix-ball}
 1-Z^*Z>0,\qquad\text{equivalently}\qquad \lVert Z\rVert_{\mathrm{op}}<1.
\end{equation}
The operator-norm unit ball in
\(\operatorname{Hom}_\C(V_+,V_-)\) is convex.  In particular the
underlying domain is connected and contractible and has complex
dimension nine.

For the holomorphic convention, use cohomological periods.  Write
\(W_\sigma\oplus W_{\bar\sigma}=H^1(Y_0,\C)\) for the two
\(K\)-eigenspaces.  The cohomological polarization induced by \(e_0\)
pairs these two spaces perfectly and vanishes on each separately.
If \(F^1\) is a compatible Lagrangian Hodge subspace, then
\(P^\vee=F^1\cap W_\sigma\) is a three-plane and
\begin{equation}
\label{p3sw:endpoint-lagrangian-period}
 F^1=P^\vee\oplus\operatorname{Ann}(P^\vee)
       \subset W_\sigma\oplus W_{\bar\sigma},
\end{equation}
where the annihilator is taken using that perfect pairing.
Conversely this formula, restricted by the open Riemann positivity
condition, gives precisely the compatible periods.  The annihilator
map between the two Grassmannians is algebraic.  We give
\(\mathscr D\) the complex structure in the coordinate \(P^\vee\);
then \eqref{p3sw:endpoint-lagrangian-period} is holomorphic into the
Siegel period domain.  This coordinate can be the conjugate of the
homological graph coordinate in \eqref{p3sw:endpoint-matrix-ball};
the connected underlying space and its complex dimension are unchanged.

\subsubsection{A fine principal-level family and the seed lattice}

Choose a rational symplectic basis
\(a_1,\ldots,a_6,b_1,\ldots,b_6\) for \(e_0\), normalized by
\(e_0(a_i,b_j)=\delta_{ij}\), and let \(\Lambda_{\mathrm{pp}}\)
be its integral span.  Such a basis follows by repeatedly choosing
a pair of pairing one and taking its symplectic orthogonal complement.
The restriction of \(e_0\) to \(\Lambda_{\mathrm{pp}}\) is integral
and unimodular.  The Riemann criterion for polarized tori therefore
makes
\begin{equation}
\label{p3sw:endpoint-reference-tori}
 A_I^{\mathrm{pp}}
       =(V_{0,\R},I)/\Lambda_{\mathrm{pp}},\qquad I\in\mathscr D,
\end{equation}
principally polarized abelian sixfolds.  The rational \(K\)-action
persists in \(\operatorname{End}^0(A_I^{\mathrm{pp}})\): every rational
\(K\)-operator has an integer multiple preserving
\(\Lambda_{\mathrm{pp}}\), and it commutes with \(I\).

We use the following precise fine-level theorem.  Over \(\C\), fix a
primitive fourth root for the symplectic Weil pairing.  The functor
of principally polarized abelian schemes of relative dimension six
with a full symplectic level-\(4\) isomorphism to the standard
\((\Z/4)^6\times\mu_4^6\) has a smooth quasi-projective fine moduli
scheme \(\mathscr A_6(4)\) and a universal abelian scheme
\(\pi:\mathscr U\to\mathscr A_6(4)\).  On the chosen analytic component
it is the quotient of the universal marked torus family on
\(\mathfrak H_6\) by the principal congruence subgroup
\(\Gamma_6(4)\); both this component and the universal total space
are quasi-projective algebraic varieties.  These are Chai's
symplectic-level and uniformization statements
\cite[Chapter~I, Theorems~1.4 and~1.7.2, and Sections~2.3--2.5,
pp.~5--13]{P3SWChai1985}.
The hypotheses here are explicit: \(e_0\) is perfect on
\(\Lambda_{\mathrm{pp}}\), the chosen basis supplies the level,
\(4\geq3\) removes stabilizers, and the ground field has
characteristic zero.

The universal \(\pi\) is smooth and projective.  Smoothness and
properness are part of being the universal abelian scheme.  To see
projectivity from the stated theorem without choosing a universal
theta line, embed the quasi-projective \(\mathscr U\) locally
closed in some \(\mathbb P^b_\C\).  Its graph over the separated
scheme \(\mathscr A_6(4)\) is a locally closed immersion
\[
 \mathscr U\longrightarrow\mathbb P^b_{\mathscr A_6(4)}.
\]
This morphism is proper: for any morphism from an
\(\mathscr A_6(4)\)-proper scheme to an
\(\mathscr A_6(4)\)-separated scheme, the graph is closed and the
second projection is a base change of a proper morphism.  A proper
immersion is a closed immersion.  The displayed map is therefore a
projective embedding over \(\mathscr A_6(4)\).

The symplectic basis identifies the space of all positive Lagrangians
with \(\mathfrak H_6\).  The compatible Lagrangians in
\eqref{p3sw:endpoint-lagrangian-period} give a holomorphic map
\(\mathscr D\to\mathfrak H_6\).  Composition with the level quotient gives
\begin{equation}
\label{p3sw:endpoint-level-period-map}
 \nu:\mathscr D\longrightarrow\mathscr A_6(4)^{\mathrm{an}}.
\end{equation}
The quoted marked uniformization identifies
\(\nu^*\mathscr U^{\mathrm{an}}\) with the family
\eqref{p3sw:endpoint-reference-tori}.  This construction uses the
ordinary principally polarized level family.  It asserts neither
that \(\mathscr D\) or its image is algebraic nor that the ambient
universal family has a \(K\)-action.

We also obtain the precise family with central fiber \(Y_0\).
Choose an integer \(c_{\mathrm{lat}}>0\) such that
\(c_{\mathrm{lat}}\Lambda_{\mathrm{pp}}\subset\Lambda_0\).  The
finite group
\[
 G_{\mathrm{lat}}=
  c_{\mathrm{lat}}^{-1}\Lambda_0/\Lambda_{\mathrm{pp}}
       \subset V_0/\Lambda_{\mathrm{pp}}
\]
consists of torsion points of uniformly bounded order.  In the
marked family on the contractible \(\mathscr D\) these are
holomorphic torsion sections.  Indeed, the kernel of multiplication
by their common order is a finite \'etale holomorphic cover of
\(\mathscr D\), and the marking identifies it with the constant
torsion lattice.  The chosen torsion classes therefore select
sections, which are holomorphic because the cover is locally
biholomorphic.  Translation by
\(G_{\mathrm{lat}}\) is free.  Taking the quotient, locally by
disjoint translated coordinate neighborhoods, gives a smooth
holomorphic group family with zero section.  It is proper because
it is the finite quotient of a proper family.  Its fiber is
\((V_{0,\R},I)/(c_{\mathrm{lat}}^{-1}\Lambda_0)\).  Multiplication
by \(c_{\mathrm{lat}}\) identifies it with
\begin{equation}
\label{p3sw:endpoint-seed-lattice-family}
 Y_I=(V_{0,\R},I)/\Lambda_0,
\end{equation}
and identifies the seed fiber with \(Y_0\).

This family is polarized in the relative sense required in
Proposition~\ref{p3sw:ordinary-conversion}.  Indeed, choose an integer
\(c_{\mathrm{pol}}>0\) so that \(c_{\mathrm{pol}}e_0\) is integral
on \(\Lambda_0\).  Let \(\mathcal P_I\) be the normalized Poincar\'e
line on \(Y_I\times\widehat Y_I\).  Mark the dual homology by
\[
 H_1(\widehat Y_I,\Z)\xrightarrow{\sim}\Lambda_0^\vee,\qquad
 \xi\longmapsto\bigl[y\longmapsto
 c_{1,\mathrm B}(\mathcal P_I)((y,0),(0,\xi))\bigr].
\]
In this marking the dual complex structure is \(-I^t\).  The constant
lattice map
\[
 F:\Lambda_0\longrightarrow\Lambda_0^\vee,\qquad
 x\longmapsto\bigl[y\longmapsto c_{\mathrm{pol}}e_0(y,x)\bigr]
\]
extends to the marked real vector spaces.  It is complex-linear from
\(I\) to the dual complex structure \(-I^t\), since
\[
 F(Ix)(y)=c_{\mathrm{pol}}e_0(y,Ix)
   =-c_{\mathrm{pol}}e_0(Iy,x)=(-I^tF(x))(y).
\]
Equivalently, its constant
complexification preserves the holomorphic Hodge subbundles in the
marked uniformizations, so it induces a holomorphic homomorphism
from the family \(Y\) to its dual family.  On each fiber the Riemann
criterion supplies an ample line \(L_I\) with alternating form
\(c_{\mathrm{pol}}e_0\).  The normalized mixed-biextension identity
\eqref{p3sw:support-polarization-biextension}, evaluated on
\((y,0),(0,x)\), gives
\[
 ((\phi_{L_I})_*x)(y)=c_{\mathrm{pol}}e_0(y,x),\qquad
 \phi_{L_I}(z)=t_z^*L_I\otimes L_I^{-1}.
\]
Homomorphisms of complex tori are determined by their maps on the
universal covers, so the family homomorphism is fiberwise this
polarization.  Its ample class is the marked transport of
\(c_{\mathrm{pol}}\theta_0\).  In particular all fibers are projective.
The local representing line, and its relative ampleness after
shrinking, are supplied by the relative Picard argument immediately
preceding Proposition~\ref{p3sw:ordinary-conversion}.  Restricting
\eqref{p3sw:endpoint-seed-lattice-family} to a small contractible
open neighborhood \(S\) of the seed gives the polarized
holomorphic abelian scheme \(a:Y\to S\) used there; \(S\) is a
smooth nine-dimensional analytic germ.

The quotient also gives, on each fiber, an isogeny
\begin{equation}
\label{p3sw:endpoint-reference-isogeny}
 j_I:A_I^{\mathrm{pp}}\longrightarrow Y_I
\end{equation}
induced by \(c_{\mathrm{lat}}\operatorname{id}_{V_0}\).  It is
\(K\)-linear in the category up to isogeny.  It is algebraic:
both tori are projective, and Serre's GAGA Proposition~15 says that
a holomorphic map from a compact algebraic variety to an algebraic
variety is regular
\cite[Section~4, no.~19, Proposition~15, pp.~29--30]{P3SWSerreGAGA1956}.
Only this fiberwise isogeny is needed; the central isogeny \(q\)
has not been extended as a relative algebraic isogeny.

\subsubsection{The flat descended character}

Proposition~\ref{p3sw:twisted-descent} supplies the determinant-trivial
twisted coherent sheaf \(\mathcal B\) of rank \(r\) on \(Y_0\),
with full twisted-derived semiregularity.  Its character comparison
\eqref{p3sw:descent-character-pullback} is the equality
\begin{equation}
\label{p3sw:endpoint-descended-character}
 \kappa_B:=\operatorname{ch}_B(\mathcal B)\in
       \bigoplus_{j=0}^6 H^{2j}(Y_0,\Q),\qquad
 q^*\kappa_B=\kappa_{M,B}.
\end{equation}
Rationality here is part of that comparison, or follows by applying
the rational trace \((\deg q)^{-1}q_*^B\) to its right side.
Because \(q_*\) is an isomorphism of rational first homology,
\(q^*\) is an isomorphism of rational cohomology rings.  It is
\(K\)-linear on first cohomology for the transported action and
therefore identifies the two Weil planes.  Consequently
\eqref{p3sw:endpoint-exceptional-component} becomes
\begin{equation}
\label{p3sw:endpoint-descended-component}
 \kappa_{B,3}=c_\theta\theta_0^3+w_0,\qquad
 0\ne w_0\in\mathcal W_K(Y_0).
\end{equation}

Use the lattice marking in \eqref{p3sw:endpoint-seed-lattice-family}
to transport \(\kappa_B\), \(\theta_0\), and \(w_0\), denoting their
values by \(\kappa_B(I)\), \(\theta_I\), and \(w_I\).
They are rational flat sections.  For each \(I\in\mathscr D\),
transporting \(I\) through \(q_*\) gives a compatible polarized
complex structure on the rational homology of \(M\).  The constant
map \(q^*\) is a Hodge map for these corresponding structures.
Proposition~\ref{p3sw:hodge-circle}, applied to the invariant
character in \eqref{p3sw:endpoint-seed-character}, thus proves
\begin{equation}
\label{p3sw:endpoint-flat-hodge}
 \kappa_{B,j}(I)\in H^{j,j}(Y_I)\cap H^{2j}(Y_I,\Q)
       \quad(0\leq j\leq6,\ I\in\mathscr D).
\end{equation}
Also \(\theta_I\) is a rational polarization class at every
period, and \eqref{p3sw:endpoint-descended-component} remains the
flat equality
\[
 \kappa_{B,3}(I)=c_\theta\theta_I^3+w_I,\qquad
 0\ne w_I\in\mathcal W_K(Y_I).
\]
This is an assertion on the marked domain.  No section on an
arithmetic quotient and no relative twisted sheaf is used to
define these flat classes.

\subsection{Local algebraicity and projective extraction}
\label{p3sw:endpoint-local-section}

\begin{proposition}[Local algebraicity of the descended character]
\label{p3sw:local-algebraicity}
There is a nonempty analytic neighborhood of the seed in
\(\mathscr D\) on which every component \(\kappa_{B,j}(I)\) is a
rational algebraic cycle class on \(Y_I\).
\end{proposition}

\begin{proof}
Take the smooth contractible nine-dimensional germ \(S\) and its
polarized abelian scheme \(a:Y\to S\) constructed above.
Apply Proposition~\ref{p3sw:ordinary-conversion} to the descended
sheaf \(\mathcal B\), and retain its notation
\[
 m=r^{12},\qquad N=6+m-1,\qquad p:Z\to Y,\qquad
 F_0=p_0^*\mathcal B\otimes\ell_0.
\]
Thus \(a\circ p\) is a proper holomorphic submersion of relative
dimension \(N\) with a relative ample line.  The sheaf \(F_0\) is
ordinary coherent, and its raw semiregularity map is injective on
the entire group \(\operatorname{Ext}^2_{Z_0}(F_0,F_0)\).
The proposition also supplies the relative anticanonical line and
its normalized Betti class
\begin{equation}
\label{p3sw:endpoint-zeta}
 H=\omega_{Z/Y}^{-1},\qquad
 \zeta_{B,s}=\frac{c_1^B(H|_{Z_s})}{m}.
\end{equation}
The class \(\zeta_B\) is rational, flat, and of type \((1,1)\),
and restricts to the hyperplane class on every \(p_s\)-fiber.
Each \(Z_s\) is projective algebraic, \(p_s\) is a projective
algebraic morphism, and \(H|_{Z_s}\) is an algebraic line bundle.

Now put
\begin{equation}
\label{p3sw:endpoint-horizontal-A}
 A_B(s)=p_s^*\kappa_B(s)e^{\zeta_{B,s}}
       =\sum_{k=0}^{N} A_{B,k}(s).
\end{equation}
Equations
\eqref{p3sw:endpoint-flat-hodge} and \eqref{p3sw:endpoint-zeta}
show that \(A_{B,k}\) is a rational flat section of
\(R^{2k}(a\circ p)_*\Q\) of type \((k,k)\).  In particular its
flat de Rham lift lies in \(F^k\).  The central character identity
\eqref{p3sw:descent-central-character} gives
\begin{equation}
\label{p3sw:endpoint-central-A}
 A_B(0)=p_0^*\operatorname{ch}_B(\mathcal B)e^{\zeta_{B,0}}
       =\operatorname{ch}_B(F_0).
\end{equation}

In the raw \(d\log\) convention, the prescribed Hodge section in
degree \(k\) is the edge of \((2\pi i)^k A_{B,k}\).  Its central
value is \(\operatorname{ch}_{H,k}(F_0)\) by
\eqref{p3sw:descent-tate} and \eqref{p3sw:endpoint-central-A}.
Apply Theorem~\ref{p3sw:projective-germ} with
\(\beta_{B,k}=A_{B,k}\) and \(I_{\mathrm{BF}}=\{1,\ldots,N\}\).
This index set retains the full semiregularity injection: every
potentially nonzero target \(H^{j+2}(Z_0,\Omega_{Z_0}^j)\) has
\(j\leq N-2\) and occurs at \(k=j+1\); the target for \(k=N\)
is zero.  All the theorem's hypotheses now hold, so after one common
shrinking of \(S\) every \(A_{B,k}(s)\), \(1\leq k\leq N\),
is rational algebraic.  The degree-zero term \(A_{B,0}=r\) is
already algebraic.

On a fiber of this smaller germ, choose rational Chow classes
representing all \(A_{B,k}(s)\), and represent \(\zeta_{B,s}\)
by \(c_1^{\mathrm{CH}}(H|_{Z_s})/m\).
Products and the finite exponential polynomial give a rational
Chow class whose Betti image is
\(\zeta_{B,s}^{m-1}e^{-\zeta_{B,s}}A_B(s)\).
Proper pushforward along the algebraic \(p_s\) has Betti class
\(\kappa_B(s)\) by \eqref{p3sw:descent-extraction}.
Every exponential is truncated in the cohomology or Chow ring
at the dimension \(N\).  Taking homogeneous components of this
proper pushforward proves that every \(\kappa_{B,j}(s)\) is
rational algebraic.
\end{proof}

\subsubsection{Two algebraic generators of the rational plane}

\begin{lemma}
\label{p3sw:endpoint-two-generators}
Let \(A\) be an abelian sixfold with a unital action of the imaginary
quadratic field \(K\).  If a nonzero vector of \(\mathcal W_K(A)\)
is a rational algebraic class, then every vector of
\(\mathcal W_K(A)\) is a rational algebraic class.
\end{lemma}

\begin{proof}
Choose \(k=1+tu\in K^\times\) with \(t\in\Q\) such that
\(\sigma(k)^6\ne\bar\sigma(k)^6\).  Such \(t\) exists because
\[
 (1+tu)^6-(1-tu)^6
\]
is a nonzero polynomial in \(t\), with linear coefficient \(12u\);
only finitely many rational \(t\) are excluded.  Over \(K\) (or
\(\overline\Q\)) the plane is the direct sum of the two conjugate
sixth exterior eigenspaces from the definition of
\(\mathcal W_K(A)\).  A nonzero rational vector \(w\) has nonzero
coordinates in both lines: scalar conjugation exchanges them, so
one zero coordinate would force the other to be zero.  On the two
lines, \(k^*=\eta(k)^*\) has respective eigenvalues
\(\sigma(k)^6\) and \(\bar\sigma(k)^6\).  The determinant formed
by \(w\) and \(k^*w\) in those coordinates is therefore nonzero.
They are linearly independent over \(\Q\), and form a basis of
\(\mathcal W_K(A)\).

Choose an integer \(n_k>0\) clearing the denominator of
\(\eta(k)\in\operatorname{End}^0(A)\).  The resulting
\(h=n_k\eta(k)\) is an algebraic endomorphism, and on \(H^6\)
one has \(k^*=n_k^{-6}h^*\).  Pullback by \(h\) preserves rational
algebraic cycle classes, so \(k^*w\) is algebraic when \(w\) is.
Both basis vectors, hence the whole rational plane, are algebraic.
\end{proof}

On the open supplied by Proposition~\ref{p3sw:local-algebraicity},
the class
\[
 w_I=\kappa_{B,3}(I)-c_\theta\theta_I^3
\]
is rational algebraic: an integer multiple of \(\theta_I\) is an
ample polarization class.  It is nonzero everywhere because it is
the transport of the single nonzero marked vector \(w_0\).
Lemma~\ref{p3sw:endpoint-two-generators} makes
\(\mathcal W_K(Y_I)\) algebraic on that open.
The isogeny \eqref{p3sw:endpoint-reference-isogeny} pulls this
plane isomorphically to \(\mathcal W_K(A_I^{\mathrm{pp}})\):
its pullback on \(H^1\) is a \(K\)-linear rational isomorphism and
thus identifies both sixth exterior eigenspaces.  Algebraic
pullback preserves their cycle classes.

Let
\[
 \mathbb W_{\mathrm{pp}}\subset
       \nu^*R^6\pi^{\mathrm{an}}_*\Q
\]
be the constant rank-two rational local subsystem defined by the
fixed marked \(K\)-eigenspaces; its fiber is
\(\mathcal W_K(A_I^{\mathrm{pp}})\).  We have proved that all its
fibers are algebraic on a nonempty analytic open of \(\mathscr D\).
Only this subsystem, not the individual character \(\kappa_B\),
will be used in the global step.

\subsection{Spreading the plane in the fine family}
\label{p3sw:endpoint-spread-section}

\begin{proposition}[Analytic Hilbert spread]
\label{p3sw:spread}
Let \(\pi:\mathcal X\to S\) be a smooth projective morphism of
relative dimension \(n\) over a smooth quasi-projective complex
variety.  Let \(\nu:T\to S^{\mathrm{an}}\) be holomorphic, where
\(T\) is a connected complex manifold.  Fix \(0\leq c\leq n\)
and a finite-rank rational local subsystem
\[
 \mathbb V\subset\nu^*R^{2c}\pi^{\mathrm{an}}_*\Q.
\]
If every vector of \(\mathbb V_t\) is a rational algebraic
codimension-\(c\) class on \(\mathcal X_{\nu(t)}\) for all \(t\)
in a nonempty analytic open subset of \(T\), then the same is
true for every \(t\in T\).
\end{proposition}

\begin{proof}
The assertion is immediate if the rank \(b\) of \(\mathbb V\) is
zero, so assume \(b>0\).  Fix a relatively very ample line for
\(\mathcal X/S\).  For each of \(b\) candidate generators choose
a finite list of rational coefficients and, for each term in
that list, a Hilbert polynomial of degree \(n-c\).  There are
only countably many such finite choices.  For one choice let
\(\mathscr H\) be the product over \(S\) of the corresponding
relative Hilbert schemes, with structure map
\(\rho:\mathscr H\to S\).

The exact representability result used here is Grothendieck's
fixed-polynomial Hilbert theorem: for a projective scheme over a
noetherian base, a fixed relatively very ample line and Hilbert
polynomial \(P\) give a projective relative Hilbert scheme
representing closed subschemes flat over a locally noetherian test
scheme with fiber polynomial \(P\), with its universal flat subscheme.
It is the case \(F=\mathcal O_{\mathcal X}\) of the fixed-polynomial
Quot theorem
\cite[Section~3, Theorem~3.2, pp.~259--260, and the Hilbert
specialization on p.~265]{P3SWGrothendieckHilbert1961}.
Here \(S\) is noetherian and \(\pi\) is projective, so the theorem
applies.  In particular \(\mathscr H\to S\) is projective.
The space \(\mathscr H\) may be singular or nonreduced; we do not
discard any of its fibers.

After analytic base change to \(T\), the parameters whose \(b\)
rational combinations span \(\mathbb V\) will form an open and
closed locus.  Properness will make its image a closed analytic
subset of \(T\).  The countably many choices cover the assumed
open algebraicity locus, so Baire's theorem and the identity
principle will force one image to be all of \(T\).  The first
step is to prove the cycle-class constancy that makes the spanning
condition locally constant, including at singular parameter points.

Put \(X_{\mathscr H}=\mathcal X\times_S\mathscr H\).
For each factor of \(\mathscr H\), let \(\mathcal F\) be the
structure sheaf of its pulled-back universal subscheme on
\(X_{\mathscr H}\).  It is coherent and flat over \(\mathscr H\).
It is also perfect on \(X_{\mathscr H}\), with a uniform local
projective-dimension bound \(n\).  To see this, at \(x\mapsto h\)
write
\[
 R=\mathcal O_{\mathscr H,h},\qquad
 B=\mathcal O_{X_{\mathscr H},x},\qquad
 M=\mathcal F_x,\qquad \overline B=B\otimes_R k(h).
\]
Both \(B\) and \(M\) are \(R\)-flat.  The local fiber ring
\(\overline B\) is regular of dimension at most \(n\), since
\(X_{\mathscr H}\to\mathscr H\) is smooth of relative dimension
\(n\).  Derived associativity and this flatness give
\begin{equation}
\label{p3sw:endpoint-hilbert-tor}
 M\otimes_B^{\mathbf L}k(x)
   \simeq (M\otimes_R k(h))
                \otimes_{\overline B}^{\mathbf L}k(x).
\end{equation}
A regular local ring of dimension \(e\) has global dimension at
most \(e\), as in the regular-local theorem
\cite[Tag~00O7, Proposition~10.110.1]{P3SWStacks}.
The right side of \eqref{p3sw:endpoint-hilbert-tor} has no Tor
homology above \(n\).  Over the noetherian local ring \(B\), build
a minimal resolution of the finite module \(M\) by choosing
minimal finite generators of each successive kernel.  All its
differentials become zero after tensoring with \(k(x)\).  The
vanishing of \(\operatorname{Tor}_{n+1}^B(M,k(x))\) therefore
forces its \((n+1)\)-st free term to vanish.  This proves
\(\operatorname{pd}_B M\leq n\).

The scheme \(X_{\mathscr H}\) is quasi-projective over \(\C\):
both \(\mathscr H\to S\) and \(X_{\mathscr H}\to\mathscr H\)
are projective and \(S\) is quasi-projective.  Choose an ample
line on \(X_{\mathscr H}\).  The ample-line resolution property
says that on a quasi-compact scheme with such a line every
finite-type quasi-coherent sheaf is a quotient of a finite direct
sum of negative powers of the line
\cite[Tag~01Q3, Proposition~28.27.13(8)]{P3SWStacks}.
Successively take these vector-bundle surjections onto
\(\mathcal F\) and its coherent kernels.  The \(n\)-th syzygy is
locally free by the local projective-dimension bound just proved
(and if \(n=0\), \(\mathcal F\) itself is locally free).
Thus \(\mathcal F\) has a finite global vector-bundle resolution
on \(X_{\mathscr H}\).  Derived restriction of that resolution
to \(h\) is an ordinary resolution of \(\mathcal F_h\): its
terms are \(\mathscr H\)-flat because \(X_{\mathscr H}\) is
\(\mathscr H\)-flat, and \(\mathcal F\) is \(\mathscr H\)-flat,
so derived tensoring with \(k(h)\) introduces no higher Tor.

Here is a topological justification of local constancy that does
not require a smooth Hilbert space.  The original
\(\pi^{\mathrm{an}}:\mathcal X^{\mathrm{an}}\to S^{\mathrm{an}}\)
is a proper smooth submersion between manifolds, because \(S\)
and \(\pi\) are smooth.  Ehresmann's fibration theorem gives a
differentiable product over a sufficiently small coordinate ball \(B\)
around \(s_0\):
\(\mathcal X_{s_0}^{\mathrm{an}}\times B\simeq
\pi^{-1}(B)\).  Pull back this trivialization of the original smooth
family to \(\mathscr H^{\mathrm{an}}\); smoothness of the Hilbert
space is not required.

Under this product, restrict one vector bundle in the global
resolution to \(F\times V\), where \(F=\mathcal X_{s_0}^{\mathrm{an}}\)
is compact and \(V\) is a neighborhood of a chosen
\(h_0\in\mathscr H^{\mathrm{an}}\) over \(B\).  Shrinking \(V\),
finitely many product neighborhoods \(U_i\times V\) trivialize
the bundle and cover \(F\times V\).  A partition of unity on
\(F\) subordinate to the \(U_i\), multiplied into these
trivializations, embeds the bundle continuously in a fixed
trivial complex vector bundle on \(F\times V\).  Let \(P(x,h)\)
be the orthogonal projection onto its image; it is continuous.
Compactness of \(F\) permits a further shrinking such that
\[
 \lVert P(x,h)-P(x,h_0)\rVert<1\qquad(x\in F,\ h\in V).
\]
The map \(P(x,h)\) from \(\operatorname{im}P(x,h_0)\) to
\(\operatorname{im}P(x,h)\) is then injective, since a nonzero
vector in its kernel would contradict this inequality; the
ranks agree, so it is an isomorphism.  These maps are continuous
bundle isomorphisms.  Thus the topological Chern character of
the bundle on each fiber is locally constant in \(h\).
Apply this to all the finitely many terms of the resolution.
Their alternating Chern character is
\(\operatorname{ch}_B(\mathcal F_h)\) by the exact fiber
restriction above.  It defines a locally constant section of
the cohomology local system pulled back to
\(\mathscr H^{\mathrm{an}}\).  This argument is valid at singular
and nonreduced parameter points, since their underlying topology
and the restricted topological bundles are all that was used.

To identify the cycle class including nonreduced multiplicities,
use Baum--Fulton--MacPherson's map
\(\tau:K_0\to\CH_*(-)_\Q\) for quasi-projective schemes over a
field.  It preserves the support-dimension filtration and sends
\([\mathcal O_V]\) on the associated graded to \([V]\); on smooth
\(X\), it sends a vector bundle \(E\) to
\(\operatorname{ch}_{\mathrm{CH}}(E)\operatorname{td}(T_X)\cap[X]\)
\cite[Chapter~II, Theorem, pp.~119--120; Chapter~III, Section~1,
Theorem and Lemma, pp.~128--129]{P3SWBaumFultonMacPherson1975}.
Put \(s=\rho(h)\).  For \(\mathcal F_h\) supported in dimension
at most \(n-c\), let \(V_i\) be its dimension-\((n-c)\) components
with generic points \(\eta_i\), and put
\(\ell_i=\operatorname{length}_{\mathcal O_{\mathcal X_s,\eta_i}}
(\mathcal F_{h,\eta_i})\).  Devissage gives
\[
 [\mathcal F_h]\equiv\sum_i\ell_i[\mathcal O_{V_i}]
       \pmod{\text{classes supported in dimension }<n-c}.
\]
Indeed, generic composition series spread to dense opens; their
discrepancies are lower-dimensional.  Thus \(\tau(\mathcal F_h)\)
has no term above dimension \(n-c\) and has top term
\(\sum_i\ell_i[V_i]\).  Additivity on a vector-bundle resolution gives
\[
 \tau(\mathcal F_h)
   =\operatorname{ch}_{\mathrm{CH}}(\mathcal F_h)
       \operatorname{td}(T_{\mathcal X_s})\cap[\mathcal X_s].
\]
Since \(\operatorname{td}_0=1\), comparison in successive
codimensions gives \(\operatorname{ch}_{\mathrm{CH},j}=0\) for
\(j<c\) and
\(\operatorname{ch}_{\mathrm{CH},c}(\mathcal F_h)\cap[\mathcal X_s]
=\sum_i\ell_i[V_i]\).  Compatibility of the cycle map with
normalized Betti Chern classes
\cite[Section~0.8, p.~106]{P3SWBaumFultonMacPherson1975}
therefore gives, for the universal fiber subscheme \(\mathscr Z_h\),
\begin{equation}
\label{p3sw:endpoint-leading-character}
 \operatorname{ch}_{B,c}(\mathcal O_{\mathscr Z_h})
       =\operatorname{cl}_B\bigl([\mathscr Z_h]_{n-c}\bigr).
\end{equation}
Here \([\mathscr Z_h]_{n-c}\) is the fundamental cycle with its
generic scheme lengths.  Lower-dimensional and embedded
contributions affect only higher character terms.  The preceding
local constancy of the character therefore proves precisely the
local constancy of these rational cycle classes on
\(\mathscr H^{\mathrm{an}}\).

Now form the analytic base change
\[
 \mathscr H_T=\mathscr H^{\mathrm{an}}
                 \times_{S^{\mathrm{an}}}T\longrightarrow T.
\]
It is proper and holomorphic: a projective embedding of
\(\mathscr H\) over \(S\) realizes this base change as a closed
analytic subspace of \(T\times\mathbb P^a\), whose projection
is proper.  The fixed rational combinations of the universal
classes in \eqref{p3sw:endpoint-leading-character} give locally
constant sections of the pulled-back cohomology local system.
In local trivializations the condition that these \(b\) sections
lie in and span \(\mathbb V\) is a condition on fixed rational
vectors.  It is locally constant.  Its good locus is open and
closed in \(\mathscr H_T\), hence is a closed analytic subspace
and remains proper over \(T\).

We use Remmert's proper-image theorem in its reduced-space form:
the image of a proper holomorphic map between reduced complex
spaces is an analytic subset.  Grauert gives this precise
statement and its derivation from the coherent direct-image
theorem in Section~7, Satz~1, p.~60
\cite{P3SWGrauert1960}.  Reduce the good locus before applying
it.  Reduction does not change its underlying proper map or
its image, while \(T\) is already reduced.  The theorem thus
applies even though the original Hilbert space could be
nonreduced or singular.  The image of the good locus is a
closed analytic subset of \(T\), closed also by properness.

The countably many choices of coefficients and polynomials
give countably many such good images.  Their union is exactly
the locus where \(\mathbb V_t\) is algebraic.  In one direction,
a good point supplies \(b\) rational algebraic classes spanning
the fiber.  Conversely, if that fiber is algebraic, choose a
rational basis and express each basis vector as a finite
rational combination of integral codimension-\(c\) subvarieties.
Their reduced subschemes have Hilbert polynomials of degree
\(n-c\) and give a good point for one of the choices.

A closed analytic subset unequal to the connected manifold \(T\)
is nowhere dense by the analytic identity principle.  The assumed
nonempty open subset of \(T\) is itself a Baire space, so its
countable covering by the good images cannot consist entirely of
such nowhere dense subsets.  Hence one good image equals \(T\),
which proves the assertion.
\end{proof}

Apply Proposition~\ref{p3sw:spread} with
\[
 S=\mathscr A_6(4),\quad \mathcal X=\mathscr U,\quad
 T=\mathscr D,\quad n=6,\quad c=3,\quad
 \mathbb V=\mathbb W_{\mathrm{pp}}.
\]
The family is smooth projective, \(\mathscr D\) is a connected
complex manifold, and the local algebraicity hypothesis was
proved above.  Therefore
\begin{equation}
\label{p3sw:endpoint-all-reference-periods}
 \mathcal W_K(A_I^{\mathrm{pp}})
     \subset\operatorname{im}\bigl(
        \CH^3(A_I^{\mathrm{pp}})_\Q
             \xrightarrow{\operatorname{cl}_B}
                 H^6(A_I^{\mathrm{pp}},\Q)\bigr)
       \qquad(I\in\mathscr D).
\end{equation}
The proof applies to the analytically pulled-back rational
subsystem.  It does not require an algebraic Weil locus or
an algebraic \(K\)-action on the ambient fine family.

\subsection{Every hyperbolic rational lattice}
\label{p3sw:all-periods}

\begin{proposition}
\label{p3sw:endpoint-all-lattices}
For the fixed \(K=\Q(u)\) with \(u^2=-d\), \(d\geq3\), let
\(A\) be any complex abelian sixfold with a unital \(K\)-action
and a compatible polarization.  If its actual homological
Hermitian form \(h_A(x,y)=e_A(x,uy)+u e_A(x,y)\) is hyperbolic,
then \(\mathcal W_K(A)\) is a rational algebraic plane.
\end{proposition}

\begin{proof}
The hyperbolic-basis assertion of Lemma~\ref{p3ab:hyperbolic}
identifies every nondegenerate rank-six hyperbolic Hermitian space
over \(K\), in the second-linear convention, with three planes
of Gram matrix
\(\left(\begin{smallmatrix}0&1\\1&0\end{smallmatrix}\right)\).
That elementary assertion was proved by splitting isotropic pairs,
independently of the Weil-plane theorem.  Since both \(h_A\) and
the actual seed form \(h_0\) are hyperbolic, it gives a \(K\)-isometry
\begin{equation}
\label{p3sw:endpoint-hermitian-isometry}
 f_A:(H_1(A,\Q),h_A)\longrightarrow(V_0,h_0).
\end{equation}

Let \(\Lambda_A=H_1(A,\Z)\) and let \(I_A\) be the complex
structure of \(A\).  Transport it to
\(I=f_AI_Af_A^{-1}\) on \(V_{0,\R}\).  The isometry is \(K\)-linear,
so \(I\) commutes with \(T_0\).  The Riemann form is recovered
without any dual convention: it is exactly the coefficient of
\(u\) when \(h(x,y)=e(x,uy)+u e(x,y)\) is expanded in the
\(\Q\)-basis \(1,u\) of \(K\).  Thus
\eqref{p3sw:endpoint-hermitian-isometry} implies
\[
 e_0(f_Ax,f_Ay)=e_A(x,y).
\]
It follows that \(I\) preserves \(e_0\) and satisfies
\(e_0(x,Ix)>0\) for \(x\ne0\).  Hence \(I\) is a point of the
full connected domain \(\mathscr D\).  Equation
\eqref{p3sw:endpoint-all-reference-periods} algebraizes the
Weil plane of \(A_I^{\mathrm{pp}}\) at this particular period.

The full rational lattices \(f_A(\Lambda_A)\) and
\(\Lambda_{\mathrm{pp}}\) in \(V_0\) are commensurable.  Choose
an integer \(n_A>0\) such that
\(n_A f_A(\Lambda_A)\subset\Lambda_{\mathrm{pp}}\).  The
invertible complex-linear map \(n_A f_A\) then induces a
surjective holomorphic homomorphism with finite kernel
\begin{equation}
\label{p3sw:endpoint-target-isogeny}
 \varphi:A=(H_1(A,\R),I_A)/\Lambda_A
       \longrightarrow
       A_I^{\mathrm{pp}}=(V_{0,\R},I)/\Lambda_{\mathrm{pp}}.
\end{equation}
Both tori are projective, so this homomorphism is algebraic by
Serre's Proposition~15 cited above.  It is a \(K\)-linear
isogeny in the category up to isogeny.  Its pullback on first
rational cohomology is therefore a \(K\)-linear isomorphism.
It maps each of the two embedding eigenspaces isomorphically,
and taking their sixth exterior powers gives
\[
 \varphi^*\mathcal W_K(A_I^{\mathrm{pp}})=\mathcal W_K(A).
\]
Pullback along the algebraic isogeny preserves rational
codimension-three cycle classes.  The left side is algebraic
by \eqref{p3sw:endpoint-all-reference-periods}, proving the
proposition.

The isogeny need not identify the integral polarizations:
its pullback scales the rational Riemann form by \(n_A^2\).
Nor was a lift of \(f_A\) to a rational spin group used.
Only the elementary hyperbolic-basis argument and a lattice
denominator were required, so extra endomorphisms or a product
decomposition of \(A\) impose no restriction.
\end{proof}

\begin{proof}[Proof of Theorem~\ref{p3sw:weil-plane}]
Every imaginary quadratic field can be written
\(K=\Q(u_0)\), with \(u_0^2=-d_0\) for a positive squarefree
integer \(d_0\).  Set \(u=2u_0\) and \(d=4d_0\geq4\).
This is the same field and is a parameter in the range of the
seed construction.  For any compatible Riemann form,
\[
 h_u(x,y)=e(x,uy)+u e(x,y)=2h_{u_0}(x,y).
\]
More generally, changing a purely imaginary generator by any
nonzero rational factor scales this Hermitian form by that
factor.  Such scaling does not change its isotropic subspaces;
in signature \((3,3)\) it does not change the signature either.
Thus the hyperbolic hypothesis of the theorem is the hyperbolic
hypothesis for the generator \(u\) used by this seed.
Proposition~\ref{p3sw:endpoint-all-lattices} applies and gives
the stated algebraicity of \(\mathcal W_K(A)\).
The full domain in \eqref{p3sw:endpoint-compatible-structures}
and the lattice-clearing isogeny include every compatible period,
with no simplicity or endomorphism-generality assumption.
\end{proof}


\begin{thebibliography}{99}
\setlength{\itemsep}{3pt plus 1pt minus 1pt}
\addcontentsline{toc}{section}{References}

\bibitem{P3SWArtin1968}
M.~Artin,
On the solutions of analytic equations,
\emph{Inventiones Mathematicae} \textbf{5} (1968), 277--291.

\bibitem{P3SWBaumFultonMacPherson1975}
P.~Baum, W.~Fulton, and R.~MacPherson,
\href{https://www.numdam.org/article/PMIHES_1975__45__101_0.pdf}
{Riemann--Roch for singular varieties},
\emph{Publications Math\'ematiques de l'IH\'ES} \textbf{45} (1975), 101--145.

\bibitem{p3swBorelSerreRR}
A.~Borel and J.-P.~Serre,
\href{https://doi.org/10.24033/bsmf.1500}{Le th\'eor\`eme de Riemann--Roch},
\emph{Bulletin de la Soci\'et\'e Math\'ematique de France}
\textbf{86} (1958), 97--136.

\bibitem{P3SWBuchweitzFlenner2003}
R.-O.~Buchweitz and H.~Flenner,
A semiregularity map for modules and applications to deformations,
\emph{Compositio Mathematica} \textbf{137} (2003), no.~2, 135--210.

\bibitem{P3HHBF}
R.-O.~Buchweitz and H.~Flenner,
The global decomposition theorem for Hochschild \mbox{(co-)homology}
of singular spaces via the Atiyah--Chern character,
\emph{Advances in Mathematics} \textbf{217} (2008), 243--281.
Numbered locators refer to
\href{https://arxiv.org/abs/math/0606730v1}{arXiv:math/0606730v1}.

\bibitem{P3HHCRVdB}
D.~Calaque, C.~A.~Rossi, and M.~Van den Bergh,
\href{https://doi.org/10.4007/annals.2012.176.2.4}
{C\u{a}ld\u{a}raru's conjecture and Tsygan's formality},
\emph{Annals of Mathematics} (2) \textbf{176} (2012), no.~2, 865--923.

\bibitem{P3HHCaldararuII}
A.~C\u{a}ld\u{a}raru,
\href{https://doi.org/10.1016/j.aim.2004.05.012}
{The Mukai pairing---II: the Hochschild--Kostant--Rosenberg isomorphism},
\emph{Advances in Mathematics} \textbf{194} (2005), 34--66.

\bibitem{P3HHCW}
A.~C\u{a}ld\u{a}raru and S.~Willerton,
The Mukai pairing, I: a categorical approach,
\emph{New York Journal of Mathematics} \textbf{16} (2010), 61--98.
\href{https://arxiv.org/abs/0707.2052}{arXiv:0707.2052}.

\bibitem{P3SWChai1985}
C.-L.~Chai,
\href{https://www.cambridge.org/core/books/compactification-of-siegel-moduli-schemes/3372406AEC90801C3DC5EB20D33C3795}
{\emph{Compactification of Siegel Moduli Schemes}},
London Mathematical Society Lecture Note Series, vol.~107,
Cambridge University Press, Cambridge, 1985.

\bibitem{ChaiOortModuli}
C.-L.~Chai and F.~Oort,
\href{https://www.claymath.org/wp-content/uploads/2022/03/cmip08c.pdf}{Moduli of abelian varieties and $p$-divisible groups},
in \emph{Arithmetic Geometry}, H.~Darmon, D.~A.~Ellwood, B.~Hassett,
and Y.~Tschinkel (eds.), Clay Mathematics Proceedings 8,
American Mathematical Society, 2009, 441--536.

\bibitem{P3SWCoelhoPacini}
J.~Coelho and M.~Pacini,
\emph{Abel maps for curves of compact type},
\texttt{arXiv:0904.0039v1} (2009).

\bibitem{P3SWDeligne1968}
P.~Deligne,
Th\'eor\`eme de Lefschetz et crit\`eres de d\'eg\'en\'erescence de suites spectrales,
\emph{Publications Math\'ematiques de l'IH\'ES} \textbf{35} (1968), 107--126.

\bibitem{P3SWDeligneGabber}
P.~Deligne,
Le lemme de Gabber,
\emph{Ast\'erisque} \textbf{127} (1985), 131--150.

\bibitem{Etingof2024Lie}
P.~Etingof,
\emph{Lie groups and Lie algebras},
\href{https://arxiv.org/abs/2201.09397v4}{arXiv:2201.09397v4} (2024).
Section locators refer to this version.

\bibitem{Floccari2026K3}
S.~Floccari,
\href{https://doi.org/10.2140/gt.2026.30.1129}{K3 surfaces associated with varieties of generalized Kummer type},
\emph{Geometry \& Topology} \textbf{30} (2026), no.~3, 1129--1154.
Theorem locators refer to
\href{https://arxiv.org/abs/2501.02315v2}{arXiv:2501.02315v2}.

\bibitem{FloccariFu2026}
S.~Floccari and L.~Fu,
\href{https://doi.org/10.1016/j.matpur.2026.103876}{The Hodge conjecture for Weil fourfolds with discriminant 1 via singular OG6-varieties},
\emph{Journal de Math\'ematiques Pures et Appliqu\'ees} \textbf{210} (2026), article~103876.
Theorem locators refer to
\href{https://arxiv.org/abs/2504.13607v2}{arXiv:2504.13607v2}.

\bibitem{GolyshevLuntsOrlov2001}
V.~Golyshev, V.~Lunts, and D.~Orlov,
Mirror symmetry for abelian varieties,
\emph{Journal of Algebraic Geometry} \textbf{10} (2001), no.~3,
433--496; version used:
\href{https://arxiv.org/abs/math/9812003v2}{arXiv:math/9812003v2}.

\bibitem{P3SWGrauert1960}
H.~Grauert,
\href{https://www.numdam.org/article/PMIHES_1960__5__5_0.pdf}
{Ein Theorem der analytischen Garbentheorie und die Modulr\"aume
komplexer Strukturen},
\emph{Publications Math\'ematiques de l'IH\'ES} \textbf{5} (1960), 5--64.

\bibitem{P3SWGrothendieckHilbert1961}
A.~Grothendieck,
\href{https://www.numdam.org/article/SB_1960-1961__6__249_0.pdf}
{Techniques de construction et th\'eor\`emes d'existence en
g\'eom\'etrie alg\'ebrique IV: les sch\'emas de Hilbert},
\emph{S\'eminaire Bourbaki}, 1960/61, expos\'e~221,
vol.~6 (1961), 249--276.

\bibitem{P3SWHirzebruch1954}
F.~Hirzebruch,
Arithmetic genera and the theorem of Riemann--Roch for algebraic varieties,
\emph{Proc. Natl. Acad. Sci. USA} \textbf{40} (1954), 110--114.
\texttt{doi:10.1073/pnas.40.2.110}.

\bibitem{P3HHHuang}
S.~Huang,
\href{https://doi.org/10.1016/j.jpaa.2021.106673}
{A note on a question of Markman},
\emph{Journal of Pure and Applied Algebra} \textbf{225} (2021),
no.~9, article 106673.
\href{https://arxiv.org/abs/2011.10207v3}{arXiv:2011.10207v3}.

\bibitem{P3SWKodaira1953}
K.~Kodaira,
On a differential-geometric method in the theory of analytic stacks,
\emph{Proc. Natl. Acad. Sci. USA} \textbf{39} (1953), 1268--1273.
\texttt{doi:10.1073/pnas.39.12.1268}.

\bibitem{LehrerZhang2015}
G.~I.~Lehrer and R.~B.~Zhang,
\href{https://doi.org/10.4171/JEMS/558}{The Brauer category and invariant theory},
\emph{Journal of the European Mathematical Society} \textbf{17} (2015), no.~9, 2311--2351.

\bibitem{LehrerZhang2017}
G.~I.~Lehrer and R.~B.~Zhang,
\href{https://doi.org/10.4171/LEM/63-1/2-6}{Invariants of the special orthogonal group and an enhanced Brauer category},
\emph{L'Enseignement Math\'ematique} \textbf{63} (2017), nos.~1--2, 181--200.

\bibitem{Markman2025Secant}
E.~Markman,
\emph{Cycles on abelian $2n$-folds of Weil type from secant sheaves on abelian $n$-folds},
\href{https://arxiv.org/abs/2502.03415v2}{arXiv:2502.03415v2} (2025).
Theorem locators refer to this version.

\bibitem{Markman2026ICM}
E.~Markman,
\href{https://doi.org/10.1137/25M1803796}{Secant sheaves and Weil classes on abelian varieties},
in \emph{Proceedings of the International Congress of Mathematicians 2026},
vol.~3, S.~Friedlander and Y.~Tschinkel (eds.), SIAM, 2026, 586--602.

\bibitem{MilneAG2022}
J.~S.~Milne,
\emph{Algebraic Groups: The Theory of Group Schemes of Finite Type over a Field},
Cambridge Studies in Advanced Mathematics 170, Cambridge University Press, 2017.
Theorem locators refer to the
\href{https://www.jmilne.org/math/Books/iAG2022.pdf}{corrected online version (2022)}.

\bibitem{MoonenZarhin1999}
B.~J.~J.~Moonen and Yu.~G.~Zarhin,
\href{https://doi.org/10.1007/s002080050333}{Hodge classes on abelian varieties of low dimension},
\emph{Mathematische Annalen} \textbf{315} (1999), no.~4, 711--733.
Theorem locators refer to
\href{https://arxiv.org/abs/math/9901113v2}{arXiv:math/9901113v2}.

\bibitem{P3SWMukai1981}
S.~Mukai,
Duality between \(D(X)\) and \(D(\widehat X)\) with its application to Picard sheaves,
\emph{Nagoya Math. J.} \textbf{81} (1981), 153--175.
\texttt{doi:10.1017/S002776300001922X}.

\bibitem{OpenAICM2026}
OpenAI,
\emph{The rational Hodge conjecture for CM abelian varieties},
OpenAI Math Release preprint
\href{https://github.com/openai/math/blob/main/preprints/The-rational-Hodge-conjecture-for-CM-abelian-varieties-October-6-2026/paper.pdf}{OAI:The-rational-Hodge-conjecture-for-CM-abelian-varieties-October-6-2026},
2026.

\bibitem{PetersSterk2024QuadraticForms}
C.~Peters and H.~Sterk,
\href{https://www-fourier.univ-grenoble-alpes.fr/~peters/Books/QuadraticForms/QuadForms.pdf}{\emph{Symmetric and Quadratic Forms, with Applications to Coding Theory, Algebraic Geometry and Topology}},
author monograph draft, version June 2024, dated 16 September 2024.
The corollary locator refers to this version.

\bibitem{P3SWPourcin1969}
G.~Pourcin,
Th\'eor\`eme de Douady au-dessus de \(S\),
\emph{Annali della Scuola Normale Superiore di Pisa, Scienze Fisiche e Matematiche},
Serie~3, \textbf{23} (1969), no.~3, 451--459.

\bibitem{Schoen1988}
C.~Schoen,
\href{https://www.numdam.org/item/CM_1988__65_1_3_0/}{Hodge classes on self-products of a variety with an automorphism},
\emph{Compositio Mathematica} \textbf{65} (1988), no.~1, 3--32.

\bibitem{Schoen1998Addendum}
C.~Schoen,
\href{https://doi.org/10.1023/A:1000566205021}{Addendum to: Hodge classes on self-products of a variety with an automorphism},
\emph{Compositio Mathematica} \textbf{114} (1998), no.~3, 329--336.

\bibitem{Schrijver2008}
A.~Schrijver,
\href{https://doi.org/10.1016/j.jalgebra.2007.10.039}{Tensor subalgebras and first fundamental theorems in invariant theory},
\emph{Journal of Algebra} \textbf{319} (2008), no.~3, 1305--1319.
The section locator refers to the
\href{https://homepages.cwi.nl/~lex/files/tensfft_ja.pdf}{author version}.

\bibitem{p3swSerreDuality}
J.-P.~Serre,
\href{https://doi.org/10.5169/seals-23275}{Un th\'eor\`eme de dualit\'e},
\emph{Commentarii Mathematici Helvetici}
\textbf{29} (1955), 9--26.

\bibitem{P3SWSerreGAGA1956}
J.-P.~Serre,
G\'eom\'etrie alg\'ebrique et g\'eom\'etrie analytique,
\emph{Annales de l'Institut Fourier} \textbf{6} (1956), 1--42.

\bibitem{Serre1979HodgeTate}
J.-P.~Serre,
\href{https://www.numdam.org/item/AST_1979__65__155_0/}{Groupes alg\'ebriques associ\'es aux modules de Hodge--Tate},
\emph{Ast\'erisque} \textbf{65} (1979), 155--188.

\bibitem{P3SWStacks}
The Stacks Project Authors,
\emph{The Stacks project},
\texttt{https://stacks.math.columbia.edu}.
The cited tag and theorem locators were consulted on 30 September 2026.

\bibitem{vanGeemenWeilNotes}
B.~van Geemen,
\href{https://doi.org/10.1007/978-3-540-49046-3_5}{An introduction to the Hodge conjecture for abelian varieties},
in \emph{Algebraic Cycles and Hodge Theory} (Torino, 1993),
Lecture Notes in Mathematics 1594, Springer, 1994, 233--252.
Lemma and section locators refer to the
\href{https://sites.unimi.it/vangeemen/cime.pdf}{author version}.

\bibitem{Weil1979}
A.~Weil,
Abelian varieties and the Hodge ring,
in \emph{Collected Papers}, vol.~III,
Springer-Verlag, New York, 1979, 421--429 (essay~1977c).

\end{thebibliography}
\end{document}
